\documentclass[11pt]{article}
\usepackage[a4paper,margin=27mm]{geometry}
\usepackage{amsmath,amssymb,amsthm,mathtools}
\usepackage[T1]{fontenc}
\usepackage{lmodern}
\usepackage{xcolor}
\newtheorem{theorem}{Theorem}
\newtheorem{lemma}[theorem]{Lemma}
\newtheorem{proposition}[theorem]{Proposition}
\newtheorem{corollary}[theorem]{Corollary}
\newcommand{\PROVED}{\textcolor{green!45!black}{\textbf{PROVED}}}
\newcommand{\OPEN}{\textcolor{red!70!black}{\textbf{OPEN}}}
\newcommand{\AUDIT}{\textcolor{orange!75!black}{\textbf{AUDIT}}}
\newcommand{\Ran}{\operatorname{Ran}}
\newcommand{\Ker}{\operatorname{Ker}}
\newcommand{\old}{\mathrm{old}}
\newcommand{\Hod}{\mathrm{Hod}}
\newcommand{\comp}{\mathrm{comp}}
\newcommand{\cC}{\mathscr C}
\newcommand{\cZ}{\mathscr Z}
\newcommand{\cL}{\mathscr L}
\title{Row (d): Gamma--Fourier--Tate amplitude audit\\
v94 (typed leakage and Gamma--arithmetic continuation audit)}
\author{}
\date{12 August 2026}

\begin{document}
\maketitle

\begin{abstract}
This note attacks the amplitude identity left by v83--v86 and reconciles the
result with the supplied files \texttt{main.tex} and
\texttt{row-d-local-analysis.tex}.  It constructs a
source-defined doubled logarithmic Tate port and proves that its divergence is
exactly $-B_Q$, thereby closing G39 at every finite threshold.  It also derives
the connection term lost when the Fourier--Gamma identity is differentiated
after a moving compression.  The calculation does not prove the required
positive energy estimate G40.  A new exact balanced factorization of the
localized primitive operator is proved below; it supplies explicit meanings
for the previously abstract factors $X_T,Y_T$.  It also shows precisely why
the local Gamma--Tate normalizations do not by themselves imply the missing
unit contraction.  New critical gcd and Archimedean Gamma GNS boundaries,
an exact Gamma ratio--scale formula deriving $m_0$, and the coherent completed
Euler--Gamma metric are constructed.  A domain theorem rules out unweighted
Gamma deconvolution on every nonzero compact-window source.  The critical
form boundary is then constructed and identified with divisor incidence, the
gcd GNS state and M\"obius whitening.  Explicit Julia, beta--Gamma and Euler
Markov colligations derive $g_\Gamma$, $m_0$ and every prime-power coefficient
as metric/defect rates.  The apparent opposition of the Euler and Gamma
Markov arrows is removed by the Bayes adjoint of the Euler conditional
expectation; its tensor product with the Gamma map is an explicit joint
Markov contraction with a conditional-variance defect.  A physical Gamma ODE realization further absorbs
$m_0$ into an explicit four-state head and turns every Gamma/shift commutator
into positive rank-one storage.  Every prime-power channel is resummed into
one Julia filter and one boundary port per prime.  The resulting physical
angular operator is explicit and is certified strictly contractive for
$0<T\le\log2$.  Any later failure is reduced to a first-loss zero mode and,
equivalently, to a homogeneous exterior-leakage equation for a clamped Sonin
potential.  Its Poisson defect is further reduced to the additive-Fourier
eigenspace opposite to the parity of the physical mode.  Equality of the two
meromorphic half-lift germs, exact differentiation of that equality, and a
finite-pole extended-S rigidity theorem are proved.  What remains is to
derive the corresponding physical boundary/reciprocal co-Poisson realization
from the zero-mode Gamma equation.  Old compatibility of the physical
continuum factors under window enlargement is proved by explicit unitary
neutral birth channels; this is distinguished from the impossible finite-port
identification.  Uniqueness for the explicit leakage/boundary-layer problem
and the positive continuum interconnection remain the final bridge; neither
condition $G$ nor row (d) is declared closed.
\end{abstract}

\section{Input from v83--v86}

The following identities are taken as the proved input of the preceding
notes.  At the unweighted endpoint let
\[
 G'=R^*\cL R-\frac12(QG+GQ),\qquad Gv=S\delta_1,
 \qquad Q\delta_d=(\log d)\delta_d .
\]
Let $P_{\old}$ be the old-coordinate projection and let $M_0$ be the old
block of $G$.  Differentiation of the harmonic normal equation gives
\begin{equation}
 P_{\old}G'v=M_0h'.                                      \tag{1}
\end{equation}
The natural Hodge current and its old analysis operator are
\[
 W_{\Hod}=\cL^{1/2}Rv,
 \qquad H_0=\cL^{1/2}RP_{\old}.
\]
Consequently
\begin{equation}
 H_0^*W_{\Hod}=M_0h'+B_Q,
 \qquad B_Q:=\frac12P_{\old}GQv.                         \tag{2}
\end{equation}
In the Zeta gauge, with $m=\cZ^{-1}\delta_1$, this is
\begin{equation}
 B_Q=\frac S2P_{\old}
 \left[-\cC\delta_1+\cZ\cC^*m\right].                  \tag{3}
\end{equation}
The first summand is the directed von--Mangoldt forcing; the second is its
adjoint return through the Zeta triangular system.

The desired archimedean cancellation is
\begin{equation}
 H_{\infty,0}^*W_\infty=-B_Q.                            \tag{6}
\end{equation}

\paragraph{Status.} Equations (1)--(3) are \PROVED{} within the stated
finite-dimensional conventions of v83--v86.  Equation (6) is proved below in
Theorem G39, equations (27e)--(27h).

\section{Exact solvability of the amplitude equation}

The next statement is elementary, but it is the correct non-circular test for
any proposed Gamma--Tate construction.

\begin{theorem}[Amplitude range theorem --- \PROVED]
Let $H:\mathcal K_0\to\mathcal K_\infty$ and
$B:E\to\mathcal K_0$ be bounded operators (all spaces may be finite
dimensional).  The equation
\[
 H^*W=-B,\qquad W:E\to\mathcal K_\infty,
\]
has a solution if and only if
\begin{equation}
 \Ran B\subseteq\Ran H^*.                                \tag{4}
\end{equation}
Equivalently, in finite dimension,
\begin{equation}
 \Ker H\subseteq\Ker B^*.                               \tag{5}
\end{equation}
When it exists, the unique solution whose range is orthogonal to $\Ker H^*$ is
\begin{equation}
 W_{\min}=-(H^*)^\dagger B,                              \tag{7}
\end{equation}
and every solution is $W_{\min}+K$ with $\Ran K\subseteq\Ker H^*$.  Moreover
\begin{equation}
 W^*W\ge W_{\min}^*W_{\min}
 =B^*(H^*H)^\dagger B.                                  \tag{8}
\end{equation}
\end{theorem}

\begin{proof}
The equation is solvable precisely when every column of $B$ lies in
$\Ran H^*$.  The orthogonal-complement identity
$(\Ran H^*)^\perp=\Ker H$ gives (5).  Moore--Penrose inversion gives (7).
The general solution differs from it by a map into $\Ker H^*$; this summand is
orthogonal to $W_{\min}$, proving (8).
\end{proof}

Applied to (6), the first genuinely new target is therefore
\begin{equation}
 \boxed{\Ran B_Q\subseteq\Ran H_{\infty,0}^*.}           \tag{9}
\end{equation}
This is an amplitude/range statement.  It is not implied by equality of the
associated quadratic-form generators.

\section{Why the known Fourier--Gamma identity does not yet give (6)}

The proved Fourier--Gamma bridge has the form
\begin{equation}
 G_\Gamma-m_0I=-(K+X).                                  \tag{10}
\end{equation}
It identifies an archimedean generator.  The desired equation (6), however,
contains an old/new matrix coefficient of a factorization of that generator.
Factorizations are not determined by their Gram operator: if
$G_\Gamma-m_0I=J^*\Sigma J$, then $UJ$ gives the same Gram identity for every
unitary $U$ commuting with the signature, while its compressed amplitudes can
change.

There is a second, independent issue: the relevant harmonic/old projection
moves under the deformation.

\begin{lemma}[Moving-compression connection --- \PROVED]
Let $U(t)$ be a differentiable unitary family, $A=U(0)^*U'(0)$, and
$P(t)=U(t)P U(t)^*$.  Then
\begin{equation}
 P'(0)=U(0)[A,P]U(0)^*.                                  \tag{11}
\end{equation}
For a differentiable operator family $L(t)$,
\begin{align}
 \frac d{dt}\bigg|_0 P(t)L(t)P(t)
  ={}&P'(0)L(0)P(0)+P(0)L'(0)P(0)\notag\\
    &+P(0)L(0)P'(0).                                    \tag{12}
\end{align}
Thus a formula for $L'(0)$ alone does not determine the differentiated
compressed block.
\end{lemma}

\begin{proof}
Differentiate $U(t)PU(t)^*$ and then apply the product rule to the compression.
\end{proof}

In the present problem the two terms containing $P'(0)$ are precisely the
kind of connection terms capable of producing the adjoint-return summand
$\cZ\cC^*m$ in (3).  Therefore a valid derivation of (6) from Tate must specify
the boundary embeddings and prove the differentiated compressed intertwining,
not only cite (10).

\section{A sharp conditional completion}

Assume a source-defined Gamma--Tate factor $H_{\infty,0}$ is given.  Then (9)
constructs, without choice,
\[
 W_{\infty,\min}=-(H_{\infty,0}^*)^\dagger B_Q .
\]
Combining it with (2) yields
\[
 (H_0^*\oplus H_{\infty,0}^*)
 (W_{\Hod}\oplus W_{\infty,\min})=M_0h'.               \tag{13}
\]
This proves the required divergence identity.  To obtain the Pythagorean
storage identity one still needs the independently derived energy budget
\begin{equation}
 S_M'\ge B_Q^*(H_{\infty,0}^*H_{\infty,0})^\dagger B_Q
       +\text{the finite-part minimum energy}.           \tag{14}
\end{equation}
Defining a residual as the square root of the difference in (14) would be
circular unless positivity of that difference has already been proved from
the source network.

\section{Source audit: Hilbertian transport is not yet unitary amplitude}

The semilocal Sonin theorem of Connes--Consani--Moscovici supplies a
hilbertian isomorphism
\[
 \theta_S:{\bf S}_\lambda(\mathbb R,e_\infty)
 \longrightarrow {\bf S}_\lambda(X_S,\alpha).
\]
This is enough to transport the underlying topological Hilbert space and to
prove stability when places are added.  It is not, as stated, an isometry for
the fixed physical metrics entering row (d).  In the same discussion the
authors explicitly note that the finite set $S$ plays a key role in fixing the
inner product of the corresponding space of entire functions.

This distinction has an exact operator consequence.  Put
\begin{equation}
 D_S:=\theta_S^*\theta_S>0.                               \tag{15}
\end{equation}
Then the canonically normalized isometry is
\begin{equation}
 U_S:=\theta_SD_S^{-1/2},\qquad U_S^*U_S=I.              \tag{16}
\end{equation}
For a differentiable deformation $S=S(t)$ one has
\begin{equation}
 U_S'=\theta_S'D_S^{-1/2}
       +\theta_S(D_S^{-1/2})'.                           \tag{17}
\end{equation}
The second summand is a metric-connection term.  It vanishes only if
$D_S'=0$, i.e. only if the raw hilbertian transport is already isometric along
the deformation.

\begin{proposition}[Metric correction identity --- \PROVED]
Let $D(t)=\theta(t)^*\theta(t)>0$ and
$U(t)=\theta(t)D(t)^{-1/2}$.  Then $U(t)$ is isometric and
\begin{equation}
 U^*U'=D^{-1/2}\theta^*\theta'D^{-1/2}
       +D^{1/2}(D^{-1/2})',                              \tag{18}
\end{equation}
whose Hermitian part is zero.  Hence the Hermitian part of the raw logarithmic
connection $D^{-1/2}\theta^*\theta'D^{-1/2}$ is cancelled exactly by the
derivative of the metric normalization.
\end{proposition}

\begin{proof}
Equation (18) follows by the product rule.  Differentiating $U^*U=I$ gives
$U'^*U+U^*U'=0$.
\end{proof}

The correction in (17) is computable without assuming that $D$ and $D'$ commute.
If $X=(D^{-1/2})'$, differentiation of
$D^{-1/2}DD^{-1/2}=I$ gives the Sylvester equation
\begin{equation}
 XD^{1/2}+D^{1/2}X=-D^{-1/2}D'D^{-1/2}.                  \tag{19}
\end{equation}
Since $D^{1/2}>0$, (19) has a unique bounded selfadjoint solution for
selfadjoint $D'$.  Thus the next calculation has no unitary gauge freedom at
the metric-correction level: once the semilocal Gram derivative $D'$ is known,
the missing connection $X$ is uniquely fixed.

Thus the published stability theorem cannot be used to infer a norm-one
storage law without explicitly carrying $D_S^{-1/2}$ and its derivative.
The missing correction is of exactly the same structural type as the moving
harmonic-lift term found in v83--v86.  This rules out the shortcut
``Sonin stability $\Rightarrow$ (6)'', while leaving open the possibility that
the fully normalized connection in (17) produces $-B_Q$.

\section{Exact CCM multiplier and its normalized connection}

The source definitions permit the metric correction to be evaluated
explicitly.  On the multiplicative Fourier line put
\begin{equation}
 a_{S,t}(s):=\prod_{p\in S_f}
 \left(1-p^{-t/2-is}\right).                             \tag{20}
\end{equation}
The dual Sonin map $\theta_{S,t}$ is multiplication by $a_{S,t}$, while the
Zeta map $\eta_{S,t}$ is multiplication by
$\overline{a_{S,t}}^{-1}$ (for real $s$).  Hence
\begin{equation}
 \theta_{S,t}^*\eta_{S,t}=I,
 \qquad
 D_{S,t}=\theta_{S,t}^*\theta_{S,t}=M_{|a_{S,t}|^2}.     \tag{21}
\end{equation}
In particular, the polar normalization is not abstract:
\begin{equation}
 U_{S,t}=\theta_{S,t}D_{S,t}^{-1/2}
 =M_{a_{S,t}/|a_{S,t}|}.                                \tag{22}
\end{equation}

Define the finite Euler logarithmic multiplier
\begin{equation}
 c_{S,t}(s):=\sum_{p\in S_f}(\log p)
 \frac{p^{-t/2-is}}{1-p^{-t/2-is}}.                     \tag{23}
\end{equation}

\begin{theorem}[Normalized Euler connection --- \PROVED]
On any compact $t$-interval on which the factors in (20) do not vanish on the
real Fourier line,
\begin{align}
 \theta_{S,t}^{-1}\partial_t\theta_{S,t}
   &=\frac12M_{c_{S,t}},                                 \tag{24}\\
 D_{S,t}^{-1}D_{S,t}'
   &=\frac12M_{c_{S,t}+\overline{c_{S,t}}},              \tag{25}\\
 U_{S,t}^*U_{S,t}'
   &=\frac14M_{c_{S,t}-\overline{c_{S,t}}}.              \tag{26}
\end{align}
Thus metric normalization cancels exactly the Hermitian Euler component and
retains exactly its skew-Hermitian component.
\end{theorem}

\begin{proof}
For one factor, with $r_p(t)=p^{-t/2}$,
\[
 \partial_t\log(1-r_p(t)p^{-is})
 =\frac{\log p}{2}\frac{p^{-t/2-is}}
 {1-p^{-t/2-is}}.
\]
Summation gives (24).  Equations (25) and (26) follow by differentiating
$|a|^2$ and $a/|a|$, respectively.
\end{proof}

This theorem sharpens the target.  The normalized CCM transport contributes
only the oriented difference $c-\bar c$.  Therefore the symmetric Hodge
energy and the amplitude $B_Q$ cannot be read from this phase connection
alone.  They can arise only after the old/new (Sonin/cutoff) compression and
the two moving boundary embeddings are differentiated.  Equivalently, the
remaining source identity is now the explicit compressed formula
\begin{equation}
 P_{\old}\left(
 j_{\old}'{}^*U^*j_{\new}
 +j_{\old}^*U^*U'j_{\new}
 +j_{\old}^*U^*j_{\new}'
 \right)
 =-B_Q,                                                  \tag{27}
\end{equation}
with $U^*U'$ fixed by (26).  There is no longer an unspecified bulk
connection in (27); the only unknown data are the differentiated Tate/cutoff
boundary maps.

\paragraph{Source boundary.}
The cited CCM theorem proves (at fixed finite set of places $S$ and fixed
cutoff $\lambda$) the stability of the Sonin space and the duality
$\theta_S^*\eta_S=I$.  It does not state a differentiable family of Tate
boundary embeddings $j_{\old}(t),j_{\new}(t)$, nor formula (27).  Consequently
those maps cannot be silently imported from that theorem.  The finite
divisibility model below constructs the required logarithmic boundary port
directly, without attributing it to the fixed-parameter CCM theorem.

\section{Closure of G39: the doubled logarithmic Tate port}

We now give a source-defined realization of the boundary amplitude.  Let
$\mathcal H_N=\operatorname{span}\{\delta_d:1\le d\le N\}$ and let
$S_n\delta_d=\delta_{nd}$ when $nd\le N$ and zero otherwise.  Put
\begin{align}
 \cZ_t&:=\sum_{n\le N}n^{-t/2}S_n,\notag\\
 \cC_t&:=\sum_{n\le N}\Lambda(n)n^{-t/2}S_n,
 \qquad Q\delta_d=(\log d)\delta_d.                     \tag{27a}
\end{align}
All sums and products are in the finite divisibility algebra, so there are no
domain or convergence issues.

\begin{lemma}[Finite logarithmic connection --- \PROVED]
One has
\begin{align}
 [Q,\cZ_t]&=\cZ_t\cC_t=\cC_t\cZ_t,                    \tag{27b}\\
 \cZ_t'&=-\frac12\cZ_t\cC_t,                           \tag{27c}\\
 (\cZ_t^{-1})'&=\frac12\cC_t\cZ_t^{-1}.                \tag{27d}
\end{align}
Thus $\cC_t=-2\cZ_t^{-1}\cZ_t'$ is the logarithmic
derivative of the Zeta transport, and $\cC_t^*$ is its adjoint-return
connection.
\end{lemma}

\begin{proof}
Since $[Q,S_n]=(\log n)S_n$, the coefficient of $S_k$ in
$[Q,\cZ_t]$ is $(\log k)k^{-t/2}$.  The coefficient of $S_k$ in
$\cZ_t\cC_t$ is
\[
 k^{-t/2}\sum_{d\mid k}\Lambda(d)
 =k^{-t/2}\log k.
\]
This proves (27b).  Direct differentiation gives (27c), and differentiating
$\cZ_t^{-1}\cZ_t=I$ gives (27d).
\end{proof}

At the critical endpoint suppress the subscript $t$, put
\[
 m=\cZ^{-1}\delta_1,\qquad
 \mathcal K_\infty^{(1)}:=\mathcal H_N\oplus\mathcal H_N,
\]
and define, before invoking any sign condition,
\begin{align}
 W_\infty
 &:=\sqrt{\frac S2}
 \binom{\cC\delta_1}{\cC^*m},                           \tag{27e}\\
 H_{\infty,0}x
 &:=\sqrt{\frac S2}
 \binom{P_{\old}^*x}{-\cZ^*P_{\old}^*x}.              \tag{27f}
\end{align}
The first component of (27e) is the differentiated Zeta/Möbius incoming
port; the second is its adjoint return.  Formula (27f) is the graph analysis
of the dual Tate termination $[I,-\cZ]$.

\begin{theorem}[G39, exact Gamma--Tate amplitude --- \PROVED]
The doubled logarithmic port (27e)--(27f) is source-defined by the finite
Euler transport and satisfies
\begin{equation}
 \boxed{H_{\infty,0}^*W_\infty=-B_Q}.                   \tag{27g}
\end{equation}
Equivalently, it supplies exactly the missing divergence in (2), and hence
\begin{equation}
 (H_0^*\oplus H_{\infty,0}^*)
 (W_{\Hod}\oplus W_\infty)=M_0h'.                       \tag{27h}
\end{equation}
No positivity of $G$, no Moore--Penrose choice, and no residual square root is
used in this construction.
\end{theorem}

\begin{proof}
Taking the adjoint of (27f) and applying it to (27e) gives
\[
 H_{\infty,0}^*W_\infty
 =\frac S2P_{\old}
 \left(\cC\delta_1-\cZ\cC^*m\right).
\]
By the already proved v84 formula (3), the right-hand side is exactly $-B_Q$.
Adding this equality to (2) proves (27h).
\end{proof}

The construction also explains the two moving-boundary contributions in
(27): they are the two orientations of the logarithmic derivative of the
dual pair $\cZ_t^{-1}$ and $\cZ_t^{-*}$.  The bulk normalized phase accounts
for their skew part, while the graph termination $[I,-\cZ]$ retains their
oriented difference.  Thus G39 is closed at every finite threshold.  This
does not assert the Hilbert-space energy bound G40.

\section{Reverse-engineered completed boundary object}

We now construct the object required by the proof without presupposing its
positivity.  This separates existence from the load-bearing sign theorem.
Let
\[
 A_0=\mathcal G_0^*\mathcal G_0\ge0,
 \qquad q=M_0h',\qquad \beta=S_M'.
\]
All operators below are understood on the finite threshold space; the same
formulas hold for closed ranges in the Hilbert-space setting.

\begin{theorem}[Canonical completed boundary colligation --- \PROVED]
Assume $q\in\Ran A_0^{1/2}$ and define
\begin{equation}
 w_0:=\mathcal G_0A_0^\dagger q,
 \qquad
 \Delta_T:=\beta-q^*A_0^\dagger q.                      \tag{28}
\end{equation}
Let $\Delta_T=J_T|\Delta_T|$ be its polar decomposition, where
$J_T=\operatorname{sgn}(\Delta_T)$ on
$\overline{\Ran|\Delta_T|}$, and put
\begin{equation}
 \mathcal K_T^{\rm bd}:=\overline{\Ran|\Delta_T|},
 \qquad r_T:=|\Delta_T|^{1/2}.                           \tag{29}
\end{equation}
Then the map
\begin{equation}
 \mathcal W_Te:=w_0e\oplus r_Te                         \tag{30}
\end{equation}
into the Krein boundary space
$\overline{\Ran\mathcal G_0}\oplus(\mathcal K_T^{\rm bd},J_T)$
satisfies the two exact identities
\begin{align}
 \mathcal G_0^*P_0\mathcal W_T&=q,                      \tag{31}\\
 \mathcal W_T^*(I\oplus J_T)\mathcal W_T&=\beta.        \tag{32}
\end{align}
It is minimal: any other realization of (31)--(32) contains this one up to a
Krein-isometric summand orthogonal to the generated boundary space.
\end{theorem}

\begin{proof}
Since $q\in\Ran A_0^{1/2}$, Moore--Penrose calculus gives
\[
 \mathcal G_0^*w_0
 =A_0A_0^\dagger q=q,
 \qquad
 w_0^*w_0=q^*A_0^\dagger q.
\]
Moreover $r_T^*J_Tr_T=\Delta_T$.  Adding the last two identities proves
(31)--(32).  Minimality follows from the minimum-norm property of $w_0$ and
the minimal polar realization of the selfadjoint form $\Delta_T$.
\end{proof}

This is the desired inverse-engineered object: it always exists once the
previously proved range condition holds, but in general its boundary metric is
Krein rather than Hilbert.

\begin{theorem}[Exact positivity threshold --- \PROVED]
The following assertions are equivalent:
\begin{enumerate}
\item the canonical boundary metric in (29) is positive, $J_T=I$;
\item $\Delta_T\ge0$;
\item there exist Hilbert-space operators $W,R$ such that
\[
 \mathcal G_0^*W=q,
 \qquad W^*W+R^*R=\beta;
\]
\item the block completion is positive,
\begin{equation}
 \begin{pmatrix}A_0&q\\q^*&\beta\end{pmatrix}\ge0;      \tag{33}
\end{equation}
\item the v85 Pythagorean identity, hence the corresponding threshold case of
condition $G$, holds.
\end{enumerate}
\end{theorem}

\begin{proof}
The equivalence of (1) and (2) is the definition of $J_T$.  If (2) holds take
$W=w_0$ and $R=\Delta_T^{1/2}$.  Conversely minimum norm gives
$W^*W\ge q^*A_0^\dagger q$, so any Hilbert realization implies
$\Delta_T\ge0$.  The equivalence with (33) is the generalized Schur-complement
criterion, including $q\in\Ran A_0^{1/2}$.  The last equivalence is precisely
the reduction proved in v85--v86.
\end{proof}

The theorem shows both the power and the limit of pure inverse engineering.
It constructs the unique missing port and its exact metric, but proving that
the metric is Hilbert rather than Krein is exactly the remaining content of
$G$.  Renaming $r_T$ a ``Tate residual'' would therefore be circular.

\subsection{The genuinely new theorem sufficient to close row (d)}

The preceding construction identifies a source-level statement which is
strictly stronger than a formal completion and is independently falsifiable:

\begin{quote}
\textbf{Tate boundary positivity theorem (OPEN).}
There is an isomorphism
\[
 \Phi_T:\mathcal K_T^{\rm bd}\longrightarrow
 \mathcal T_T^{\rm out}
\]
onto a boundary space obtained from the two terminated Tate evaluations and
the Gamma storage tail, constructed before using $\Delta_T$, such that
\begin{equation}
 \langle\Phi_Tx,\Phi_Ty\rangle_{\mathcal T_T^{\rm out}}
 =\langle J_Tx,y\rangle.                                \tag{34}
\end{equation}
The target carries its ordinary positive $L^2$/Gamma-storage inner product.
\end{quote}

If (34) is proved from Fourier--Gamma--Tate, then its left-hand side is
positive, hence $J_T=I$ on the minimal boundary space.  The preceding theorem
then gives $G$ immediately.  Unlike defining
$R=\Delta_T^{1/2}$, statement (34) is non-circular only when
$\mathcal T_T^{\rm out}$ and $\Phi_T$ are given by explicit source formulas
which do not use $\Delta_T$ or its sign.

The proved amplitude identity (27g) is the divergence part of (34).  The missing
norm identity is
\begin{equation}
 \Phi_T^*\Phi_T=J_T                                    \tag{35}
\end{equation}
on the generated boundary space.  Thus (27g) alone is necessary but not
sufficient; (35) is exactly what prevents an indefinite completion from being
mistaken for a conservative Hilbert-space network.

\section{Intermediate status ledger (superseded by the final ledger below)}

\subsection*{PROVED}
\begin{itemize}
\item The v83 Möbius/von--Mangoldt forcing and the v84 excess divergence
      formula (2)--(3), subject to their stated finite Zeta-gauge definitions.
\item The range/kernel criterion (4)--(5), the canonical current (7), and its
      exact minimum energy (8).
\item The moving-compression connection formula (11)--(12).
\item The metric-normalization correction (15)--(18).
\item The exact CCM multiplier identities (20)--(26).
\item G39: the source-defined doubled logarithmic port (27e)--(27f), the exact
      amplitude identity (27g), and the completed divergence (27h).
\item The canonical Krein boundary colligation (28)--(32) and the exact
      equivalence between its positivity, block positivity (33), the v85
      storage identity, and $G$.
\item The amplitude range inclusion (9) for the constructed port, since
      $-B_Q=H_{\infty,0}^*W_\infty$, and the composite divergence (13)/(27h).
\end{itemize}

\subsection*{OPEN}
\begin{itemize}
\item The non-circular source energy balance (14), hence the v85 Pythagorean
      identity, condition $G$, and row (d).
\item The Tate boundary positivity theorem (34)--(35), with a source-defined
      positive target and an intertwiner independent of $\Delta_T$.
\end{itemize}

\subsection*{AUDIT / forbidden shortcuts}
\begin{itemize}
\item Equation (10) is an operator-generator identity, not an amplitude
      identity; it does not by itself imply (6).
\item Möbius/Zeta adjoint-inverse triangularity does not determine the moving
      old/new connection block.
\item A hilbertian Sonin isomorphism is not automatically an isometry in the
      row-(d) physical metric.  The derivative of $D_S^{-1/2}$ cannot be
      discarded.
\item The parameter deformation and the differentiated Tate boundary maps are
      not supplied by the fixed-$S$, fixed-$\lambda$ Sonin stability theorem.
\item Choosing $W_\infty$ by Moore--Penrose inversion is legitimate only after
      (9); claiming the desired budget for that choice is equivalent to the
      remaining inequality and cannot serve as its proof.
\item The inverse-engineered boundary space exists canonically, but its metric
      is a priori Krein.  Declaring it positive or calling
      $|\Delta_T|^{1/2}$ a Hilbert residual would assume $G$.
\item No statement in this note declares $G$ or row (d) proved.
\end{itemize}

\section{Next exact calculation}
With G39 closed, the next calculation is exclusively the G40 energy theorem.
For the explicitly constructed port it asks whether its source norm, together
with the finite Hodge current and the Gamma storage tail, telescopes to
\[
 S_M'=(M_0h')^*A_0^\dagger(M_0h')+R_T^*R_T
\]
with $R_T$ constructed before assuming the sign of the difference.  The
divergence is no longer an unknown.  The remaining task is to identify the
positive Gamma--Tate storage realization (34)--(35); this is G40, not G39.

\section{G40 bridge audit: KYP completion and the circularity barrier}

The completion of G39 supplies an exact current, but amplitude and passivity
are logically different.  The following theorem determines, without guesswork,
the least energy compatible with the proved amplitude.

Put $H:=H_0\oplus H_{\infty,0}$ and
$q:=M_0h'$.  By (27h), $q\in\Ran H^*$.  Define
\begin{equation}
 W_{\min}:=(H^*)^\dagger q,\qquad
 E_{\min}:=q^*(H^*H)^\dagger q.                         \tag{36}
\end{equation}

\begin{theorem}[Minimal passive completion test --- \PROVED]
For the source budget $\beta=S_M'$, the following are equivalent:
\begin{enumerate}
\item there are a Hilbert-space current $W$ and residual $R$ satisfying
\[
 H^*W=q,\qquad \beta=W^*W+R^*R;
\]
\item $\beta-E_{\min}\ge0$;
\item the KYP block
\begin{equation}
 \mathfrak K_T:=
 \begin{pmatrix}H^*H&q\\q^*&\beta\end{pmatrix}          \tag{37}
\end{equation}
is positive.
\end{enumerate}
If these conditions hold, the canonical solution is
\begin{equation}
 W=W_{\min},\qquad R=(\beta-E_{\min})^{1/2}.             \tag{38}
\end{equation}
\end{theorem}

\begin{proof}
Moore--Penrose minimality gives $W^*W\ge E_{\min}$ for every solution of
$H^*W=q$.  Hence a passive completion implies (2).  Conversely (2) gives
(38).  Equivalence with (3) is the generalized Schur-complement theorem,
using $q\in\Ran H^*$, which is already supplied by G39.
\end{proof}

This theorem is a useful advance: G39 has discharged the entire range part of
the generalized Schur criterion.  G40 consists only of the remaining scalar
operator inequality
\begin{equation}
 \boxed{S_M'-q^*(H^*H)^\dagger q\ge0.}                  \tag{39}
\end{equation}
However, inverse engineering cannot prove (39) by defining the residual in
(38): the square root exists as a Hilbert operator if and only if (39) is
already true.

\begin{theorem}[No-free-bridge theorem --- \PROVED]
Suppose a proposed G40 bridge is constructed only from the already fixed data
$(H,q,\beta)$ and is required to satisfy the exact amplitude and energy
identities.  Then its existence as a Hilbert-space colligation is equivalent
to (39).  In particular, polar decomposition, Julia completion, graph
normalization, Moore--Penrose lifting, or adjoining a formal residual cannot
establish G40 unless one independently proves positivity of (39).
\end{theorem}

\begin{proof}
Every such bridge gives a pair $(W,R)$ in item (1) of the preceding theorem,
so it implies (39).  Conversely (39) constructs the bridge by (38).  Each of
the listed procedures is a realization of this same completion and therefore
cannot change the equivalence.
\end{proof}

Thus the non-circular content still required from Gamma--Tate is an
\emph{independently positive} source formula for the left side of (39), for
example
\begin{equation}
 S_M'-q^*(H^*H)^\dagger q
 =\int_{\Omega_T}K_T(\omega)^*K_T(\omega)\,d\mu_T(\omega)
  +E_{0,T}^*E_{0,T}+E_{1,T}^*E_{1,T},                  \tag{40}
\end{equation}
where $K_T,E_{0,T},E_{1,T}$ must be given explicitly from the Gamma tail and
the two Tate evaluations without using the left side or its spectral
decomposition.  Formula (40), not a formal completion, is the precise new
mathematical theorem which would prove G40.

\subsection{Compatibility with the fixed old factor}

There is one indispensable bookkeeping point.  The analysis operator $H$ in
(36) is the explicit amplitude realization of G39.  The old factor used in
the threshold induction is fixed independently by
\[
 A_0=\mathcal G_0^*\mathcal G_0.
\]
One may replace $\mathcal G_0$ by $H$ in the storage argument only after
proving the Gram compatibility
\begin{equation}
H^*H=A_0,                                               \tag{41}
\end{equation}
or, more generally, constructing a source-defined isometry between their
minimal factorizations.  G39 proves $H^*W=q$; it does not by itself prove
(41).

The notation must be reconciled with the supplied row (d).  There
$R_0=X_0^*X_0$ is the \emph{positive reference} and
$L_0=Y_0^*Y_0$ is the opposite channel; neither $X_0$ alone nor $Y_0$ alone
factors the signed old block.  Under the old-core hypothesis set
\begin{equation}
 T_0=R_0^{\dagger/2}L_0R_0^{\dagger/2},
 \qquad D_0=I-T_0.                                    \tag{41a}
\end{equation}
On the support of $R_0$, one has the exact form identity
\begin{equation}
 A_0=R_0-L_0
 =R_0^{1/2}D_0R_0^{1/2}.                              \tag{41b}
\end{equation}
Consequently, if and only if $D_0\ge0$, a minimal source factor of the signed
old block is
\begin{equation}
 \mathcal G_0=D_0^{1/2}R_0^{1/2},
 \qquad \mathcal G_0^*\mathcal G_0=A_0,               \tag{41c}
\end{equation}
up to the standard support projections.  Thus the actual compatibility test
for the G39 current is (41) with $\mathcal G_0$ from (41c), not
$H^*H=X_0^*X_0$.

\begin{proof}
Since $L_0=R_0^{1/2}T_0R_0^{1/2}$ on
$\overline{\Ran R_0}$, subtraction gives (41b).  The old-core hypothesis is
$0\le T_0\le I$, hence $D_0\ge0$ and functional calculus gives (41c).
Conversely existence of a Hilbert factor of (41b) forces the signed old block
to be nonnegative.
\end{proof}

This correction also locates the return-load condition in the current
language.  A normalized old load can be lifted through $\mathcal G_0^*$
exactly when it lies in $\Ran D_0^{1/2}$ after reference normalization.  This
is the range condition explicitly left open before (and used conditionally
in) \texttt{eq:returnsum} of the supplied row (d).  Hence replacing
$\mathcal G_0$ by the reference factor $X_0$ would erase precisely the
capacity defect that G40 is required to control.

\begin{theorem}[Exact reference-harmonic capacity identity --- \PROVED]
Write the old/new columns of the authoritative factors as
\[
 X=(X_0,X_E),\qquad Y=(Y_0,Y_E),
\]
and put
\[
 R_0=X_0^*X_0,\quad L_0=Y_0^*Y_0,\quad
 r=X_0^*X_E,\quad \ell=Y_0^*Y_E.
\]
Let $H=R_0^\dagger r$ and assume the usual reference range condition
$r\in\Ran R_0$.  Define
\begin{align}
 S_E&:=X_E^*X_E-r^*R_0^\dagger r,                     \tag{41d}\\
 Z_E&:=Y_E-Y_0H,                                      \tag{41e}\\
 b_E&:=L_0H-\ell.                                     \tag{41f}
\end{align}
Then the reference-harmonic column $v_E=(-H,I)^t$ satisfies
\begin{align}
 v_E^*(X^*X-Y^*Y)v_E&=S_E-Z_E^*Z_E,                  \tag{41g}\\
 P_{\rm old}(X^*X-Y^*Y)v_E&=b_E.                     \tag{41h}
\end{align}
If $A_0=R_0-L_0\ge0$, the full old/new signed block is nonnegative if and
only if
\begin{equation}
 b_E\in\Ran A_0^{1/2},\qquad
 \boxed{S_E-Z_E^*Z_E-b_E^*A_0^\dagger b_E\ge0.}       \tag{41i}
\end{equation}
\end{theorem}

\begin{proof}
The normal equation $R_0H=r$ gives
$X_0^*(X_E-X_0H)=0$ and
$(X_E-X_0H)^*(X_E-X_0H)=S_E$.  This proves the reference part of
(41g); its opposite part is $Z_E^*Z_E$ by definition.  The old cross of the
signed block against $v_E$ is
\[
 (r-\ell)-(R_0-L_0)H=L_0H-\ell=b_E,
\]
which proves (41h).  Congruence by the invertible triangular change from the
coordinate newborn column to $v_E$ leaves the old block equal to $A_0$ and
changes the other entries to (41h) and (41g).  The generalized Schur
criterion gives exactly (41i).
\end{proof}

Formula (41i) is the authoritative version of the G40 storage target.  It
also proves why the derivative calculation in v83 found a moving-lift square:
$v_E$ is harmonic for $R_0$, but its signed old divergence is the generally
nonzero $b_E$.  The Moore--Penrose term in (41i) is the energy required to
remove that divergence.  Any proposed Hodge residual must be identified with
the entire left side of (41i), not merely with $S_E-Z_E^*Z_E$.

On supported vectors for which $S_E^{-1/2}$ is defined, the normalized
variables of the supplied row (d) obey the exact translation
\begin{equation}
 u_e=D_0h_e-q_e
 =-R_0^{\dagger/2}b_ES_E^{\dagger/2}e.                \tag{41j}
\end{equation}
Indeed, substitute $R_0H=r$ into the definitions of $D_0,h_e,q_e$.  Thus the
apparently different range obstruction in \texttt{eq:returnsum} is precisely
the normalized form of the Feshbach range condition in (41i), and its return
energy is the normalized Moore--Penrose term there.  This closes the
bookkeeping equivalence between the threshold-capacity and balanced-factor
languages; it does not prove their common sign.

\begin{theorem}[Physical old-compatible amplitude port --- \PROVED]
Assume the inductive old-core hypothesis $A_0\ge0$ and the range condition
$b_E\in\Ran A_0^{1/2}$.  With $\mathcal G_0$ as in (41c), define
\begin{equation}
 W_{0,E}:=\mathcal G_0A_0^\dagger b_E.                 \tag{41k}
\end{equation}
Then
\begin{align}
 \mathcal G_0^*W_{0,E}&=b_E,                           \tag{41l}\\
 W_{0,E}^*W_{0,E}&=b_E^*A_0^\dagger b_E.              \tag{41m}
\end{align}
Moreover (41m) is the least Hilbert-current energy among all solutions of
$\mathcal G_0^*W=b_E$.
\end{theorem}

\begin{proof}
Moore--Penrose calculus and the range hypothesis give
$A_0A_0^\dagger b_E=b_E$, proving (41l).  Using
$\mathcal G_0^*\mathcal G_0=A_0$ gives (41m).  Every other solution differs
from (41k) by a vector in $\ker\mathcal G_0^*$, orthogonal to the minimal
solution, proving minimality.
\end{proof}

Thus amplitude and old-factor compatibility are completely realized inside
the authoritative threshold model once the inductive old sign and its range
condition are available.  The remaining newborn energy after paying this
minimal old current is exactly the boxed form (41i).  Identifying the finite
divisibility G39 port with (41k) would require a source-defined unitary (or
isometry plus null channels) preserving both the old Gram and the cross; the
present files provide no such comparison map.

\begin{theorem}[Joint-Gram identification criterion --- \PROVED]
Let $(K,W):E_0\oplus E\to\mathcal K$ and
$(\widetilde K,\widetilde W):E_0\oplus E\to\widetilde{\mathcal K}$ be two
old/new analysis pairs.  There is a unique unitary
\begin{equation}
 U:\overline{\Ran(K,W)}\longrightarrow
   \overline{\Ran(\widetilde K,\widetilde W)},
 \qquad UK=\widetilde K,\quad UW=\widetilde W,          \tag{41n}
\end{equation}
if and only if their complete joint Grams agree:
\begin{equation}
 \begin{pmatrix}K^*K&K^*W\\W^*K&W^*W\end{pmatrix}
 =
 \begin{pmatrix}
 \widetilde K^*\widetilde K&\widetilde K^*\widetilde W\\
 \widetilde W^*\widetilde K&\widetilde W^*\widetilde W
 \end{pmatrix}.                                       \tag{41o}
\end{equation}
\end{theorem}

\begin{proof}
Necessity follows from $U^*U=I$.  Conversely define
$U(Kx+We)=\widetilde Kx+\widetilde We$.  Equality (41o) shows that this rule
is well defined and preserves every inner product.  It extends to an
isometry of the closed generated ranges and is onto by construction.
\end{proof}

For G39 versus the physical threshold port one must first construct source
maps $\iota_0:E_0^{\rm phys}\to E_0^{\rm div}$ and
$\iota_E:E^{\rm phys}\to E^{\rm div}$.  After these have been specified,
(41o) requires simultaneously
\begin{align}
 \iota_0^*H_{\infty,0}^*H_{\infty,0}\iota_0&=A_0,   \tag{41p}\\
 \iota_0^*(-B_Q)\iota_E&=b_E,                         \tag{41q}\\
 \iota_E^*W_\infty^*W_\infty\iota_E
 &=b_E^*A_0^\dagger b_E.                               \tag{41r}
\end{align}
The adjoint cross is automatic once (41q) holds.  Equation (27g) establishes
the left cross in its finite divisibility model, but the source maps and none
of (41p)--(41r) are constructed in the supplied continuum row (d): the two
models use different source spaces and different deformation variables.  In
particular, (42) expands the untransported G39 old Gram, while (41a)--(41c)
expand the physical old Gram; their equality after transport is not a
normalization convention.  This joint-Gram test is
the exact, correctly typed interface required before G39 may be used as the
physical G40 current.

\begin{theorem}[Finite-port/full-Hodge dimension obstruction --- \PROVED]
Suppose the physical old block $A_0$ has infinite rank (in particular, suppose
it is strictly positive on an infinite-dimensional old primitive core).  No
map $\iota_0:E_0^{\rm phys}\to\mathcal H_N$ can satisfy (41p), because
$H_{\infty,0}$ in (27f) has finite-dimensional codomain.  Consequently the
finite divisibility G39 port cannot be unitarily identified with the complete
physical old/Hodge factor of the supplied row (d).
\end{theorem}

\begin{proof}
The rank of
$\iota_0^*H_{\infty,0}^*H_{\infty,0}\iota_0$ is at most
$\dim\mathcal K_\infty^{(1)}=2N$, whereas $A_0$ has infinite rank.  Thus
(41p) is impossible.  Under strict positivity the contradiction can also be
seen from injectivity: (41p) would force $\iota_0$ to be injective into a
finite-dimensional space.
\end{proof}

Hence Theorem G39 remains a valid finite arithmetic amplitude identity, but
it is not a factorization of the full continuum row-(d) current.  It may be
used for compatible finite compressions only.  A full G40 proof requires an
infinite-dimensional continuum port constructed directly from (81)--(84),
not an isometric embedding of (27e)--(27f).

\begin{theorem}[Old compatibility by neutral birth channels --- \PROVED]
Let $0<T_0<T_1$ and embed $L^2(I_{T_0})$ in $L^2(I_{T_1})$ by zero
extension.  For a channel born between the two thresholds, put
$a=\log n\ge2T_0$.  Then, on the old source,
\begin{equation}
 \|J_{a,-}F\|^2=\|J_{a,+}F\|^2=\|F\|^2,             \tag{41s}
\end{equation}
and there is a source-defined unitary $U_{a,T_0}$ on the two disjoint support
copies such that
\begin{equation}
             U_{a,T_0}J_{a,-}F=J_{a,+}F.             \tag{41t}
\end{equation}
Consequently the physical balanced factors $\mathcal X_T,\mathcal Y_T$ in
(83)--(84) restrict from $T_1$ to $T_0$ as the old factors at $T_0$ plus
orthogonal neutral channels on which the angular observation is unitary.
In particular,
\begin{equation}
 \bigl(\mathcal X_{T_1}^*\mathcal X_{T_1}
       -\mathcal Y_{T_1}^*\mathcal Y_{T_1}\bigr)\big|_{I_{T_0}}
 =\mathcal X_{T_0}^*\mathcal X_{T_0}
       -\mathcal Y_{T_0}^*\mathcal Y_{T_0}.          \tag{41u}
\end{equation}
Thus old-factor compatibility of the \emph{physical continuum
factorization} is exact and does not require (41p)--(41r).
\end{theorem}

\begin{proof}
For $F$ supported in $I_{T_0}$ and $a\ge2T_0$, the supports of $F$ and
$\tau_aF$ are disjoint up to a null endpoint.  Hence
\[
 \|\tau_aF\mp F\|^2=2\|F\|^2,
\]
which proves (41s).  On
$E_{T_0}L^2\oplus\tau_aE_{T_0}L^2$ define $U_{a,T_0}$ to be minus the
identity on the first summand and the identity on the translated summand.
It is unitary and sends
$(\tau_aF-F)/\sqrt2$ to $(\tau_aF+F)/\sqrt2$, proving (41t).

Every channel already present at $T_0$ is intertwined exactly by ambient
zero extension, as is the Gamma difference feature (82) and the constant
channel.  Each newly born label is orthogonal in the direct-sum analysis
space and supplies the unitary pair just constructed.  Its two Grams are
equal, so it cancels from the signed difference.  Summing over the finitely
many new labels proves (41u) and the restriction assertion.
\end{proof}

\begin{corollary}[Correct scope of the joint-Gram obstruction --- \PROVED]
Equations (41p)--(41r) test the stronger and, for an infinite-rank old block,
impossible identification of the \emph{finite G39 port alone} with the full
physical Hodge factor.  They are not conditions for old compatibility of
the continuum factors.  The latter is supplied by (41s)--(41u).  The role of
G39 in a continuum proof is therefore restricted to its cross-amplitude
identity (27g)--(27h); its positive normalization must be absorbed into the
infinite Gamma--Poisson colligation rather than matched by a finite unitary.
\end{corollary}

\begin{proof}
The finite-rank obstruction proves the first statement, while Theorem
(41s)--(41u) proves the second.  The last statement follows because adjoining
the finite G39 analysis as an orthogonal old factor would change the old
Gram, whereas a source interconnection inside the continuum colligation can
change amplitudes while retaining the neutral restriction law.
\end{proof}

For the doubled logarithmic port its additional old Gram is explicitly
\begin{equation}
 H_{\infty,0}^*H_{\infty,0}
 =\frac S2P_{\old}(I+\cZ\cZ^*)P_{\old}^*.               \tag{42}
\end{equation}
Thus (41) is a concrete, falsifiable operator identity, not a normalization
convention.

The complete G40 bridge can be stated without subdivision as the single
source factorization
\begin{equation}
 \boxed{
 \begin{pmatrix}A_0&q\\q^*&S_M'\end{pmatrix}
 =\mathfrak T_T^*\mathfrak T_T,}                        \tag{43}
\end{equation}
where the old/new columns of $\mathfrak T_T$ must restrict to the G39
amplitude realization and the remaining rows must be explicit Gamma-tail and
Tate-terminal outputs.  Equality (43) simultaneously enforces the correct old
Gram, the proved cross amplitude, and the positive source budget.

Reverse engineering always produces a Krein factorization of the left side of
(43), but it produces a Hilbert factorization if and only if that block is
positive.  By the generalized Schur theorem this is equivalent to
$S_M'-q^*A_0^\dagger q\ge0$, namely G42/G.  Therefore (43) proves G40 only if
$\mathfrak T_T$ is constructed independently from Gamma--Tate formulas; taking
a square root or polar factor of the block would be circular.

\section{The unitary midpoint bridge}

There is a non-circular mechanism capable of proving (43): the midpoint split
of a source-defined unitary Poisson--Fourier operator.  It avoids defining a
Schur residual.

\begin{theorem}[Unitary midpoint storage --- \PROVED]
Let $\mathcal U$ be unitary on a Hilbert space $\mathcal K$, let $W$ be a
current, and define
\begin{equation}
 W_+:=\frac12(I+\mathcal U)W,
 \qquad
 W_-:=\frac12(I-\mathcal U)W.                           \tag{44}
\end{equation}
Then
\begin{equation}
 W^*W=W_+^*W_++W_-^*W_-.                               \tag{45}
\end{equation}
If an analysis operator $H$ satisfies
\begin{equation}
 H^*W=q+B,\qquad H^*W_-=B,                             \tag{46}
\end{equation}
then
\begin{equation}
 H^*W_+=q                                               \tag{47}
\end{equation}
and (45) is the required non-circular positive storage identity.
\end{theorem}

\begin{proof}
Expanding the two squares in (44), the mixed terms cancel and
\[
 W_+^*W_++W_-^*W_-
 =\frac12W^*W+\frac12W^*\mathcal U^*\mathcal UW=W^*W.
\]
Equation (47) follows by subtracting the two equations in (46).
\end{proof}

For the present problem take $W=W_{\Hod}$, so that $W^*W=S_M'$.  G39 has
constructed a current with divergence $-B_Q$.  Consequently the single bridge
identity which would finish G40 is
\begin{equation}
 \boxed{
 W_\infty=-W_-
 =-\frac12(I-\mathcal U_{\rm PF})W_{\Hod},}             \tag{48}
\end{equation}
where $\mathcal U_{\rm PF}$ must be the normalized Poisson--Fourier transport
in the same current Hilbert space.  Indeed, (2) and G39 give
\[
 H^*W_{\Hod}=q+B_Q,\qquad H^*(-W_\infty)=B_Q.
\]
Thus (44)--(48) yield
\begin{align}
 H^*W_+&=q,\notag\\
 S_M'&=W_+^*W_++W_-^*W_-,                              \tag{49}
\end{align}
which is exactly the desired Hilbert storage theorem; $W_-$ is the positive
Tate residual and no square root of $\Delta_T$ has been introduced.

The differentiated Poisson identity from the preceding row-(c) work,
\begin{equation}
 (C+C^\sharp-K-X)f=(\mathcal U_{\rm PF}-I)Kf,           \tag{50}
\end{equation}
has precisely the algebraic form required by (48): the two Euler orientations
give the doubled G39 port and $K+X$ is converted to Gamma by the proved
Fourier--Gamma bridge.  Hence (50) proves (48) at the level of unnormalized
source coefficients.

The remaining referee-level requirement is metric: the map carrying $Kf$ to
$W_{\Hod}$ and the left side of (50) to the doubled current (27e) must be shown
to be isometric for the physical current norm.  The fixed-parameter CCM
theorem gives a hilbertian isomorphism, not this isometry; its normalization
derivative is (25)--(26).  Therefore the exact final identity to check is
\begin{equation}
 \mathcal I_T(\mathcal U_{\rm PF}-I)Kf
 =(\mathcal U_{\rm cur}-I)\mathcal I_TKf,\qquad
 \mathcal I_T^*\mathcal I_T=I,                          \tag{51}
\end{equation}
with $\mathcal I_TKf=W_{\Hod}$.  If (51) is established from the source
normalizations, (48)--(49) prove G40 completely.  Without the second equality
in (51), replacing a hilbertian transport by an isometry would assume the
missing energy statement.

\section{Restricted polar normalization on the Sonin space}

The metric part of (51) can in fact be obtained directly from the available
CCM normalizations, provided the polar normalization is performed after
restriction to the Sonin space.

Let $\mathcal S_T$ be the archimedean Sonin space at the fixed cutoff and
$\mathcal S_T^S$ its semilocal counterpart.  Write
\[
 \Theta_T:=\theta_S|_{\mathcal S_T}:\mathcal S_T
 \longrightarrow\mathcal S_T^S.
\]
The CCM stability theorem says that $\Theta_T$ is bounded and bijective with
bounded inverse, and the Fourier intertwining says
\begin{equation}
 \mathcal F_S\Theta_T=\Theta_T\mathcal F_T.              \tag{52}
\end{equation}
Define the restricted Gram and its polar normalization by
\begin{equation}
 D_T:=\Theta_T^*\Theta_T>0,\qquad
 \mathcal I_T:=\Theta_TD_T^{-1/2}.                       \tag{53}
\end{equation}

\begin{theorem}[Sonin polar isometry --- \PROVED]
The operator $D_T$ commutes with the restricted Fourier symmetry
$\mathcal F_T$, the operator $D_T^{-1/2}$ preserves $\mathcal S_T$, and
\begin{align}
 \mathcal I_T^*\mathcal I_T&=I_{\mathcal S_T},           \tag{54}\\
 \mathcal F_S\mathcal I_T&=\mathcal I_T\mathcal F_T.    \tag{55}
\end{align}
Thus $\mathcal I_T$ is a source-defined unitary from $\mathcal S_T$ onto
$\mathcal S_T^S$; no positivity statement from row (d) is used.
\end{theorem}

\begin{proof}
Both Fourier transforms are unitary and (52) holds.  Hence
\[
 \mathcal F_T^*D_T\mathcal F_T
 =\mathcal F_T^*\Theta_T^*\Theta_T\mathcal F_T
 =\Theta_T^*\mathcal F_S^*\mathcal F_S\Theta_T
 =D_T.
\]
Therefore $D_T$ and every bounded Borel function of $D_T$, in particular
$D_T^{-1/2}$, commute with $\mathcal F_T$.  Since $D_T$ is an operator on the
already restricted space $\mathcal S_T$, its functional calculus preserves
that space.  Equation (54) follows directly:
\[
 \mathcal I_T^*\mathcal I_T
 =D_T^{-1/2}\Theta_T^*\Theta_TD_T^{-1/2}=I.
\]
Finally, using (52) and commutation with $D_T^{-1/2}$,
\[
 \mathcal F_S\mathcal I_T
 =\mathcal F_S\Theta_TD_T^{-1/2}
 =\Theta_T\mathcal F_TD_T^{-1/2}
 =\mathcal I_T\mathcal F_T.
\]
Surjectivity follows from the invertibility of $\Theta_T$ and $D_T^{-1/2}$.
\end{proof}

Consequently the second equality in (51) and the unitary intertwining portion
of the first equality are proved.  The normalization is canonical: it is the
unitary factor in the polar decomposition of the restricted CCM map.

\paragraph{Remaining identification.}
To apply (54)--(55) to (48), one must still prove that the source tangent vector
$Kf$ and the physical Hodge current use these restricted Sonin coordinates:
\begin{equation}
 \boxed{\mathcal I_TKf=W_{\Hod}.}                        \tag{56}
\end{equation}
Equation (56) is not the normalization identity; it is the observation
identity between the Poisson tangent feature and the finite-threshold Hodge
feature.  The multiplier/connection calculations (24)--(27) and G39 prove its
logarithmic boundary component, but equality of the complete feature vectors
must be checked from the definitions of $K$, $R$, and $\cL$.  Thus the
normalization requested in (51) is closed, while the application to the
specific Hodge current retains the explicit observation check (56).

\section{Type audit of the final observation}

As written, (56) suppresses two different codomains.  The polar map has type
\[
 \mathcal I_T:\mathcal S_T\longrightarrow\mathcal S_T^S,
\]
whereas the Hodge feature has type
\[
 W_{\Hod}:E_T\longrightarrow
 \mathcal K_{\Hod,T}:=\overline{\Ran\cL^{1/2}R}.
\]
Thus a referee-level statement requires an additional map
\begin{equation}
 J_T:\mathcal S_T^S\longrightarrow\mathcal K_{\Hod,T}  \tag{57}
\end{equation}
and the correctly typed observation is
\begin{equation}
 J_T\mathcal I_TK_T=W_{\Hod}.                            \tag{58}
\end{equation}

\begin{theorem}[Unitary observation criterion --- \PROVED]
Let $A:E\to\mathcal H_1$ and $B:E\to\mathcal H_2$ be bounded.  There is an
isometry $J:\overline{\Ran A}\to\overline{\Ran B}$ satisfying $JA=B$ if and
only if
\begin{equation}
 A^*A=B^*B.                                              \tag{59}
\end{equation}
In that case $J$ is uniquely determined on $\overline{\Ran A}$ and is unitary
onto $\overline{\Ran B}$.
\end{theorem}

\begin{proof}
Necessity follows from $A^*J^*JA=A^*A$.  Conversely, define $J(Ae)=Be$.
Equation (59) shows first that this is well defined and then that it preserves
inner products (by polarization).  It therefore extends uniquely to an
isometry of the closed ranges, and its range is dense in and closed in
$\overline{\Ran B}$.
\end{proof}

Apply the theorem with $A=\mathcal I_TK_T$ and $B=W_{\Hod}$.  Since
$\mathcal I_T$ is isometric, (59) becomes
\begin{equation}
 \boxed{K_T^*K_T=W_{\Hod}^*W_{\Hod}=S_M'.}              \tag{60}
\end{equation}
Consequently the CCM polar normalization proves the isometry inside the Sonin
transport, but it cannot by itself prove the observation into the Hodge
current space.  The latter exists isometrically precisely when the source
tangent Gram is the physical Hodge Gram.

Equation (60) is a concrete coefficient identity to be proved from the actual
definitions of $K_T$, $R$, and $\cL$:
\begin{equation}
 \boxed{K_T^*K_T=v^*R^*\cL Rv.}                         \tag{61}
\end{equation}
The current local dossier specifies the right side and the logarithmic
boundary of the left side, but it does not contain a definition of $K_T$ as an
operator into the restricted Sonin space sufficient to expand the left side.
Therefore identifying $J_T$ by inverse engineering before proving (61) would
simply define an isometry whose existence is equivalent to the desired Gram
identity.

\section{Audit of the symbol $K$: the midpoint candidate fails}

The preceding row-(c) notes define
\begin{equation}
 K=\mathcal F^{-1}X\mathcal F,                           \tag{62}
\end{equation}
where $X$ is the position/scaling-coordinate operator.  Thus $K$ is a
first-order unbounded operator; it was not defined as a Kolmogorov factor of
the Hodge quadratic form.  Consequently the use of $Kf$ as though it were the
source energy feature in (56) requires the Gram identity (60), which is not a
formal consequence of differentiated Poisson.

In fact the identity fails on the full natural domain.

\begin{theorem}[Order mismatch no-go --- \PROVED]
Let $K=\mathcal F^{-1}X\mathcal F$ on $L^2(\mathbb R)$ and let
$G_\Gamma$ be the archimedean Weil operator whose Fourier multiplier has
logarithmic growth.  Then no isometry $J$ can satisfy
\begin{equation}
 J Kf=G_\Gamma^{1/2}f                              \tag{63}
\end{equation}
for all functions in a common invariant core.  In particular
$K^*K\ne G_\Gamma$ and the proposed observation (56) cannot hold with the
row-(c) operator $K$.
\end{theorem}

\begin{proof}
By unitarity of Fourier,
\[
 K^*K=\mathcal F^{-1}X^2\mathcal F.
\]
Choose normalized Fourier-localized functions $f_n$ supported where
$|x-n|\le1$.  Then
\[
 \langle K^*Kf_n,f_n\rangle=n^2+O(n).
\]
The Gamma multiplier obtained from the logarithmic derivative of the Gamma
factor is $O(\log(2+|x|))$; hence
\[
 \langle G_\Gamma f_n,f_n\rangle=O(\log n).
\]
The two quadratic forms cannot agree.  If (63) held with $J$ isometric, their
norm squares would agree for every $f_n$, a contradiction.
\end{proof}

The finite Euler terms do not repair the order mismatch: at a fixed threshold
they are a finite sum of bounded translations, hence contribute $O(1)$ on the
sequence above.  Therefore the full finite-threshold Hodge energy still cannot
be $K^*K$.

\paragraph{Consequence for G40.}
The unitary midpoint theorem (44)--(47) is correct, but its application
(48)--(51) with the row-(c) vector $Kf$ is \AUDIT{} and cannot be used as a
proof of G40.  Differentiated Poisson proves a connection/amplitude identity,
not equality of source and Hodge energy features.  A valid midpoint proof
would require an independently constructed feature $B_T$ satisfying
\[
 B_T^*B_T=S_M'
\]
and a source-defined unitary acting on its range whose difference port is the
G39 current.  Taking $B_T=(S_M')^{1/2}$ without deriving that unitary from the
Gamma--Tate network would merely restate the missing theorem.

\section{Unitary-orbit audit of the G39 auxiliary port}

For vectors/operators $W,Z$ in one Hilbert current space, a unitary $U$ can
satisfy
\begin{equation}
 Z=-\frac12(I-U)W                                      \tag{64}
\end{equation}
only if
\begin{equation}
 W^*Z+Z^*W+2Z^*Z=0.                                    \tag{65}
\end{equation}
Indeed $UW=W+2Z$, and (65) is exactly
$(W+2Z)^*(W+2Z)=W^*W$.

\begin{proposition}[Orthogonal-port no-go --- \PROVED]
If $Z\perp W$ and $Z\ne0$, no unitary $U$ satisfies (64).
\end{proposition}

\begin{proof}
Under orthogonality, (65) reduces to $2Z^*Z=0$, a contradiction.
\end{proof}

The G39 construction (27e) was adjoined as an auxiliary direct-sum port to the
Hodge current.  In that realization it is orthogonal to $W_{\Hod}$, so it
cannot be the midpoint residual asserted in (48).  A successful G40 network
must instead construct a nonorthogonal embedding
\[
 \jmath_T:\mathcal K_\infty^{(1)}\to\mathcal K_{\Hod,T}
\]
such that
\begin{align}
 H^*\jmath_TW_\infty&=-B_Q,\notag\\
 W_{\Hod}^*\jmath_TW_\infty
 +(\jmath_TW_\infty)^*W_{\Hod}
 +2(\jmath_TW_\infty)^*\jmath_TW_\infty&=0.            \tag{66}
\end{align}
The first line is amplitude; G39 proves it only before this embedding.  The
second line is the exact unitary-orbit energy constraint.  Constructing
$\jmath_T$ from Gamma--Tate, independently of (66), would give the desired
midpoint bridge.  Choosing it by solving (66) is equivalent to the missing
passivity theorem.

\section{Oriented-delay realization of the positive source metric}

Let $E_T:L^2(-T,T)\to L^2(\mathbb R)$ be extension by zero, let $S_u$ be
the unitary translation group on the \emph{ambient} space $L^2(\mathbb R)$,
and let $\mathfrak F$ be an involutive unitary satisfying
\begin{equation}
 \mathfrak F S_u\mathfrak F=S_{-u}.                     \tag{67}
\end{equation}
For $u>0$ define the oriented-delay feature
\begin{equation}
 \mathcal B_uF:=\frac1{\sqrt2}
 \binom{(I-S_u)F}{(I-S_{-u})F}.                         \tag{68}
\end{equation}

\begin{theorem}[Oriented-delay Fourier colligation --- \PROVED]
Define $\widehat{\mathfrak F}$ on the doubled current space by
\begin{equation}
 \widehat{\mathfrak F}\binom{x}{y}
 :=\binom{\mathfrak Fy}{\mathfrak Fx}.                 \tag{69}
\end{equation}
Then $\widehat{\mathfrak F}$ is an involutive unitary,
\begin{align}
 \mathcal B_u^*\mathcal B_uF
 &\text{ has quadratic form }\|(I-S_u)F\|^2,           \tag{70}\\
 \widehat{\mathfrak F}\mathcal B_u
 &=\mathcal B_u\mathfrak F.                            \tag{71}
\end{align}
Consequently the midpoint currents
\begin{equation}
 \mathcal B_{u,\pm}:=\frac12
 (I\pm\widehat{\mathfrak F})\mathcal B_u              \tag{72}
\end{equation}
satisfy the exact positive conservation law
\begin{equation}
 \mathcal B_u^*\mathcal B_u
 =\mathcal B_{u,+}^*\mathcal B_{u,+}
 +\mathcal B_{u,-}^*\mathcal B_{u,-}.                  \tag{73}
\end{equation}
\end{theorem}

\begin{proof}
Equation (69) is plainly an involutive unitary.  Since $S_u$ is unitary,
$\|(I-S_{-u})F\|=\|(I-S_u)F\|$, which proves (70).
From (67),
\[
 \mathfrak F(I-S_{-u})=(I-S_u)\mathfrak F,
 \qquad
 \mathfrak F(I-S_u)=(I-S_{-u})\mathfrak F,
\]
giving (71).  Equation (73) is the unitary midpoint identity.
\end{proof}

Let
\[
 w_\Gamma(u)=\frac{e^{-u/2}}{1-e^{-2u}}\qquad(u>0),
 \qquad c_n=\frac{\Lambda(n)}{\sqrt n}.
\]
The completed finite-threshold feature is the direct integral/sum
\begin{equation}
 \mathcal B_TF:=
 \left(
   \{w_\Gamma(u)^{1/2}\mathcal B_uF\}_{u>0},
   \{c_n^{1/2}\mathcal B_{\log n}F\}_{n<e^{2T}}
 \right).                                               \tag{74}
\end{equation}
Using the already proved Gamma delay identity, (70) gives
\begin{equation}
 \|\mathcal B_TF\|^2
 =\int_0^\infty w_\Gamma(u)\|(I-S_u)F\|^2du
 +\sum_{n<e^{2T}}c_n\|(I-S_{\log n})F\|^2.             \tag{75}
\end{equation}
Here the argument of every $\mathcal B_u$ is $E_TF$.  Thus $\mathcal B_T$
is a source-defined Kolmogorov factor of the positive Euler--Gamma reference
metric.  The direct-integral unitary obtained from (69)
intertwines $\mathcal B_T$ with $\mathfrak F$ and gives
\begin{equation}
 \mathcal B_T^*\mathcal B_T
 =\mathcal B_{T,+}^*\mathcal B_{T,+}
 +\mathcal B_{T,-}^*\mathcal B_{T,-}.                  \tag{76}
\end{equation}
This construction has the correct logarithmic spectral order and replaces the
invalid use of $K=\mathcal F^{-1}X\mathcal F$.

The preceding display proves a Gram identity for the positive reference.
It does \emph{not} identify that Gram with the independently defined Hodge
operator $\cL$ or with $S_M'$.  Defining $\cL:=\mathcal B_T^*\mathcal B_T$
would merely rename the reference form and would not prove compatibility
with the old factor used by the threshold induction.  The former equations
(76a)--(76b) are therefore withdrawn and classified as \AUDIT{}.

\paragraph{Interface still to prove.}
For G40 it remains to restrict the difference port to the harmonic newborn
input and verify the exact identity
\begin{equation}
 H^*\mathcal B_{T,-}Rv=B_Q.                              \tag{78}
\end{equation}
Equation (78) is the required identification of the
oriented difference port with the algebraic G39 port.  Unlike the former
$K$-bridge, (78) has the correct logarithmic order.  It remains only a
candidate until both its type map and its compatibility with the independently
defined Hodge/old Gram have been constructed.

\section{Authoritative balanced factorization of the supplied row (d)}

We now use the definitions in the supplied
\texttt{row-d-local-analysis.tex}.  Let $I_T=(-T,T)$, let $E_T$ denote
extension by zero to $\mathbb R$, and put
\[
 S_a:=E_T^*\tau_aE_T,
 \qquad (\tau_aF)(t)=F(t+a).
\]
Then $S_a^*=S_{-a}$.  For $a>0$ define the two \emph{ambient}
zero-extension channels from $L^2(I_T)$ to $L^2(\mathbb R)$
\begin{equation}
 J_{a,-}F:=\frac{\tau_aE_TF-E_TF}{\sqrt2},
 \qquad J_{a,+}F:=\frac{\tau_aE_TF+E_TF}{\sqrt2}.     \tag{79}
\end{equation}
They obey the exact polarization identity
\begin{equation}
 J_{a,-}^*J_{a,-}-J_{a,+}^*J_{a,+}
 =-(S_a+S_{-a}).                                      \tag{80}
\end{equation}
Indeed $\tau_a$ is unitary on the ambient line and
$E_T^*\tau_aE_T=S_a$.  This is the exact zero-extension factorization used
by the supplied row (d).  Both signs receive the same diagonal energy.

The Gamma multiplier in the supplied file is
\[
 g_\Gamma(\tau)=\Re\psi\left(\frac14+\frac{i\tau}{2}\right)
                  -\psi\left(\frac14\right).
\]
The integral formula for the digamma function and the substitution $t=2u$
give
\begin{equation}
 g_\Gamma(\tau)
 =2\int_0^\infty
 \frac{e^{-u/2}}{1-e^{-2u}}(1-\cos(\tau u))\,du.       \tag{81}
\end{equation}
Equivalently, on the natural form core,
\begin{equation}
 \langle G_{\Gamma,T}F,F\rangle
 =\int_0^\infty\frac{e^{-u/2}}{1-e^{-2u}}
       \|(I-\tau_u)E_TF\|_{L^2(\mathbb R)}^2\,du.     \tag{82}
\end{equation}
Thus $G_{\Gamma,T}\ge0$ and (82) is a source-defined Gamma Kolmogorov
factor, not a square root chosen from the desired signed operator.

Recall
\[
 m_0=\log\pi+\gamma+\frac\pi2+3\log2
     =\log\pi-\psi\left(\frac14\right)>0,
 \qquad w_n=\frac{\Lambda(n)}{\sqrt n}.
\]
Let $\mathcal X_T$ be the direct-sum analysis map whose components are the
Gamma feature in (82) and $w_n^{1/2}J_{\log n,-}$, and let
$\mathcal Y_T$ have components $m_0^{1/2}I$ and
$w_n^{1/2}J_{\log n,+}$, for $2\le n<e^{2T}$.

\begin{theorem}[Exact Gamma--Tate balanced factorization --- \PROVED]
On the form domain of the localized primitive operator in the supplied row
(d),
\begin{equation}
 A_T=\mathcal X_T^*\mathcal X_T-
     \mathcal Y_T^*\mathcal Y_T.                       \tag{83}
\end{equation}
After imposing the two Tate moments, with $P_T$ the orthogonal projection
onto $\mathcal P_T=\ker M_-\cap\ker M_+$,
\begin{equation}
 P_TA_TP_T=(\mathcal X_TP_T)^*(\mathcal X_TP_T)
           -(\mathcal Y_TP_T)^*(\mathcal Y_TP_T).      \tag{84}
\end{equation}
In particular this gives explicit, source-defined meanings to the abstract
factors $X_T,Y_T$ used in the return calculation of the supplied file.
\end{theorem}

\begin{proof}
The Gamma and constant pieces of (83) are respectively
$G_{\Gamma,T}$ and $-m_0I$.  For every prime power, (80), multiplied by
$w_n$, is exactly
$-w_n(S_{\log n}+S_{-\log n})$.  Summing gives the displayed definition of
$A_T$.  Compression by $P_T$ proves (84).
\end{proof}

The factorization also isolates the missing theorem without any notation
from v83--v86.

\begin{theorem}[Exact Douglas form of the remaining gate --- \PROVED]
The following are equivalent:
\begin{enumerate}
\item $P_TA_TP_T\ge0$;
\item $\|\mathcal Y_TP_TF\|\le\|\mathcal X_TP_TF\|$ for every $F$;
\item there is a contraction
\[
 C_T:\overline{\Ran(\mathcal X_TP_T)}
       \longrightarrow\overline{\Ran(\mathcal Y_TP_T)}
 \]
such that $C_T\mathcal X_TP_T=\mathcal Y_TP_T$.
\end{enumerate}
\end{theorem}

\begin{proof}
The equivalence of the first two statements is (84).  The inequality makes
$C_T(\mathcal X_TP_TF):=\mathcal Y_TP_TF$ well-defined and contractive on
the range; continuity extends it to the closure.  Conversely such a
contraction gives the norm inequality.
\end{proof}

There is a useful refinement: the source angular object exists before its
contractivity is known.

\begin{theorem}[Canonical source angular operator --- \PROVED]
The Gamma feature in $\mathcal X_T$ is injective.  Hence the rule
\begin{equation}
 \mathfrak C_T^{\rm src}(\mathcal X_TP_TF)
 :=\mathcal Y_TP_TF                                    \tag{84a}
\end{equation}
defines a linear operator on $\Ran(\mathcal X_TP_T)$ without assuming row
(d).  In the sense of sesquilinear forms on source vectors, its defect pulls
back to the primitive operator:
\begin{equation}
 \langle P_TA_TP_TF,G\rangle
 =\langle\mathcal X_TP_TF,\mathcal X_TP_TG\rangle
 -\langle\mathfrak C_T^{\rm src}\mathcal X_TP_TF,
          \mathfrak C_T^{\rm src}\mathcal X_TP_TG\rangle. \tag{84b}
\end{equation}
For every old/new decomposition of the primitive source, restriction of
(84a) gives the same old columns as (83); thus its compatibility with the
authoritative old factor is automatic.  It extends to a contraction on the
closed range if and only if row (d) holds at $T$.
\end{theorem}

\begin{proof}
By (81), $g_\Gamma(\tau)>0$ for every $\tau\ne0$.  If the Gamma feature of
$F$ vanishes, Plancherel gives
$g_\Gamma(\tau)|\widehat{E_TF}(\tau)|^2=0$ almost
everywhere.  An $L^2$ function supported at the singleton $\{0\}$ is zero,
so $E_TF=0$ and $F=0$.  Therefore $\mathcal X_T$ and its restriction to
$\mathcal P_T$ are injective, proving that (84a) is well defined.

For source vectors $F,G$, the sesquilinear form of the right side of (84b)
is
\[
 \langle\mathcal X_TP_TF,\mathcal X_TP_TG\rangle
 -\langle\mathcal Y_TP_TF,\mathcal Y_TP_TG\rangle,
\]
which is (84) by polarization.  Restriction to old/new columns commutes with
the source definition, proving compatibility.  Finally, bounded extension
with norm at most one is equivalent to the norm inequality in the preceding
theorem.
\end{proof}

Thus inverse engineering has produced a correctly typed, source-defined
bridge rather than a square root of the desired Schur remainder.  Its
unproved property is now a single quantitative statement,
\begin{equation}
 \boxed{\|\mathfrak C_T^{\rm src}\|\le1
        \quad\hbox{for every }T>0.}                    \tag{84c}
\end{equation}
Within the authoritative balanced factorization, no normalization or
old-column ambiguity remains in (84a)--(84b); the missing content is precisely
the global passivity bound (84c).  Identification of these columns with the
separately introduced v83--v86 Hodge factor remains subject to the Gram
compatibility audit (41).

\begin{theorem}[Correctly typed Poisson--Sonin observation --- \PROVED]
Let $U_T^X$ be any source-defined isometry from
$\overline{\Ran(\mathcal X_TP_T)}$ into a proposed normalized Sonin current
space.  Define on $U_T^X\Ran(\mathcal X_TP_T)$
\begin{equation}
 \Pi_T^{\rm src}\,U_T^X\mathcal X_TP_TF
 :=\mathcal Y_TP_TF.                                   \tag{86}
\end{equation}
Then $\Pi_T^{\rm src}$ is well defined, and
\begin{align}
 &\langle P_TA_TP_TF,G\rangle \notag\\
 &\quad=
 \langle U_T^X\mathcal X_TP_TF,U_T^X\mathcal X_TP_TG\rangle
 -\langle\Pi_T^{\rm src}U_T^X\mathcal X_TP_TF,
          \Pi_T^{\rm src}U_T^X\mathcal X_TP_TG\rangle. \tag{87}
\end{align}
The observation $\Pi_T^{\rm src}$ extends contractively to the closed current
space if and only if row (d) holds at $T$.  Requiring it to be isometric would
instead force $P_TA_TP_T=0$.
\end{theorem}

\begin{proof}
Injectivity of $\mathcal X_TP_T$ and isometry of $U_T^X$ make (86) well
defined.  Isometry and (84b) give (87).  The contractive-extension criterion
is exactly (84c).  If $\Pi_T^{\rm src}$ were isometric, the two terms on the
right of (87) would agree for every $F,G$, so the compressed primitive form
would vanish identically.
\end{proof}

Thus the normalized Sonin map $\mathcal I_T$ from (53) can serve as the
isometry $U_T^X$ only after one supplies a correctly typed source isometry
\begin{equation}
 V_T:\overline{\Ran(\mathcal X_TP_T)}\longrightarrow\mathcal S_T,
 \qquad V_T^*V_T=I,                                    \tag{88}
\end{equation}
derived from the Gamma--Tate Gram.  The actual normalized transport is then
$U_T^X=\mathcal I_TV_T$, and (86), not an isometric observation, is the
correct Poisson--Sonin--Gamma--Tate diagram.  Equation (88) is the precise
Gram/type interface absent from the supplied files.  Even after (88), the
contractivity of (86) remains the global sign theorem.

\paragraph{Audit against the Sonin metric in \texttt{main.tex}.}
The later section ``A derived metric: the Sonin compression'' defines instead
the positive Hilbert--Schmidt Gram
\[
 M(f)_{mn}=\operatorname{Tr}
 \bigl(\vartheta(\psi_m^f)\mathbf S
       \vartheta(\psi_n^f)^*\bigr).
\]
That positivity is unconditional, but the supplied manuscript explicitly
states that comparison with the Weil form is proved only in the prime-free
window and that the semilocal extension is not carried out.  For general
windows it leaves open the quantitative lower-frame estimate labelled
\texttt{eq:framebound}.  No equality in the supplied files identifies this
$M(f)$ with $(\mathcal X_TP_T)^*(\mathcal X_TP_T)$, so it does not construct
$V_T$ in (88).  Treating two positive Grams as identical merely because both
come from Sonin/Gamma data would be a Gram substitution and is forbidden.

\begin{theorem}[Global equivalence audit --- \PROVED]
The following assertions are equivalent:
\begin{enumerate}
\item $\|\mathfrak C_T^{\rm src}\|\le1$ for every $T>0$;
\item $P_TA_TP_T\ge0$ for every $T>0$;
\item the Weil quadratic form has the required sign on every compactly
supported test function satisfying the two Tate moment conditions;
\item the Riemann hypothesis holds.
\end{enumerate}
\end{theorem}

\begin{proof}
The equivalence of the first two assertions is (84b)--(84c).  Extension by
zero identifies the union over $T$ of the primitive spaces with the compactly
supported two-moment test core, so the second and third assertions are
equivalent.  The equivalence of the third and fourth assertions is Weil's
positivity criterion, with the sign convention fixed by
$A_T=-B_{\rm nuc}$ in the supplied manuscript.
\end{proof}

This theorem is not a proof of the final sign.  It records that any derivation
of (84c) from Poisson--Sonin data is itself a proof of RH and therefore must
contain a new global estimate beyond the currently available local
normalizations and intertwining identities.

\begin{theorem}[No channelwise Gamma--Tate contraction --- \PROVED]
Fix $a>0$.  For all sufficiently large $T$ there is an
$F\in\mathcal P_T$ such that
\begin{equation}
 \|J_{a,+}F\|^2-\|J_{a,-}F\|^2
 =2\Re\langle S_aF,F\rangle>0.                         \tag{89}
\end{equation}
Consequently a contraction implementing the global identity in the preceding
theorem cannot be a direct sum of contractions which send each arithmetic
difference channel $J_{a,-}$ to the corresponding sum channel $J_{a,+}$.
It must mix different places and/or the Gamma and Tate-terminal channels.
\end{theorem}

\begin{proof}
Choose a nonzero nonnegative smooth bump $\varphi$ supported in an interval
of length smaller than $a/4$.  For $T$ large, choose three numbers $x_j$,
separated from one another by more than $3a$, so that all six translates
supporting
\[
 F_j(t):=\varphi(t-x_j)+\varphi(t-x_j-a),\qquad 1\le j\le3,
\]
lie in $I_T$ and are mutually disjoint.  The two Tate moments of the three
$F_j$ form a $2\times3$ matrix.  Hence there is a nonzero real vector
$c=(c_1,c_2,c_3)$ such that
$F=\sum_jc_jF_j$ satisfies $M_-F=M_+F=0$.

By the disjointness of the three pairs, cross terms between different $j$
vanish.  Inside each pair exactly one overlap survives after translation by
$a$, and therefore
\[
 \Re\langle S_aF,F\rangle
 =\|\varphi\|^2\sum_{j=1}^3c_j^2>0.
\]
Expanding (79) gives (89).  A channelwise contraction would imply the
opposite weak inequality for every primitive $F$, a contradiction.
\end{proof}

This no-go result is important for the Poisson--Sonin proposal.  Polar
normalization of each local multiplier, even when isometric, cannot be the
missing $C_T$.  Any successful source construction must contain a genuinely
global off-diagonal mixing law whose Gram includes the centered atomic--
continuous cross $Q_c$.

\paragraph{Normalization audit.}
Equations (79)--(83) use all available local normalizations: the unique
self-dual coefficient $w_n=\Lambda(n)/\sqrt n$, the constant $m_0$, the
complete Gamma place, and the two exact Tate moments.  They determine the
two Grams and their signed difference, but they do not supply the contraction
$C_T$.  Producing $C_T$ from (83) by polar decomposition is legitimate only
after proving the norm inequality, and is therefore circular.  In the
threshold notation of the supplied file, the same missing assertion is the
factorization of the centered cross $Q_c$ through $D_0^{1/2}$ with unit
shorted-capacity budget.  Hence (83) is a genuine completion of the balanced
factorization bookkeeping, but not a proof of G40.

\paragraph{Notation audit.}
The symbol $\mathcal H_{5/4}$ occurs once in the supplied row (d), in the
definition of $\widehat{\mathcal R}_T$, and is not defined in any of the five
supplied files.  Formula (82) provides an explicit positive Gamma reference.
Identifying $\mathcal H_{5/4}$ with that reference requires either a stated
definition or a separately proved equality of forms; it must not be used
silently in the G40 argument.

The labels G39--G42 likewise do not occur in any of the five supplied files.
In this audit they refer only to the v83--v87 decomposition.  The authoritative
row-(d) target is (84), or equivalently the unit Douglas contraction above;
no status assigned to a G-label may override that operator inequality.

\section{Primitive-potential reduction and the continuous defect kernel}

The two Tate moments admit an exact differential parametrization which
removes the moving two-moment projection from the global estimate.  Put
\[
 L:=\partial_t^2-\frac14,
 \qquad H_0^2(I_T):=\{\phi\in H^2(I_T):
 \phi(\pm T)=\phi'(\pm T)=0\}.
\]

\begin{theorem}[Tate potential isomorphism --- \PROVED]
The map
\begin{equation}
 L:H_0^2(I_T)\longrightarrow
 \mathcal P_T=\ker M_-\cap\ker M_+                 \tag{90}
\end{equation}
is a continuous bijection onto the $L^2$ primitive space.  Its inverse is
the causal Green formula
\begin{equation}
 \phi(t)=\int_{-T}^{t}2\sinh\!\left(\frac{t-u}{2}\right)F(u)\,du.
                                                               \tag{91}
\end{equation}
After extension by zero, $L$ commutes distributionally with every ambient
translation $\tau_a$.
\end{theorem}

\begin{proof}
The kernel $k(x)=2\sinh(x/2)\mathbf1_{x\ge0}$ satisfies $Lk=\delta_0$.
Thus (91) gives $L\phi=F$ and the left endpoint conditions.  For $t\ge T$ it
equals
\[
 e^{t/2}\int_{-T}^{T}e^{-u/2}F(u)\,du
 -e^{-t/2}\int_{-T}^{T}e^{u/2}F(u)\,du,
\]
which vanishes together with its derivative exactly when $M_-F=M_+F=0$.
Hence its zero extension lies in $H^2(\mathbb R)$ and gives the right endpoint
conditions.  Conversely, integration by parts twice shows that $L\phi$ has
both moments zero for every clamped $\phi$.  The homogeneous equation
$L\phi=0$ and the two left boundary conditions give injectivity.  Standard
one-dimensional Green estimates give continuity of the inverse.  Constant-
coefficient differentiation commutes with translations on distributions.
\end{proof}

Let
\begin{equation}
 a_T(\tau):=g_\Gamma(\tau)-m_0
 -2\sum_{2\le n<e^{2T}}\frac{\Lambda(n)}{\sqrt n}
        \cos(\tau\log n).                              \tag{92}
\end{equation}

\begin{theorem}[Exact continuous defect representation --- \PROVED]
For $F=L\phi\in\mathcal P_T$,
\begin{align}
 \langle A_TF,F\rangle
 &=\frac1{2\pi}\int_{\mathbb R}
       a_T(\tau)(\tau^2+\tfrac14)^2
       |\widehat\phi(\tau)|^2\,d\tau                 \tag{93}\\
 &=\int_0^\infty w_\Gamma(u)
       \|(I-\tau_u)L\phi\|^2\,du
   -m_0\|L\phi\|^2 \notag\\
 &\quad-\sum_{2\le n<e^{2T}}\frac{\Lambda(n)}{\sqrt n}
       \bigl(\langle\tau_{\log n}L\phi,L\phi\rangle
             +\langle\tau_{-\log n}L\phi,L\phi\rangle\bigr). \tag{94}
\end{align}
Here every function is extended by zero before translation.  Thus row (d)
is equivalent to nonnegativity of the explicit scalar kernel (93) on Fourier
transforms of clamped potentials; no moving projection or unspecified Hodge
operator remains.
\end{theorem}

\begin{proof}
On the ambient line, Plancherel turns $S_a+S_{-a}$ into multiplication by
$2\cos(a\tau)$, and (81)--(82) give the Gamma term.  Since
$\widehat{L\phi}(\tau)=-(\tau^2+1/4)\widehat\phi(\tau)$, this proves (93).
Expanding the same terms before Fourier transformation gives (94).
\end{proof}

Formula (93) is not a pointwise multiplier test: $a_T$ changes sign, while
$\widehat\phi$ belongs to the Paley--Wiener class forced by support and four
clamped boundary conditions.  Replacing that constrained class by arbitrary
$L^2$ Fourier data would destroy the problem.  Any new Poisson--Tate square
decomposition must therefore act on this Paley--Wiener subspace and include
cross-place terms; theorem (89) rules out a placewise decomposition.

\subsection{The global differentiated Poisson current}

On Meyer's common nuclear Fr\'echet core let $Z$ be the Zeta convolution
operator, let $\partial$ denote multiplication by $\log x$ (the derivation of
multiplicative convolution), and put
\begin{equation}
 C:=Z\partial(Z^{-1})=\sum_{n\ge2}\Lambda(n)U_n.       \tag{95}
\end{equation}
Let $\mathscr Jf(x)=x^{-1}f(x^{-1})$ and let $\mathcal F$ be the self-dual
additive Fourier transform.

\begin{theorem}[Differentiated Poisson mixing identity --- \PROVED]
For every $f$ in the Poisson core
$\mathcal H_\cap=\{f:f(0)=\mathcal Ff(0)=0\}$,
\begin{align}
 Z\partial f+\mathscr JZ\partial\mathcal Ff
 &=CZf+\mathscr JCZ\mathcal Ff                         \tag{96}\\
 &=(C+\mathscr JC\mathscr J)Zf.                        \tag{97}
\end{align}
Thus all Euler places and both orientations occur in one source-defined
Poisson current before any positivity or spectral decomposition is invoked.
\end{theorem}

\begin{proof}
The derivation rule and $C=Z\partial(Z^{-1})$ give
$\partial Z=-CZ$.  Differentiate the Poisson identity
$Zf=\mathscr JZ\mathcal Ff$.  Since
$\partial\mathscr J=-\mathscr J\partial$, one obtains
\[
 -CZf+Z\partial f
 =\mathscr JCZ\mathcal Ff-\mathscr JZ\partial\mathcal Ff,
\]
which is (96).  Applying $\mathscr J$ to Poisson gives
$Z\mathcal Ff=\mathscr JZf$ and proves (97).
\end{proof}

Under the self-dual critical normalization, conjugation by $\mathscr J$
turns $U_n$ into the opposite weighted translation; consequently the right
side of (97) is the global source representative of the symmetric arithmetic
operator appearing in (92).  Unlike (27e), it acts on the full continuum
Poisson core and contains cross-place mixing through $Z$.

Define the unnormalized Poisson current
\begin{equation}
 \mathcal W_{\rm P}f:=Z\partial f+\mathscr JZ\partial\mathcal Ff. \tag{98}
\end{equation}
Equations (96)--(97) prove its divergence, but not its physical Gram.  The
next required identity is the continuum metric comparison
\begin{equation}
 \langle\mathcal W_{\rm P}f,
          (ZZ^*)^\dagger\mathcal W_{\rm P}f\rangle
 \stackrel{?}{=}
 \int_{\mathbb R}a_T(\tau)|\widehat f(\tau)|^2
       \frac{d\tau}{2\pi}+\text{explicit Tate terminals},       \tag{99}
\end{equation}
with all compressions and domains made precise.  Formula (99) is recorded as
\OPEN{}, not as an identity: $Z$ is a topological isomorphism on Meyer's
weighted space, not a bounded invertible operator on the physical critical
$L^2$ space, so $(ZZ^*)^\dagger$ cannot be imported without a closed-form
domain construction.  The restricted Sonin polar normalization (53) is the
available mechanism for such a construction, but equality with the
Gamma--Tate Gram remains to be derived.

\begin{proposition}[Poisson-current Gram homogeneity obstruction --- \PROVED]
The divergence identity (97) cannot be converted into the signed arithmetic
term in $A_T$ by a parameter-independent quadratic normalization of the
Euler current alone.  More precisely, scale the contact operator in the
right side of (97) and set
$\mathcal W_{\rm E}(\lambda):=\lambda(C+\mathscr JC\mathscr J)Zf$.
Its divergence is linear in $\lambda$, whereas every quadratic Gram
\begin{equation}
 \mathcal W_{\rm E}(\lambda)^*M\mathcal W_{\rm E}(\lambda) \tag{100}
\end{equation}
with $M$ independent of $\lambda$ is quadratic in $\lambda$.  Hence no
contact-scale-stable identity of this form can equal the affine family
$G_{\Gamma,T}-m_0I-\lambda(C+C^*)$.
\end{proposition}

\begin{proof}
By definition
$\mathcal W_{\rm E}(\lambda)=\lambda\mathcal W_{\rm E}(1)$.
Therefore (100) is multiplied by $\lambda^2$, while the arithmetic
piece of (92) is multiplied by $\lambda$.  Equality on an interval of
$\lambda$ is impossible unless the arithmetic operator vanishes.
\end{proof}

This homogeneity argument does not rule out an isolated numerical coincidence
at $\lambda=1$; it rules out deriving the desired linear form as a structural
Gram identity of the differentiated Euler current.  A valid construction
must use an affine background/cross term or, equivalently, a metric variation.

The correct linear object is the first variation of the Zeta metric.  For
the weighted Euler family
\[
 Z_t=\sum_{n\ge1}n^{-t/2}U_n,\qquad
 C_t=\sum_{n\ge2}\Lambda(n)n^{-t/2}U_n,\qquad
 G_t=Z_t^*Z_t,
\]
one has on the common core
\begin{align}
 Z_t'&=-\frac12Z_tC_t,                                 \tag{101}\\
 G_t'&=-\frac12(C_t^*G_t+G_tC_t),                     \tag{102}\\
 -2\frac d{dt}\|Z_tf\|^2
 &=\langle(C_t^*G_t+G_tC_t)f,f\rangle.                \tag{103}
\end{align}
These identities are algebraic consequences of the Euler logarithmic
derivative and require no sign assumption.

Thus the viable global theorem is not positivity of a current norm, but
monotonicity of a \emph{normalized} Poisson--Sonin metric after the complete
Gamma and Tate boundary corrections have been included.  Writing
$\widetilde G_t$ for that still-to-be-constructed physical normalized Gram,
the required source theorem has the form
\begin{equation}
 -2P_T\widetilde G_t'P_T
 =\mathcal R_t^*\mathcal R_t,
 \qquad t=1,                                           \tag{104}
\end{equation}
with $\widetilde G_t$ and $\mathcal R_t$ defined independently of the left
side.  Equation (104) would supply the missing global square.  The raw
$G_t=Z_t^*Z_t$ does not suffice: its derivative (102) has no automatic sign,
and restricted Sonin polar normalization makes its internal Gram constant,
moving the whole derivative into the cutoff/Tate embeddings.  Computing
those boundary derivatives is therefore the remaining load-bearing step.

\subsection{Exact motion and energy of the two Tate terminals}

The Tate part of that boundary motion can be computed explicitly.  On
$H_T=L^2(-T,T)$ define, for $t>0$,
\begin{equation}
 B_tF:=\binom{\displaystyle\int_{-T}^TF(x)e^{-tx/2}\,dx}
                 {\displaystyle\int_{-T}^TF(x)e^{tx/2}\,dx},
 \qquad P_t:=I-B_t^*K_t^{-1}B_t,                       \tag{105}
\end{equation}
where
\begin{equation}
 K_t:=B_tB_t^*
 =\begin{pmatrix}
 2\sinh(tT)/t&2T\\2T&2\sinh(tT)/t
 \end{pmatrix}.                                       \tag{106}
\end{equation}
For $t>0$ the eigenvalues of $K_t$ are
$2\sinh(tT)/t\pm2T>0$, so the inverse is ordinary, not generalized.

\begin{theorem}[Tate boundary velocity and terminal square --- \PROVED]
For $F\in\ker B_t$,
\begin{equation}
 P_t'F=-B_t^*K_t^{-1}B_t'F.                            \tag{107}
\end{equation}
This vector is the unique minimum-norm normal velocity $\nu_t(F)$ satisfying
the differentiated primitive constraint
\begin{equation}
 B_t\nu_t(F)+B_t'F=0.                                  \tag{108}
\end{equation}
Its energy is the explicit positive terminal form
\begin{equation}
 \boxed{
 \|\nu_t(F)\|^2
 =\langle B_t'F,K_t^{-1}B_t'F\rangle_{\mathbb C^2}.}   \tag{109}
\end{equation}
At $t=1$, $B_t'F$ consists exactly of the two logarithmic Tate moments
\begin{equation}
 \frac12\binom{-\int xF(x)e^{-x/2}\,dx}
                    {\int xF(x)e^{x/2}\,dx}.           \tag{110}
\end{equation}
\end{theorem}

\begin{proof}
Differentiate $P_t=I-B_t^*K_t^{-1}B_t$ and apply the result to a vector with
$B_tF=0$.  Every term containing $B_tF$ vanishes, leaving (107).  Applying
$B_t$ proves (108).  The normal space is $\Ran B_t^*$; hence any other
solution of (108) differs from (107) by an element of $\ker B_t$, orthogonal
to it.  Finally
\[
 \|B_t^*K_t^{-1}B_t'F\|^2
 =\langle B_t'F,K_t^{-1}B_tB_t^*K_t^{-1}B_t'F\rangle,
\]
which is (109).  Direct differentiation of (105) gives (110).
\end{proof}

Writing $a_t=2\sinh(tT)/t$, $b=2T$, and
$B_t'F=(u_-,u_+)^t$, the terminal square is completely scalar:
\begin{equation}
 \|\nu_t(F)\|^2
 =\frac{|u_-+u_+|^2}{2(a_t+b)}
  +\frac{|u_--u_+|^2}{2(a_t-b)}.                       \tag{111}
\end{equation}
This separates the even and odd Tate boundary modes and proves that no
undetermined normalization constant remains in the terminal contribution.

For the potential parametrization $F=L\phi$ of (90), integration by parts
and the clamped boundary conditions give
\begin{equation}
 \boxed{
 B_1'(L\phi)=\frac12
 \binom{\displaystyle\int_{-T}^T\phi(x)e^{-x/2}\,dx}
       {\displaystyle\int_{-T}^T\phi(x)e^{x/2}\,dx}.} \tag{111a}
\end{equation}
Indeed
$L(xe^{x/2})=e^{x/2}$ and
$L(xe^{-x/2})=-e^{-x/2}$.  Thus the two moving Tate terminals of the
primitive vector are exactly the two ordinary exponential moments of its
potential.  This gives a direct type-compatible link between the
Paley--Wiener reduction (93) and the terminal Gram (109).

\begin{theorem}[Kato transport of the primitive metric --- \PROVED]
Let $M_t=M_t^*$ be a differentiable ambient metric and set
\begin{equation}
 \Omega_t:=[P_t',P_t].                                 \tag{112}
\end{equation}
Then $\Omega_t^*=-\Omega_t$.  If $U_t$ solves
$U_t'=\Omega_tU_t$, $U_{t_0}=I$, it is unitary and transports the primitive
spaces:
\begin{equation}
 P_tU_t=U_tP_{t_0}.                                    \tag{113}
\end{equation}
On the fixed source $P_{t_0}H_T$, the transported compression
\begin{equation}
 \widehat M_t:=P_{t_0}U_t^*M_tU_tP_{t_0}              \tag{114}
\end{equation}
satisfies
\begin{equation}
 \widehat M_t'
 =P_{t_0}U_t^*P_t
   \bigl(M_t'+[M_t,\Omega_t]\bigr)
   P_tU_tP_{t_0}.                                      \tag{115}
\end{equation}
\end{theorem}

\begin{proof}
The commutator of two selfadjoint operators is skew-adjoint, so the evolution
is unitary.  Differentiating $P_t^2=P_t$ gives
$P_t'=P_tP_t'+P_t'P_t$ and $P_tP_t'P_t=0$.  These imply
$P_t'+[P_t,\Omega_t]=0$, which proves (113) by differentiation.  Finally,
differentiate (114), use $U_t'=\Omega_tU_t$ and
$(U_t^*)'=-U_t^*\Omega_t$, and insert (113); this gives (115).
\end{proof}

For the raw Zeta metric $M_t=G_t=Z_t^*Z_t$, equations (102) and (115) give
the completely explicit primitive derivative
\begin{equation}
 \widehat G_t'
 =P_{t_0}U_t^*P_t
 \left[-\frac12(C_t^*G_t+G_tC_t)+[G_t,\Omega_t]\right]
 P_tU_tP_{t_0}.                                       \tag{116}
\end{equation}
Here $P_t'$ and hence $\Omega_t$ are given by (105)--(107), so (116) has no
unknown connection.  The remaining issue is not differentiation but the
identification of the complete Gamma-corrected physical metric whose
derivative equals the kernel (93), followed by a proof that the resulting
operator in (115) is negative semidefinite.

\subsection{The completed Euler--Gamma logarithmic metric}

The ambient metric can be identified before taking its critical boundary
value.  Put
\begin{equation}
 \Lambda_{\mathbb Q}(s):=\pi^{-s/2}\Gamma(s/2)\zeta(s),
 \qquad s=\frac t2+i\tau,                              \tag{117}
\end{equation}
and, away from its zeros and poles,
$\mathfrak M_t(\tau):=|\Lambda_{\mathbb Q}(t/2+i\tau)|^2$.

\begin{theorem}[Completed logarithmic-metric identity --- \PROVED]
For $t>2$ one has the absolutely convergent identity
\begin{align}
 2\partial_t\log\mathfrak M_t(\tau)
 &=2\Re\frac{\Lambda_{\mathbb Q}'}
                  {\Lambda_{\mathbb Q}}(t/2+i\tau)    \tag{118}\\
 &=\Re\psi\left(\frac t4+\frac{i\tau}{2}\right)
   -\log\pi
   -2\sum_{n\ge2}\Lambda(n)n^{-t/2}
          \cos(\tau\log n).                           \tag{119}
\end{align}
At $t=1$, the Gamma part of (119) is exactly
$g_\Gamma(\tau)-m_0$.  When (119) is continued distributionally to $t=1$
and paired with $|\widehat F(\tau)|^2$ for
$F\in\mathcal P_T$, the polar terms at $s=0,1$ vanish by the two Tate
moments, and every arithmetic term with $\log n\ge2T$ pairs to zero by
support.  The resulting pairing is precisely
$\langle A_TF,F\rangle$ in (93).
\end{theorem}

\begin{proof}
Differentiate the logarithm of (117).  Since $ds/dt=1/2$, differentiating
$\log|\Lambda_{\mathbb Q}(s)|^2$ gives
$\Re(\Lambda_{\mathbb Q}'/\Lambda_{\mathbb Q})(s)$;
the extra factor two yields (118).  The logarithmic derivative of the Gamma
factor is
$-\tfrac12\log\pi+\tfrac12\psi(s/2)$, while
$\zeta'/\zeta(s)=-\sum_{n\ge2}\Lambda(n)n^{-s}$ for
$\Re s>1$.  Taking twice the real part proves (119).

At $t=1$,
$\Re\psi(1/4+i\tau/2)-\log\pi=g_\Gamma(\tau)-m_0$ by
the definitions in the supplied row (d).  The distributional continuation
is the standard explicit-formula continuation already realized in row (c).
The Mellin values at $0$ and $1$ are the two moments $M_-,M_+$, so their
polar contributions vanish.  Finally
$\langle\tau_{\log n}F,F\rangle=0$ when $\log n\ge2T$ (endpoint equality is
null), which truncates the arithmetic sum and gives (92)--(93).
\end{proof}

This theorem supplies the previously unspecified complete physical
Gamma--Euler metric variation.  It does not yet supply a positive Hilbert
boundary metric at $t=1$: $\mathfrak M_1$ vanishes at critical-line zeros,
and continuation from $t>2$ crosses possible zeros with real part greater
than $1/2$.  A proof of (104) must therefore construct a source boundary
space whose positive defect kernel includes exactly the residues encountered
in this continuation, rather than applying an inverse to
$\mathfrak M_1$ pointwise.

\begin{proposition}[Three-lines inverse-metric obstruction --- \PROVED]
Let $H$ be a Hilbert-valued holomorphic function in a vertical strip, with
the usual Hardy decay.  The function
\begin{equation}
 N_H(\sigma):=\int_{\mathbb R}|H(\sigma+i\tau)|^2\,d\tau \tag{120}
\end{equation}
is log-convex.  If it is symmetric about $\sigma=1/2$, its right derivative
at the center is nonnegative.  Applied to
$H(s)=\Lambda_{\mathbb Q}(s)\Phi(s)$, however, this controls a derivative
weighted by $|\Lambda_{\mathbb Q}|^2$ and does not imply (93).

To obtain the unweighted physical logarithmic derivative from (120), the
boundary datum would have to be divided by $\Lambda_{\mathbb Q}$ (up to an
outer factor of the same modulus).  Such a division is holomorphic across a
domain only if the datum cancels every zero there; requiring this for the
full primitive test space is equivalent to requiring the relevant domain to
be zero-free.  Therefore a bare three-lines/functional-equation argument
cannot prove row (d) without importing the missing zero information.
\end{proposition}

\begin{proof}
Log-convexity of Hardy $L^2$ norms is the Hilbert-valued three-lines theorem;
symmetry makes the center a minimum.  Differentiating
$|\Lambda_{\mathbb Q}\Phi|^2$ produces the logarithmic derivative multiplied
by $|\Lambda_{\mathbb Q}Phi|^2$, together with the derivative of $\Phi$.
The target (93) instead has the arbitrary Paley--Wiener weight
$|\widehat F|^2$.  Matching the weights requires boundary modulus
$|\Phi|=|\widehat F|/|\Lambda_{\mathbb Q}|$.  A holomorphic quotient with
this prescribed cancellation cannot exist for arbitrary $F$ at a zero of
$\Lambda_{\mathbb Q}$.  Hence the cancellation step assumes the missing
zero-free assertion.
\end{proof}

\subsection{Residue boundary space and its exact signature}

For $F\in\mathcal P_T$ use the entire Mellin--Fourier transform
\begin{equation}
 \Phi_F(z):=\int_{-T}^TF(t)e^{izt}\,dt,                \tag{121}
\end{equation}
and for a nontrivial zero $\rho=\beta+i\gamma$ put
$z_\rho=(\rho-1/2)/i=\gamma-i(\beta-1/2)$.  The involution
$\rho\mapsto1-\bar\rho$ becomes $z\mapsto\bar z$.

\begin{theorem}[Exact residue Krein boundary --- \PROVED]
After the two polar modes have been killed by $F\in\mathcal P_T$, the
distributional boundary value in (118)--(119) has the source realization
\begin{equation}
 \langle A_TF,F\rangle
 =\sum_{\mathcal O}\langle
     \mathcal E_{\mathcal O}F,
     J_{\mathcal O}\mathcal E_{\mathcal O}F\rangle,   \tag{122}
\end{equation}
where the sum is over functional-equation orbits of nontrivial zeros,
with multiplicity.  For a critical-line orbit $\rho=1/2+i\gamma$,
\begin{equation}
 \mathcal K_{\mathcal O}=\mathbb C,\qquad
 J_{\mathcal O}=1,\qquad
 \mathcal E_{\mathcal O}F=\Phi_F(\gamma).             \tag{123}
\end{equation}
For an off-line two-point orbit
$\{\rho,1-\bar\rho\}$,
\begin{equation}
 \mathcal K_{\mathcal O}=\mathbb C^2,\qquad
 J_{\mathcal O}=\begin{pmatrix}0&1\\1&0\end{pmatrix},\qquad
 \mathcal E_{\mathcal O}F=
 \binom{\Phi_F(z_\rho)}{\Phi_F(\bar z_\rho)}.         \tag{124}
\end{equation}
In the eigenbasis of $J_{\mathcal O}$ its contribution is
\begin{equation}
 \frac12|\Phi_F(z_\rho)+\Phi_F(\bar z_\rho)|^2
 -\frac12|\Phi_F(z_\rho)-\Phi_F(\bar z_\rho)|^2.     \tag{125}
\end{equation}
Thus the source residue boundary is canonically Krein; its metric is Hilbert
on every residue fiber if and only if every nontrivial zero lies on the
critical line.
\end{theorem}

\begin{proof}
The nuclear Lefschetz identity in the supplied row (c), equivalently the
contour shift of (118), identifies the primitive explicit form with the sum
of the autocorrelation transform at the nontrivial zeros.  The entire
continuation of the real-line identity
$\widehat{F*F^*}(\tau)=|\Phi_F(\tau)|^2$ is
\[
 \Phi_F(z)\overline{\Phi_F(\bar z)}.
\]
For a fixed critical zero this is $|\Phi_F(\gamma)|^2$, giving (123).  The
two members of an off-line functional-equation orbit contribute
\[
 \Phi_F(z)\overline{\Phi_F(\bar z)}
 +\Phi_F(\bar z)\overline{\Phi_F(z)},
\]
which is exactly the Gram in (124).  Diagonalizing the exchange matrix gives
(125).  Absolute convergence for compactly supported autocorrelations is the
convergence statement used by the supplied nuclear explicit formula.
\end{proof}

The negative vector in (125) is not an artifact of coordinates.  If
$z\ne\bar z$ and neither point is one of the two polar evaluation points,
the map
\begin{equation}
 F\longmapsto\bigl(M_-F,M_+F,Phi_F(z),\Phi_F(\bar z)\bigr) \tag{126}
\end{equation}
from $L^2(-T,T)$ onto $\mathbb C^4$ is surjective: its four representing
exponentials have distinct exponents and are linearly independent on every
nonempty interval.  Consequently the restricted evaluation
$F\in\mathcal P_T\mapsto(\Phi_F(z),\Phi_F(\bar z))$ is onto
$\mathbb C^2$.  Hence each off-line residue fiber contains a genuine negative
direction already on the primitive source.  A positive Hodge realization
must eliminate that fiber structurally; it cannot turn (124) positive by a
change of basis.

\subsection{A new global Hodge object: the unitarized Poisson quotient}

Let $\pi_0(u)$ denote the centrally normalized scaling action on the
source-defined quotient $\mathcal H_-^0=\mathcal H_-/Z\mathcal H_\cap$ of
the supplied row (c); central normalization means that the spectral
character attached to a zero $\rho$ is $e^{u(\rho-1/2)}$.

\begin{theorem}[Amenable unitarization Hodge bridge --- \PROVED]
Assume that $\mathcal H_-^0$ admits a Hilbert completion
$\mathscr H_{\rm P}$ satisfying the following source conditions:
\begin{enumerate}
\item the quotient map and the nuclear test-action of row (c) extend
continuously to a dense subspace of $\mathscr H_{\rm P}$;
\item $\pi_0(u)$ extends to a strongly continuous group on
$\mathscr H_{\rm P}$;
\item there is a source-derived constant $M<\infty$ such that
\begin{equation}
 \sup_{u\in\mathbb R}\|\pi_0(u)\|\le M.               \tag{127}
\end{equation}
\end{enumerate}
Then there is an equivalent positive Hilbert metric
$\langle\cdot,\cdot\rangle_{\rm Hod}$ on $\mathscr H_{\rm P}$ for which
every $\pi_0(u)$ is unitary.  Consequently the generator of the centered
scaling group is skew-adjoint, every transpose eigencharacter supplied by
row (c) has exponent in $i\mathbb R$, and all nontrivial zeros satisfy
$\Re\rho=1/2$.
\end{theorem}

\begin{proof}
Let $m$ be a translation-invariant mean on $L^\infty(\mathbb R)$ and define
\begin{equation}
 \langle x,y\rangle_{\rm Hod}
 :=m_u\,\langle\pi_0(u)x,\pi_0(u)y\rangle_{\mathscr H_{\rm P}}. \tag{128}
\end{equation}
Uniform boundedness gives the upper estimate
$\|x\|_{\rm Hod}\le M\|x\|$.  Since (127) also applies to $-u$,
$\|\pi_0(u)x\|\ge M^{-1}\|x\|$, giving the lower estimate
$\|x\|_{\rm Hod}\ge M^{-1}\|x\|$.  Thus (128) is a positive equivalent
Hilbert norm.  Translation invariance of $m$ makes the group unitary in this
metric.  Stone's theorem makes its generator skew-adjoint.  A character
$e^{u(\rho-1/2)}$ belonging to the transpose spectrum of a unitary group is
bounded for both signs of $u$, hence $\Re(\rho-1/2)=0$.
\end{proof}

The new object is (128): it is constructed without selecting spectral
subspaces and, once (127) is proved, supplies the missing positive Hodge
metric directly.  The load-bearing analytic theorem has therefore been
reduced to the source estimate (127), together with construction of the
Hilbert completion.  Sonin compression is a natural candidate norm, but the
supplied files only prove positivity of each compressed trace; they do not
prove that the centered scaling action descends to the quotient as a
uniformly bounded group.

There is an equivalent closed-range formulation.  If the image
$Z\mathcal H_\cap$ is a closed invariant subspace in a critical Hilbert
realization on which centered scaling is unitary, then the Hilbert quotient
is itself unitary and (127) holds with $M=1$.  Conversely, failure of closed
range is exactly where an off-line hyperbolic residue plane can enter the
nuclear quotient without contradicting unitarity of the ambient scaling
action.  Thus a source proof of critical closed range for the Poisson map is
a concrete alternative formulation of the missing Hodge theorem.

\begin{theorem}[Local Witt Hilbert completion and global Zeta obstruction
--- \PROVED]
Give the algebraic Witt orbit the Hilbert norm for which
$\{\phi_r\}_{r\ge1}$ is orthonormal, and give the Dirichlet coefficient
space the norm of $\ell^2(\mathbb N)$.  Then the row-(b) intertwiner
\begin{equation}
 J_2\phi_r=\delta_r                                      \tag{129}
\end{equation}
extends unitarily, and every $V_n$ (equivalently, convolution by $\delta_n$)
is an isometry.  However Meyer's global Zeta sum has no bounded extension to
this Hilbert completion, since
\begin{equation}
 Z\delta_1=\sum_{n\ge1}\delta_n\notin\ell^2(\mathbb N). \tag{130}
\end{equation}
Thus the positive Hilbert structure available in row (b) does not by itself
construct the Hilbert Poisson quotient required in (127).
\end{theorem}

\begin{proof}
Equation (129) maps one orthonormal basis to the other.  The rule
$V_n\phi_r=\phi_{nr}$ maps the basis injectively to an orthonormal subset,
so it is isometric; the same holds for
$\delta_r\mapsto\delta_{nr}$.  Formula (130) follows directly from the
definition $Z=\sum_nU_n$, and its coefficient sequence has infinite squared
norm.  A bounded operator must send $\delta_1$ to a Hilbert vector, proving
the obstruction.
\end{proof}

For $t>1$, the weighted cyclic vector
\begin{equation}
 Z_t\delta_1=\sum_{n\ge1}n^{-t/2}\delta_n             \tag{131}
\end{equation}
does lie in $\ell^2$, with norm squared $\zeta(t)$.  Its divergence as
$t\downarrow1$ is exactly the critical loss of the naive Hilbert completion.
Any successful completion must renormalize this pole while retaining the
two-dimensional Tate quotient and the nontrivial residue boundary; simply
discarding the divergent cyclic direction would also discard the Poisson
extension data.

\begin{theorem}[Critical gcd boundary state --- \PROVED]
For $t>1$ put
\begin{equation}
 e_t:=\zeta(t)^{-1/2}\sum_{n\ge1}n^{-t/2}\delta_n
 \in\ell^2(\mathbb N).                                 \tag{132}
\end{equation}
Then, for all $m,n\ge1$ and $g=(m,n)$,
\begin{equation}
 \langle V_ne_t,V_me_t\rangle
 =\left(\frac{mn}{g^2}\right)^{-t/2}.                 \tag{133}
\end{equation}
Consequently the critical kernel
\begin{equation}
 K(m,n):=\lim_{t\downarrow1}\langle V_ne_t,V_me_t\rangle
 =\frac{(m,n)}{\sqrt{mn}}                              \tag{134}
\end{equation}
is positive definite.  It extends to the group positive-definite function
on $\mathbb Q_+^\times$
\begin{equation}
 \varphi(q)=\prod_p p^{-|v_p(q)|/2}.                   \tag{135}
\end{equation}
Its GNS representation is a unitary representation of
$\mathbb Q_+^\times$, carries a cyclic vector $\Omega$, and satisfies
$\langle U_n\Omega,U_m\Omega\rangle=K(m,n)$.
\end{theorem}

\begin{proof}
The supports of $V_ne_t$ and $V_me_t$ meet when $nr=ms$.  Writing
$m=ga$, $n=gb$ with $(a,b)=1$, all solutions are $r=ak$, $s=bk$.  Therefore
\[
 \langle V_ne_t,V_me_t\rangle
 =\frac1{\zeta(t)}\sum_{k\ge1}(ak)^{-t/2}(bk)^{-t/2}
 =(ab)^{-t/2},
\]
which is (133) and tends to (134).  A pointwise limit of positive-definite
kernels is positive definite.  Since
$K(m,n)=\varphi(m^{-1}n)$ and (135) is the corresponding reduced-prime
factorization, the GNS theorem gives the unitary group representation.
\end{proof}

This boundary is fully explicit.  For each prime let $r_p=p^{-1/2}$ and put
\begin{equation}
 d\mu_p(e^{i\theta})
 :=\frac{1-r_p^2}{|1-r_pe^{i\theta}|^2}\frac{d\theta}{2\pi}. \tag{136}
\end{equation}
Its Fourier moments are
$\int e^{ik\theta}d\mu_p=r_p^{|k|}$.  Hence the GNS space in the theorem is
realized by the infinite product probability space
\begin{equation}
 \mathscr H_{\gcd}:=L^2\left(\prod_p\mathbb T,
                              \bigotimes_p\mu_p\right), \tag{137}
\end{equation}
with $U_q$ multiplication by
$\prod_pz_p^{v_p(q)}$ and $\Omega=1$.  Inversion
$q\mapsto q^{-1}$ is unitary because every $\mu_p$ is invariant under
$z_p\mapsto\bar z_p$.

The gcd boundary supplies exactly the cross-place structure absent from a
direct sum of local ports and realizes the central coefficient
$\langle U_n\Omega,\Omega\rangle=n^{-1/2}$.  It does not yet include the
archimedean Gamma channel or prove that the physical observation from
clamped potentials into $\mathscr H_{\gcd}$ has Gram (93).  That observation
is now a concrete map to construct, rather than an unspecified positive
space.

\begin{theorem}[Gcd--Gamma reference observation --- \PROVED]
Let $\mathcal A_T=\{n:2\le n<e^{2T},\ \Lambda(n)>0\}$ and let
$K_T=(K(m,n))_{m,n\in\mathcal A_T}$.  This matrix is strictly positive.
Writing $\chi_n=U_n\Omega\in\mathscr H_{\gcd}$, define
\begin{equation}
 \eta_n:=\sum_{m\in\mathcal A_T}
          (K_T^{-1/2})_{mn}\chi_m.                    \tag{138}
\end{equation}
Then $\{\eta_n:n\in\mathcal A_T\}$ is orthonormal.  If
$\mathscr K_{\Gamma,T}$ denotes the Gamma delay space in (82), the map
\begin{equation}
 \mathcal V_T^X(\mathcal X_TF):=
 \left(
   \{w_\Gamma(u)^{1/2}(I-\tau_u)E_TF\}_{u>0},
   \sum_{n\in\mathcal A_T}\eta_n\otimes
        w_n^{1/2}J_{\log n,-}F
 \right)                                               \tag{139}
\end{equation}
is a source-defined isometry from
$\overline{\Ran\mathcal X_T}$ onto its generated range in
\begin{equation}
 \mathscr K_T^X:=\mathscr K_{\Gamma,T}
 \oplus(\mathscr H_{\gcd}\otimes L^2(\mathbb R)).      \tag{140}
\end{equation}
In particular
\begin{equation}
 (\mathcal V_T^X\mathcal X_T)^*
 (\mathcal V_T^X\mathcal X_T)=\mathcal X_T^*\mathcal X_T. \tag{141}
\end{equation}
\end{theorem}

\begin{proof}
The vectors $\chi_n$ are distinct characters on the product torus (137).
Because the product Poisson measure has full support, a finite linear
combination of them which vanishes in $L^2$ is the zero trigonometric
polynomial.  Thus their Gram $K_T$ is strictly positive.  Equation (138)
gives
\[
 (\langle\eta_m,\eta_n\rangle)_{m,n}
 =K_T^{-1/2}K_TK_T^{-1/2}=I.
\]
Orthogonality makes the squared norm of the arithmetic component of (139)
equal to
$\sum_nw_n\|J_{\log n,-}F\|^2$.  Equation (82) gives equality for the Gamma
component.  Their direct sum is precisely the defining Gram of
$\mathcal X_T$, proving (141).
\end{proof}

This closes the Gram/type interface (88) after replacing the previously
unspecified Sonin target by the new source-derived gcd--Gamma Hodge reference
space (140).  It is compatible with restriction to old contact sets because
all finite Grams are compressions of the single global kernel (134), although
the positive square roots $K_T^{-1/2}$ themselves need not be nested.  The
canonical object is the unorthogonalized GNS span; (138) is only a coordinate
choice inside it.

The opposite observation is defined on the same arithmetic labels by
replacing $J_{\log n,-}$ with $J_{\log n,+}$ and adjoining the constant
channel $m_0^{1/2}F$.  The resulting angular map is exactly
$\mathfrak C_T^{\rm src}$ transported by $\mathcal V_T^X$.  Its
contractivity is still (84c); the GNS construction proves that its domain
Gram and cross-place type are correct, but does not yet prove the norm-one
bound.

\begin{theorem}[Global gcd-oriented midpoint colligation --- \PROVED]
Let $I_T=\mathcal A_T\times\{+,-\}$ and put
$\chi_{n,+}=U_n\Omega$, $\chi_{n,-}=U_{n^{-1}}\Omega$.  Let
$\mathbf K_T$ be their Gram matrix and define the orthonormal family
\begin{equation}
 \eta_{n,\epsilon}:=\sum_{(m,\delta)\in I_T}
 (\mathbf K_T^{-1/2})_{(m,\delta),(n,\epsilon)}
 \chi_{m,\delta}.                                      \tag{142}
\end{equation}
The inversion unitary $\mathfrak J_{\gcd}$ satisfies
\begin{equation}
 \mathfrak J_{\gcd}\eta_{n,+}=\eta_{n,-},
 \qquad \mathfrak J_{\gcd}\eta_{n,-}=\eta_{n,+}.     \tag{143}
\end{equation}
Let $\mathfrak F$ be any involutive unitary on $L^2(\mathbb R)$ with
$\mathfrak F\tau_a\mathfrak F=\tau_{-a}$, and define
\begin{equation}
 \mathbb U_T:=\mathfrak J_{\gcd}\otimes\mathfrak F.   \tag{144}
\end{equation}
The arithmetic feature
\begin{equation}
 \mathbb B_T^{\rm ar}F
 :=\sum_{n\in\mathcal A_T}\frac{w_n^{1/2}}2
 \left[
  \eta_{n,+}\otimes(I-\tau_{\log n})E_TF
 +\eta_{n,-}\otimes(I-\tau_{-\log n})E_TF
 \right]                                               \tag{145}
\end{equation}
satisfies
\begin{align}
 (\mathbb B_T^{\rm ar})^*\mathbb B_T^{\rm ar}
 &=\sum_{n\in\mathcal A_T}w_nJ_{\log n,-}^*J_{\log n,-}, \tag{146}\\
 \mathbb U_T\mathbb B_T^{\rm ar}
 &=\mathbb B_T^{\rm ar}\mathfrak F.                  \tag{147}
\end{align}
Consequently
\begin{equation}
 \mathbb B_T^{\rm ar}
 =\mathbb B_{T,+}^{\rm ar}+\mathbb B_{T,-}^{\rm ar},
 \qquad
 (\mathbb B_T^{\rm ar})^*\mathbb B_T^{\rm ar}
 =\sum_{\epsilon=\pm}
   (\mathbb B_{T,\epsilon}^{\rm ar})^*
    \mathbb B_{T,\epsilon}^{\rm ar},                  \tag{148}
\end{equation}
where
$\mathbb B_{T,\pm}^{\rm ar}=\tfrac12(I\pm\mathbb U_T)
\mathbb B_T^{\rm ar}$.
\end{theorem}

\begin{proof}
The characters indexed by $I_T$ are distinct, so the proof used for (138)
shows $\mathbf K_T>0$.  Inversion permutes the characters by the swap
$P(n,+)=(n,-)$, hence $P\mathbf K_TP=\mathbf K_T$.
Functional calculus gives $P\mathbf K_T^{-1/2}P=\mathbf K_T^{-1/2}$,
proving (143).  Orthogonality of the $\eta_{n,\epsilon}$ and equality of the
two translation-difference norms give
\[
 \|\mathbb B_T^{\rm ar}F\|^2
 =\sum_n\frac{w_n}{4}
  \left(\|(I-\tau_{\log n})E_TF\|^2
       +\|(I-\tau_{-\log n})E_TF\|^2\right)
 =\sum_nw_n\|J_{\log n,-}F\|^2,
\]
which is (146).  Equations (143)--(144) and
$\mathfrak F(I-\tau_a)=(I-\tau_{-a})\mathfrak F$ give (147).
The unitary midpoint identity proves (148).
\end{proof}

Adding the Gamma-oriented feature (68), integrated with weight
$w_\Gamma(u)$, gives a single source-defined reference current
$\mathbb B_T^X$ with
\begin{equation}
 (\mathbb B_T^X)^*\mathbb B_T^X=\mathcal X_T^*\mathcal X_T, \tag{149}
\end{equation}
and an exact positive midpoint conservation law.  Therefore normalization,
global cross-place typing, reference Gram and conservation are closed in this
new gcd--Gamma colligation.  The remaining load-bearing identity is to show
that its appropriate midpoint observation equals the opposite factor
$\mathcal Y_T$ (including the constant and Tate-terminal channels) with no
extra energy.  That equality would prove (84c).

\begin{theorem}[M\"obius whitening of the gcd boundary --- \PROVED]
For a finite set of primes $S$, let $m_S$ be Haar probability measure on
$\mathbb T^S$, let $r_p=p^{-1/2}$, and put
\begin{equation}
 d\mu_S(z)=\prod_{p\in S}
 \frac{1-r_p^2}{|1-r_pz_p|^2}\,dm_S(z),
 \qquad
 b_S(z)=\prod_{p\in S}
 \frac{1-r_pz_p}{\sqrt{1-r_p^2}}.                     \tag{150}
\end{equation}
Then multiplication by $b_S$ is a unitary
\begin{equation}
 M_{b_S}:L^2(\mathbb T^S,m_S)longrightarrow
          L^2(\mathbb T^S,\mu_S),                     \tag{151}
\end{equation}
and
\begin{equation}
 |b_S(z)|^2d\mu_S(z)=dm_S(z).                         \tag{152}
\end{equation}
The numerator of $b_S$ is the finite M\"obius Euler multiplier and its
denominator is the unique positive normalization making it isometric.
\end{theorem}

\begin{proof}
For each prime the Poisson density in (136) gives
\[
 \frac{|1-r_pz_p|^2}{1-r_p^2}
 \frac{1-r_p^2}{|1-r_pz_p|^2}=1.
\]
Multiplying these identities proves (152), and integration proves (151).
\end{proof}

For the deformation $r_p(t)=p^{-t/2}$ the same construction gives
$b_{S,t}$.  Its unnormalized numerator is exactly $a_{S,t}$ in (20), while
the derivative of the positive denominator is real scalar.  Passing to the
polar phase therefore yields
\begin{equation}
 (M_{b_{S,t}}^{\rm polar})^*
 \partial_tM_{b_{S,t}}^{\rm polar}
 =\frac14M_{c_{S,t}-\overline{c_{S,t}}},               \tag{153}
\end{equation}
which agrees with (26).  Thus the gcd boundary, finite Euler polar
normalization and CCM Sonin normalization are three realizations of the same
source isometry at finite $S$.

What survives after whitening is not an Euler norm defect---that Hermitian
part is exactly absorbed by the change from $m_S$ to $\mu_S$---but the motion
of the Poisson/cutoff observation inside the whitened space.  This confirms
from a second construction that the remaining G40 term must be a boundary
second fundamental form, now combining (107)--(116) with the finite-product
measure variation in (150).

\begin{theorem}[Centered Euler score and exact Fisher budget --- \PROVED]
Let $r_p(t)=p^{-t/2}$ and let $\mu_{p,t}$ be the Poisson measure (136) with
$r_p$ replaced by $r_p(t)$.  Its logarithmic score is
\begin{equation}
 \sigma_{p,t}(z):=\partial_t\log\frac{d\mu_{p,t}}{dm}(z)
 =(\log p)\left[
 \frac{r_p(t)^2}{1-r_p(t)^2}
 -\Re\frac{r_p(t)z}{1-r_p(t)z}
 \right].                                              \tag{154}
\end{equation}
It is centered and has the exact Fisher energy
\begin{equation}
 \int_{\mathbb T}\sigma_{p,t}\,d\mu_{p,t}=0,
 \qquad
 \int_{\mathbb T}|\sigma_{p,t}|^2d\mu_{p,t}
 =\frac{(\log p)^2r_p(t)^2}
        {2(1-r_p(t)^2)^2}.                             \tag{155}
\end{equation}
For a finite set $S$, the product score
$\sigma_{S,t}=\sum_{p\in S}\sigma_{p,t}$ satisfies
\begin{equation}
 \|\sigma_{S,t}\|_{L^2(\mu_{S,t})}^2
 =\sum_{p\in S}\frac{(\log p)^2r_p(t)^2}
        {2(1-r_p(t)^2)^2}.                             \tag{156}
\end{equation}
Moreover, for every differentiable observable $H_t$,
\begin{equation}
 \partial_t\int H_t\,d\mu_{S,t}
 =\int\partial_tH_t\,d\mu_{S,t}
  +\int(H_t-\mathbb E_tH_t)\sigma_{S,t}\,d\mu_{S,t}. \tag{157}
\end{equation}
\end{theorem}

\begin{proof}
Differentiate the density in (150), using
$r_p'=-(\log p)r_p/2$, to obtain (154).  The Poisson moments are
$\int z^k d\mu_{p,t}=r_p(t)^k$ for $k\ge0$, so the expression is centered.
Put $r=r_p(t)$ and $W=rz/(1-rz)=\sum_{k\ge1}r^kz^k$.
The same moments give
\[
 \mathbb EW=\frac{r^2}{1-r^2},\qquad
 \operatorname{Var}(\Re W)=\frac{r^2}{2(1-r^2)^2}.
\]
This proves (155).  Independence makes distinct prime scores orthogonal and
gives (156).  Differentiation under the finite product integral gives (157);
centering the second term uses $\mathbb E_t\sigma_{S,t}=0$.
\end{proof}

Expanding (154),
\begin{equation}
 \sigma_{p,t}(z)
 =(\log p)\sum_{k\ge1}r_p(t)^k
       \bigl(r_p(t)^k-\Re z^k\bigr),                  \tag{158}
\end{equation}
so it contains every prime power with the central coefficient and subtracts
exactly its Poisson expectation.  Formula (157) is therefore a probability-
space realization of the centered atomic cross denoted $Q_c$ in the supplied
threshold analysis.  Cauchy--Schwarz gives the sharp covariance estimate
\begin{equation}
 |\operatorname{Cov}_t(H,\sigma_{S,t})|^2
 \le\operatorname{Var}_t(H)\,
      \|\sigma_{S,t}\|_2^2.                           \tag{159}
\end{equation}
To finish G40 by this route one must identify the old defect form $D_0$ with
the relevant conditional variance and show that the born/reference
normalization makes the Fisher factor in (159) exactly the available unit
budget.  The score and its complete prime-power budget are now explicit;
that final variance identification remains \OPEN{}.

\begin{theorem}[Exact arithmetic innovation budget --- \PROVED]
Order the finite contact characters by their integer label.  Let
$\mathscr E_{N-1}$ be the span of the characters $\chi_m=U_m\Omega$ with
$\Lambda(m)>0$ and $m<N$, and suppose $\Lambda(N)>0$.  Write $K_{N-1}$ for
their Gram matrix and
\begin{equation}
 k_N=(\langle\chi_m,\chi_N\rangle)_{m<N},
 \qquad
 \kappa_N:=1-k_N^*K_{N-1}^{-1}k_N.                    \tag{160}
\end{equation}
Then
\begin{equation}
 0<\kappa_N\le1,
 \qquad
 \|P_{\mathscr E_{N-1}}\chi_N\|^2
 =k_N^*K_{N-1}^{-1}k_N\le1,                           \tag{161}
\end{equation}
and the normalized innovation
\begin{equation}
 r_N:=\kappa_N^{-1/2}
 \left(\chi_N-\sum_{m<N}(K_{N-1}^{-1}k_N)_m\chi_m\right) \tag{162}
\end{equation}
is a unit vector orthogonal to $\mathscr E_{N-1}$.
\end{theorem}

\begin{proof}
The enlarged character family is linearly independent in the full-support
product measure, so its Gram matrix
\[
 \begin{pmatrix}K_{N-1}&k_N\\k_N^*&1\end{pmatrix}
\]
is strictly positive.  Its Schur complement is $\kappa_N>0$; since the
subtracted term is nonnegative, $\kappa_N\le1$.  The normal equations for
orthogonal projection give the coefficients $K_{N-1}^{-1}k_N$, and direct
expansion gives (161)--(162).
\end{proof}

This theorem pays the exact unit Schur budget for every purely arithmetic
birth in the global gcd metric, with all prime-power correlations retained.
It strengthens mere all-depth summability: no estimate is used.  To transfer
(160)--(162) to the physical threshold capacity (41i), one must prove that
the Gamma--continuous and Tate-terminal columns enlarge
$\mathscr E_{N-1}$ so that the physical centered load $b_E$ is precisely the
projection load in this augmented Hilbert space.  That is the remaining
mixed observation identity.

\begin{corollary}[Closed prime-power innovation --- \PROVED]
Augment the old contact space by the Tate cyclic vector,
\begin{equation}
 \mathscr E_{N-1}^{\rm aug}
 :=\operatorname{span}\bigl(\Omega,
       \{\chi_m:\Lambda(m)>0,\ m<N\}\bigr).           \tag{163}
\end{equation}
If $N=p^k$, with the convention $\chi_{p^0}=\Omega$, then
\begin{align}
 P_{\mathscr E_{N-1}^{\rm aug}}\chi_{p^k}
 &=p^{-1/2}\chi_{p^{k-1}},                             \tag{164}\\
 \left\|\chi_{p^k}-
 P_{\mathscr E_{N-1}^{\rm aug}}\chi_{p^k}\right\|^2
 &=1-p^{-1}.                                           \tag{165}
\end{align}
The normalized innovations
\begin{equation}
 e_{p,k}(z):=\frac{(z_p-p^{-1/2})z_p^{k-1}}
                   {\sqrt{1-p^{-1}}}                  \tag{166}
\end{equation}
are orthonormal as $(p,k)$ varies.
\end{corollary}

\begin{proof}
For one Poisson factor with parameter $r=p^{-1/2}$,
\begin{equation}
 |z-r|^2d\mu_p(z)=(1-r^2)dm(z).                        \tag{167}
\end{equation}
Therefore the functions
$(z-r)z^{k-1}/\sqrt{1-r^2}$ are orthonormal and have mean zero.  Independence
makes functions belonging to different primes orthogonal.  Since
\[
 z_p^k-rz_p^{k-1}=(z_p-r)z_p^{k-1},
\]
the right side is orthogonal to the augmented old span, proving
(164)--(166).
\end{proof}

Thus the entire prime-power return chain is an exact AR(1)-type orthogonal
innovation process in the gcd boundary.  The factor $1-p^{-1}$ is the same
local denominator appearing in the M\"obius whitening (150).  This provides
an inverse-free, all-depth arithmetic capacity factorization; no Witt-moment
majorant is needed for the arithmetic part.

\begin{theorem}[Archimedean Gamma GNS boundary --- \PROVED]
Put $a_j=2j+\tfrac12$ and let
\begin{equation}
 d\nu_{a_j}(x):=\frac{a_j}{2}e^{-a_j|x|}\,dx
 \qquad(x\in\mathbb R).                                \tag{168}
\end{equation}
On $L^2(\nu_{a_j})$ let $U_\tau$ be multiplication by $e^{i\tau x}$ and
$\Omega_j=1$.  Then
\begin{align}
 \langle U_\tau\Omega_j,\Omega_j\rangle
 &=\frac{a_j^2}{a_j^2+\tau^2},                         \tag{169}\\
 \frac1{a_j}\|(U_\tau-I)\Omega_j\|^2
 &=\frac{2\tau^2}{a_j(a_j^2+\tau^2)}.                 \tag{170}
\end{align}
Consequently
\begin{equation}
 g_\Gamma(\tau)
 =\sum_{j\ge0}\frac1{a_j}\|(U_\tau-I)\Omega_j\|^2,  \tag{171}
\end{equation}
and the infinite product boundary has characteristic function
\begin{equation}
 \prod_{j\ge0}\frac{a_j^2}{a_j^2+\tau^2}
 =\frac{|\Gamma(\frac14+\frac{i\tau}{2})|^2}
        {\Gamma(\frac14)^2}.                           \tag{172}
\end{equation}
\end{theorem}

\begin{proof}
The Fourier transform of the two-sided exponential density (168) is
$a_j^2/(a_j^2+\tau^2)$, proving (169); expanding the difference norm proves
(170).  The digamma series with
$a_j=2(j+1/4)$ is exactly (171).  Finally the Weierstrass product for the
ratio
$\Gamma(1/4+i\tau/2)\Gamma(1/4-i\tau/2)/\Gamma(1/4)^2$
is the product in (172).
\end{proof}

For the metric deformation set $a_j(t)=2j+t/2$ and let $\nu_{a_j(t)}$ be
as in (168).  Its centered score and Fisher energy are
\begin{equation}
 \sigma_{\Gamma,j,t}(x)
 =\frac12\left(\frac1{a_j(t)}-|x|\right),
 \qquad
 \|\sigma_{\Gamma,j,t}\|_2^2=\frac1{4a_j(t)^2}.       \tag{173}
\end{equation}
Indeed $|x|$ is exponentially distributed with rate $a_j(t)$.  The scores
are independent in the product boundary and
$\sum_j1/(4a_j(t)^2)<\infty$, so the complete Gamma score is an honest
$L^2$ random variable.  Combining it with the finite Euler score (154)
gives a single centered Gamma--Euler score whose Fisher energies add.

The constant $m_0$ is not an external correction to this boundary.  It is
the scale coordinate paired with the Gamma ratio coordinate, as follows.

\begin{theorem}[Gamma ratio--scale normalization --- \PROVED]
Let $R_+,R_-$ be independent Gamma random variables with shape
$a=\tfrac14$ and unit scale, and put
\begin{equation}
 X:=\frac12\log\frac{R_+}{R_-},
 \qquad
 Y:=\log\pi-\frac12\log(R_+R_-).                    \tag{174}
\end{equation}
Then
\begin{align}
 \phi_\Gamma(\tau):=\mathbb E e^{i\tau X}
 &=\frac{|\Gamma(\frac14+\frac{i\tau}{2})|^2}
         {\Gamma(\frac14)^2},                       \tag{175}\\
 \mathbb EY&=\log\pi-\psi\left(\frac14\right)=m_0,\tag{176}\\
 \mathbb E\!\left(Ye^{i\tau X}\right)
 &=\bigl(m_0-g_\Gamma(\tau)\bigr)\phi_\Gamma(\tau).
                                                               \tag{177}
\end{align}
Moreover the conditional scale is the explicit positive function
\begin{equation}
 h_\Gamma(x):=\mathbb E(Y\mid X=x)
 =\log\pi-\psi\left(\frac12\right)+\log(2\cosh x)>0. \tag{178}
\end{equation}
Thus the finite positive measure $h_\Gamma\,d\mu_X$ has Fourier transform
$\phi_\Gamma(\tau)(m_0-g_\Gamma(\tau))$.
The ratio law itself is explicit:
\begin{equation}
 d\mu_X(x)=\frac{\sqrt2}{B(\frac14,\frac14)}
             (\cosh x)^{-1/2}\,dx.                  \tag{178a}
\end{equation}
\end{theorem}

\begin{proof}
For a unit-scale Gamma variable $R$ of shape $a$,
$\mathbb E R^z=\Gamma(a+z)/\Gamma(a)$.  Independence with
$z=\pm i\tau/2$ proves (175), and
$\mathbb E\log R=\psi(a)$ proves (176).  Differentiating the two Mellin
moments gives
\[
 \mathbb E\!\left(e^{i\tau X}\log R_+\right)
 =\phi_\Gamma(\tau)\psi\left(a+\frac{i\tau}{2}\right),
 \quad
 \mathbb E\!\left(e^{i\tau X}\log R_-\right)
 =\phi_\Gamma(\tau)\psi\left(a-\frac{i\tau}{2}\right).
\]
Substitution in (174), followed by
$g_\Gamma(\tau)-m_0=\Re\psi(a+i\tau/2)-\log\pi$, proves (177).

For (178), write $S=R_++R_-$ and $U=R_+/S$.  The beta--Gamma algebra says
that $S\sim\mathrm{Gamma}(2a,1)$, $U\sim\mathrm{Beta}(a,a)$, and $S,U$
are independent.  Also
$X=\tfrac12\log(U/(1-U))$ and
\[
 Y=\log\pi-\log S-\frac12\log(U(1-U)).
\]
For $a=1/4$, $U(1-U)=(4\cosh^2X)^{-1}$ and
$\mathbb E\log S=\psi(1/2)$, giving (178).  Positivity follows at
$x=0$ from $\psi(1/2)=-\gamma-2\log2$ and then from $\cosh x\ge1$.
The change of variables $u=e^{2x}/(1+e^{2x})$ in the beta density gives
(178a).
\end{proof}

Equation (177) is the exact Gamma--Tate normalization identity sought in
the preceding constructions: the Gamma multiplier and the constant channel
are respectively the ratio and scale matrix coefficients of one canonical
Mellin probability space.  It is stronger than merely checking the numerical
value of $m_0$.

It does not, by itself, prove (84c).  Passing from the positive Fourier
transform in (177) to the unweighted multiplier $m_0-g_\Gamma$ requires
division by $\phi_\Gamma$.  Stirling's formula gives
\begin{equation}
 \phi_\Gamma(\tau)^{-1/2}
 \asymp e^{\pi|\tau|/4}|\tau|^{1/4}
 \qquad(|\tau|\to\infty),                            \tag{179}
\end{equation}
so this deconvolution is unbounded on $L^2(\mathbb R)$ and on the general
Paley--Wiener image of compactly supported $L^2$ sources.  Consequently
(177) may be used as a weighted boundary identity, but treating its
deconvolution as a Hilbert-space contraction would require exponential
regularity absent from the ambient compact-window source.

In fact the two Tate constraints cannot supply that regularity.

\begin{proposition}[Gamma deconvolution domain no-go --- \PROVED]
Let $F\in L^2(I_T)$ and let $E_TF$ be its zero extension.  If
\begin{equation}
 \phi_\Gamma^{-1/2}\widehat{E_TF}\in L^2(\mathbb R), \tag{180}
\end{equation}
then $F=0$.  This remains true after restriction to the primitive space
$\mathcal P_T$.  Therefore no nonzero physical compact-window source can be
obtained by applying the unbounded Gamma deconvolution in (179), even after
the Tate moment conditions.
\end{proposition}

\begin{proof}
By Stirling's formula, (180) implies
$e^{b|\tau|}\widehat{E_TF}(\tau)\in L^2(\mathbb R)$ for every
$0<b<\pi/4$ (a polynomial factor is harmless after slightly decreasing
$b$).  Fourier inversion and Cauchy--Schwarz then extend $E_TF$ to a
holomorphic function on the strip $|\Im z|<b$.  Its boundary value on the
real axis vanishes on each open half-line outside $[-T,T]$.  The identity
theorem therefore gives $E_TF=0$.  The primitive space is a subspace of
$L^2(I_T)$, so the same conclusion holds there.
\end{proof}

Thus the weighted matrix coefficient (177) is a valid normalization theorem,
but deconvolution is not the missing G40 observation.  A successful bridge
must retain the Gamma ratio law inside its Hilbert metric and combine it with
the Euler/Tate channels before taking the physical compression.

This combined metric is exact in the absolutely convergent half-plane.

\begin{theorem}[Coherent completed Euler--Gamma metric --- \PROVED]
For $t>2$ put $r_p(t)=p^{-t/2}$ and evaluate the finite/infinite product
Poisson density on the scaling orbit $z_p(\tau)=p^{-i\tau}$:
\begin{equation}
 D_{E,t}(\tau):=prod_p
 \frac{1-r_p(t)^2}{|1-r_p(t)z_p(\tau)|^2}
 =\frac{|\zeta(\frac t2+i\tau)|^2}{\zeta(t)}.        \tag{181}
\end{equation}
Let
\begin{equation}
 \phi_{\Gamma,t}(\tau):=
 \frac{|\Gamma(\frac t4+\frac{i\tau}{2})|^2}
      {\Gamma(\frac t4)^2},
 \qquad
 c_t:=\pi^{-t/2}\Gamma\left(\frac t4\right)^2\zeta(t).
                                                               \tag{182}
\end{equation}
Then the completed metric has the exact factorization
\begin{equation}
 \left|\Lambda_{\mathbb Q}\left(\frac t2+i\tau\right)\right|^2
 =c_tD_{E,t}(\tau)\phi_{\Gamma,t}(\tau).             \tag{183}
\end{equation}
Thus the Euler Poisson boundary and the Gamma ratio boundary are not merely
analogous: their normalized coherent-orbit density is exactly the completed
Zeta metric, up to the scalar $c_t$.
\end{theorem}

\begin{proof}
Absolute convergence for $t>2$ permits multiplication of the Euler factors:
\[
 \prod_p(1-p^{-t})=\zeta(t)^{-1},\qquad
 \prod_p|1-p^{-t/2-i\tau}|^{-2}
 =|\zeta(t/2+i\tau)|^2,
\]
which proves (181).  Multiplication by (182) and the definition
$\Lambda_{\mathbb Q}(s)=\pi^{-s/2}\Gamma(s/2)\zeta(s)$ prove (183).
\end{proof}

The identity also locates the critical difficulty without hiding it in a
formal product.  As $t\downarrow1$, $c_t$ has the pole of $\zeta(t)$ while
the normalized Euler density loses its ordinary product-probability meaning;
only their product in (183) has the meromorphic continuation used in (119).
Consequently a critical Hilbert compression must renormalize the pair
$c_tD_{E,t}$ jointly.  Substituting the pointwise continuation of (183) for
such a renormalized boundary operator would be circular.

The required renormalization cannot be a vector limit in the original
$\ell^2$ completion.

\begin{theorem}[Euler coherent-orbit escape of mass --- \PROVED]
For $t>1$ let $e_t$ be the unit vector (131), and let
$W_\tau\delta_n=n^{-i\tau}\delta_n$ on $\ell^2(\mathbb N)$.  Then
\begin{equation}
 \left\langle W_\tau e_t,W_\sigma e_t\right\rangle
 =\frac{\zeta(t+i(\tau-\sigma))}{\zeta(t)}.           \tag{184}
\end{equation}
For every fixed $\tau$, $W_\tau e_t\rightharpoonup0$ as $t\downarrow1$,
although its norm is one.  Moreover
\begin{equation}
 \lim_{t\downarrow1}
 \left\langle W_\tau e_t,W_\sigma e_t\right\rangle
 =\begin{cases}1,&\tau=\sigma,\\0,&\tau\ne\sigma.
 \end{cases}                                          \tag{185}
\end{equation}
For $t>2$, the formal boundary-sum functional is absolutely defined on this
vector and satisfies
\begin{equation}
 \left|\sum_{n\ge1}(W_\tau e_t)_n\right|^2
 =\frac{|\zeta(\frac t2+i\tau)|^2}{\zeta(t)}
 =D_{E,t}(\tau).                                      \tag{186}
\end{equation}
Thus (181) is the squared boundary evaluation of the coherent vector, but
that evaluation and the vector itself have no ordinary critical limit.
The kernel (185) also cannot be realized by nonzero vectors indexed by all
$\tau\in\mathbb R$ in a separable Hilbert space.
\end{theorem}

\begin{proof}
Expanding in the standard basis gives
\[
 \langle W_\tau e_t,W_\sigma e_t\rangle
 =\zeta(t)^{-1}\sum_{n\ge1}n^{-t-i(\tau-\sigma)},
\]
which is (184).  Every fixed coordinate tends to zero because
$\zeta(t)^{-1/2}\to0$; boundedness of the norms and density of finitely
supported vectors imply weak convergence to zero.  If $\tau\ne\sigma$,
$\zeta(t+i(\tau-\sigma))$ has a finite limit while $\zeta(t)\to\infty$,
proving (185).  Equation (186) follows by summing the coordinates when
$t/2>1$.  Finally, a separable Hilbert space contains at most a countable
orthonormal set, whereas the limiting kernel in (185) indexes an uncountable
one.
\end{proof}

The gcd limit (134) uses the multiplicative isometries $V_n$, not the
continuous diagonal orbit $W_\tau$.  The two constructions therefore capture
different matrix coefficients of the same precritical vector $e_t$.  A valid
Poisson--Tate bridge must join them after removing the escaping polar
direction at the level of quadratic forms.  Identifying either limit with the
other without this quotient is an \AUDIT{} error.

There is nevertheless a canonical critical limit at the level of forms.

\begin{theorem}[Rigged critical Euler form --- \PROVED]
For $t>1$ set $k_t(u)=\zeta(t+iu)$.  If $f\in\mathcal S(\mathbb R)$ and
\begin{equation}
 \widehat f(x):=\int_{\mathbb R}f(\tau)e^{-i\tau x}\,d\tau,
\end{equation}
then
\begin{align}
 \mathfrak q_t(f)
 &:=\iint_{\mathbb R^2}f(\tau)\overline{f(\sigma)}
       k_t(\tau-\sigma)\,d\tau d\sigma                 \notag\\
 &=\sum_{n\ge1}n^{-t}|\widehat f(\log n)|^2\ge0.      \tag{187}
\end{align}
Moreover
\begin{equation}
 \lim_{t\downarrow1}\mathfrak q_t(f)
 =\mathfrak q_1(f)
 :=\sum_{n\ge1}\frac1n|\widehat f(\log n)|^2.         \tag{188}
\end{equation}
Hence $k_t\to\zeta(1+i\,\cdot)$ in the sense of positive tempered
distributions, where the boundary value on the right is the Fourier transform
of the positive locally finite measure
\begin{equation}
 \mu_{E,\mathrm{crit}}:=\sum_{n\ge1}\frac1n\delta_{\log n}. \tag{189}
\end{equation}
The completion of the form (188) is
$L^2(\mu_{E,\mathrm{crit}})$ (modulo the null space of the sampling map).
It is the canonical form-level renormalization of the escaping coherent
orbit.
\end{theorem}

\begin{proof}
For $t>1$, absolute convergence of the Dirichlet series and Fubini give
(187).  For every $N>1$, Schwartz decay gives
$|\widehat f(\log n)|\ll_N(1+\log n)^{-N}$, and
\[
 \sum_{n\ge2}\frac1{n(1+\log n)^{2N}}<\infty.
\]
Dominated convergence proves (188).  Equations (187)--(188) identify the
distributional Fourier transform and prove positivity.  The last assertion
is the standard completion of a sampled $L^2$ form.
\end{proof}

This critical form produces the arithmetic reference weights without an
Euler product or a sign assumption.  On the vector-valued space
$L^2(\mu_{E,\mathrm{crit}};L^2(\mathbb R))$, define for a fixed window
\begin{equation}
 (D_Tg)(\log n):=
 \mathbf1_{\{2\le n<e^{2T}\}}\mathbf1_{\{\Lambda(n)>0\}}
 (\Lambda(n)\sqrt n)^{1/2}g(\log n).                 \tag{190}
\end{equation}
Then, exactly,
\[
 \|D_Tg\|^2
 =\sum_{\substack{2\le n<e^{2T}\\\Lambda(n)>0}}
   \frac{\Lambda(n)}{\sqrt n}\|g(\log n)\|^2.
\]
Taking $g(\log n)=J_{\log n,-}F$ recovers the complete arithmetic part of
$\|\mathcal X_TF\|^2$.  Thus the arithmetic reference port is now obtained
as a diagonal observation of a source-defined critical Euler form, rather
than inserted as an unrelated direct sum.

What (187)--(190) do not yet provide is the opposite $J_{\log n,+}$ channel
or its global mixing with Gamma and the two Tate terminals.  That is precisely
the observation, rather than the reference Gram, which remains in G40.

The Euler critical form and the Gamma ratio law admit a positive joint
completion.

\begin{theorem}[Joint critical Euler--Gamma rigged boundary --- \PROVED]
For $t>1$ let
\begin{equation}
 \phi_{\Gamma,t}(u)=
 \frac{|\Gamma(\frac t4+\frac{iu}{2})|^2}
      {\Gamma(\frac t4)^2},
 \qquad
 K_t^{E\Gamma}(u):=\zeta(t+iu)\phi_{\Gamma,t}(u).     \tag{191}
\end{equation}
Then $K_t^{E\Gamma}$ is a positive-definite tempered distribution.  For
every $f\in\mathcal S(\mathbb R)$ its quadratic form has a finite critical
limit
\begin{equation}
 \mathfrak q_1^{E\Gamma}(f)
 =\int_{\mathbb R}|\widehat f(x)|^2
       d(\mu_{E,\mathrm{crit}}*\mu_{\Gamma,1})(x)\ge0, \tag{192}
\end{equation}
where $\mu_{\Gamma,1}$ is the probability law of the Gamma ratio $X$ in
(174).  Equivalently,
\begin{equation}
 K_t^{E\Gamma}\longrightarrow
 K_1^{E\Gamma}:=\zeta(1+i\,\cdot)\phi_{\Gamma,1}
 \quad\hbox{as positive tempered distributions}.     \tag{193}
\end{equation}
Consequently
\begin{equation}
 \mathscr H_{E\Gamma}^{\mathrm{crit}}
 :=L^2(\mu_{E,\mathrm{crit}}*\mu_{\Gamma,1})          \tag{194}
\end{equation}
is a canonical positive critical Euler--Gamma boundary constructed without
$\Delta_T$, the Weil sign, or a zero-free hypothesis.
\end{theorem}

\begin{proof}
The first factor in (191) is the Fourier transform of
$\sum_n n^{-t}\delta_{\log n}$ and the second is the characteristic function
of the Gamma ratio.  Their product is therefore the Fourier transform of the
convolution of two positive measures.  The density (178a), with $1/4$
replaced by $t/4$, has exponential tails uniformly for $t$ in a compact
right-neighbourhood of one.  Combined with the Schwartz decay of
$\widehat f$ and the summability argument in the proof of (188), this gives
an integrable majorant for the convolved sampling sum.  Dominated convergence
proves (192)--(193).  Positivity and the Hilbert completion (194) follow.
\end{proof}

This theorem constructs the positive target space required in the open
Tate-boundary statement (34) at the level of the global critical
Euler--Gamma form.  To prove (35), one must still construct the actual
intertwiner $\Phi_T$ into the windowed subspace of (194) and prove that its
compressed norm equals the Schur boundary form $J_T$.  The target is now
source-defined independently of that sign; the remaining equality is no
longer allowed to choose its codomain ad hoc.

The relation between this form boundary and the gcd boundary is exact, not
heuristic.

\begin{theorem}[Divisor-incidence realization of the gcd boundary --- \PROVED]
Let
\begin{equation}
 \mathscr E_N:=\ell^2\left(\{1,\ldots,N\},\frac1k\right),
 \qquad H_N:=\sum_{k\le N}\frac1k,
\end{equation}
and, for $n\le N$, define the truncated Zeta-incidence column
\begin{equation}
 z_n^{(N)}(k):=\sqrt n\,\mathbf1_{\{n\mid k\}}.
                                                               \tag{195}
\end{equation}
Then
\begin{equation}
 \frac1{H_N}\langle z_m^{(N)},z_n^{(N)}\rangle_{\mathscr E_N}
 =\frac{(m,n)}{\sqrt{mn}}
   \frac{H_{\lfloor N/[m,n]\rfloor}}{H_N},           \tag{196}
\end{equation}
where $[m,n]$ denotes the least common multiple.  Consequently, for fixed
$m,n$,
\begin{equation}
 \lim_{N\to\infty}\frac1{H_N}
 \langle z_m^{(N)},z_n^{(N)}\rangle
 =\frac{(m,n)}{\sqrt{mn}}=K(m,n).                    \tag{197}
\end{equation}
The inverse incidence relation is the exact finite M\"obius formula
\begin{equation}
 \mathbf1_{\{k=n\}}
 =\sum_{d\le N/n}\mu(d)\mathbf1_{\{nd\mid k\}}
 =\sum_{d\le N/n}\frac{\mu(d)}{\sqrt{nd}}
       z_{nd}^{(N)}(k).                               \tag{198}
\end{equation}
Thus the critical gcd GNS state is the harmonic-density renormalization of
the Zeta divisor-incidence columns in the rigged Euler boundary, and its
whitening is precisely M\"obius inversion.
\end{theorem}

\begin{proof}
The simultaneous support of the two columns consists of multiples of
$[m,n]$.  Hence
\[
 \langle z_m^{(N)},z_n^{(N)}\rangle
 =\sqrt{mn}\sum_{j\le N/[m,n]}\frac1{j[m,n]}
 =\frac{(m,n)}{\sqrt{mn}}H_{\lfloor N/[m,n]\rfloor},
\]
which proves (196).  Since $H_{\lfloor N/c\rfloor}/H_N\to1$ for fixed $c$,
(197) follows.  Finally, if $n\mid k$, the last sum in (198) is
$\sum_{d\mid k/n}\mu(d)$, equal to one only when $k=n$; if $n\nmid k$ it
is zero.  This proves (198).
\end{proof}

\begin{corollary}[Deformed incidence/GNS equivalence --- \PROVED]
For $t>1$ equip sequences with weight $k^{-t}$ and put
$z_{n,t}(k)=n^{t/2}\mathbf1_{\{n\mid k\}}$.  With
$H_M^{(t)}=\sum_{j\le M}j^{-t}$,
\begin{equation}
 \langle z_{m,t}^{(N)},z_{n,t}^{(N)}\rangle
 =\left(\frac{(m,n)}{\sqrt{mn}}\right)^t
   H_{\lfloor N/[m,n]\rfloor}^{(t)}.                \tag{199}
\end{equation}
After $N\to\infty$ and normalization by $\zeta(t)$,
\begin{equation}
 \frac1{\zeta(t)}\langle z_{m,t},z_{n,t}\rangle
 =\left(\frac{(m,n)}{\sqrt{mn}}\right)^t
 =\langle V_ne_t,V_me_t\rangle.                     \tag{200}
\end{equation}
Thus the cyclic spans of normalized incidence columns and of the shifted
Zeta vector $e_t$ are canonically isometric.  Their common critical limit is
the gcd GNS space (134)--(137).
If $A_t\delta_k=k^{t/2}\delta_k$, regarded as a unitary from the standard
finite $\ell^2$ space to the space with weight $k^{-t}$, then the Zeta
transport (27a) satisfies the exact identity
\begin{equation}
 z_{n,t}^{(N)}=A_t\mathscr Z_t\delta_n.               \tag{201}
\end{equation}
Hence the incidence model is unitarily the same finite Zeta transport used
in G39, before passing to either critical limit.
\end{corollary}

\begin{proof}
The calculation proving (196), with $k^{-1}$ replaced by $k^{-t}$, gives
(199).  Letting $N\to\infty$ gives the factor $\zeta(t)$; (133) gives the
last equality in (200).  Equality of all Gram matrices defines the canonical
isometry of the generated cyclic spans, and their pointwise Gram limit is
(134).  Finally, at $k=nd$ the $k$-th coordinate of
$A_t\mathscr Z_t\delta_n$ is
$(nd)^{t/2}d^{-t/2}=n^{t/2}$, and it vanishes when $n\nmid k$; this is
(201).
\end{proof}

Equations (195)--(198) close the previously open compatibility between the
finite Zeta gauge used in G39 and the global gcd GNS space: both arise from
one divisor-incidence system, and the normalizing scalar is forced by the
escaping harmonic mass.  In particular, the unorthogonalized incidence
columns are nested under $N\mapsto N'$, solving the nesting defect of the
coordinate square roots $K_T^{-1/2}$ noted after (141).

This does not yet identify the Gamma/Tate opposite observation.  It does,
however, place its arithmetic divergence (27e)--(27h), its reference Gram
(139)--(149), and the M\"obius normalization (150)--(153) in a single
source-derived boundary model.

There is a necessary affine component at the Gamma place.

\begin{theorem}[Linear/affine Gamma boundary separation --- \PROVED]
Let $\mathscr K_\Gamma=\bigoplus_{j\ge0}L^2(\nu_{a_j})$ and let
$\mathbf U_\tau=\bigoplus_jU_{j,\tau}$.  The vector-valued function
\begin{equation}
 b_\Gamma(\tau):=\bigoplus_{j\ge0}
 a_j^{-1/2}(U_{j,\tau}-I)\Omega_j                  \tag{202}
\end{equation}
is a continuous $1$-cocycle:
\begin{equation}
 b_\Gamma(\tau+\sigma)
 =b_\Gamma(\tau)+\mathbf U_\tau b_\Gamma(\sigma),
 \qquad
 \|b_\Gamma(\tau)\|^2=g_\Gamma(\tau).              \tag{203}
\end{equation}
In contrast, in the infinite tensor-product Gamma-ratio representation,
the coherent orbit $\Omega_\Gamma(\tau)$ satisfies
\begin{equation}
 \langle\Omega_\Gamma(\tau),\Omega_\Gamma(0)\rangle
 =\phi_\Gamma(\tau),
 \qquad
 \|\Omega_\Gamma(\tau)-\Omega_\Gamma(0)\|^2
 =2(1-\phi_\Gamma(\tau))\le2.                         \tag{204}
\end{equation}
There is no bounded operator $B$ such that
$B(\Omega_\Gamma(\tau)-\Omega_\Gamma(0))=b_\Gamma(\tau)$ for all $\tau$.
Thus the Gamma coherent metric and the Gamma energy feature are respectively
the linear and affine parts of the boundary and cannot be identified by a
bounded change of coordinates.
\end{theorem}

\begin{proof}
The cocycle identity follows from
$U_{\tau+\sigma}-I=(U_\tau-I)+U_\tau(U_\sigma-I)$.
Equations (170)--(171) show that (202) belongs to the direct sum and give its
norm in (203).  The tensor-product coefficient is (172), and the difference
norm gives (204).  Finally $g_\Gamma(\tau)\to\infty$ like $\log|\tau|$ by
the digamma asymptotic, whereas the norm in (204) is bounded.  A bounded
$B$ would contradict these two estimates.
\end{proof}

Consequently the source-defined positive target obtained so far must be
enlarged from (194) to an affine Hodge boundary carrying both
$\mathscr H_{E\Gamma}^{\mathrm{crit}}$ and the cocycle space
$\mathscr K_\Gamma$.  The Gamma channel of $\mathcal X_T$ lands in the
second component by (202)--(203); the scale vector (177) belongs to the first.
The missing G40 intertwiner is therefore a coupling between these two
components and the moving Tate terminal, not an isometry between either pair
alone.

The arithmetic component of that coupling has an exact lossless realization.

\begin{theorem}[Prime Julia colligation and M\"obius innovation --- \PROVED]
For a prime $p$ put $r=p^{-1/2}$, $s=(1-r^2)^{1/2}$, and
\begin{equation}
 \mathcal J_p:=
 \begin{pmatrix}r&s\\s&-r\end{pmatrix}.              \tag{205}
\end{equation}
Then $\mathcal J_p$ is a selfadjoint unitary colligation and its transfer
function is the Blaschke factor
\begin{equation}
 \Theta_p(z)=-r+zs(1-zr)^{-1}s
 =\frac{z-r}{1-rz},\qquad |\Theta_p(z)|=1\quad(|z|=1). \tag{206}
\end{equation}
With the M\"obius whitening multiplier
$b_p(z)=(1-rz)/s$, the prime-power innovations (166) obey
\begin{equation}
 e_{p,k}(z)=b_p(z)\Theta_p(z)z^{k-1}.                 \tag{207}
\end{equation}
For every finite set of primes $S$, the tensor product
$\bigotimes_{p\in S}\mathcal J_p$ is unitary and has transfer multiplier
$\Theta_S=\prod_{p\in S}\Theta_p$.  Hence the complete finite arithmetic
whitening/innovation system is a source-defined lossless scattering network.
\end{theorem}

\begin{proof}
The rows of (205) are orthonormal, so $\mathcal J_p^*\mathcal J_p=I$.
The standard transfer formula $D+zC(I-zA)^{-1}B$, with
$A=r$, $B=C=s$, $D=-r$, gives (206).  Its boundary modulus is one because
$|z-r|=|1-rz|$ on the unit circle.  Finally
\[
 b_p(z)\Theta_p(z)z^{k-1}
 =\frac{(z-r)z^{k-1}}s=e_{p,k}(z),
\]
which proves (207).  Finite tensor products of unitaries are unitary, and
scalar transfer functions multiply under the corresponding cascade.
\end{proof}

Together, (195)--(207) prove amplitude typing, nested Gram, M\"obius
normalization, prime-power innovation and positive conservation for the
arithmetic \emph{label} network.  The remaining physical cell must identify
its lossless output with the $J_{\log n,+}$ delay observation while coupling
the scale vector (177), the Gamma cocycle (202), and the moving two-moment
square (109).  Only after that joint affine Gamma--Tate termination is proved
does the network output equal the full factor $\mathcal Y_T$.

The intrinsic normalization of the Archimedean cell can now be closed.

\begin{theorem}[Exact affine Gamma--Tate termination identity --- \PROVED]
For $t>0$ set $a_j(t)=2j+t/2$ and
\begin{align}
 g_{\Gamma,t}(\tau)
 &:=\sum_{j\ge0}\frac{2\tau^2}
 {a_j(t)(a_j(t)^2+\tau^2)},                            \tag{208}\\
 c_\Gamma(t)&:=\pi^{-t/2}\Gamma\left(\frac t4\right)^2.
\end{align}
With $\phi_{\Gamma,t}$ from (191), one has the exact logarithmic metric
identity
\begin{equation}
 2\partial_t\log\!\left(c_\Gamma(t)\phi_{\Gamma,t}(\tau)\right)
 =g_{\Gamma,t}(\tau)+\psi\left(\frac t4\right)-\log\pi. \tag{209}
\end{equation}
At the critical endpoint this becomes
\begin{equation}
 \left.2\partial_t\log\!\left(
 c_\Gamma(t)\phi_{\Gamma,t}(\tau)\right)\right|_{t=1}
 =g_\Gamma(\tau)-m_0.                                 \tag{210}
\end{equation}
Thus the positive Gamma coherent metric, its affine energy cocycle and the
constant Tate scale channel form one differentiable Hilbert metric cell; the
normalization $m_0$ is forced by its scalar metric derivative.
\end{theorem}

\begin{proof}
For one factor,
\[
 \partial_t\log\frac{a_j(t)^2}{a_j(t)^2+\tau^2}
 =\frac1{a_j(t)}-\frac{a_j(t)}{a_j(t)^2+\tau^2}
 =\frac{\tau^2}{a_j(t)(a_j(t)^2+\tau^2)}.
\]
The series and its derivative converge locally uniformly, so summing gives
$2\partial_t\log\phi_{\Gamma,t}=g_{\Gamma,t}$.  Also
\[
 2\partial_t\log c_\Gamma(t)
 =-\log\pi+\psi(t/4).
\]
Adding proves (209).  At $t=1$, (171) identifies
$g_{\Gamma,1}=g_\Gamma$ and
$\psi(1/4)-\log\pi=-m_0$, proving (210).
\end{proof}

Equation (210) is the non-deconvolved form of (177).  It supplies the exact
normalization of the affine Gamma--Tate termination in the moving-metric
picture: the coherent metric must be transported, not cancelled.  The
remaining G40 identity is therefore the equality between the Kato-transported
physical compression of this cell, joined to the Julia arithmetic cascade,
and the opposite delay block $\mathcal Y_T$, including the already computed
second fundamental form (109).

The Euler cell has the complementary exact normalization.

\begin{theorem}[Exact coherent Euler--Gamma bulk normalization --- \PROVED]
For $t>2$, with $s=t/2+i\tau$, one has
\begin{align}
 2\partial_t\log\!\left(\zeta(t)D_{E,t}(\tau)\right)
 &=2\Re\frac{\zeta'}{\zeta}(s)                       \notag\\
 &=-2\sum_{n\ge2}\Lambda(n)n^{-t/2}
                  \cos(\tau\log n),                  \tag{211}\\
 2\partial_t\log\!\left(
 c_\Gamma(t)\zeta(t)D_{E,t}(\tau)
 \phi_{\Gamma,t}(\tau)\right)
 &=\Re\psi\left(\frac t4+\frac{i\tau}{2}\right)
   -\log\pi
   -2\sum_{n\ge2}\Lambda(n)n^{-t/2}
                  \cos(\tau\log n).                 \tag{212}
\end{align}
The product inside the logarithm in (212) is exactly
$|\Lambda_{\mathbb Q}(s)|^2$.  Its distributional continuation to $t=1$,
paired with a source supported in $I_T$, is
\begin{equation}
 g_\Gamma(\tau)-m_0
 -2\sum_{2\le n<e^{2T}}\frac{\Lambda(n)}{\sqrt n}
       \cos(\tau\log n)=a_T(\tau).                  \tag{213}
\end{equation}
Thus the entire Gamma, constant and prime-power bulk symbol of the
authoritative row-(d) operator is the covariant logarithmic derivative of the
single positive coherent Euler--Gamma metric, with no fitted scalar and no
omitted prime powers.
\end{theorem}

\begin{proof}
Equation (181) gives
$\zeta(t)D_{E,t}(\tau)=|\zeta(s)|^2$.  Differentiation yields the first line
of (211), and the absolutely convergent logarithmic Euler derivative gives
the series.  Adding (209), or differentiating (183) directly, proves (212).
The distributional continuation and compact-support localization were proved
in (93) and (117)--(119): a translation by $\log n\ge2T$ has zero quadratic
pairing with the zero-extended source, and the endpoint contact has null
overlap.  At $t=1$, (210) supplies $g_\Gamma-m_0$, proving (213).
\end{proof}

Equations (208)--(213) complete the normalization identity of the proposed
G40 bridge.  They do not assert that the derivative of the compressed metric
is positive: a differentiable family of positive metrics need not be
monotone.  The last load-bearing step is now solely the Kato/Julia compression
identity which must turn that covariant derivative, after removal of (109),
into the squared norm of the physical output residual.

The moving Gamma boundary embeddings themselves are explicit.

\begin{theorem}[Beta--Gamma Markov transport --- \PROVED]
Let $0<t<t'$, put $a=t/4$, $b=t'/4$, and let $X_t$ have the Gamma-ratio law
with characteristic function $\phi_{\Gamma,t}$.  If $B_+,B_-$ are independent
$\mathrm{Beta}(a,b-a)$ variables and
\begin{equation}
 Z_{t,t'}:=\frac12\log\frac{B_+}{B_-},                \tag{214}
\end{equation}
then there is a coupling with
\begin{equation}
 X_t=X_{t'}+Z_{t,t'},\qquad X_{t'}\perp Z_{t,t'},     \tag{215}
\end{equation}
and
\begin{equation}
 \mathbb E e^{i\tau Z_{t,t'}}
 =\frac{\phi_{\Gamma,t}(\tau)}{\phi_{\Gamma,t'}(\tau)}
 =\left|
 \frac{\Gamma(a+\frac{i\tau}{2})\Gamma(b)}
      {\Gamma(a)\Gamma(b+\frac{i\tau}{2})}
 \right|^2.                                          \tag{216}
\end{equation}
Consequently
\begin{equation}
 (\mathcal C_{t,t'}f)(x)
 :=\mathbb E\bigl[f(x+Z_{t,t'})\bigr]                \tag{217}
\end{equation}
defines a Markov contraction
$L^2(\mu_{\Gamma,t})\to L^2(\mu_{\Gamma,t'})$ under the coupling (215).
Its action on characters is multiplication by the explicit coefficient
(216).  These maps give a source-defined differentiable transport of the
Gamma boundary; no spectral sign or choice of $\Delta_T$ occurs.
\end{theorem}

\begin{proof}
Take independent Gamma variables $R_{a,\pm}$ of shape $a$ and
$G_{b-a,\pm}$ of shape $b-a$.  The beta--Gamma algebra gives
$R_{b,\pm}=R_{a,\pm}+G_{b-a,\pm}$ and
$B_\pm=R_{a,\pm}/R_{b,\pm}$, with $B_\pm$ independent of
$R_{b,\pm}$.  Taking half the difference of logarithms proves (215).
The beta Mellin formula
$\mathbb EB^z=\Gamma(a+z)\Gamma(b)/(\Gamma(a)\Gamma(b+z))$
proves (216).  Finally (217) is conditional expectation in the joint model,
and is therefore positivity preserving, unital and contractive on $L^2$.
\end{proof}

The infinitesimal logarithmic derivative of the character coefficient in
(216) is $-g_{\Gamma,t}(\tau)/2$, while the scalar metric derivative is the
remaining term in (209).  Hence (214)--(217) give the previously missing
boundary motion whose covariant generator is the affine Gamma--Tate cell.
What remains is to tensor this Markov transport with the finite Julia cascade
and prove that compression by the moving two-moment projector agrees with the
authoritative old/new decomposition.

The Euler transport exists as well, but with the opposite orientation.

\begin{theorem}[Oppositely oriented Euler and Gamma Markov transports
--- \PROVED]
Let $0<t<t'$ and, for every prime, put
$r_{p,t}=p^{-t/2}$ and $q_p=p^{-(t'-t)/2}$.  If
$Z_{p,t}$ has Poisson law $\mu_{r_{p,t}}$ and $Y_p$ independently has
Poisson law $\mu_{q_p}$, then
\begin{equation}
 Z_{p,t'}\ \buildrel d\over=\ Z_{p,t}Y_p,             \tag{218}
\end{equation}
because Poisson measures on the circle satisfy
$\mu_r*\mu_q=\mu_{rq}$.  Hence, on every finite prime window,
\begin{equation}
 (\mathcal C^E_{t,t'}f)(z)
 :=\mathbb E_Y f(zY)                                  \tag{219}
\end{equation}
is a Markov contraction
$L^2(\mu_{E,t'})\to L^2(\mu_{E,t})$.
The Gamma contraction (217), however, has direction
\begin{equation}
 \mathcal C^\Gamma_{t,t'}:
 L^2(\mu_{\Gamma,t})\longrightarrow
 L^2(\mu_{\Gamma,t'}).                               \tag{220}
\end{equation}
Thus the canonical Euler and Gamma conditional expectations point in
opposite directions.  There is no tensor product of these two
\emph{forward} conditional expectations implementing both metric motions
with the same time orientation.
\end{theorem}

\begin{proof}
The $k$-th Fourier coefficient of $\mu_r*\mu_q$ is
$r^{|k|}q^{|k|}=(rq)^{|k|}$, proving (218).  Jensen's inequality in the
coupling $Z_{t'}=Z_tY$ proves (219).  Equation (220) is (217).  Since the
domain of one canonical conditional expectation is the codomain of the other
with the time indices reversed, their tensor product is not a Markov map in
either common direction.
\end{proof}

The forward-arrow mismatch does not obstruct use of the Hilbert adjoint of
the Euler map.  That adjoint supplies the reverse arrow canonically.

\begin{theorem}[Bayes-reversed joint Euler--Gamma transport --- \PROVED]
Let $\mathcal C^E_{t,t'}$ and $\mathcal C^\Gamma_{t,t'}$ be (219) and
(217), with the Euler product restricted to a finite prime set $S$.  Define
\begin{equation}
 \mathcal R^E_{t,t'}:=(\mathcal C^E_{t,t'})^*:
 L^2(\mu_{E,t})\longrightarrow L^2(\mu_{E,t'}).      \tag{220a}
\end{equation}
Then $\mathcal R^E_{t,t'}$ is a positivity-preserving unital contraction
with the Bayes formula
\begin{equation}
 (\mathcal R^E_{t,t'}h)(z')
   =\mathbb E[h(Z_t)\mid Z_{t'}=z'],
 \qquad Z_{t'}=Z_tY.                                 \tag{220b}
\end{equation}
Consequently
\begin{equation}
 \boxed{\mathcal M_{t,t'}:=
 \mathcal R^E_{t,t'}\otimes\mathcal C^\Gamma_{t,t'}:
 L^2(\mu_{E,t}\otimes\mu_{\Gamma,t})\to
 L^2(\mu_{E,t'}\otimes\mu_{\Gamma,t'})}            \tag{220c}
\end{equation}
is a source-defined Markov contraction in one common direction.  In the
independent product coupling it is joint conditional expectation, and
\begin{equation}
 \|H\|_t^2-\|\mathcal M_{t,t'}H\|_{t'}^2
 =\mathbb E\!\left[
    \operatorname{Var}(H(Z_t,X_t)\mid Z_{t'},X_{t'})
   \right]\ge0.                                     \tag{220d}
\end{equation}
No polar factor, Weil sign or $\Delta_T$ enters this construction.
\end{theorem}

\begin{proof}
For $h\in L^2(\mu_{E,t})$ and $f\in L^2(\mu_{E,t'})$, (218) gives
\[
 \langle h,\mathcal C^E_{t,t'}f\rangle_{\mu_{E,t}}
 =\mathbb E[\overline{h(Z_t)}f(Z_{t'})]
 =\langle\mathbb E[h(Z_t)\mid Z_{t'}],f\rangle_{\mu_{E,t'}},
\]
with the harmless conjugation changed if the opposite inner-product
convention is used.  This proves (220b).  An adjoint of a contraction is a
contraction; (220b) proves positivity and preservation of the constant.

The beta--Gamma coupling gives $X_t=X_{t'}+Z_{t,t'}$ with the added variable
independent of $X_{t'}$.  Taking it independently of the Euler coupling
shows that (220c) is
$\mathbb E[H(Z_t,X_t)\mid Z_{t'},X_{t'}]$.  Orthogonal projection in
$L^2$ gives Pythagoras and proves (220d).
\end{proof}

\begin{theorem}[Orthogonal joint innovation Gram --- \PROVED]
Write $R=\mathcal R^E_{t,t'}$, $G=\mathcal C^\Gamma_{t,t'}$ and
$M=R\otimes G$.  Then
\begin{align}
 I-M^*M
 &=(I-R^*R)\otimes I
   +R^*R\otimes(I-G^*G)                               \tag{220f}\
 &=I\otimes(I-G^*G)
   +(I-R^*R)\otimes G^*G.                             \tag{220g}
\end{align}
In particular, with $D_R=(I-R^*R)^{1/2}$ and
$D_G=(I-G^*G)^{1/2}$, the analysis map
\begin{equation}
 \mathcal D_{E\Gamma}H:=
 \bigl((D_R\otimes I)H, (R\otimes D_G)H\bigr)       \tag{220h}
\end{equation}
satisfies
\begin{equation}
 \|H\|^2-\|MH\|^2=\|\mathcal D_{E\Gamma}H\|^2.
                                                               \tag{220i}
\end{equation}
Thus the ambient Euler and Gamma losses have an exact positive joint Gram
with no cross term and no double counting.
\end{theorem}

\begin{proof}
Since $M^*M=R^*R\otimes G^*G$, add and subtract $R^*R\otimes I$ to obtain
(220f), or add and subtract $I\otimes G^*G$ to obtain (220g).  Both summands
in (220f) are positive because $R$ and $G$ are contractions.  Taking their
canonical square roots in the displayed order gives (220h), and its squared
norm is (220f), proving (220i).
\end{proof}

\begin{theorem}[Honest critical joint boundary at every finite window ---
\PROVED]
For fixed $T$, put
\begin{equation}
 S_T:=\{p\text{ prime}:\log p<2T\}.                  \tag{220k}
\end{equation}
Then $S_T$ is finite, the product Poisson probability $\mu_{S_T,t}$ and the
Gamma-ratio probability $\mu_{\Gamma,t}$ are defined at $t=1$, and all maps
(220a)--(220i) exist as ordinary Hilbert-space contractions for
$t=1$, $t'=1+\varepsilon$.  Moreover, every power $p^k$ with
$k\log p\ge2T$ has zero compressed shift on $I_T$ and therefore contributes
nothing to the physical quadratic form.  Thus no infinite-Euler analytic
continuation is needed to define the \emph{ambient} critical Bayes defect for
a fixed window.
\end{theorem}

\begin{proof}
The inequality $p<e^{2T}$ leaves finitely many primes.  A finite product of
the Poisson measures (150) is a probability for every $t>0$ because
$0<p^{-t/2}<1$.  The beta--Gamma construction (214)--(217) is valid for every
$0<t<t'$, hence in particular at $t=1$.  The conditional expectations,
their adjoints, tensor product and defect operators are consequently bounded
ordinary Hilbert-space maps there.

If $a=k\log p\ge2T$, the zero-extension compression
$S_a=E_T^*\tau_aE_T$ is zero: the translated and original open intervals do
not overlap except possibly at a null endpoint.  Hence all such prime powers
drop from (83), (92) and (213).  This proves the assertion.
\end{proof}

\begin{corollary}[Location of the critical obstruction --- \PROVED]
For a fixed window, the word ``renormalized'' in (220j) concerns only the
physical source embedding: its M\"obius whitening, moving cutoff/primitive
projection, and affine scalar/Tate terminals.  It does not concern existence
or positivity of the finite-prime ambient conditional expectation
$\mathcal M_{1,1+\varepsilon}$ or its joint defect Gram; those are already
ordinary positive Hilbert objects.
\end{corollary}

\begin{proof}
This is the separation between Theorem (220k) and the source pullback in
(220j).  The former constructs the ambient maps at $t=1$; the latter asks
whether the physical embedding intertwines their innovation ports with the
authoritative balanced factors.
\end{proof}

\begin{corollary}[Bayes pullback criterion for G40 --- \PROVED]
For $t>2$, let $\mathcal E_{T,t}$ denote the coherent physical embedding
constructed in (181)--(194), augmented by M\"obius whitening and the affine
Tate channel (105)--(111).  Suppose it admits a critical renormalized tangent
realization $\mathcal E_{T,1}^{\rm ren}$ for which
\begin{equation}
 \boxed{
 \operatorname*{form-lim}_{\varepsilon\downarrow0}
 \varepsilon^{-1}
 (\mathcal D_{E\Gamma}^{(1,1+\varepsilon)}
       \mathcal E_{T,1}^{\rm ren})^*
 (\mathcal D_{E\Gamma}^{(1,1+\varepsilon)}
       \mathcal E_{T,1}^{\rm ren})
 =P_T(\mathcal X_T^*\mathcal X_T
             -\mathcal Y_T^*\mathcal Y_T)P_T.}       \tag{220j}
\end{equation}
Then G40 holds.  For (220j) to be a source identity rather than a restatement,
the realization must be obtained as the form limit of the displayed coherent
embeddings and its projection derivative must be the already computed Tate
square (109), transported by (112)--(116); it may not be defined from the
right side of (220j).
\end{corollary}

\begin{proof}
The left side of (220j) is positive by (220f)--(220i).  Hence (220j) gives
$P_TA_TP_T\ge0$, which is (84c); the exact positivity-threshold theorem gives
G40.  The restrictions on the construction prevent defining the left side as
a square root of $A_T$ and therefore state the required non-circular scope.
\end{proof}

\paragraph{Status of (220j).}
The ambient operator $\mathcal D_{E\Gamma}$ is now explicit and positive.
The physical signed Gram and Tate correction are also explicit.  What is not
yet proved is that $\mathcal E_{T,t}$ has a critical renormalized tangent with a
form derivative whose Euler defect lands in the physical plus-delay port and
whose Gamma defect lands in the four-state-head/tail port.  This is the same
boundary-layer matching isolated in (343)--(345), now written as one joint
Gram identity.  Equation (220j) remains \OPEN{} and is not used as a lemma.

\begin{corollary}[Canonical Julia dilation of the Bayes reversal --- \PROVED]
The Bayes contraction $\mathcal R^E_{t,t'}$ has the canonical Julia unitary
\begin{equation}
 \mathbf J(\mathcal R^E)=
 \begin{pmatrix}
  \mathcal R^E&(I-\mathcal R^E(\mathcal R^E)^*)^{1/2}\\
  (I-(\mathcal R^E)^*\mathcal R^E)^{1/2}&-(\mathcal R^E)^*
 \end{pmatrix}.                                      \tag{220e}
\end{equation}
On the forward Euler character $z_p^k$, the forward-map norm defect is
$1-p^{-k(t'-t)}$ and its infinitesimal normalization is (221)--(223); in
(220e) the forward and adjoint defect spaces occupy complementary Julia
ports.
Thus reversing the Euler arrow introduces no sign hypothesis.  Identification
of this abstract Bayes-Julia dilation with the separate fixed-critical
whitening colligation (205)--(207) is not asserted here.
\end{corollary}

\begin{proof}
The standard Julia-operator multiplication, using
$\mathcal R^E(I-(\mathcal R^E)^*\mathcal R^E)^{1/2}
=(I-\mathcal R^E(\mathcal R^E)^*)^{1/2}\mathcal R^E$,
shows that (220e) is unitary.  The forward map sends $z_p^k$ to
$p^{-k(t'-t)/2}z_p^k$ by (221), and both characters have norm one in their
respective probability spaces.  Hence its norm defect on that vector is
$1-p^{-k(t'-t)}$.  The Julia matrix contains the two defect operators
$D_{\mathcal C^E}$ and $D_{(\mathcal C^E)^*}=D_{\mathcal R^E}$ in its
complementary off-diagonal ports, while (222)--(223) give the stated
forward-character limit.  The final sentence records
the exact scope: (205)--(207) encode a different fixed-critical whitening
realization and require an additional intertwining proof.
\end{proof}

The final G40 bridge is now the concrete intertwining statement that the
joint conditional expectation (220c), after M\"obius whitening, affine
Gamma normalization and Kato compression by $P_t$, has output feature
exactly $\mathcal Y_TP_T$ and orthogonal variance equal to the terminal square
(109) plus the old-compatible residual.  Unlike a postulated contraction,
the ambient positive map and its defect have now been constructed.

The infinitesimal Euler defects recover the authoritative prime-power weights
exactly.

\begin{theorem}[Prime-power Markov defect normalization --- \PROVED]
Let $\varepsilon=t'-t>0$.  On the $k$-th character of the $p$-th Poisson
factor, the Euler Markov contraction (219) acts by
\begin{equation}
 \mathcal C^E_{t,t'}z_p^k
 =p^{-k\varepsilon/2}z_p^k.                           \tag{221}
\end{equation}
Put $r_p=p^{-1/2}$.  The Julia defect of this scalar contraction, amplified
in an independent port for every prime power, has normalized squared energy
\begin{equation}
 d_{p,k}(\varepsilon)^2
 :=\frac{r_p^k}{k\varepsilon}
       \left(1-p^{-k\varepsilon}\right).             \tag{222}
\end{equation}
Then
\begin{equation}
 \lim_{\varepsilon\downarrow0}d_{p,k}(\varepsilon)^2
 =\frac{\log p}{p^{k/2}}
 =\frac{\Lambda(p^k)}{\sqrt{p^k}}.                  \tag{223}
\end{equation}
Thus every arithmetic coefficient in (83), including all prime powers, is
the positive infinitesimal defect rate of a source-defined Euler Markov/Julia
cell.  The factor $1/k$ is forced by the logarithmic Euler expansion and is
exactly cancelled by the character decay rate $k\log p$.
\end{theorem}

\begin{proof}
Equation (221) follows because the $k$-th moment of the independent Poisson
noise in (218) is $q_p^k=p^{-k\varepsilon/2}$.  A scalar contraction of
modulus $q_p^k$ has squared Julia defect $1-q_p^{2k}=1-p^{-k\varepsilon}$,
which gives (222) after the logarithmic Euler weight and the critical central
amplitude are inserted.  Finally
$(1-e^{-k\varepsilon\log p})/\varepsilon\to k\log p$,
proving (223).
\end{proof}

\begin{theorem}[Midpoint Julia two-port rotation --- \PROVED]
Let $U$ be a unitary, set
\begin{equation}
 J_-F:=\frac{UF-F}{\sqrt2},\qquad
 J_+F:=\frac{UF+F}{\sqrt2},                          \tag{223a}
\end{equation}
and let $0\le q\le1$, $d=(1-q^2)^{1/2}$.  With the Hadamard unitary
$H=2^{-1/2}\left(\begin{smallmatrix}1&1\\1&-1\end{smallmatrix}\right)$,
\begin{equation}
 H\begin{pmatrix}q&d\\d&-q\end{pmatrix}H
 =\begin{pmatrix}d&q\\q&-d\end{pmatrix}.           \tag{223b}
\end{equation}
Consequently the two physical midpoint ports have the lossless rotation
\begin{equation}
 \binom{O_+}{O_-}
 :=\begin{pmatrix}d&q\\q&-d\end{pmatrix}
       \binom{J_-F}{J_+F}
 =\binom{dJ_-F+qJ_+F}{qJ_-F-dJ_+F}.                 \tag{223c}
\end{equation}
They satisfy
\begin{align}
 \|O_+\|^2&=d^2\|J_-F\|^2+q^2\|J_+F\|^2,\notag\\
 \|O_-\|^2&=q^2\|J_-F\|^2+d^2\|J_+F\|^2,          \tag{223d}\\
 \|O_+\|^2+\|O_-\|^2
 &=\|J_-F\|^2+\|J_+F\|^2=2\|F\|^2.                \tag{223e}
\end{align}
For $q=p^{-k\varepsilon/2}$, $O_+\to J_+F$ and
$O_-\to J_-F$ as $\varepsilon\downarrow0$, while
$d^2=1-p^{-k\varepsilon}$ has precisely the infinitesimal weight in
(222)--(223).  Thus the local Julia dilation really produces the physical
two-port interpolation and the correct defect scale.  This statement alone
does not construct $J_+F$ from $J_-F$, because $J_+F$ is already the second
input in (223c).
\end{theorem}

\begin{proof}
Direct multiplication gives (223b), and its matrix is orthogonal/unitary.
Moreover
\[
 \langle J_-F,J_+F\rangle
 =\frac12\bigl(\langle UF,F\rangle
                -\langle F,UF\rangle\bigr)
\]
is purely imaginary.  Since $q,d$ are real, its contribution to each squared
norm in (223c) has zero real part.  This proves the two identities in (223d).
Their sum, or unitarity of (223b), gives (223e); the last equality uses
$\|UF\|=\|F\|$.  The limiting assertions and the formula for $d^2$ are
immediate.
\end{proof}

\begin{proposition}[Scope of the midpoint rotation --- \AUDIT]
For every active prime power $n=p^k$, take $U=\tau_{\log n}$ on the ambient
zero-extension line and $q=p^{-k\varepsilon/2}$.  The first component in
(223c) converges strongly to $J_{\log n,+}F$, and its $dJ_-$ coupling, after
multiplication by $p^{-k/2}/(k\varepsilon)$ at the squared-norm level,
converges to
\begin{equation}
 \frac{\Lambda(p^k)}{\sqrt{p^k}}\,
       \|J_{\log n,-}F\|^2.                          \tag{223f}
\end{equation}
This verifies the coefficient and the lossless two-port normalization, but
not the Bayes pullback: the input vector in (223c) already contains
$J_{\log n,+}F$.  To prove (220j), one must construct that second input from
an independent Gamma--Tate/old storage channel or show that it is the actual
Bayes main output.  Relabelling the pre-existing $J_+$ entry as a generated
output would be circular.
\end{proposition}

\begin{proof}
Ambient translation is unitary even though its compression is not.  Apply
Theorem (223a)--(223e) before compressing back to $I_T$.  Strong convergence
follows from $q\to1$, $d\to0$.  Multiplying
$d^2=1-p^{-k\varepsilon}$ by the stated central normalization and using
(223) proves (223f).  The displayed computation does not provide the second
input of (223c), which proves the scope assertion.
\end{proof}

\begin{theorem}[Zeroth-order scattering is necessary --- \PROVED]
Let $\mathcal M_\varepsilon:\mathscr K_1\to\mathscr K_{1+\varepsilon}$ be
the Bayes joint transports (220c), identified in a common continuous Hilbert
field so that $\mathcal M_\varepsilon\to I$ strongly.  Let
$E_X,E_Y:H\to\mathscr K_1$ be bounded source embeddings.  If
\begin{equation}
             \mathcal M_\varepsilon E_XF\longrightarrow E_YF
             \quad\hbox{for every }F\in H,           \tag{223g}
\end{equation}
then $E_X=E_Y$.  In particular a strongly continuous Markov transport alone
cannot convert the physical minus port $J_{a,-}$ into $J_{a,+}$ on a window
where the translated copies overlap.
\end{theorem}

\begin{proof}
Strong convergence gives
$\mathcal M_\varepsilon E_XF\to E_XF$ for every fixed $F$.  Uniqueness of
limits and (223g) imply $E_XF=E_YF$ for all $F$.

For the physical ports,
\begin{equation}
 J_{a,+}^*J_{a,+}-J_{a,-}^*J_{a,-}=S_a+S_{-a}.       \tag{223h}
\end{equation}
When $0<a<2T$, the compressed shift $S_a$ is nonzero, so the two Grams and
hence the two embeddings cannot be equal.  This proves the last assertion.
\end{proof}

\begin{corollary}[Required order-zero Poisson--Fourier scattering ---
\PROVED]
A non-circular Bayes completion of (220j) must contain a source-defined
unitary (or isometry with explicitly controlled orthogonal ports)
\begin{equation}
 \mathscr S_{T,t}:\mathscr K_{T,t}^{\rm in}
          \longrightarrow\mathscr K_{T,t}^{\rm out} \tag{223i}
\end{equation}
such that its zeroth-order physical intertwining is
\begin{equation}
 \mathscr S_{T,1}E_XF=E_YF,                           \tag{223j}
\end{equation}
and the actual evolution is
\begin{equation}
 \mathscr S_{T,1+\varepsilon}\mathcal M_{1,1+\varepsilon}. \tag{223k}
\end{equation}
The Markov factor supplies the positive first-order innovation Gram; the
scattering factor supplies the order-zero change from minus to plus.  The
global gcd--Fourier involution $\mathbb U_T$ in (144)--(149) is an explicit
candidate for $\mathscr S_{T,1}$, but its physical intertwining (223j), after
Gamma affine and Tate compression, must be checked rather than assumed.
\end{corollary}

\begin{proof}
The preceding theorem excludes (223g) whenever the input and output ports
are genuinely different.  Therefore an order-zero map is necessary.  If it
is unitary, composition with $\mathcal M$ preserves the conditional-variance
defect (220d)--(220i), so the division of roles stated after (223k) is exact.
Equations (144)--(149) prove that $\mathbb U_T$ is source-defined, unitary and
has the required label/reflection type.  They do not prove (223j).  The next
proposition shows that the canonical reflection choice does not supply it.
\end{proof}

\begin{proposition}[Gcd inversion/reflection does not supply the scattering
--- \PROVED]
Choose the physical involution in (144) to be reflection
$(\mathcal RF)(t)=F(-t)$, the canonical unitary satisfying
$\mathcal R\tau_a\mathcal R=\tau_{-a}$.  Then (147) gives
\begin{equation}
 \frac12(I\pm\mathbb U_T)\mathbb B_T^{\rm ar}
 =\mathbb B_T^{\rm ar}\frac12(I\pm\mathcal R).       \tag{223l}
\end{equation}
Thus the midpoint outputs of $\mathbb U_T$ are merely the even and odd
source features.  They do not map the minus-delay factor to the plus-delay
factor.  In particular this canonical $\mathbb U_T$ cannot be the scattering
$\mathscr S_{T,1}$ required in (223j), except on a degenerate subspace where
the two physical Grams already coincide.
\end{proposition}

\begin{proof}
Equation (223l) follows by adding and subtracting
$\mathbb U_T\mathbb B_T^{\rm ar}=\mathbb B_T^{\rm ar}\mathcal R$ from
$\mathbb B_T^{\rm ar}$.  Its right side is the feature of the parity
projection of $F$.

If the same unitary mapped the complete minus factor to the plus factor on a
source subspace, their Grams would agree there.  But (80) gives
\begin{equation}
 J_{a,+}^*J_{a,+}-J_{a,-}^*J_{a,-}=S_a+S_{-a},       \tag{223m}
\end{equation}
which is nonzero on an overlapping channel.  Hence no such identification
follows from (223l).
\end{proof}

\begin{corollary}[Cross-place character of the missing scattering ---
\PROVED]
The order-zero factor $\mathscr S_{T,1}$ cannot be built solely from label
inversion and physical reflection, nor as a direct sum of independent
prime-channel unitaries.  It must mix the arithmetic delay ports with the
Gamma affine and/or moving Tate/old channels.  Any channelwise construction
which claims to produce $J_+$ from $J_-$ without such mixing either assumes
the second port as in (223c), or contradicts the Gram difference (223m).
\end{corollary}

\begin{proof}
The first assertion is the proposition.  A direct sum of prime-channel
unitaries would preserve each channel Gram separately, contradicting (223m)
on every overlapping shift.  Therefore compensation must occur across
places or terminal channels, as also required abstractly by (89).
\end{proof}

\begin{proposition}[Scalar completed scattering loses the Euler ports ---
\PROVED]
On the critical line, away from the null set of zeros, define
\begin{equation}
 S_\zeta(s):=\frac{\zeta(1-s)}{\zeta(s)},
 \qquad \Re s=\frac12.                               \tag{223n}
\end{equation}
Then $|S_\zeta(s)|=1$ almost everywhere, but the functional equation gives
the exact identity
\begin{equation}
 S_\zeta(s)
 =\pi^{\frac12-s}
   \frac{\Gamma(s/2)}{\Gamma((1-s)/2)}               \tag{223o}
\end{equation}
up to the fixed convention for the completed zeta factor.  Hence the scalar
scattering is entirely Archimedean after continuation.  Compressing its
unitary multiplier to a time window gives a contraction, but cannot produce
the independent finite Euler shifts
$\{S_{\log n}+S_{-\log n}\}$ in (83).
\end{proposition}

\begin{proof}
For $s=1/2+i\tau$, real Dirichlet coefficients give
$\zeta(1-s)=\overline{\zeta(s)}$, so (223n) has modulus one wherever it is
defined.  Solving the functional equation
$\pi^{-s/2}\Gamma(s/2)\zeta(s)
=\pi^{-(1-s)/2}\Gamma((1-s)/2)\zeta(1-s)$
for the ratio gives (223o).  Thus its logarithmic derivative contains only
the continued Gamma factor; the separate Euler orientations have cancelled.
A compression of a scalar multiplier can only produce the convolution kernel
of that multiplier.  It has no source-labelled direct-sum ports capable of
recovering the locally finite prime translations in (83).  This proves the
last assertion.
\end{proof}

\begin{corollary}[Two-contour jump scattering target --- \PROVED]
The order-zero map in (223i) must be a matrix/port scattering
\begin{equation}
 \mathscr S_T^{\rm jump}:
 \mathscr K_T^{>}\oplus\mathscr K_T^{<}\oplus
       \mathscr K_{\Gamma,T}^{\rm aff}\oplus\mathbb C^2_{\rm Tate}
 \longrightarrow
 \mathscr K_T^{>}\oplus\mathscr K_T^{<}\oplus
       \mathscr K_{\Gamma,T}^{\rm aff}\oplus\mathbb C^2_{\rm Tate},
                                                               \tag{223p}
\end{equation}
defined before identifying the two Mellin continuations.  Its supertrace may
reduce to the scalar identity (223o), but its off-diagonal ports must retain
the contour jump (275), equivalently the source operator (279a).  A proof of
unitarity/passivity for (223p), together with the physical intertwining
(223j), would prove (220j).  Replacing (223p) by (223n) would erase the
load-bearing arithmetic carrier and is invalid.
\end{corollary}

\begin{proof}
The proposition excludes a scalar continued multiplier.  Equations
(270)--(279a) prove that the locally finite row-(d) source is precisely the
difference of the two contour orientations before their meromorphic
cancellation.  Therefore retaining those two boundary components and the
Gamma/Tate terminals is necessary.  If the resulting source-defined map is
unitary or passive and satisfies (223j), composition with the Bayes
contraction preserves the positive defect and the Bayes pullback criterion
gives G40.  This proves the conditional sufficiency and the no-replacement
statement.
\end{proof}

Equations (208)--(223) now derive all scalar normalizations of the proposed
bridge from positive metric motions: $g_\Gamma$, $m_0$ and every
$\Lambda(n)/\sqrt n$.  The remaining statement is not a scalar
normalization; it is the operator-level assertion that the amplified Julia
defect ports, after tensoring with the physical translated source and Kato
compressing by $P_t$, assemble with the plus orientation
$J_{\log n,+}$ while respecting the neutral old-restriction law
(41s)--(41u).  The finite-port identities (41p)--(41r) are not required.

The complete reference factor is the infinitesimal defect of these explicit
contractions.

\begin{theorem}[Markov-defect realization of $\mathcal X_T$ --- \PROVED]
Let
\begin{equation}
 \beta_\varepsilon(\tau):=
 \frac{\phi_{\Gamma,1}(\tau)}
      {\phi_{\Gamma,1+\varepsilon}(\tau)},
 \qquad
 d_{\Gamma,\varepsilon}(\tau)^2
 :=\frac{1-|\beta_\varepsilon(\tau)|^2}{\varepsilon}. \tag{224}
\end{equation}
For a prime power $n=p^k$ put
\begin{equation}
 d_{n,\varepsilon}^2
 :=\frac{p^{-k/2}}{k\varepsilon}
       (1-p^{-k\varepsilon}).                         \tag{225}
\end{equation}
On the form domain of $\mathcal X_T$, define the explicit defect analysis
map
\begin{equation}
 \mathcal D_{T,\varepsilon}F:=
 \left(
   \mathcal F^{-1}d_{\Gamma,\varepsilon}\widehat{E_TF},
   \{d_{n,\varepsilon}J_{\log n,-}F\}_{n\in\mathcal A_T}
 \right).                                             \tag{226}
\end{equation}
Then
\begin{equation}
 \lim_{\varepsilon\downarrow0}
 \|\mathcal D_{T,\varepsilon}F\|^2
 =\langle G_{\Gamma,T}F,F\rangle
  +\sum_{n\in\mathcal A_T}\frac{\Lambda(n)}{\sqrt n}
       \|J_{\log n,-}F\|^2
 =\|\mathcal X_TF\|^2.                               \tag{227}
\end{equation}
Every finite-$\varepsilon$ summand in (226) is a Julia defect of one of the
source-defined Markov contractions (217), (219).  Thus the full reference
Gram is a positive infinitesimal defect rate independently of row (d).
\end{theorem}

\begin{proof}
By (216) and (209),
\[
 -\log|\beta_\varepsilon(\tau)|^2
 =\int_1^{1+\varepsilon}g_{\Gamma,t}(\tau)\,dt.
\]
Hence $d_{\Gamma,\varepsilon}^2\to g_\Gamma$ pointwise.
Since $g_{\Gamma,t}\le g_{\Gamma,1}$ for $t\ge1$ and
$1-e^{-x}\le x$, the defect multipliers are dominated by $g_\Gamma$.
Dominated convergence on the Gamma form domain and (82) give the first
summand in (227).  Equations (222)--(223) give convergence of every
arithmetic coefficient; the contact set $\mathcal A_T$ is finite.  Summing
proves (227).  The defect interpretation is the definition of the Julia
defect operator $D_C=(I-C^*C)^{1/2}$ applied to the character eigenchannels.
\end{proof}

Therefore the input/reference half of the G40 network is now completely
constructed from positive contractions.  The remaining output construction
and identity are
\begin{equation}
 \boxed{\quad
 \operatorname*{form-lim}_{\varepsilon\downarrow0}
 \mathcal O_{T,\varepsilon}P_T
 =\mathcal Y_TP_T,
 \quad}                                                \tag{228}
\end{equation}
where $\mathcal O_{T,\varepsilon}$ must be constructed as the main output of
the coupled Julia dilations of the Euler contraction (219) and beta--Gamma
contraction (217), with the Euler ports reversed as required by (218)--(220),
and then transported by (112)--(116).  The admissible construction is thus
restricted to explicit source contractions, but its precise physical port
interconnection has not yet been proved.  Constructing it and proving (228),
with its orthogonal defect equal to (109) plus the old residual, would give
$\|\mathcal Y_TP_TF\|\le\|\mathcal X_TP_TF\|$ and close G40 without
circularity.

There is also a physical state-space realization of the same Archimedean
termination, which is better adapted to the windowed shifts.

\begin{theorem}[Gamma ODE storage and finite-head termination --- \PROVED]
Put $A_j=j+\tfrac14$.  On an interval $[a,b]$, let $y$ solve
\begin{equation}
 \frac12y'+A_jy=\sqrt{A_j}\,F.                        \tag{229}
\end{equation}
Then, for arbitrary initial state,
\begin{equation}
 \frac1{A_j}\int_a^b|F|^2+\frac{|y(a)|^2}{2A_j}
 =\int_a^b|y|^2+
  \int_a^b\left|\frac{y'}{2A_j}\right|^2
  +\frac{|y(b)|^2}{2A_j}.                             \tag{230}
\end{equation}
For the causal zero-state filter on $I_T$,
\begin{equation}
 (C_j^+F)(t)=2\sqrt{A_j}\int_{-T}^t
 e^{-2A_j(t-s)}F(s)\,ds,                              \tag{231}
\end{equation}
and the compressed shift $V_aF(t)=\widetilde F(t+a)$, one has, for
$t\le T-a$,
\begin{equation}
 (V_aC_j^+-C_j^+V_a)F(t)
 =e^{-2A_j(t+T)}(C_j^+F)(-T+a).                       \tag{232}
\end{equation}
Thus every Gamma/shift commutator is a rank-one free response whose energy is
the positive boundary storage in (230).

On the whole-line spectral realization write $C_jF=y_j$ and
$D_jF=y_j'/(2A_j)$.  Then
\begin{equation}
 C_j^*C_j+D_j^*D_j=\frac1{A_j}I,
 \qquad
 G_\Gamma=\sum_{j\ge0}D_j^*D_j.                      \tag{233}
\end{equation}
Choose finitely supported numbers $0\le\theta_j\le1$ with
\begin{equation}
 \sum_j\frac{\theta_j}{A_j}=m_0,                     \tag{234}
\end{equation}
which is possible because $\sum_jA_j^{-1}=\infty$.  Define
\begin{align}
 C_H&:=\bigoplus_{j:\theta_j>0}\sqrt{\theta_j}\,C_j,
 &D_H&:=\bigoplus_{j:\theta_j>0}\sqrt{\theta_j}\,D_j,\notag\\
 D_{\rm tail}&:=\left(\bigoplus_{j:\theta_j>0}
                   \sqrt{1-\theta_j}\,D_j\right)
       \oplus\left(\bigoplus_{j:\theta_j=0}D_j\right), \tag{235}
\end{align}
where in the last display each index is included exactly once: a head index
uses its complementary coefficient and every index beyond the finite head
uses coefficient one.  Then
\begin{equation}
 C_H^*C_H+D_H^*D_H=m_0I,
 \qquad
 D_H^*D_H+D_{\rm tail}^*D_{\rm tail}=G_\Gamma,
\end{equation}
and therefore
\begin{equation}
 \boxed{G_\Gamma-m_0I=D_{\rm tail}^*D_{\rm tail}-C_H^*C_H.} \tag{236}
\end{equation}
The negative Archimedean output is finite, the Gamma tail is a positive loss
reservoir, and all internal cell states telescope by (230).  On the primitive
source the only external states are the two Tate moments, hence vanish.
\end{theorem}

\begin{proof}
Substitute
$F=(\frac12y'+A_jy)/\sqrt{A_j}$, expand the square, and integrate
$\Re(y'\overline y)=(|y|^2)'/2$ to obtain (230).  Formula (232) follows by
writing the two integrals in (231), translating the second variable, and
observing that their difference is precisely the integral over
$[-T,-T+a]$ propagated freely to $t$.

The Fourier multipliers of $C_j$ and $D_j$ are respectively
\[
 \frac{2\sqrt{A_j}}{2A_j+i\tau},
 \qquad
 \frac{i\tau}{\sqrt{A_j}(2A_j+i\tau)}.
\]
Their squared moduli sum to $1/A_j$, while the second is
$\tau^2/[A_j(4A_j^2+\tau^2)]$, the $j$-th summand of (171) because
$a_j=2A_j$.  This proves (233).  Take the shortest initial segment for which
the reciprocal sum exceeds $m_0$ and fractionally weight its last member to
obtain (234).  Weighted summation of (233) proves the two identities before
(236), and subtraction proves (236).  Summing (230) over adjacent cells
cancels every internal terminal with the next initial storage.  For the
Gamma index $j=0$, $A_{j}=1/4$, and the two outer oriented states
are scalar multiples of $M_+F$ and
$M_-F$; the primitive conditions annihilate them.  All $j\ge1$ outer
remainders are stable positive storages retained in $D_{\rm tail}$.
\end{proof}

\begin{corollary}[Explicit four-state Gamma head --- \PROVED]
In (234) one may take
\begin{equation}
 \theta_0=\theta_1=\theta_2=1,
 \qquad
 \theta_3=\frac{13}{4}
 \left(m_0-\frac{236}{45}\right),
 \qquad
 \theta_j=0\quad(j\ge4).                              \tag{239}
\end{equation}
Moreover $0<\theta_3<1$ (numerically
$\theta_3=0.4151516680\ldots$).  Thus the negative Gamma head $C_H$ has
exactly four scalar ODE states.
\end{corollary}

\begin{proof}
Since $A_0=1/4$, $A_1=5/4$, $A_2=9/4$, and $A_3=13/4$,
\[
 \frac1{A_0}+\frac1{A_1}+\frac1{A_2}
 =4+\frac45+\frac49=\frac{236}{45}.
\]
The elementary bounds
$1.14<\log\pi<1.15$, $0.57<\gamma<0.58$,
$1.57<\pi/2<1.572$, and $0.69<\log2<0.694$ give
\[
 \frac{236}{45}<m_0<\frac{236}{45}+\frac4{13}.
\]
(Each displayed logarithmic bound follows, for example, from the alternating
series for $\log(1+x)$ after a rational range reduction.)  Hence (239) lies
strictly between zero and one, and its contribution is
$\theta_3/A_3=m_0-236/45$, proving (234).
\end{proof}

Replacing the Gamma and constant components in (83) by (236) gives the
equivalent source-defined balanced factorization
\begin{equation}
 \mathcal X_T^\sharp F=
 \left(D_{\rm tail}F,\{w_n^{1/2}J_{\log n,-}F\}_{n\in\mathcal A_T}\right),
 \quad
 \mathcal Y_T^\sharp F=
 \left(C_HF,\{w_n^{1/2}J_{\log n,+}F\}_{n\in\mathcal A_T}\right), \tag{237}
\end{equation}
with
\begin{equation}
 A_T=(\mathcal X_T^\sharp)^*\mathcal X_T^\sharp
     -(\mathcal Y_T^\sharp)^*\mathcal Y_T^\sharp.    \tag{238}
\end{equation}
This removes the separate $m_0^{1/2}I$ output and makes the opposite Gamma
head finite without changing the signed physical operator.

\begin{theorem}[Gamma common-channel cancellation --- \PROVED]
Let $G_\Gamma^{1/2}$ denote the canonical positive spectral factor only for
the already positive Gamma operator.  On its generated range define
\begin{equation}
 \mathcal V_\Gamma(G_\Gamma^{1/2}F)
 :=(D_HF,D_{\rm tail}F),                              \tag{238a}
\end{equation}
and on the scalar affine channel define
\begin{equation}
 \mathcal U_H(m_0^{1/2}F):=(D_HF,C_HF).              \tag{238b}
\end{equation}
Then $\mathcal V_\Gamma$ and $\mathcal U_H$ are source-defined isometries
onto their generated ranges.  Their first output is the identical channel
$D_HF$.  Cancelling this common internal channel gives exactly
\begin{equation}
 \|G_\Gamma^{1/2}F\|^2-m_0\|F\|^2
 =\|D_{\rm tail}F\|^2-\|C_HF\|^2,                   \tag{238c}
\end{equation}
and hence the Gamma components of
$\mathcal X_T^\sharp,\mathcal Y_T^\sharp$ in (237).
No square root of $G_\Gamma-m_0I$ and no sign assumption is used.
\end{theorem}

\begin{proof}
The two identities preceding (236) give
\[
 \|D_HF\|^2+\|D_{\rm tail}F\|^2
 =\langle G_\Gamma F,F\rangle
 =\|G_\Gamma^{1/2}F\|^2
\]
and
\[
 \|D_HF\|^2+\|C_HF\|^2=m_0\|F\|^2.
\]
They prove that (238a)--(238b) preserve all polarized inner products, so
they extend uniquely to isometries of the generated ranges.  Subtracting the
second equality from the first cancels the common $D_H$ component and gives
(238c).  This is exactly (236) at the quadratic and polarized Gram levels.
\end{proof}

\begin{corollary}[Gamma part of the Bayes pullback --- \PROVED]
The infinitesimal beta--Gamma Markov defect in (224)--(227), whose Gram is
$G_\Gamma$, may be transported by $\mathcal V_\Gamma$ to
$D_H\oplus D_{\rm tail}$.  The affine normalization $m_0$ may be transported
by $\mathcal U_H$ to $D_H\oplus C_H$.  After feedback cancellation of the
common $D_H$ port, the input/output pair is precisely
\begin{equation}
             D_{\rm tail}F\quad/\quad C_HF.          \tag{238d}
\end{equation}
Thus the Gamma defect, constant normalization and four-state-head/tail
orientation required in (220j) are completely matched.  The still open part
of (220j) is the coupling of these already matched bulk ports to the moving
Tate projection and old storage.
\end{corollary}

\begin{proof}
Equation (227) identifies the infinitesimal Gamma defect Gram with
$G_\Gamma$; equality of Grams permits the isometry (238a).  Equation (210)
identifies the affine scalar normalization and (238b) realizes its source
Gram.  Theorem (238a)--(238c) performs the common-port cancellation, proving
(238d) and the claim.
\end{proof}

The prime-power outputs can likewise be collected into one lossless cell per
prime, including the exact zero-extension boundary storage.

\begin{theorem}[Compressed Euler--Julia resolvent identity --- \PROVED]
Fix a prime $p$, put $r=p^{-1/2}$, $s=(1-r^2)^{1/2}$, and let
$S_p=S_{\log p}$ be the positive zero-extension shift on $L^2(I_T)$.  It is
nilpotent and $S_p^k=S_{k\log p}$.  Define
\begin{align}
 C_{p,T}&:=rS_p(I-rS_p)^{-1}
          =\sum_{k\ge1}r^kS_p^k,                     \tag{240}\\
 K_{p,T}&:=s(I-rS_p)^{-1},\notag\\
 B_{p,T}&:=(I-S_p^*S_p)^{1/2}(I-rS_p)^{-1}.          \tag{241}
\end{align}
All sums are finite.  Then
\begin{equation}
 \boxed{I+C_{p,T}+C_{p,T}^*
 =K_{p,T}^*K_{p,T}+r^2B_{p,T}^*B_{p,T}.}             \tag{242}
\end{equation}
Consequently the complete contribution of all active powers of $p$ is
\begin{equation}
 -\sum_{k\ge1}\frac{\Lambda(p^k)}{p^{k/2}}
   (S_{k\log p}+S_{-k\log p})
 = (\log p)\bigl(I-K_{p,T}^*K_{p,T}
                    -r^2B_{p,T}^*B_{p,T}\bigr).      \tag{243}
\end{equation}
The boundary term vanishes for a bilateral unitary shift and is positive for
the compressed physical shift.
\end{theorem}

\begin{proof}
Nilpotence and convexity of the interval give $S_p^k=S_{k\log p}$ and the
finite Neumann series (240).  Since
$I+C_{p,T}=(I-rS_p)^{-1}$, multiplication on the left by
$I-rS_p^*$ and on the right by $I-rS_p$ gives
\[
 (I-rS_p^*)(I+C_{p,T}+C_{p,T}^*)(I-rS_p)
 =I-r^2S_p^*S_p.
\]
But
\[
 I-r^2S_p^*S_p=s^2I+r^2(I-S_p^*S_p).
\]
Conjugating back proves (242).  Finally
$\Lambda(p^k)=\log p$ and (240) identify the left side of (243) with
$-(\log p)(C_{p,T}+C_{p,T}^*)$; subtracting (242) from $I$ proves (243).
\end{proof}

Combining (236) and (243) gives a second, polar-free balanced factorization
of the authoritative operator:
\begin{align}
 \mathcal X_T^{\rm free}F
 &:=\left(D_{\rm tail}F,
          \{(\log p)^{1/2}F\}_{p<e^{2T}}\right),      \tag{244}\\
 \mathcal Y_T^{\rm fin}F
 &:=\left(C_HF,
   \{(\log p)^{1/2}K_{p,T}F\}_{p<e^{2T}},
   \{(\log p)^{1/2}p^{-1/2}B_{p,T}F\}_{p<e^{2T}}
          \right),                                   \tag{245}\\
 A_T&=(\mathcal X_T^{\rm free})^*\mathcal X_T^{\rm free}
     -(\mathcal Y_T^{\rm fin})^*\mathcal Y_T^{\rm fin}. \tag{246}
\end{align}
Only primes with $\log p<2T$ occur.  Formula (246) contains the four-state
Gamma head, one rational Julia filter and one positive boundary-storage port
per active prime; every prime power is already resummed in (240).

\begin{theorem}[Finite-head physical angular operator --- \PROVED]
The positive Gamma-tail feature in $\mathcal X_T^{\rm free}$ is injective.
Consequently
\begin{equation}
 \mathfrak C_T^{\rm fin}
   (\mathcal X_T^{\rm free}P_TF)
 :=\mathcal Y_T^{\rm fin}P_TF                         \tag{247}
\end{equation}
defines a source operator on
$\Ran(\mathcal X_T^{\rm free}P_T)$ without assuming a sign.  Its exact
defect law is
\begin{equation}
 \|\mathcal X_T^{\rm free}P_TF\|^2
 -\|\mathfrak C_T^{\rm fin}\mathcal X_T^{\rm free}P_TF\|^2
 =\langle A_TP_TF,P_TF\rangle.                       \tag{248}
\end{equation}
It extends to a contraction on the closed generated range if and only if
\begin{equation}
 \boxed{\|\mathfrak C_T^{\rm fin}\|\le1,}             \tag{249}
\end{equation}
equivalently if and only if row (d) holds at the window $T$.
\end{theorem}

\begin{proof}
The squared Fourier multiplier of $D_{\rm tail}$ is a sum of positive
$D_j$ multipliers and is strictly positive for every $\tau\ne0$, because
infinitely many full tail indices remain after removing the four-state head.
If $D_{\rm tail}E_TF=0$, the Fourier transform of $E_TF$ is supported at the
singleton $\{0\}$ and therefore vanishes in $L^2$; hence $F=0$.  This proves
well-definedness of (247).  Equation (248) is (246) after compression by
$P_T$.  The norm inequality in (249) is therefore equivalent to
$P_TA_TP_T\ge0$.
\end{proof}

\begin{corollary}[Certified initial passivity --- \PROVED]
For every $0<T\le\log2$, the operator (247) extends to a strict contraction
on its generated primitive range:
\begin{equation}
 \|\mathfrak C_T^{\rm fin}
       \mathcal X_T^{\rm free}P_TF\|
 <\|\mathcal X_T^{\rm free}P_TF\|\qquad(F\ne0).     \tag{250}
\end{equation}
\end{corollary}

\begin{proof}
The certified-initial-interval theorem in the supplied
\texttt{row-d-local-analysis.tex} proves
$\langle A_TF,F\rangle>0$ for every nonzero $F\in\mathcal P_T$ and every
$0<T\le\log2$.  Substitute this strict inequality in the exact defect law
(248).
\end{proof}

\begin{theorem}[First-loss zero-mode reduction --- \PROVED]
For each $T>0$, the primitive Friedrichs operator
\begin{equation}
 H_T:=P_TA_TP_T\big|_{\mathcal P_T}                  \tag{251}
\end{equation}
is selfadjoint, bounded below and has compact resolvent.  After transporting
$L^2(I_T)$ to a fixed interval by dilation, its closed quadratic forms depend
continuously on $T$; the same is true at contact thresholds because a newborn
translated overlap has measure zero.  Consequently, if row (d) fails at some
finite window, there is a first $T_*>\log2$ and a nonzero
$F_*\in\mathcal P_{T_*}$ such that
\begin{equation}
 H_{T_*}F_*=0.                                        \tag{252}
\end{equation}
Equivalently, for some scalars $\alpha,\beta$,
\begin{equation}
 A_{T_*}F_*
 =\alpha e^{-t/2}+\beta e^{t/2}\quad\hbox{on }I_{T_*}. \tag{253}
\end{equation}
Thus global row (d) follows once nonzero solutions of (252)--(253) are
excluded at every finite window.
\end{theorem}

\begin{proof}
The Gamma multiplier satisfies
$g_\Gamma(\tau)=\log(2+|\tau|)+O(1)$, while the constant and the finitely
many active zero-extension shifts are bounded.  Hence the closed Gamma form
domain, equipped with its graph norm, embeds compactly into $L^2(I_T)$: in a
Fourier basis its coercive weights tend to infinity.  Bounded perturbation and
finite-codimensional compression preserve selfadjointness, semiboundedness
and compact resolvent.

Conjugate by the unitary dilation from $I_T$ to $(-1,1)$.  On the common form
core of smooth compactly supported functions, the Gamma kernel, moment
vectors and each active translated overlap vary continuously with $T$.
Uniform lower bounds on compact $T$-intervals and form-core approximation
give norm-resolvent, hence eigenvalue, continuity by the standard theorem for
continuous closed forms.  At $2T=\log n$ the new overlap is supported on a
null set, so its form enters continuously from zero.

The lowest primitive eigenvalue is strictly positive for
$T\le\log2$ by the supplied certified theorem.  If it is negative at a later
finite window, zero extension from $I_T$ to $I_{T'}$ preserves both Tate
moments and the global quadratic value; hence the lowest eigenvalue is
nonincreasing under enlargement.  Continuity gives a first zero $T_*$ and
compact resolvent gives an eigenvector $F_*$, proving (252).  Since
$\mathcal P_T=\ker M_-\cap\ker M_+$, its orthogonal complement in $L^2(I_T)$
is the span of $e^{-t/2},e^{t/2}$.  Equation (252) is therefore equivalent to
(253).
\end{proof}

\begin{corollary}[Parity reduction of a first-loss mode --- \PROVED]
The operator $A_T$, the primitive space and the finite-head angular system
commute with reflection $(\mathcal RF)(t)=F(-t)$.  Hence a nonzero first-loss
mode may be chosen either even or odd.  In the even sector (252)--(253) reduce
to
\begin{equation}
 A_TF=c\cosh(t/2),
 \qquad \int_{-T}^TF(t)\cosh(t/2)\,dt=0,              \tag{254}
\end{equation}
while in the odd sector they reduce to
\begin{equation}
 A_TF=c\sinh(t/2),
 \qquad \int_{-T}^TF(t)\sinh(t/2)\,dt=0.              \tag{255}
\end{equation}
The other Tate moment follows automatically from parity.  Thus it suffices to
exclude the two scalar equality systems (254)--(255).
\end{corollary}

\begin{proof}
The Gamma kernel is even, reflection interchanges $S_a$ and $S_{-a}$, and
$A_T$ contains their symmetric sum.  Reflection also interchanges $M_+$ and
$M_-$, so it preserves their joint kernel.  The kernel of $H_T$ therefore
splits into its even and odd parts.  In (253), an even right side requires
$\alpha=\beta$ and is a multiple of $\cosh(t/2)$; an odd right side requires
$\alpha=-\beta$ and is a multiple of $\sinh(t/2)$.  The two primitive moments
then coincide up to sign and give the displayed scalar conditions.
\end{proof}

\begin{theorem}[Homogeneous exterior-leakage form of a zero mode --- \PROVED]
For $h\in\mathcal E'(\mathbb R)$ define the arithmetic part of the global
explicit-form operator by the locally finite distributional sum
\begin{equation}
 \mathcal E_{\rm ar}h
 :=-\sum_{n\ge2}\frac{\Lambda(n)}{\sqrt n}
       (\tau_{\log n}+\tau_{-\log n})h.               \tag{256}
\end{equation}
Thus
$\mathcal A_\infty h:=G_\Gamma h-m_0h+\mathcal E_{\rm ar}h$
is a well-defined element of $\mathcal D'(\mathbb R)$ on the compactly
supported core.  Let $L=\partial_t^2-\tfrac14$.  For
$\phi\in H_0^2(I_T)$ put $F=L\phi$.  Then the primitive zero-mode equation
(252)--(253) is equivalent to the homogeneous exterior-leakage condition
\begin{equation}
 \boxed{P_{I_T}\mathcal A_\infty L^2E_T\phi=0.}       \tag{257}
\end{equation}
Equivalently, the distribution
$\mathcal A_\infty L^2E_T\phi$ is supported in
$\mathbb R\setminus I_T$.  Thus a first-loss mode is exactly a nonzero
clamped potential whose differentiated Euler--Gamma current has complete
interior cancellation.
\end{theorem}

\begin{proof}
The sum (256) is locally finite: for compact sets $K$ and
$\operatorname{supp}h$, only finitely many $n$ can make
$K\cap(\operatorname{supp}h\pm\log n)$ nonempty.  Hence all commutations
below are legitimate in $\mathcal D'$.  Because $\phi$ and $\phi'$ vanish at
both endpoints, distributional differentiation creates no endpoint deltas
and $L(E_T\phi)=E_T(L\phi)=E_TF$.  The global operator commutes with $L$ on
this compactly supported core.  For $t\in I_T$, translations by
$\log n\ge2T$ do not meet the support, so the restriction of
$\mathcal A_\infty F$ is $A_TF$.  Applying $L$ to (253) kills both
$e^{-t/2}$ and $e^{t/2}$ and gives (257).

Conversely, (257) says that $L(\mathcal A_\infty F)=0$ distributionally on
$I_T$.  Every distributional solution of $Lg=0$ on an interval is a linear
combination of $e^{-t/2}$ and $e^{t/2}$, so (253) follows.  Since
$F=L\phi$ with clamped boundary data, integration by parts gives
$M_-F=M_+F=0$, hence $F\in\mathcal P_T$.
\end{proof}

\begin{proposition}[Fourier-multiplier withdrawal and correct analytic
typing --- \PROVED]
The formal expression
\begin{equation}
 g_\Gamma(\tau)-m_0
 -2\sum_{n\ge2}\frac{\Lambda(n)}{\sqrt n}
       \cos(\tau\log n)                               \tag{258}
\end{equation}
is not, from the presently proved estimates, a tempered-distribution
multiplier on the real Fourier line.  Formula (256), as a locally finite
map $\mathcal E'\to\mathcal D'$, is the authoritative meaning used in
(257).  The two one-sided Dirichlet series attached to (258) live initially
in disjoint half-planes and may be joined only through the completed
explicit-form continuation.  Accordingly, multiplying (258) by
$(\tau^2+\tfrac14)^2\widehat\phi(\tau)$ is not used as a Fourier-space proof
of (257).
\end{proposition}

\begin{proof}
For the positive translation orientation one obtains a Dirichlet series
containing $n^{-1/2-i\tau}$, while the negative orientation contains
$n^{-1/2+i\tau}$.  Absolute convergence of the first requires one open
half-plane and that of the second the reflected half-plane; there is no
common initial strip.  In the time variable, (256) is nevertheless locally
finite on compactly supported inputs, which proves the asserted
$\mathcal E'\to\mathcal D'$ typing.  Promoting the resulting output to a
tempered distribution requires an additional global growth theorem not
contained in the local row-(d) estimates.  The completed continuation in
(117)--(119) may be paired with compactly supported test functions, but that
does not turn the separate formal series in (258) into an ordinary real-line
multiplier.  This proves the withdrawal.
\end{proof}

\begin{proposition}[Exact scope of logarithmic unique continuation ---
\AUDIT]
Let $L_\Delta$ be the one-dimensional logarithmic Laplacian.  Its established
weak unique-continuation theorem has the form
\begin{equation}
 u=0\ \hbox{and}\ L_\Delta u=0\ \hbox{on one nonempty open set}
 \quad\Longrightarrow\quad u=0.                      \tag{259}
\end{equation}
It cannot be applied directly to (257).  On $I_T$, the unknown
$L^2E_T\phi$ need not vanish, and the Gamma output is forced by the nonlocal
arithmetic term (256), rather than vanishing separately.  Therefore the
remaining statement is the genuinely coupled implication
\begin{equation}
 \boxed{\quad
 \phi\in H_0^2(I_T),\quad
 P_{I_T}\mathcal A_\infty L^2E_T\phi=0
 \quad\Longrightarrow\quad\phi=0 .\quad}             \tag{260}
\end{equation}
Proposition (259) proves (260) only after an independent argument produces
simultaneous local vanishing of an input channel and its logarithmic
Laplacian channel.  No such argument has yet been proved here.
\end{proposition}

\begin{proof}
The hypothesis of (259) contains two local equalities for the same input.
Equation (257) contains only the cancellation of the sum of its Gamma,
constant and translated arithmetic outputs.  Neither equality required by
(259) follows by separating summands of a zero sum.  Such a separation would
already assert the missing orthogonality/passivity.  Hence a direct appeal to
logarithmic unique continuation would be circular, while (260) is exactly
the new Gamma--arithmetic continuation theorem still required.
\end{proof}

\begin{theorem}[Gamma exterior antilocality --- \PROVED]
Let $h\in\mathcal E'(\mathbb R)$ be supported in $[-T,T]$.  If
\begin{equation}
 h=0\quad\hbox{and}\quad G_\Gamma h=0                \tag{261}
\end{equation}
on a nonempty open subinterval of either exterior component
$(T,\infty)$ or $(-\infty,-T)$, then $h=0$ on $\mathbb R$.
\end{theorem}

\begin{proof}
Put $a_j=2j+\tfrac12$ and let $R_j$ be convolution with
$(a_j/2)e^{-a_j|t|}$.  From (171), on the compactly supported core,
\begin{equation}
 G_\Gamma=\sum_{j\ge0}\frac2{a_j}(I-R_j).            \tag{262}
\end{equation}
For $t>T$, the local term vanishes and the series is normally convergent on
every closed exterior subinterval; explicitly,
\begin{equation}
 (G_\Gamma h)(t)
 =-\sum_{j\ge0}e^{-a_jt}M_j^+,
 \qquad
 M_j^+:=\int_{-T}^{T}e^{a_js}h(s)\,ds.               \tag{263}
\end{equation}
Here the integral notation means the distributional pairing
$M_j^+=\langle h,e^{a_j\cdot}\rangle$.  A compactly supported distribution
of order $N$ satisfies $|M_j^+|\le C(1+a_j)^Ne^{a_jT}$, so the series and all
its derivatives converge normally for $t\ge T+\delta$.  The right side is real analytic on
$(T,\infty)$.  If it vanishes on one open interval, it vanishes on that
whole component.  Multiplying successively by $e^{a_0t}$ and sending
$t\to\infty$ gives $M_0^+=0$; subtracting the first term and iterating gives
$M_j^+=0$ for every $j$.

Set $x=e^{2s}$.  The equalities just obtained say
\begin{equation}
 \int_{-T}^{T}x^j e^{s/2}h(s)\,ds=0
 \qquad(j=0,1,2,\ldots).                              \tag{264}
\end{equation}
After multiplying $h$ by $e^{s/2}$ and pushing it forward under the
diffeomorphism $s\mapsto e^{2s}$, (264) gives a compactly supported
distribution $v$ which annihilates every polynomial.  The entire function
$z\mapsto\langle v,e^{zx}\rangle$ has all derivatives at $z=0$ equal to
zero; hence it vanishes identically.  Injectivity of the Fourier transform
on compactly supported distributions gives $v=0$, and therefore $h=0$.
The left exterior component is identical with $e^{-a_js}$ in place of
$e^{a_js}$.
\end{proof}

\begin{theorem}[Exact infinite Gamma--arithmetic local system --- \PROVED]
For compactly supported $h$ define
\begin{equation}
 r_j:=R_jh,
 \qquad q_j:=\frac2{a_j}(h-r_j),
 \qquad a_j=2j+\tfrac12.                              \tag{265}
\end{equation}
Then, distributionally on the line,
\begin{align}
 (-\partial_t^2+a_j^2)r_j&=a_j^2h,                    \tag{266}\\
 G_\Gamma h&=\sum_{j\ge0}q_j.                         \tag{267}
\end{align}
The series in (267) converges on every smooth compactly supported test
function.  Hence (257), with $h=L^2E_T\phi$, is precisely the local-in-$t$
infinite-state system
\begin{equation}
 \boxed{
 \sum_{j\ge0}q_j-m_0h
 -\sum_{2\le n<e^{2T}}\frac{\Lambda(n)}{\sqrt n}
       (\tau_{\log n}+\tau_{-\log n})h=0
 \quad\hbox{on }I_T,}                                \tag{268}
\end{equation}
together with (265)--(266) and the clamped representation
$h=L^2E_T\phi$.  All arithmetic couplings in (268) are finite; the only
infinite state is the positive Gamma resolvent ladder.
\end{theorem}

\begin{proof}
The Fourier multiplier of $R_j$ is $a_j^2/(a_j^2+\tau^2)$, proving (266).
The multiplier of $q_j$ is
$2\tau^2/[a_j(a_j^2+\tau^2)]$, and summing gives the digamma series (171),
which proves (267).  Against a fixed test function the large-$j$ multiplier
is $O(a_j^{-3}\tau^2)$ after moving derivatives to the test function, so
the series converges because $\sum_j a_j^{-3}<\infty$.  Finally, when both
the argument and observation lie in $I_T$, a translate by $\log n\ge2T$
has no overlap.  Substitution of (267) and (256) in (257) therefore gives
(268).
\end{proof}

\paragraph{Sharpened remaining bridge.}
Theorem (261) shows that no Gamma state can hide once both its input and
Gamma output vanish in an exterior open set.  Thus (260) will follow from a
non-circular \emph{arithmetic separation lemma}: the coupled cancellation
(268), the two Tate moments and the clamped boundary conditions must force
such a simultaneous exterior vanishing interval.  Establishing that lemma
is now the sole continuation step; assuming it is equivalent to assuming
the missing passivity and is therefore retained as \OPEN{}.

\begin{theorem}[Completed logarithmic-derivative collapse --- \PROVED]
Let
\begin{equation}
 \xi(s)=\frac12s(s-1)\pi^{-s/2}\Gamma(s/2)\zeta(s),
 \qquad D_\xi(s)=\frac{\xi'(s)}{\xi(s)}.              \tag{269}
\end{equation}
Away from the zeros of $\xi$, one has the meromorphic identity
\begin{align}
 0&=D_\xi(s)+D_\xi(1-s)                              \tag{270}\\
 &=\frac12\psi(s/2)+\frac12\psi((1-s)/2)-\log\pi
   +\frac{\zeta'}{\zeta}(s)
   +\frac{\zeta'}{\zeta}(1-s).                       \tag{271}
\end{align}
The apparent poles at every zero cancel between the two terms in (270).
On the central line the formal substitution of the two prime Dirichlet
series into (271) is exactly the scalar expression (258).  Hence a common
meromorphic continuation of the two orientations makes that expression
identically zero; it does \emph{not} produce the row-(d) operator.
\end{theorem}

\begin{proof}
Differentiate the functional equation $\xi(s)=\xi(1-s)$ to obtain
$D_\xi(s)=-D_\xi(1-s)$, which is (270).  Logarithmic differentiation of
(269) gives
\[
 D_\xi(s)=\frac1s+\frac1{s-1}-\frac12\log\pi
 +\frac12\psi(s/2)+\frac{\zeta'}{\zeta}(s).
\]
Adding the expression at $1-s$ cancels the four rational polar terms and
gives (271).  In the half-plane of its Euler expansion,
$\zeta'/\zeta(s)=-\sum\Lambda(n)n^{-s}$.  At
$s=\tfrac12+i\tau$, the two orientations formally become the two cosine
terms of (258), and the two digamma terms become
$g_\Gamma(\tau)-m_0$.  The expansions, however, have disjoint convergence
half-planes, as proved after (258).  Therefore only the formal central-line
expression agrees; their common meromorphic continuation is the zero in
(270).
\end{proof}

\begin{corollary}[The missing object is a jump carrier, not a scalar
symbol --- \AUDIT]
The nonzero row-(c)/(d) distribution is carried by retaining the two
Poisson orientations before contour deformation and before taking their
supercharacter.  Equivalently, it is the contour-jump/residue datum lost in
the meromorphic cancellation (270).  Consequently no proof of (260) may
replace the source-defined locally finite operator (256) by multiplication
with a single meromorphic function on the central line.  Such a replacement
annihilates the very carrier whose sign is at issue.

The lateral form of the remaining construction is therefore precise: build
a Hilbert/Krein observation
\begin{equation}
 \mathfrak O_T: H_0^2(I_T)\longrightarrow
 \mathscr K_T^+\,[\dotplus],\mathscr K_T^-           \tag{272}
\end{equation}
from the two unsupertraced Poisson orientations and the Gamma resolvent
ladder, such that its Krein square is the left side of (257), while its
Hilbert defect is the sum of the exterior Gamma moments in (263) and
nonnegative prime Julia boundary storages.  Strictness of that defect would
give (260).  Existence and strictness of (272), not its formal scalar
supertrace, remain \OPEN{}.
\end{corollary}

\begin{theorem}[Support saturation of a first-loss mode --- \PROVED]
Let $T_*$ be the first-loss window in (252) and let $F_*$ be a corresponding
nonzero primitive zero mode.  Then
\begin{equation}
 \operatorname{diam}(\operatorname{ess,supp}F_*)=2T_* . \tag{273}
\end{equation}
In particular the essential support reaches both endpoints of $I_{T_*}$.
The same conclusion holds for either parity representative in
(254)--(255).
\end{theorem}

\begin{proof}
Suppose the essential support is contained in an interval $[a,b]$ with
$b-a<2T_*$.  Choose a translation $c$ which centres that interval and put
$T'=(b-a)/2+\varepsilon<T_*$.  Translation invariance of the global form
gives
\[
 \langle\mathcal A_\infty\tau_cF_*,\tau_cF_*\rangle
 =\langle\mathcal A_\infty F_*,F_*\rangle=0.
\]
Moreover
\[
 M_\pm(\tau_cF_*)=e^{\mp c/2}M_\pm(F_*)=0,
\]
with the sign depending only on the translation convention.  Hence the
translated function is a nonzero member of $\mathcal P_{T'}$ with zero
quadratic value.  This contradicts strict positivity at every
$T'<T_*$, which follows from the definition of the first-loss window.
Thus the support diameter is $2T_*$.  Since it is contained in
$[-T_*,T_*]$, both endpoints belong to its essential closed hull.  Reflection
does not change this conclusion.
\end{proof}

\paragraph{Consequence for lateral strategies.}
The arithmetic separation lemma cannot be obtained by finding a smaller
support interval and applying Theorem (261) in the resulting empty collar:
Theorem (273) rules out such a collar for a first-loss mode.  A successful
argument must instead show that the two unsupertraced orientations in (272)
cannot cancel while their common input saturates both boundaries.

\begin{theorem}[Finite-energy jump rigidity --- \AUDIT]
Let $h\in\mathcal E'(\mathbb R)$ and put
\begin{equation}
 H(s):=\langle h(t),e^{(s-1/2)t}\rangle.              \tag{274}
\end{equation}
Assume $H(0)=H(1)=0$.  For $c>1$ and
$\eta\in C_c^\infty(\mathbb R)$ put
\begin{equation}
 E_\eta(s):=\int_{\mathbb R}\eta(t)e^{(s-1/2)t}\,dt
\end{equation}
and define the source contour jump by
\begin{equation}
 \langle\mathfrak J_\xi h,\eta\rangle
 :=\frac1{2\pi i}\left(
 \int_{c-i\infty}^{c+i\infty}
 -\int_{1-c-i\infty}^{1-c+i\infty}\right)
 D_\xi(s)H(s)E_\eta(s)\,ds.                           \tag{275}
\end{equation}
Both lines are oriented upwards.  This definition is independent of $c$
and retains the two contour orientations before the cancellation (270).  If
\begin{equation}
 \mathfrak J_\xi h\in L^2(\mathbb R),                 \tag{276}
\end{equation}
then $h=0$.
\end{theorem}

\begin{proof}
On vertical lines $D_\xi(s)=O(\log(2+|\Im s|))$ away from bounded
neighbourhoods of its poles, $H$ has polynomial growth, and $E_\eta$ has
rapid decay of every order.  Thus (275) converges distributionally.
Horizontal sides tend to zero, so the residue theorem gives, in
Fourier-hyperfunction notation,
\begin{equation}
 \mathfrak J_\xi h
 =\sum_{\rho}m_\rho H(\rho)
       e^{(\rho-1/2)t}.                               \tag{277}
\end{equation}
The completed function $\xi$ is entire and nonzero at $0,1$, so the poles
crossed are precisely its nontrivial zeros.  Formula (277) is a spectral description of the
source-defined contour jump, not its definition.

An $L^2$ function is a tempered distribution.  Uniqueness of the carrier
decomposition of Fourier hyperfunctions therefore forces the coefficient
of every point off the real Fourier axis to vanish: a term with
$\Re\rho\ne1/2$ has nonzero real exponential type in one time direction
and cannot occur in a tempered carrier.  The remaining terms have
$\rho=1/2+i\gamma$.  Their Fourier transform is the pure-point distribution
\begin{equation}
 2\pi\sum_{\Re\rho=1/2}m_\rho H(\rho)\delta_\gamma.   \tag{278}
\end{equation}
By Plancherel, the Fourier transform of (276) is represented by an $L^2$
function.  An $L^2$ function supported on a discrete set is zero; hence all
coefficients in (278) vanish as well.  Thus $H(\rho)=0$ at every nontrivial
zero of $\xi$.

Since $h$ is compactly supported, Paley--Wiener--Schwartz makes $H$ an
entire function of exponential type and finite polynomial order on
horizontal strips.  A nonzero such function has zero counting function
$O(R)$ by Jensen's formula.  Here multiplicities cannot simply be discarded
from Riemann--von Mangoldt.  We use instead the unconditional distinct-zero
theorem of X.~Wu (\emph{Q. J. Math.} 66 (2015), 759--771,
doi:10.1093/qmath/hav014): a fixed positive proportion of the nontrivial
zeros are distinct.  Hence their distinct counting function is
$\gg R\log R$.  Therefore an exponential-type entire function vanishing at
every nontrivial zero of $\xi$ must vanish identically.  Injectivity of the
bilateral Laplace transform on $\mathcal E'$ gives $h=0$.
\end{proof}

\paragraph{Scope of the preceding carrier argument --- \AUDIT.}
The passage from the infinite residue expansion (277) to coefficientwise
vanishing of every off-axis exponential uses uniqueness of the
Fourier-hyperfunction carrier decomposition with the required convergence
and growth topology.  That theorem has not been supplied in the present
manuscript.  Accordingly (274)--(278) is retained as an audited route, not
as load-bearing evidence for G.  The Hardy proof (283)--(292), which tests
the two half-planes separately and uses ordinary $H^2$ pole removal, avoids
this issue and is authoritative for the later closure reductions.

\begin{proposition}[Source identification of the jump --- \PROVED]
With the convolution/reflection conventions fixed in the supplied row (c),
the contour jump (275) and the locally finite operator (256) satisfy
\begin{equation}
 \boxed{\mathfrak J_\xi h=\mathcal A_\infty h}         \tag{279a}
\end{equation}
in $\mathcal D'(\mathbb R)$ whenever $H(0)=H(1)=0$.
This is an equality of source-defined distributions; it does not define
either side by the zeros.
\end{proposition}

\begin{proof}
Pair both sides with $\eta\in C_c^\infty(\mathbb R)$.  On the multiplicative
line the resulting test function is the convolution of the tests associated
with $h$ and $\eta^\vee$.  The arithmetic nuclear Lefschetz identity in the
supplied row (c) evaluates its source supercharacter as the two locally
finite prime orientations plus the Fourier-normalized archimedean term.
After central metric conjugation these are respectively
$\mathcal E_{\rm ar}h$ and $G_\Gamma h-m_0h$, hence their sum is the
right side of (279a).  The same row-(c) theorem evaluates that
supercharacter by shifting the two Mellin contours.  Its odd quotient gives
the residues of $D_\xi$, while the even quotient gives the characters at
$s=0,1$.  The latter are $H(0)E_\eta(0)$ and $H(1)E_\eta(1)$ and vanish by
hypothesis.  The remaining contour expression is exactly (275), with the
orientation chosen there.  Equality for every $\eta$ proves (279a).
\end{proof}

\begin{corollary}[Finite-energy closure criterion for G --- \PROVED]
Let $F\in\mathcal P_T$ be a primitive zero mode, so
\begin{equation}
 P_TA_TF=0,
 \qquad
 A_TF=\alpha e^{-t/2}+\beta e^{t/2}\quad(t\in I_T).  \tag{279}
\end{equation}
Extend $F$ by zero and put
$H_F(s)=\langle E_TF,e^{(s-1/2)t}\rangle$.  Then
$H_F(0)=H_F(1)=0$.  Consequently G40 and condition $G$ follow from the
single exterior estimate
\begin{equation}
 \boxed{\quad
 P_TA_TF=0\quad\Longrightarrow\quad
 1_{\mathbb R\setminus I_T}\mathcal A_\infty E_TF
       \in L^2(\mathbb R).\quad}                      \tag{280}
\end{equation}
It is enough to construct the unsupertraced observation (272) so that its
outgoing port is this exterior leakage and its equality-state storage
controls its $L^2$ norm.
\end{corollary}

\begin{proof}
The two equalities for $H_F$ are precisely the two primitive Tate moments.
On the right exterior the opposite orientation $Y_-$ and the compact
constant channel vanish, so (280) and the Gamma bound (281) imply
$Y_{+,T}^{\rm out}\in L^2(T,\infty)$.  Reflection gives
$Y_{-,T}^{\rm out}\in L^2(-\infty,-T)$.  By (286) these are exactly the two
Hardy memberships, and Hardy pole rigidity (292) gives $F=0$.
Thus no nonzero first-loss mode exists.  The first-loss theorem proves
$P_TA_TP_T\ge0$ for every $T$, and (248)--(249) give G40 and $G$.
\end{proof}

\paragraph{New minimal open bridge.}
Equation (280) is strictly weaker in appearance than directly proving the
pointwise arithmetic separation in (268): only finite exterior energy of
the undifferentiated $L^2$ mode is required.  It is also source typed.  G39
supplies the divergence/amplitude, the Gamma ODE and prime Julia systems
supply conservative local storages, and (272) must show that, when the
interior primitive divergence vanishes, the uncancelled exterior jump is
their $L^2$ outgoing port.  This exact finite-energy estimate remains
\OPEN{}; defining its norm from the desired sign would be circular.

\begin{theorem}[Exterior Gamma energy is automatic --- \PROVED]
For every $T>0$ there is $C_T<\infty$ such that
\begin{equation}
 \|1_{\mathbb R\setminus I_T}G_\Gamma E_TF\|_2
 \le C_T\|F\|_{L^2(I_T)}                              \tag{281}
\end{equation}
for all $F\in L^2(I_T)$.  Thus the open estimate (280) concerns only the
two oriented arithmetic synthesis tails.
\end{theorem}

\begin{proof}
Put $k(u)=e^{-u/2}/(1-e^{-2u})$.  Formula (81) gives, outside the support,
\begin{equation}
 (G_\Gamma E_TF)(t)
 =-\int_0^\infty k(u)
       \bigl(E_TF(t-u)+E_TF(t+u)\bigr)\,du.           \tag{282}
\end{equation}
For $t=T+x$, $x>0$, only the first summand occurs.  With $s=T-y$ it has
kernel $k(x+y)$ on $x>0$, $0<y<2T$.  Since
$k(v)=(2v)^{-1}+O(1)$ for $0<v\le1$ and
$k(v)=O(e^{-v/2})$ for $v\ge1$, this operator is the sum of a truncated
Carleman operator and a Hilbert--Schmidt operator.  Carleman's inequality
\[
 \int_0^\infty\left|\int_0^\infty
       \frac{f(y)}{x+y}\,dy\right|^2dx
 \le\pi^2\|f\|_2^2
\]
proves the right-exterior bound.  Reflection gives the left-exterior bound,
and their sum proves (281).
\end{proof}

\begin{theorem}[Exact Hardy form of the remaining arithmetic bridge ---
\PROVED]
For $F\in L^2(I_T)$ define
\begin{align}
 K_F(z)&:=\int_{-T}^{T}F(u)e^{-zu}\,du,\notag\\
 Y_+(t)&:=\sum_{n\ge2}\frac{\Lambda(n)}{\sqrt n}
                 E_TF(t-\log n).                     \tag{283}
\end{align}
The sum is locally finite as an $L^2_{\rm loc}$ function.  Put
$Y_{+,T}^{\rm out}=1_{(T,\infty)}Y_+$ and
$Y_{+,T}^{\rm in}=1_{(-\infty,T]}Y_+$.  The latter has compact support in
$[\log2-T,T]$.  If
\begin{equation}
 Q_{+,T}(z):=\int_{\log2-T}^{T}
               Y_+(t)e^{-zt}\,dt,                    \tag{284}
\end{equation}
then $Q_{+,T}$ is entire and, initially for $\Re z>1/2$,
\begin{equation}
 \int_T^\infty Y_+(t)e^{-zt}\,dt
 =-\frac{\zeta'}{\zeta}\left(\frac12+z\right)K_F(z)
  -Q_{+,T}(z).                                        \tag{285}
\end{equation}
Consequently
\begin{equation}
 Y_{+,T}^{\rm out}\in L^2(T,\infty)
 \Longleftrightarrow
 e^{Tz}\left[-\frac{\zeta'}{\zeta}
     \left(\frac12+z\right)K_F(z)-Q_{+,T}(z)\right]
 \in H^2(\mathbb C_+).                                \tag{286}
\end{equation}
The reflected formulas, with $F(u)$ replaced by $F(-u)$, characterize the
left arithmetic tail.  Therefore (280) is equivalent to the pair of
explicit Hardy memberships (286).
\end{theorem}

\begin{proof}
Local finiteness follows because on a fixed compact $t$-set only finitely
many translates of $[-T,T]$ occur.  For $\Re z>1/2$, Tonelli followed by
$t=u+\log n$ gives
\begin{align*}
 \int_{\mathbb R}Y_+(t)e^{-zt}\,dt
 &=K_F(z)\sum_{n\ge2}\frac{\Lambda(n)}{n^{z+1/2}}\\
 &=-\frac{\zeta'}{\zeta}\left(\frac12+z\right)K_F(z).
\end{align*}
The support lower bound and the split at $T$ give (284)--(285).  The
Paley--Wiener theorem for the Laplace transform identifies
$L^2(T,\infty)$ isometrically, after multiplication by $e^{Tz}$, with
$H^2(\mathbb C_+)$, proving (286).  Reflection proves the other orientation.
Finally (281) removes the Gamma tail, while $m_0E_TF$ has no exterior part;
hence the two Hardy conditions are exactly (280).
\end{proof}

\paragraph{Hardy target for the jump carrier.}
The remaining task is now to derive (286), for a first-loss solution (279),
from the conservative Gamma ODE/prime Julia storage system.  Analytic
continuation alone is insufficient: $H^2$ requires the uniform boundary
estimate
\begin{equation}
 \sup_{\sigma>0}\int_{\mathbb R}
 \left|e^{T(\sigma+i\tau)}\left[
 -\frac{\zeta'}{\zeta}(\tfrac12+\sigma+i\tau)
 K_F(\sigma+i\tau)-Q_{+,T}(\sigma+i\tau)\right]
 \right|^2d\tau<\infty,                               \tag{287}
\end{equation}
and its reflection.  Proving (287) by telescoping the already constructed
positive cell storages is the exact non-circular G40 bridge still \OPEN{}.

\begin{theorem}[The interior Wiener--Hopf polynomial does not alter the
critical Abel defect --- \PROVED]
Let
\begin{equation}
 \mathcal H_{+,T}F(z):=
 -\frac{\zeta'}{\zeta}(\tfrac12+z)K_F(z)
 -Q_{+,T}(z),                                       \tag{695}
\end{equation}
the meromorphic continuation of the right-tail transform in (285).
For every compact interval \(J\) whose endpoints are not critical-zero
ordinates,
\begin{equation}
 \frac{\sigma}{\pi}
 |\mathcal H_{+,T}F(\sigma+i\tau)|^2\,d\tau
 \buildrel *\over\rightharpoonup
 \sum_{\substack{\rho=1/2+i\gamma\\\gamma\in J}}
 m_\rho^2|K_F(i\gamma)|^2\delta_\gamma             \tag{696}
\end{equation}
on \(J\) as \(\sigma\downarrow0\).  Thus the compact prehistory
\(Q_{+,T}\) changes only the \(O(1)\) part of the local Abel energy; it
cannot cancel a critical \(1/\sigma\) atom.
\end{theorem}

\begin{proof}
The function \(Q_{+,T}\) is entire and uniformly bounded on a compact
neighbourhood of \(\{i\tau:\tau\in J\}\).  In disjoint discs around the
finitely many critical zeros over \(J\), write the first term in (695) as
\[
 -\frac{m_\rho K_F(i\gamma)}
        {\sigma+i(\tau-\gamma)}+O(1).
\]
The square of the bounded term, multiplied by \(\sigma\), tends to zero in
\(L^1(J)\).  The cross term is bounded after integration by
\[
 C\sigma\int_{-\delta}^{\delta}
 \frac{dx}{\sqrt{\sigma^2+x^2}}
 =O(\sigma\log(1/\sigma))\longrightarrow0.
\]
The Cauchy-kernel square gives the atoms exactly as in (495)--(496).
Away from the discs every term is bounded and disappears after
multiplication by \(\sigma\).  This proves (696).
\end{proof}

\begin{corollary}[For parity sources, the one-sided estimate is an exact
zero theorem --- \PROVED]
Let \(F\in\mathcal P_T\) have reflection parity.  Then
\begin{equation}
 R_{+,T}\in L^2(T,\infty)\quad\Longleftrightarrow\quad F=0.       \tag{697}
\end{equation}
In particular, restricted to parity first-loss states, proving (417) is
logically identical to excluding the state; the causal estimate contains
no weaker residual regularity gate.
\end{corollary}

\begin{proof}
The reverse implication is immediate.  If \(R_{+,T}\in L^2\), the
Paley--Wiener theorem gives a uniform \(H^2(\mathbb C_+)\) norm for
\(e^{Tz}\mathcal H_{+,T}F(z)\).  Equation (696) forces
\(K_F(i\gamma)=0\) at every critical-line zero.  The unconditional family
of simple critical zeros has counting function \(\gg R\log R\), whereas a
nonzero Cartwright function of type \(T\) has only \(O_T(R)\) zeros.
Hence \(K_F\equiv0\), and \(F=0\).  Parity is retained here because it is
the form in which (417) is used in the first-loss reduction, although the
critical-zero argument itself already proves the implication.
\end{proof}

\begin{theorem}[Causal Hardy closure and negative-port closure coincide on
the zero-mode space --- \PROVED]
Let
\[
 (\mathcal C_T^{\rm simp}F)_\gamma:=K_F(i\gamma),
\]
where \(\frac12+i\gamma\) ranges over the simple zeros of \(\zeta\) on the
critical line.  The sampling theorem used in (549) makes
\(\mathcal C_T^{\rm simp}\) a bounded observation on the Gamma form domain.
For every compact \(J\) as in (696),
\begin{equation}
 \lim_{\sigma\downarrow0}\frac{\sigma}{\pi}
 \int_J|\mathcal H_{+,T}F(\sigma+i\tau)|^2\,d\tau
 \ge
 \sum_{\substack{\gamma\in J\\
       1/2+i\gamma\ {\rm simple}}}
 |(\mathcal C_T^{\rm simp}F)_\gamma|^2 .            \tag{698}
\end{equation}
If in addition \(F\in\ker H_T\), then
\begin{equation}
 \|\mathcal C_T^{\rm simp}F\|^2
 \le \|P_+\mathcal E_TF\|^2
 =\|\mathcal N_TF\|^2.                              \tag{699}
\end{equation}
On \(\ker H_T\) the following four statements are equivalent:
\begin{enumerate}
\item the right causal current belongs to \(L^2(T,\infty)\), i.e. (417);
\item \(\mathcal C_T^{\rm simp}F=0\);
\item \(\mathcal N_TF=0\), i.e. the equality-state closure (513);
\item \(F=0\).
\end{enumerate}
Thus (417) and the Hermitian-connection route do not leave two independent
terminal gaps: on a first-contact space they are two boundary realizations
of the same zero theorem.
\end{theorem}

\begin{proof}
Equation (698) is the restriction of the positive atomic limit (696) to
the simple critical zeros.  In the residue decomposition, these coordinates
are an orthogonal subfamily of the positive port; hence
\(\|\mathcal C_T^{\rm simp}F\|^2\le
\|P_+\mathcal E_TF\|^2\).  For \(F\in\ker H_T\), the Krein balance (512)
gives the equality in (699).

The implication (1)\(\Rightarrow\)(4) is (697), and
(2)\(\Rightarrow\)(4) is the same unconditional
Heath--Brown--Levinson/Cartwright rigidity argument: a nonzero entire
function of exponential type \(T\) cannot vanish at
\(\gg R\log R\) distinct simple critical zeros.  If (3) holds, (512)
gives \(P_+\mathcal E_TF=0\), hence (2), and therefore (4).  All reverse
implications follow from \(F=0\).  This proves the equivalence.
\end{proof}

\paragraph{Single surviving implication --- \OPEN.}
The exact remaining assertion may now be written without choosing between
the two constructions:
\begin{equation}
 \boxed{F\in\ker H_T\ \Longrightarrow\
        \mathcal C_T^{\rm simp}F=0.}                 \tag{700}
\end{equation}
The zero-mode balance proves only the non-strict estimate (699).  Replacing
it by equality with zero requires either the causal source estimate (417)
or the bounded Hermitian product rule (661).  In particular, discarding the
negative off-line port in (699) would assume the assertion to be proved.

\begin{corollary}[The compensated causal current has trivial physical
\(L^2\)-domain --- \PROVED]
Regard the locally finite right current in (567) as a linear map
\[
 \mathscr R_{+,T}:\mathcal P_T\longrightarrow
 L^2_{\rm loc}(T,\infty).
\]
Then
\begin{equation}
 \operatorname{Dom}_{L^2}(\mathscr R_{+,T})
 :=\{F\in\mathcal P_T:\mathscr R_{+,T}F\in L^2(T,\infty)\}
 =\{0\}.                                             \tag{701}
\end{equation}
Consequently no bounded Hilbert-output operator on a nonzero source
subspace can agree with the physical current (567).  In particular, a
telescoping prime-cell construction cannot prove (417) as an ordinary
bounded-input/bounded-output estimate valid on a dense source core.
It can only prove the equality-state implication (700), in which case its
conclusion immediately annihilates that state.
\end{corollary}

\begin{proof}
The current (567) equals the right arithmetic tail after the two primitive
Tate moments remove the mean term, by (288)--(291).  The implication
\(\mathscr R_{+,T}F\in L^2\Rightarrow F=0\) is (697); the zero source gives
the converse.  This proves (701).  If a Hilbert-output operator agreed with
\(\mathscr R_{+,T}\) on a nonzero source subspace, every vector in that
subspace would belong to the left side of (701), a contradiction.
\end{proof}

\paragraph{Causal-route verdict --- \AUDIT.}
Equation (701) does not disprove the conditional implication (417) for a
first-loss state; it proves that such an implication cannot arise by first
extending the displayed current to a bounded causal observation on ordinary
sources.  The viable calculation must use the zero-mode equation before
taking the exterior norm, or equivalently construct the Hermitian
equality-state product rule (661).  Treating the cell cascade as a generic
\(L^2\) synthesis would contradict (701).

\begin{theorem}[Exact Tate cancellation of the prime mean --- \PROVED]
Let
\begin{equation}
 \Psi(x)=\sum_{n\le x}\Lambda(n),
 \qquad R_\Psi(x)=\Psi(x)-x.                          \tag{288}
\end{equation}
For $F\in L^2(I_T)$ and almost every $t>T$ the right arithmetic tail in
(283) has the Stieltjes representation
\begin{align}
 Y_+(t)
 &=\int_{e^{t-T}}^{e^{t+T}}
 x^{-1/2}F(t-\log x)\,d\Psi(x)                       \tag{289}\\
 &=e^{t/2}M_-F+
 \int_{e^{t-T}}^{e^{t+T}}
 x^{-1/2}F(t-\log x)\,dR_\Psi(x).                    \tag{290}
\end{align}
Hence on the primitive space
\begin{equation}
 \boxed{Y_+(t)=
 \int_{e^{t-T}}^{e^{t+T}}
 x^{-1/2}F(t-\log x)\,dR_\Psi(x).}                  \tag{291}
\end{equation}
The reflected left tail has the identical formula with $M_+F$, which also
vanishes.  Thus both outgoing arithmetic ports see only the global
Chebyshev fluctuation, with their full mean removed exactly by the Tate
terminals.
\end{theorem}

\begin{proof}
The jumps of $\Psi$ at $n$ are $\Lambda(n)$, so (289) is exactly (283).
Splitting $d\Psi=dx+dR_\Psi$ and substituting $u=t-\log x$ gives
\[
 \int_{e^{t-T}}^{e^{t+T}}x^{-1/2}F(t-\log x)\,dx
 =e^{t/2}\int_{-T}^{T}F(u)e^{-u/2}\,du,
\]
which proves (290)--(291).  Approximation by smooth compactly supported
functions justifies the identities for $L^2$ inputs.  Reflection replaces
$e^{-u/2}$ by $e^{u/2}$ and proves the final assertion.
\end{proof}

\begin{theorem}[Hardy pole rigidity --- \PROVED]
Let $F\in\mathcal P_T$ and let the two oriented functions in (286) belong
to their reflected Hardy spaces.  Then
\begin{equation}
 K_F(\rho-\tfrac12)=0                                \tag{292}
\end{equation}
for every distinct nontrivial zero $\rho$ of $\zeta$.  Consequently
$F=0$.
\end{theorem}

\begin{proof}
The right-hand expression in (285) has a meromorphic continuation to
$\mathbb C_+$.  The pole of $\zeta'/\zeta$ at $s=1$ is cancelled by
$K_F(1/2)=M_-F=0$.  At a zero $\rho$ with $\Re\rho>1/2$, holomorphy of an
$H^2(\mathbb C_+)$ function forces the residue to vanish and gives (292).
If $\rho=1/2+i\gamma$, a nonzero residue would produce locally a multiple
of $(z-i\gamma)^{-1}$; its squared boundary norm on $\Re z=\sigma$ grows
like $\sigma^{-1}$, contradicting the uniform Hardy norm.  Hence (292)
also holds on the central line.  Applying the reflected Hardy condition
gives (292) for $\Re\rho<1/2$.

The function $z\mapsto K_F(z)$ is entire of exponential type.  By the
unconditional distinct-zero theorem of X.~Wu
(\emph{Q. J. Math.} 66 (2015), 759--771,
doi:10.1093/qmath/hav014), the set in (292) has counting function
$\gg R\log R$.
Jensen's formula permits only $O(R)$ zeros for a nonzero exponential-type
entire function.  Thus $K_F\equiv0$, and injectivity of the Fourier/Laplace
transform gives $F=0$.
\end{proof}

\paragraph{Interpretation.}
The two Hardy estimates do not need to estimate the prime mean: (291)
shows that it is identically absent.  They must control the Chebyshev
fluctuation against the very special first-loss mode.  Hardy pole rigidity
then converts that control into vanishing at every zeta zero and forces the
mode itself to vanish.  Invoking a general estimate for $R_\Psi$ strong
enough to imply (287) would merely import RH; the required estimate must
come from the equality-state Poisson--Julia storage relation.

\begin{theorem}[Exact Wiener--Hopf equation of a zero mode --- \PROVED]
Define the opposite arithmetic orientation
\begin{equation}
 Y_-(t):=\sum_{n\ge2}\frac{\Lambda(n)}{\sqrt n}
                    E_TF(t+\log n)                   \tag{293}
\end{equation}
and its interior transform
\begin{equation}
 Q_{-,T}(z):=\int_{-T}^{T}Y_-(t)e^{-zt}\,dt.          \tag{294}
\end{equation}
Both $Q_{+,T}$ and $Q_{-,T}$ are finite sums and hence entire.  Put
\begin{align}
 \Gamma_T(z)&:=\int_{-T}^{T}(G_\Gamma E_TF)(t)e^{-zt}\,dt,\notag\\
 P_{\alpha,\beta,T}(z)&:=\int_{-T}^{T}
   (\alpha e^{-t/2}+\beta e^{t/2})e^{-zt}\,dt.        \tag{295}
\end{align}
Then the zero-mode equation (279) is exactly
\begin{equation}
 \boxed{
 Q_{+,T}(z)+Q_{-,T}(z)
 =\Gamma_T(z)-m_0K_F(z)-P_{\alpha,\beta,T}(z)}        \tag{296}
\end{equation}
for every $z\in\mathbb C$.  If $F(-t)=\varepsilon F(t)$ with
$\varepsilon\in\{1,-1\}$, then
\begin{equation}
 Q_{-,T}(z)=\varepsilon Q_{+,T}(-z),
 \qquad K_F(-z)=\varepsilon K_F(z),                  \tag{297}
\end{equation}
and therefore
\begin{equation}
 Q_{+,T}(z)+\varepsilon Q_{+,T}(-z)
 =\Gamma_T(z)-m_0K_F(z)-P_{\alpha,\beta,T}(z).       \tag{298}
\end{equation}
Thus G39/zero-mode divergence fixes the reflected symmetric part of the
oriented Poisson transform.  The complementary reflected antisymmetric
part, equivalently its causal Hardy choice in (286), is the sole datum not
fixed by (298).
\end{theorem}

\begin{proof}
Only $n<e^{2T}$ can contribute when both the input and observation variable
belong to $I_T$, so (284) and (294) are finite sums.  Multiply (279) by
$e^{-zt}$, integrate on $I_T$, and use
$A_TF=G_{\Gamma,T}F-m_0F-Y_+-Y_-$ to obtain (296) first on the real line
and then everywhere by entireness.  Reflection gives
$Y_-(t)=\varepsilon Y_+(-t)$, which proves (297) by changing $t$ to $-t$.
Substitution gives (298).  Finally a function and its reflection are the
sum of their reflected symmetric and antisymmetric parts, so (298) fixes
only the former.  Selecting the latter by its right/left support is exactly
the Wiener--Hopf/Hardy separation (286).
\end{proof}

\paragraph{Engineering specification for the missing object.}
The new bridge may now be required to do one sharply delimited job.  It must
lift the supertraced equality (298) to the two causal Poisson ports and prove
that the lifted right port and its reflection have the $H^2$ bounds (287).
It need not reconstruct G39, Gamma normalization or the Tate terminals;
those already constitute the right side of (298).  A proposed object which
only reproduces (298) but supplies no bounded causal splitting has not
advanced G40.

\begin{theorem}[Equality-mode Poisson lift closes G --- \PROVED]
Let $\mathcal U_c$ denote the fixed critical logarithmic conjugation
\begin{equation}
 (\mathcal U_cF)(x)=x^{-1/2}F(\log x),
 \qquad x>0.                                          \tag{299}
\end{equation}
Suppose that for every primitive zero mode $F$ there is
$f_F\in\mathcal H_\cap$ such that
\begin{equation}
 Zf_F=\mathcal U_cE_TF.                               \tag{300}
\end{equation}
Then the exterior estimate (280) holds and consequently condition $G$ is
true.  No isometry or uniform Gram identity is required in (300); existence
of the lift on the zero-mode space is sufficient.
\end{theorem}

\begin{proof}
For $f_F\in\mathcal H_\cap$, differentiated Poisson (96)--(98) gives
\begin{equation}
 (C+\mathscr JC\mathscr J)Zf_F
 =Z\partial f_F+\mathscr JZ\partial\mathcal Ff_F.     \tag{301}
\end{equation}
By the defining weighted-Schwartz properties of $\mathcal H_\cap$ and
Meyer's Poisson theorem, the right side belongs to $\mathcal H_-$ and hence,
after critical logarithmic conjugation, to $L^2(\mathbb R)$.  The left side
is exactly the sum of the two oriented arithmetic synthesis outputs acting
on $E_TF$.  The exterior Gamma output is in $L^2$ by (281), and the constant
term has compact support.  Therefore
$\mathcal A_\infty E_TF\in L^2(\mathbb R)$, proving (280).  Finite-energy
closure then excludes every nonzero first-loss mode and proves $G$.
\end{proof}

\begin{proposition}[Spectral meaning and non-circular construction target ---
\PROVED]
For a compactly supported primitive $F$, a lift (300), if it exists, is
unique because $Z$ is injective on the relevant Meyer source space.  In
Mellin coordinates it necessarily satisfies
\begin{equation}
 \widehat f_F(s)=\frac{H_F(s)}{\zeta(s)}               \tag{302}
\end{equation}
in the initial half-plane.  Thus existence of (300) forces the apparent
poles in (302) to be removable.  Conversely, merely defining (302) by
meromorphic division is not a construction: proving that the quotient is
the Mellin transform of an element of $\mathcal H_\cap$ is exactly the
closed-range/Hardy assertion (286)--(287).

The fixed-parameter CCM Sonin isomorphism is bounded with bounded inverse on
its own restricted Sonin spaces, so it can transport a lift after (300) is
constructed.  It does not supply the source map from the physical zero mode
to $\mathcal H_\cap$; the missing map (88) was precisely that type/Gram
interface.  For the first-loss contradiction the metric requirement can now
be weakened: one needs a source-defined equality-mode lift (300), not an
all-source isometry (88).
\end{proposition}

\begin{proof}
Injectivity is part of Meyer's closed-range theorem for $Z$ on the source
space.  Mellin transformation changes Zeta convolution into multiplication
by $\zeta(s)$, proving (302).  If a lift lies in $\mathcal H_\cap$, its Mellin
transform has the prescribed weighted-Schwartz/Poisson boundary behaviour,
so all poles are removable.  The converse demands precisely those topology
and boundary estimates and cannot follow from formal division.  The final
claim follows from the domain statement of the CCM stability theorem: its
bounded map begins on the Sonin source, whereas $F$ initially belongs only
to the physical primitive $L^2$ space.
\end{proof}

\paragraph{Smallest presently sufficient new object.}
One may therefore replace the ambitious global observation isometry by the
partial operator
\begin{equation}
 \mathcal L_T:\ker(P_TA_TP_T)\longrightarrow\mathcal H_\cap,
 \qquad Z\mathcal L_T=\mathcal U_cE_T.                \tag{303}
\end{equation}
The kernel is finite dimensional by compact resolvent, so boundedness is
automatic once the lift exists.  Proving existence of (303) from the
zero-mode Wiener--Hopf equation (298), without dividing by unverified zeros,
would close G.  This equality-mode Poisson lifting theorem remains \OPEN{}.

\begin{theorem}[Exact source criterion for the equality-mode Poisson lift ---
\PROVED]
Let $g\in\mathcal H_-$ be compactly supported and define its two half-lifts
on Meyer's common source domain by
\begin{equation}
 f_>(g):=Z^{-1}g,
 \qquad
 f_<(g):=\mathcal FZ^{-1}\mathscr Jg.                 \tag{304}
\end{equation}
Then
\begin{equation}
 \boxed{\quad
 g\in Z\mathcal H_\cap
 \quad\Longleftrightarrow\quad
 f_>(g)=f_<(g).\quad}                                \tag{305}
\end{equation}
When these conditions hold, their common value is the unique lift
$f=Z^{-1}g\in\mathcal H_\cap$.  Consequently, for a zero mode in the
common source domain, the equality-mode lift (303) exists if and only if
the source-defined Poisson defect
\begin{equation}
 \mathfrak d_{\rm P,T}F
 :=Z^{-1}\mathcal U_cE_TF
  -\mathcal FZ^{-1}\mathscr J\mathcal U_cE_TF         \tag{306}
\end{equation}
vanishes.
\end{theorem}

\begin{proof}
If $g=Zf$ with $f\in\mathcal H_\cap$, the first equality in (304) gives
$f_>(g)=f$.  Poisson summation on $\mathcal H_\cap$ is
$Zf=\mathscr JZ\mathcal Ff$.  Applying $\mathscr J$, then $Z^{-1}$ and
$\mathcal F$ (with $\mathcal F^2=I$ on the even source) gives
$f_<(g)=f$.

Conversely suppose the two half-lifts agree and call their common value
$f$.  Then $Zf=g$ by the first definition.  The second definition gives
$\mathcal Ff=Z^{-1}\mathscr Jg$, hence
$Z\mathcal Ff=\mathscr Jg$ and
$\mathscr JZ\mathcal Ff=g=Zf$.  The estimates in the proof of Meyer's
Theorem~3.3 show that simultaneous membership of these two halves places
$f$ in $\mathcal H_+$; moreover $Zf=g\in\mathcal H_-$, so the last assertion
of that theorem forces $f\in\mathcal H_\cap$.  Uniqueness follows from
injectivity of $Z$.  Substituting $g=\mathcal U_cE_TF$ proves (306).
\end{proof}

\begin{proposition}[Domain gate for physical zero modes --- \AUDIT]
The first-loss theorem initially gives $F\in L^2(I_T)$ in the form domain of
the logarithmic Gamma operator.  It does not by itself give
$g=\mathcal U_cE_TF\in\mathcal H_-$, whose elements are nuclear
weighted-Schwartz vectors.  Pointwise local finiteness defines
\begin{equation}
 Z^{-1}g(x)=\sum_{n\ge1}\mu(n)g(nx)                   \tag{306a}
\end{equation}
as a local distribution for compactly supported $g$, but membership in the
Meyer source topology, and the action of
$\mathcal FZ^{-1}\mathscr J$, still require a closed Hilbert/rigged
extension.  Thus applying Theorem (304)--(306) to a physical zero mode
requires one of the following independently proved facts:
\begin{enumerate}
\item zero-mode regularity places $\mathcal U_cE_TF$ in the common Meyer
      source domain; or
\item the two half-lifts extend to a physical completion and equality in
      (305) places their common value in a completion on which (301) has
      $L^2$ output.
\end{enumerate}
Neither fact follows merely from compact support or logarithmic form
regularity.  This domain gate is part of the \OPEN{} lifting theorem.
\end{proposition}

\begin{theorem}[Compact Poisson-range rigidity --- \PROVED]
One has
\begin{equation}
 Z\mathcal H_\cap\cap C_c^\infty(\mathbb R_+^\times)=\{0\}.    \tag{306b}
\end{equation}
More generally, if a compactly supported source $g$ admits the two equal
half-lifts in (304) as tempered $L^2$ functions, then $g=0$.
\end{theorem}

\begin{proof}
Suppose $g=Zf$ is supported in a compact multiplicative interval
$[a,b]\subset(0,\infty)$ and $f\in\mathcal H_\cap$.  Möbius inversion gives
\[
 f(x)=Z^{-1}g(x)=\sum_{n\ge1}\mu(n)g(nx).
\]
For $x>b$ every summand vanishes; after the even extension defining
$\mathcal H_+$, $f$ has compact additive support in $[-b,b]$.  Poisson
summation gives $\mathscr Jg=Z\mathcal Ff$, hence
\[
 \mathcal Ff=Z^{-1}\mathscr Jg.
\]
Since $\mathscr Jg$ is supported in $[b^{-1},a^{-1}]$, the same argument
shows that $\mathcal Ff$ has compact additive support.  The Fourier
uncertainty theorem (or the identity theorem applied to the entire
Paley--Wiener continuation of $\mathcal Ff$) forces $f=0$, hence $g=0$.
The identical argument applies in the stated tempered $L^2$ completion.
\end{proof}

\begin{proposition}[Meyer near-section audit --- \AUDIT]
Let $\iota_+f=(Zf,\mathscr JZ\mathcal Ff)$ and
$\iota_-g=(g,g)$.  Meyer's cutoff operator
\begin{equation}
 \pi_+(g_1,g_2)=M_\varphi Z^{-1}g_1
  +\mathcal F M_\varphi Z^{-1}\mathscr Jg_2          \tag{307}
\end{equation}
is continuous and is sufficient for his nuclear trace computation, but it
is not a section of $\iota_+$.  Indeed
\begin{equation}
 \pi_+\iota_+
 =I+M_\varphi-\mathcal F\mathscr JM_\varphi
                  \mathscr J\mathcal F,              \tag{308}
\end{equation}
and Meyer explicitly leaves open whether $\operatorname{Ran}\iota_+$ is
complemented.  Therefore neither
$\iota_+\pi_+\iota_-g$ nor its difference from $\iota_-g$ is an exact range
projection or an exact norm-square defect.  Using it as one would silently
assume the missing G40 splitting.  The authoritative obstruction is the
half-lift difference (306).
\end{proposition}

\paragraph{New zero-free target.}
The final equality-mode statement can now be written without spectral
division, Hardy boundary values or an unspecified projection:
\begin{equation}
 \boxed{\quad
 P_TA_TF=0,\quad F\in\mathcal P_T
 \quad\Longrightarrow\quad
 \mathfrak d_{\rm P,T}F=0
 \ \hbox{in a physical Poisson completion satisfying (301).}\quad}
                                                               \tag{309}
\end{equation}
By (305), (309) constructs (303); by (299)--(301), it proves finite exterior
energy; compact Poisson-range rigidity already excludes the mode directly
(and (274)--(280) supplies the independent exterior-energy route).  What
remains is to derive (309) from the zero-mode Wiener--Hopf identity (298) and
the Gamma--Tate boundary normalization.  This implication remains \OPEN{}.

\begin{theorem}[Exact transport equation for the Poisson defect --- \PROVED]
On Meyer's common source core let
\begin{equation}
 f=Z^{-1}g,\qquad
 r_{\rm P}(g):=g-\mathscr JZ\mathcal Ff.              \tag{310}
\end{equation}
Then
\begin{equation}
 \boxed{
 (C+\mathscr JC\mathscr J)g
 =\mathcal W_{\rm P}f
  -(\partial-\mathscr JC\mathscr J)r_{\rm P}(g),}     \tag{311}
\end{equation}
where
$\mathcal W_{\rm P}f=Z\partial f+\mathscr JZ\partial\mathcal Ff$
is (98).  Thus differentiated Poisson (97) is the special case
$r_{\rm P}(g)=0$.
\end{theorem}

\begin{proof}
The convolution derivation and
$C=Z\partial Z^{-1}$ give
\begin{equation}
 \partial g=-Cg+Z\partial f.                          \tag{312}
\end{equation}
Using $\partial\mathscr J=-\mathscr J\partial$ in the second term of (310)
gives
\[
 \partial(\mathscr JZ\mathcal Ff)
 =\mathscr JCZ\mathcal Ff-\mathscr JZ\partial\mathcal Ff.
\]
Since $\mathscr JZ\mathcal Ff=g-r_{\rm P}(g)$, one has
$Z\mathcal Ff=\mathscr Jg-\mathscr Jr_{\rm P}(g)$.
Substitute this equality, subtract the last display from (312), and collect
the two Euler orientations.  The result is exactly (311).
\end{proof}

\begin{corollary}[Forced defect equation for a zero mode --- \PROVED]
After critical conjugation, a primitive zero mode satisfies on $I_T$
\begin{align}
 (\partial-\mathscr JC\mathscr J)r_{\rm P}(g)
 ={}&\alpha e^{-t/2}+\beta e^{t/2}
      -(G_\Gamma-m_0)g+\mathcal W_{\rm P}Z^{-1}g,
                                                               \tag{313}
\end{align}
Consequently (313) becomes a homogeneous uniqueness equation for
$r_{\rm P}$ only after proving the source identity
\begin{equation}
 \mathcal W_{\rm P}Z^{-1}g
 =(G_\Gamma-m_0)g-\alpha e^{-t/2}-\beta e^{t/2}
 \quad\hbox{on }I_T.                                  \tag{314}
\end{equation}
Equation (314) is precisely the Gamma--Tate terminal/observation
identification, not a consequence of the algebraic differentiated Poisson
identity.  If (314) and zero incoming data for (313) are proved in the
physical completion, causal uniqueness gives $r_{\rm P}=0$ and hence (309).
\end{corollary}

\begin{proof}
Write $g=\mathcal U_cE_TF$.  Directly from (299),
\[
 \mathcal U_cU_n\mathcal U_c^{-1}g(t)
 =n^{-1/2}g(t+\log n),
 \qquad
 \mathcal U_c\mathscr J\mathcal U_c^{-1}g(t)=g(-t).
\]
Consequently $C+\mathscr JC\mathscr J$ becomes exactly
\[
 \sum_{n\ge2}\frac{\Lambda(n)}{\sqrt n}
   (\tau_{\log n}+\tau_{-\log n}),
\]
with the same positive orientation as the arithmetic sum subtracted in
$A_T$.  The zero-mode equation (279) is therefore
\[
 (C+\mathscr JC\mathscr J)g
 =(G_\Gamma-m_0)g-\alpha e^{-t/2}-\beta e^{t/2}
 \quad\hbox{on }I_T.
\]
Insert this exact equality in (311) and rearrange to obtain (313), with no
remaining sign convention.  Its right side vanishes exactly under (314),
proving the remaining claims.
\end{proof}

\begin{theorem}[Causal uniqueness for the Poisson-defect transport ---
\PROVED]
After critical logarithmic conjugation, the homogeneous equation associated
with (313) is
\begin{equation}
 t\,r(t)=\sum_{n\ge2}\frac{\Lambda(n)}{\sqrt n}
                   r(t-\log n).                       \tag{315}
\end{equation}
Let $r\in L^2_{\rm loc}(\mathbb R)$ satisfy (315) almost everywhere.  If
$r=0$ almost everywhere on $(-\infty,a]$ for some $a$, then $r=0$ almost
everywhere on the whole line.
\end{theorem}

\begin{proof}
For $t\in(a,a+\log2)$ every delayed argument $t-\log n$ lies below $a$.
Thus the right side of (315) vanishes.  Hence $t\,r(t)=0$, which gives
$r=0$ almost everywhere on that interval; the possible singleton $t=0$ is
irrelevant for an $L^2_{\rm loc}$ function.  Inductively, if $r$ vanishes on
$(-\infty,a+k\log2]$, the same argument gives its vanishing on the next
interval of length $\log2$.  Induction proves the assertion.
\end{proof}

\begin{corollary}[Two concrete gates for the defect route --- \PROVED]
Condition (309), and hence G, follows once the physical completion proves
for every first-loss mode:
\begin{enumerate}
\item the Gamma--Tate observation identity (314); and
\item the incoming support condition
\begin{equation}
 r_{\rm P}(\mathcal U_cE_TF)=0
 \quad\hbox{on }(-\infty,a]\quad\hbox{for some }a.     \tag{316}
\end{equation}
\end{enumerate}
Indeed (314) makes (313) homogeneous, and (315)--(316) give
$r_{\rm P}=0$.
\end{corollary}

\begin{proof}
This is direct substitution followed by the preceding causal-uniqueness
theorem and compact Poisson-range rigidity.
\end{proof}

\begin{theorem}[Laplace rigidity of the homogeneous defect --- \PROVED]
Let $r$ solve (315) in a class carrying a translation-covariant bilateral
Laplace transform
\begin{equation}
 R(z)=\langle r,e^{-z\cdot}\rangle,                    \tag{317}
\end{equation}
defined for $\Re z>1/2$, such that
\[
 \mathcal L(\tau_a r)(z)=e^{-az}R(z),
 \qquad
 \mathcal L(tr)(z)=-R'(z).
\]
Suppose that $R$ extends holomorphically to a connected domain meeting that
half-plane and containing $z=1/2$.  Then $r=0$.  A one-sided causal integral
is covered only when $r$ has zero prehistory, so that no lower-limit terms
occur under translation.
\end{theorem}

\begin{proof}
On $\Re z>1/2$, transformation of (315) is legitimate and gives
\[
 -R'(z)=
 \left(\sum_{n\ge2}\frac{\Lambda(n)}{n^{z+1/2}}\right)R(z)
 =-\frac{\zeta'}{\zeta}\left(\frac12+z\right)R(z).
\]
Hence
\begin{equation}
 R'(z)=\frac{\zeta'}{\zeta}\left(\frac12+z\right)R(z),
 \qquad
 R(z)=c\,\zeta\left(\frac12+z\right)                  \tag{318}
\end{equation}
on each connected component, by uniqueness for a scalar holomorphic ODE.
The right side has a pole at $z=1/2$ unless $c=0$, while $R$ is holomorphic
there by hypothesis.  Thus $c=0$, analytic continuation gives $R=0$, and
injectivity of the Laplace transform gives $r=0$.
\end{proof}

\begin{corollary}[Polar alternative to incoming support --- \PROVED]
In the defect route, condition (316) may be replaced by a
translation-covariant bilateral transform realization having the holomorphic
continuation stated in Theorem (317)--(318).  Therefore G follows from the
zero-mode observation identity (314) plus this transform-realization and
polar-regularity condition.  This is not claimed to be weaker for an
arbitrary causal transform: without zero prehistory, delay shifts create
additional history terms.
\end{corollary}

\begin{proof}
Equation (314) makes (313) homogeneous.  Apply Laplace rigidity and then
compact Poisson-range rigidity.
\end{proof}

\paragraph{Tate relevance.}
The exceptional point $z=1/2$ is $s=1$ in the Mellin variable.  The
primitive condition $H_F(1)=0$ removes the corresponding polar character
for the original compact source.  The next theorem proves that the
cancellation survives at the level of common meromorphic germs.  The
remaining issue is strictly a boundary-domain realization question: whether
the physical defect transform in (317) is represented by that germ.

\begin{theorem}[Common Mellin germ of the two Poisson half-lifts --- \PROVED]
Let $g$ be a compactly supported smooth multiplicative source, put
\[
 G(s):=\int_0^\infty g(x)x^s\,d^*x,
\]
and define the two half-lifts by (304) on their respective initial domains.
Their meromorphic Mellin continuations coincide:
\begin{equation}
 \boxed{
  \mathcal M f_>(g)(s)=\mathcal M f_<(g)(s)
       =\frac{G(s)}{\zeta(s)}.}                       \tag{319}
\end{equation}
Equivalently, the Poisson defect (310) has identically zero common
meromorphic Mellin germ.  At $s=1$ the germ in (319) is holomorphic and has
a zero of order at least one; if $G(1)=0$, it has a zero of order at least
two.  Thus the polar regularity still requested after (318) contains no
additional scalar cancellation at the Tate pole: its only remaining content
is that the physical boundary realization of the defect is represented by
this common germ in a connected transform domain.
\end{theorem}

\begin{proof}
Mellin transformation of Zeta convolution gives
\[
 \mathcal M f_>(g)(s)=\frac{G(s)}{\zeta(s)}
\]
in its initial half-plane.  The Mellin form of the self-dual Fourier
functional equation is
\begin{equation}
 \zeta(s)\,\mathcal M(\mathcal Fh)(s)
 =\zeta(1-s)\,\mathcal Mh(1-s).                       \tag{320}
\end{equation}
For $h=Z^{-1}\mathscr Jg$ one has
\[
 \mathcal Mh(1-s)
 =\frac{\mathcal M(\mathscr Jg)(1-s)}{\zeta(1-s)}
 =\frac{G(s)}{\zeta(1-s)}.
\]
Substitution in (320) proves the second equality in (319) wherever the
initial expressions converge, and meromorphic continuation proves it
globally.  Applying the same calculation before the last Fourier transform
gives
\[
 \mathcal M(\mathscr JZ\mathcal FZ^{-1}g)(s)=G(s),
\]
so the common meromorphic germ of $r_{\rm P}(g)$ is zero.

Finally $\zeta$ has a simple pole at $s=1$, hence $1/\zeta$ has a simple
zero there.  Since $G$ is entire, $G/\zeta$ is holomorphic and vanishes at
least simply at $1$; the Tate condition $G(1)=0$ raises the order by at
least one.  Notice that equality of meromorphic germs does not identify two
boundary values taken in disjoint convergence regions.  That last
topological identification is exactly the physical-completion assertion
left open above, so no Poisson range or positivity statement has been
assumed.
\end{proof}

\begin{theorem}[Meromorphic differentiated observation identity --- \PROVED]
Under the hypotheses of Theorem (319), let $f=Z^{-1}g$ and
$S_{\rm ar}:=C+\mathscr JC\mathscr J$.  The two sides of the desired
differentiated Poisson observation have the same meromorphic Mellin germ:
\begin{equation}
 \boxed{
 \mathcal M(\mathcal W_{\rm P}f)(s)
 =-\left(
       \frac{\zeta'}{\zeta}(s)
       +\frac{\zeta'}{\zeta}(1-s)
    \right)G(s)
 =\mathcal M(S_{\rm ar}g)(s).}                       \tag{321}
\end{equation}
Consequently, for a primitive zero mode, (314) is already true as an
identity of common meromorphic germs.  Its remaining content is exactly the
equality of the physical boundary representatives on $I_T$.
\end{theorem}

\begin{proof}
Write $F(s)=\mathcal Mf(s)=G(s)/\zeta(s)$.  Since multiplication by
$\log x$ differentiates the Mellin transform,
\[
 \mathcal M(Z\partial f)(s)=\zeta(s)F'(s).
\]
Put $Q(u)=\mathcal M(\mathcal Ff)(u)$.  Equation (320) reads
\[
 Q(u)=\frac{\zeta(1-u)}{\zeta(u)}F(1-u).
\]
Using $\mathcal M(\mathscr Jh)(s)=\mathcal Mh(1-s)$ and differentiating
the last display gives
\[
 \mathcal M(\mathscr JZ\partial\mathcal Ff)(s)
 =-\zeta(s)F'(s)
  -\left(\zeta'(s)
        +\zeta(s)\frac{\zeta'(1-s)}{\zeta(1-s)}\right)F(s).
\]
Adding the two current orientations cancels the derivatives of $F$ and,
because $\zeta(s)F(s)=G(s)$, gives the first equality in (321).
On the other hand $C=\sum_{n\ge2}\Lambda(n)U_n$ has Mellin multiplier
$-\zeta'/\zeta(s)$, while conjugation by $\mathscr J$ replaces $s$ by
$1-s$.  This proves the second equality.

For a zero mode the exact critical conjugation in the proof of (313) gives
\[
 S_{\rm ar}g=(G_\Gamma-m_0)g
       -\alpha e^{-t/2}-\beta e^{t/2}\quad\hbox{on }I_T.
\]
Combining this local equality with (321) proves the last assertion at the
level of common germs, but does not identify their boundary realizations.
\end{proof}

\begin{definition}[Physical Gamma--Tate boundary carrier]
For a physical zero mode for which the indicated terms are defined in a
common distributional completion, set
\begin{equation}
 \mathfrak B_TF
 :=P_{I_T}\left[
   \mathcal W_{\rm P}Z^{-1}g-S_{\rm ar}g\right]
 =P_{I_T}(\partial-\mathscr JC\mathscr J)r_{\rm P}(g),
 \qquad g=\mathcal U_cE_TF.                           \tag{322}
\end{equation}
The second equality is the exact transport identity (311).  Theorem (321)
proves that $\mathfrak B_TF$ is a boundary carrier with zero meromorphic
Mellin germ; it is not defined by a scalar symbol.
\end{definition}

\begin{corollary}[Single boundary-carrier target --- \PROVED]
Condition $G$ follows if every first-loss zero mode belongs to a physical
completion in which
\begin{equation}
 \boxed{\mathfrak B_TF=0}                              \tag{323}
\end{equation}
and the corresponding defect transform realizes the common germ (319) on
the connected domain in Theorem (317)--(318), with the
translation-covariance required there.  In particular, the former
two scalar-looking gates---Gamma--Tate normalization (314) and cancellation
at $s=1$---are two boundary aspects of one realization theorem, not two new
arithmetic identities.
\end{corollary}

\begin{proof}
Equation (323) is (314), by the zero-mode equality used in the proof of
(321), and makes the defect transport homogeneous.  The covariant
common-germ realization supplies the holomorphy at $s=1$ without delay
boundary terms.  Laplace rigidity then gives $r_{\rm P}=0$, and Theorems
(299)--(306b) give the first-loss contradiction and hence $G$.
\end{proof}

\begin{corollary}[Zero-free form of the remaining polar gate --- \PROVED]
Suppose a physical completion for a first-loss mode has the following two
properties:
\begin{enumerate}
\item the defect satisfies the homogeneous transport equation, equivalently
      the observation identity (314); and
\item it has a translation-covariant bilateral Laplace--Mellin realization
      representing the common germ (319) on one connected domain meeting
      $\Re s>1$ and a neighbourhood of $s=1$.
\end{enumerate}
Then the defect vanishes and condition $G$ follows.
\end{corollary}

\begin{proof}
The second property supplies both the covariance and the holomorphic
continuation at $s=1$, or $z=1/2$, required in Theorem (317)--(318); in fact
the represented defect germ is identically zero.  The first property permits
application of that theorem.  Hence $r_{\rm P}=0$, and (305),
(299)--(301), and compact Poisson-range rigidity close the first-loss
contradiction.
\end{proof}

\paragraph{Boundary-carrier audit.}
Theorem (319) explains why scalar meromorphic symmetrization cannot prove
G: it sees both half-lifts as the same germ and therefore sees no defect at
all.  A nonzero defect can only be a mismatch of boundary realizations
across the disjoint initial Mellin domains.  The new theorem removes the
putative scalar pole at $s=1$, but it does not prove that this boundary
carrier belongs to the physical Laplace class of (317).  Establishing that
realization, together with (314), remains \OPEN{}.

\begin{theorem}[Parity--Fourier reduction of the boundary carrier --- \PROVED]
Let $g$ be a source of parity
\begin{equation}
 \mathscr Jg=\varepsilon g,
 \qquad \varepsilon\in\{1,-1\},                       \tag{324}
\end{equation}
and put $f=Z^{-1}g$ on the common source calculus.  Then the half-lift defect
in (306) is
\begin{equation}
 d_{\varepsilon}(g):=f-\varepsilon\mathcal Ff,
 \qquad
 \mathcal Fd_{\varepsilon}(g)=-\varepsilon
             d_{\varepsilon}(g).                     \tag{325}
\end{equation}
Moreover
\begin{align}
 r_{\rm P}(g)&=\varepsilon\mathscr JZ d_{\varepsilon}(g),\notag\\
 \mathcal W_{\rm P}f-S_{\rm ar}g
 &=-\varepsilon\mathscr JZ\partial d_{\varepsilon}(g).
                                                               \tag{326}
\end{align}
Consequently the entire missing boundary carrier lies in the Fourier
eigenspace with eigenvalue opposite to the physical parity.  For a first-loss
mode, (323) is exactly
\[
 P_{I_T}\mathscr JZ\partial d_{\varepsilon}(g)=0.
\]
\end{theorem}

\begin{proof}
From $\mathscr Jg=\varepsilon g$ one gets
\[
 f_<(g)=\mathcal FZ^{-1}\mathscr Jg
       =\varepsilon\mathcal Ff,
\]
which proves the first formula in (325).  Since $\mathcal F^2=I$ on the
even additive source,
\[
 \mathcal F(f-\varepsilon\mathcal Ff)
 =\mathcal Ff-\varepsilon f
 =-\varepsilon(f-\varepsilon\mathcal Ff).
\]
Next,
\[
 \mathscr JZ\mathcal Fd_{\varepsilon}
 =\mathscr JZ\mathcal Ff-\varepsilon\mathscr JZf
 =\mathscr JZ\mathcal Ff-g=-r_{\rm P}(g),
\]
because $\mathscr JZf=\mathscr Jg=\varepsilon g$.  Using
$\mathcal Fd_{\varepsilon}=-\varepsilon d_{\varepsilon}$ gives the first
line of (326).

For the second line, differentiate the parity identity
$\mathscr JZf=\varepsilon Zf$.  The rules
$\partial\mathscr J=-\mathscr J\partial$ and
$\partial Z=-CZ+Z\partial$ give
\[
 \mathscr JZ\partial f
 =\varepsilon S_{\rm ar}g-\varepsilon Z\partial f.
\]
Since $\mathcal Ff=\varepsilon(f-d_{\varepsilon})$, substitution in
$\mathcal W_{\rm P}f=Z\partial f+\mathscr JZ\partial\mathcal Ff$
gives
\[
 \mathcal W_{\rm P}f
 =S_{\rm ar}g-\varepsilon\mathscr JZ\partial d_{\varepsilon},
\]
which is (326).
\end{proof}

\begin{theorem}[Exact polar defect on Meyer's nuclear source --- \PROVED]
For $f\in\mathcal H_+$, ordinary Poisson summation gives
\begin{equation}
 Zf-\mathscr JZ\mathcal Ff
 =\frac12\left(x^{-1}\mathcal Ff(0)-f(0)\right).       \tag{327}
\end{equation}
If in addition $Zf\in\mathcal H_-$, then $f\in\mathcal H_\cap$ and both
coefficients in (327) vanish.  Hence a physical first-loss mode would be
excluded immediately once both its M\"obius half-lift $Z^{-1}g\in\mathcal
H_+$ and the required physical-to-$\mathcal H_-$ regularity of $g$ are
proved.
\end{theorem}

\begin{proof}
Apply the self-dual Poisson formula to the full even sums:
\[
 f(0)+2Zf(x)
 =x^{-1}\bigl(\mathcal Ff(0)+2Z\mathcal Ff(x^{-1})\bigr).
\]
Rearrangement is (327).  Meyer's closed-range theorem states that
$Zf\in\mathcal H_-$ for $f\in\mathcal H_+$ if and only if
$f\in\mathcal H_\cap$, whose defining conditions are
$f(0)=\mathcal Ff(0)=0$.  The last assertion follows from (305) and compact
Poisson-range rigidity.
\end{proof}

\begin{proposition}[Co-Poisson/Sonin range audit --- \AUDIT]
The co-Poisson construction does not, by itself, prove the missing inclusion
$Z^{-1}g\in\mathcal H_+$.  In Burnol's extended Sonin space $L_a$ for
$0<a<1$, the closed co-Poisson subspace is the orthogonal complement of the
evaluation vectors associated with the nontrivial zeros of $\zeta$, counted
with multiplicity.  Thus placing the physical zero mode in that subspace
already requires the zero cancellations which finite-energy rigidity would
use to close G.  The semilocal Sonin isomorphism may transport such a vector
after its range membership is established, but cannot establish that
membership from compact support alone.

Likewise, multiplicative smoothing cannot supply a universal repair.  If
$g$ is replaced by $\eta*g$, then its Mellin transform is replaced by
$E(s)G(s)$, where $E$ is entire of exponential type for compactly supported
$\eta$.  A nonzero such $E$ cannot vanish at all distinct zeta zeros: their
unconditional counting function is $\gg R\log R$, whereas the zero count of
an exponential-type entire function is $O(R)$.  Therefore smoothing cannot
remove all possible poles of $G/\zeta$ without being the zero operator.
\end{proposition}

\begin{theorem}[Finite-pole extended-S rigidity --- \PROVED]
Let $g$ be a compactly supported multiplicative distribution whose Mellin
transform $G$ is entire of exponential type and finite polynomial order on
horizontal strips.  Suppose that its M\"obius half-lift admits a tempered
extended-S realization $f$ in the sense that
\begin{equation}
 \mathcal Mf(s)=\frac{G(s)}{\zeta(s)}                 \tag{328}
\end{equation}
is a moderate meromorphic Mellin transform with only finitely many poles at
nontrivial zeros of $\zeta$.  In particular, Burnol's extended-S class,
whose polar part is finite, satisfies this hypothesis.  Then $g=0$.
\end{theorem}

\begin{proof}
At every nontrivial zero $\rho$ of $\zeta$ outside the finite pole set of
(328), holomorphy of $G/\zeta$ forces $G(\rho)=0$; multiplicities only
strengthen this conclusion.  Unconditionally, the number of distinct
nontrivial zeta zeros of height at most $R$ is $\gg R\log R$.  On the other
hand Jensen's formula gives $O(R)$ zeros for a nonzero entire function of
exponential type and finite strip order.  Therefore $G$ must vanish
identically.  Injectivity of the Mellin transform on compactly supported
distributions gives $g=0$.
\end{proof}

\begin{corollary}[Extended-S completion criterion for G --- \PROVED]
Condition $G$ follows if every first-loss zero mode has a M\"obius half-lift
$Z^{-1}\mathcal U_cE_TF$ with the extended-S realization in (328).  More
weakly, it is enough that the opposite Fourier component
$d_\varepsilon$ in (325) have finite polar data and that the complementary
$\varepsilon$ component have an ordinary Poisson realization with no
nontrivial-zeta polar data:
\begin{equation}
 \boxed{\text{$d_\varepsilon$ has finite nontrivial polar data and
 the complementary component has none.}}              \tag{329}
\end{equation}
\end{corollary}

\begin{proof}
The first assertion is immediate from Theorem (328): the compact source of a
first-loss mode would have to vanish.  For the second, finite polar data for
$d_\varepsilon$ leave only finitely many possible uncancelled residues at
nontrivial zeta zeros, while the ordinary Poisson component has none there.
Thus $G/\zeta$ has only finitely many nontrivial-zero poles and Theorem (328)
applies.  The primitive conditions separately remove the two Tate characters
at $0$ and $1$.
\end{proof}

\paragraph{Sharpened physical target.}
The preceding theorems reduce the completion problem to one parity sector:
prove that a first-loss zero mode has a M\"obius half-lift whose
$-\varepsilon$ Fourier component $d_{\varepsilon}$ belongs to a physical
boundary class in which its zero common Mellin germ and
$P_{I_T}\mathscr JZ\partial d_{\varepsilon}=0$ force
$d_{\varepsilon}=0$.  Merely assuming $Z^{-1}g\in\mathcal H_+$ would already
give the conclusion by (327), and is therefore retained as an equivalent
\OPEN{} range theorem rather than inserted as a regularity lemma.  The
finite-pole alternative (328)--(329) is weaker in topology: it suffices to
derive extended-S behavior, equivalently finite quasi-homogeneous polar data,
for the opposite Fourier component.

\begin{theorem}[Exact Gamma--Tate bi-gap reformulation --- \PROVED]
For $T>0$ define the source-typed jump--Sonin space
\begin{equation}
 \mathscr S_T^\xi:=\left\{h=L^2E_T\phi:\
 \begin{array}{l}
   \phi\in H_0^2(I_T),\\
   \operatorname{supp}h\subset[-T,T],\\
   \operatorname{supp}\mathfrak J_\xi h
        \subset\mathbb R\setminus I_T
 \end{array}\right\}.                               \tag{330}
\end{equation}
Here $\mathfrak J_\xi$ is the two-contour source jump (275), not the
cancelled scalar continuation (270).  If
\begin{equation}
 \Phi(s):=\langle E_T\phi,e^{(s-1/2)t}\rangle,
 \qquad H_h(s):=\langle h,e^{(s-1/2)t}\rangle,
\end{equation}
then every $h\in\mathscr S_T^\xi$ satisfies the exact Tate factorization
\begin{equation}
 \boxed{H_h(s)=s^2(s-1)^2\Phi(s).}                   \tag{331}
\end{equation}
In particular $H_h$ has zeros of order at least two at both Tate terminals
$s=0,1$, and the spectral description of its exterior carrier is
\begin{equation}
 \mathfrak J_\xi h
 =\sum_\rho m_\rho\,[\rho(\rho-1)]^2\Phi(\rho)
             e^{(\rho-1/2)t}.                       \tag{332}
\end{equation}
Finally, the Gamma--Tate uniqueness statement (260) is equivalent to
\begin{equation}
 \boxed{\mathscr S_T^\xi=\{0\}\quad\hbox{for every }T>0.} \tag{333}
\end{equation}
Consequently (333) implies G40 and condition $G$ through the first-loss
theorem, without a scalar multiplier interpretation of $\mathcal A_\infty$.
\end{theorem}

\begin{proof}
The clamped conditions imply
$L(E_T\phi)=E_T(L\phi)$ in distributions, so
$h=L^2E_T\phi$ is supported in $[-T,T]$.  Since
\[
 L e^{(s-1/2)t}=\bigl((s-\tfrac12)^2-\tfrac14\bigr)
                    e^{(s-1/2)t}
                 =s(s-1)e^{(s-1/2)t},
\]
distributional integration by parts twice gives (331).  Formula (332) is
(277) with (331) substituted.  Source identification (279a) gives
$\mathfrak J_\xi h=\mathcal A_\infty h$ because (331) annihilates both
Tate terminals.  Hence the last support condition in (330) is exactly
$P_{I_T}\mathcal A_\infty L^2E_T\phi=0$.

If (260) holds, any element of $\mathscr S_T^\xi$ has $\phi=0$ and hence
$h=0$.  Conversely, if (333) holds, a solution of (260) gives
$h\in\mathscr S_T^\xi$ and therefore $h=0$.  Thus
$L^2E_T\phi=0$.  Its transform and (331) give
$s^2(s-1)^2\Phi(s)=0$ identically; since the polynomial is nonzero,
$\Phi=0$, and injectivity of the transform gives $\phi=0$.  This proves the
equivalence and the stated first-loss consequence.
\end{proof}

\begin{proposition}[Why the classical Sonin gap is not yet the bi-gap ---
\AUDIT]
Under the critical logarithmic change of variables $x=e^t$, the observation
window $I_T=(-T,T)$ becomes the multiplicative annulus
\begin{equation}
             e^{-T}<x<e^T.                          \tag{334}
\end{equation}
The classical Sonin property $S$ instead requires a function and its
additive Fourier transform to vanish on a neighbourhood $(0,a)$ of the
additive origin.  Neither a single multiplicative dilation nor reflection
turns the annular gap (334) into $(0,a)$.  Therefore the support statement
in (330), or the carrier cancellation
$P_{I_T}\mathscr JZ\partial d_\varepsilon=0$, does not by itself place the
mode in a classical or extended Sonin space.  A valid Sonin completion must
first prove a propagation theorem from the annular bi-gap to an origin gap,
or construct directly a jump-adapted space $\mathscr S_T^\xi$.
\end{proposition}

\begin{proof}
The image of $(-T,T)$ under $t\mapsto e^t$ is exactly (334).  A
multiplicative dilation sends it to $(\lambda e^{-T},\lambda e^T)$, whose
left endpoint remains positive for every $\lambda>0$; reflection
$x\mapsto x^{-1}$ merely interchanges its endpoints.  Thus these operations
cannot produce an interval beginning at zero.  The support hypotheses in
the definitions of the two spaces are consequently different, and no
inclusion follows without an additional propagation statement.
\end{proof}

\paragraph{New reverse-engineering target.}
Equations (330)--(333) identify the smallest new object which could close
the proof: a jump-adapted Sonin theorem asserting that a compact source with
the double Tate divisor $s^2(s-1)^2$, whose \emph{two-contour} $\xi$-jump
has the complementary annular support, is zero.  The density of the zeta
zeros alone does not prove this: after scalar continuation the two contour
orientations cancel, and completeness of an overcomplete exponential
system does not imply uniqueness of its expansion coefficients.  The
required theorem must retain the two Hardy boundary orientations (or an
equivalent positive Hilbert defect) before taking their supertrace.  This
bi-gap theorem remains \OPEN{}.

\begin{theorem}[Translation-jet exclusion at first contact --- \PROVED]
Let $T_*$ be a first-loss window and let $F\in\ker H_{T_*}$, with $H_T$
as in (251).  If the zero extension of $F$ belongs to
$C_c^\infty(\mathbb R)$, then $F=0$.  More generally, the same conclusion
holds whenever all distributional derivatives $D^kE_{T_*}F$ belong to the
global form domain, are represented by $L^2$ functions supported in
$[-T_*,T_*]$, and differentiation of the global form along translations is
valid.  In short,
\begin{equation}
 \boxed{F\in\ker H_{T_*}\cap C_c^\infty(\mathbb R)
        \quad\Longrightarrow\quad F=0.}              \tag{335}
\end{equation}
\end{theorem}

\begin{proof}
Write
\[
 \mathfrak q(U,V):=\langle\mathcal A_\infty U,V\rangle
\]
for the polarized global source form.  On compactly supported smooth
vectors its arithmetic sum is finite and its Gamma part is a Fourier
multiplier form, so translation invariance gives
\begin{equation}
 \mathfrak q(\tau_cU,\tau_cV)=\mathfrak q(U,V).       \tag{336}
\end{equation}
Every derivative of a primitive compact source remains primitive, since
\[
 \langle D^kF,e^{\pm t/2}\rangle
   =(-1)^k(\pm\tfrac12)^k
       \langle F,e^{\pm t/2}\rangle=0.               \tag{337}
\]
At first contact $H_{T_*}\ge0$ and $F\in\ker H_{T_*}$.  Differentiate
$\mathfrak q(\tau_cF,\tau_cF)=\mathfrak q(F,F)=0$
twice at $c=0$ to obtain
\begin{equation}
 0=2\mathfrak q(DF,DF)
     +2\operatorname{Re}\mathfrak q(F,D^2F).         \tag{338}
\end{equation}
The last term is zero: $D^2F$ is primitive by (337), while the compressed
zero-mode equation says that $\mathfrak q(F,K)=0$ for every primitive form
vector $K$.  Thus $\mathfrak q(DF,DF)=0$.  Positivity of $H_{T_*}$ implies
$DF\in\ker H_{T_*}$.

Apply the same argument to $DF$, and iterate.  It follows that
\begin{equation}
              D^kF\in\ker H_{T_*}\qquad(k\ge0).     \tag{339}
\end{equation}
The kernel is finite-dimensional because $H_{T_*}$ has compact resolvent.
Hence the translation jet is linearly dependent: for some nonzero
polynomial $p$, $p(D)F=0$.  A distributional solution of a constant-
coefficient ODE is an exponential polynomial; a compactly supported one is
zero.  Therefore $F=0$.
\end{proof}

\begin{corollary}[Smooth-core regularity criterion for G --- \PROVED]
Condition $G$ follows if every first-loss zero mode is a smooth vector for
the translation group in the sense of (335).  Equivalently, any nonzero
first-loss mode must satisfy
\begin{equation}
 \boxed{\text{loss of translation-generator regularity at the moving
 boundary or at some finite derivative order.}}       \tag{340}
\end{equation}
\end{corollary}

\begin{proof}
The first assertion follows from Theorem (335) and the first-loss theorem.
The second is its contrapositive.
\end{proof}

\begin{proposition}[Logarithmic ellipticity audit --- \AUDIT]
The zero-mode equation does not presently supply the hypothesis of (335).
The Gamma symbol grows only as $\log(2+|\tau|)$; solving once for the Gamma
channel therefore yields logarithmic, not positive Sobolev, regularity.
Iterating such estimates can give membership in every power of a logarithmic
weight without implying $D^kF\in L^2$ for even one prescribed $k$.  Moreover
zero extension creates boundary distributions unless the corresponding
traces vanish.  Thus replacing this missing boundary bootstrap by the phrase
``Gamma ellipticity'' would be invalid.  A proof of either
\begin{equation}
 E_{T_*}F\in C_c^\infty(\mathbb R),
 \quad\hbox{or directly the translation-jet hypotheses in (335),}    \tag{341}
\end{equation}
from (253)/(268) would, however, close G non-circularly.
\end{proposition}

\begin{theorem}[Prime-oriented boundary collars --- \PROVED]
Let $F$ be a primitive zero mode on $I_T$, put $u=E_TF$, and set
\begin{equation}
 C_+(T):=(T-\log2,T),\qquad
 C_-(T):=(-T,-T+\log2).                              \tag{342}
\end{equation}
Then the two-sided zero-mode equation (253) becomes the pair of one-way
Gamma--renewal equations
\begin{align}
 (G_\Gamma-m_0)u(t)
 &=\alpha e^{-t/2}+\beta e^{t/2}
   +\sum_{n\ge2}\frac{\Lambda(n)}{\sqrt n}
                         u(t-\log n),
       &&t\in C_+(T),                                \tag{343}\\
 (G_\Gamma-m_0)u(t)
 &=\alpha e^{-t/2}+\beta e^{t/2}
   +\sum_{n\ge2}\frac{\Lambda(n)}{\sqrt n}
                         u(t+\log n),
       &&t\in C_-(T).                                \tag{344}
\end{align}
Both sums are locally finite.  If $F(-t)=\varepsilon F(t)$, the two collar
equations are carried into one another by reflection.  Thus every
first-contact obstruction to the smooth-core criterion (341) is a pair of
reflected one-sided Gamma boundary layers; no outgoing prime orientation is
present in either collar.
\end{theorem}

\begin{proof}
If $t\in C_+(T)$ and $n\ge2$, then
$t+\log n>T$, hence $u(t+\log n)=0$.  Substitution in (253), with the sign
of the arithmetic part fixed by (256), gives (343).  Likewise, for
$t\in C_-(T)$ one has $t-\log n<-T$, which gives (344).  For fixed $t$ only
those $n$ for which the remaining argument lies in $[-T,T]$ contribute, so
the sums are finite.  Reflection interchanges the two translations, the two
collars and the two Tate characters; (254)--(255) give the final assertion.
\end{proof}

\begin{corollary}[Boundary-layer completion target --- \PROVED]
Condition $G$ follows if the one-way problems (343)--(344), together with
the two Tate moments and first-contact positivity, imply that the zero
extension is a smooth translation vector.  More economically, it is enough
to prove directly
\begin{equation}
 \boxed{\text{both collar boundary layers belong to the common
 translation-jet form domain in (335).}}             \tag{345}
\end{equation}
\end{corollary}

\begin{proof}
Either conclusion supplies (335); translation-jet exclusion then rules out
the first-loss mode, and the first-loss theorem gives $G$.
\end{proof}

\paragraph{Boundary-layer audit.}
Equations (343)--(344) do not themselves imply flat boundary traces.  The
Gamma channel has logarithmic order, and zero-exterior Dirichlet problems
for logarithmic Laplacians have genuine, optimally weak boundary layers.
The inward prime sums in (343)--(344) are smooth only after suitable
interior regularity is proved and in any case act as forcing terms.  The
missing G40 bridge can therefore be stated as a positive matching theorem
for these two reflected Gamma boundary layers, with the affine Tate square
(109) and the prime Julia storages supplying their terminal energy.  This
matching theorem remains \OPEN{}.

\begin{proposition}[Optimal logarithmic boundary scale does not supply the
translation jet --- \AUDIT]
Put, for $0<r<1/10$,
\begin{equation}
 \ell(r)=\frac1{|\log r|},
 \qquad b(r)=\sqrt{\ell(r)}.                         \tag{418}
\end{equation}
Then $b\in L^2(0,1/10)$ but $b\notin H^1(0,1/10)$, since
\[
 |b'(r)|^2=\frac1{4r^2|\log r|^3}
 \quad\hbox{and}\quad
 \int_0^{1/10}\frac{dr}{r^2|\log r|^3}=+\infty.
\]
The scale $b(\operatorname{dist}(t,\partial I_T))$ is the optimal generic
boundary scale for bounded-data Dirichlet problems for the logarithmic
Laplacian (Hern\'andez--Santamar\'ia, L\'opez R\'ios and Salda\~na,
\texttt{arXiv:2401.18033}).  Hence standard logarithmic boundary
regularity, even at its optimal exponent, cannot imply the first
translation derivative required in (335).  A successful boundary bootstrap
for a first-loss mode must exploit an additional Gamma--Euler equality-state
cancellation; it cannot be cited from the generic Dirichlet theory.
\end{proposition}

\begin{proof}
The displayed derivative calculation proves the Sobolev assertion.  The
cited optimal-boundary theorem applies to the logarithmic Laplacian, while
$g_\Gamma$ has the same logarithmic high-frequency order and the interior
arithmetic part consists of finitely many bounded shifts.  This explains the
relevance of the scale but is not a claimed perturbation theorem for the
present sign-changing system, nor does it assert that every zero mode
attains the barrier.  The logical
conclusion used here is only negative: generic boundary regularity at that
scale is insufficient for (335), so any stronger conclusion needs a new
equality-state argument.
\end{proof}

\paragraph{Support audit.}
Compact support of $g$ does not automatically imply (316): the second
Poisson half $\mathscr JZ\mathcal FZ^{-1}g$ can have a noncompact incoming
tail.  Thus (316) must be supplied by the Sonin cutoff/range construction or
replaced by an $L^2$ incoming boundary condition strong enough to run the
same Volterra argument.  It is retained as \OPEN{}.

\paragraph{Audit consequence.}
The lateral differentiated-Poisson-defect idea does not by itself close G:
it exposes (314), the same typed observation which earlier appeared as
(56)--(61), now required only on the zero-mode space and only with enough
norm control for causal uniqueness.  This is weaker than an all-source
isometry but stronger than equality of scalar generators.

Equation (257) is the precise Poisson--Sonin equality case still to exclude.
It is stronger and better typed than the withdrawn fixed-vector shortcut:
both the interior projection and the global distributional operator are
explicit.  A proof that the only clamped potential satisfying (257) is zero
would make the first-loss theorem contradict itself and establish (249) for
all windows.

This replaces the unspecified output symbol in (228) by an exact physical
source formula.  The four-state Gamma ODE realization and the prime Julia
resolvents specify every row of (247).  What remains is solely to prove its
passivity (249) by interconnecting those source systems; defining a defect
operator from the left side of (248) would still be circular.

This probabilistic boundary realizes the positive metric
$\mathfrak M_t$ of (117) before critical continuation: the finite Euler
Poisson factors supply the Zeta modulus, while (172) supplies the Gamma
modulus.  The remaining constant/polar normalization and the two Tate
terminal motions are already explicit in (105)--(111).  The final G40 task
is therefore reduced to showing that the physical clamped-potential
observation is the conditional-expectation martingale associated with this
product boundary; then (157) and orthogonal innovations would give the
required unit contraction.

Thus the two ``explicit Tate terminals'' in (99) are no longer unspecified:
their canonical combined square is (109).  What remains in (104) is the
continuum Gamma--Poisson bulk comparison after subtracting this finite-rank
second-fundamental-form contribution.  The positivity of (109) alone cannot
fix the sign of that bulk term.

\section{Panoramic negative certificate and change of attack}

The preceding local constructions can now be viewed from the opposite
direction.  Rather than adding another candidate positive port, assume that
condition G fails and record everything which the first failing state would
have to violate simultaneously.  This produces a single boundary object,
not a new normalization problem.

\begin{theorem}[Complete negative certificate for failure of G --- \PROVED]
If condition G fails, then there are $T_*>\log2$ and a nonzero
$F_*\in\mathcal P_{T_*}$, which may be chosen with a fixed parity, such that
all of the following hold:
\begin{align}
 P_{T_*}A_{T_*}F_*&=0,\qquad
 A_{T_*}F_*=\alpha e^{-t/2}+\beta e^{t/2}
       \quad(t\in I_{T_*}),                         \tag{346}\\
 \operatorname{diam}(\operatorname{ess,supp}F_*)&=2T_*,              \tag{347}\\
 E_{T_*}F_*&\notin\bigcap_{k\ge0}\operatorname{Dom}D^k
       \quad\hbox{in the translation-form sense of (335)},          \tag{348}\\
 \nexists f\in\mathcal H_\cap\quad Zf&=\mathcal U_cE_{T_*}F_*.       \tag{349}
\end{align}
Moreover the exterior Gamma contribution belongs to $L^2$, but at least one
of the two oriented arithmetic tails does not:
\begin{equation}
 1_{\mathbb R\setminus I_{T_*}}G_\Gamma E_{T_*}F_*\in L^2,
 \qquad
 Y_{+,T_*}^{\rm out}\notin L^2
 \quad\hbox{or}\quad
 Y_{-,T_*}^{\rm out}\notin L^2.                    \tag{350}
\end{equation}
Equivalently, at least one of the two reflected expressions (286) fails its
Hardy $H^2$ condition.  In Mellin coordinates the obstruction is therefore
exactly one of the following two mutually exhaustive phenomena:
\begin{enumerate}
\item $K_{F_*}(z)/\zeta(\tfrac12+z)$ has a nonremovable polar or boundary-pole
      contribution at a nontrivial zeta zero in one causal orientation;
\item all such local singularities cancel in that orientation, but the
      resulting holomorphic quotient fails the uniform Hardy boundary norm
      (287), equivalently it carries non-$L^2$ incoming/outgoing history.
\end{enumerate}
Thus a counterexample cannot be located in the Gamma normalization, in the
two Tate moments, in G39, or in a finite collection of prime channels.  It is
necessarily a nonsmooth, support-saturating, globally causal completion
obstruction.
\end{theorem}

\begin{proof}
The first-loss theorem and parity reduction give (346).  Support saturation
is (273), proving (347).  If (348) failed, translation-jet exclusion (335)
would give $F_*=0$.  If a lift in (349) existed, the equality-mode Poisson
lift theorem (299)--(302) would give the exterior estimate (280), and
the finite-energy closure criterion would give $F_*=0$.

The exterior Gamma estimate (281) holds for every compact-window $L^2$
source.  If both arithmetic tails were in $L^2$, then the constant channel,
which has compact support, and the Gamma tail would show that
$\mathcal A_\infty E_{T_*}F_*\in L^2$.  The finite-energy closure criterion
would force $F_*=0$.  Hence (350).  The exact Laplace/Hardy
equivalence (283)--(287) gives the final reformulation.  A failed Hardy
membership is caused either by a nonremovable interior/boundary pole or, once
all poles are removable, by failure of the uniform Hardy norm; these cases
are exhaustive.  Finally the Tate mean cancels exactly by (288)--(291), so
the non-$L^2$ part is the Chebyshev fluctuation rather than a missing polar
normalization.
\end{proof}

\begin{proposition}[Why interval shape differentiation cannot close G by
itself --- \PROVED]
Let $\lambda_1(T)$ be the lowest primitive eigenvalue of $H_T$.  Then
$\lambda_1(T)$ is nonincreasing.  Consequently, at a first contact
$\lambda_1(T_*)=0$, every left or right Dini derivative which exists is
nonpositive.  A Hadamard or Riccati formula for motion of the endpoints can
therefore describe the boundary leakage in (350), but its sign alone cannot
exclude first contact.
\end{proposition}

\begin{proof}
Zero extension maps $\mathcal P_T$ isometrically into
$\mathcal P_{T'}$ for $T<T'$ and preserves the global quadratic form.
The min--max principle gives $\lambda_1(T')\le\lambda_1(T)$.  The assertion
about Dini derivatives follows.  Thus a positive transversality sign would
contradict domain nesting; any successful shape argument must contain an
additional rigidity theorem which already rules out the kernel vector.
\end{proof}

\begin{theorem}[Reciprocal Sonin gap amplification --- \PROVED]
Let $L_a$ be Burnol's extended cosine--Sonin space: $u\in L^2(0,\infty)$
belongs to $L_a$ when both $u$ and its cosine transform are constant on
$(0,a)$.  For $0<a<1$, let $P_a\subset L_a$ be the closed co-Poisson
subspace generated by sources supported in $[a,a^{-1}]$.  Then
\begin{equation}
 \boxed{P_a\cap L_{a^{-1}}=\{0\}.}                  \tag{351}
\end{equation}
This statement is unconditional and makes no assumption about the location
of the zeta zeros.
\end{theorem}

\begin{proof}
Burnol's co-Poisson orthogonality theorem identifies $P_a$ with the
orthogonal complement in $L_a$ of all evaluation vectors
$Y^a_{\rho,k}$ attached to the nontrivial zeros $\rho$ of $\zeta$, with
$0\le k<m_\rho$.  Hence, if $u\in P_a$, its completed Mellin transform and
the required derivatives vanish at every $\rho$.

The Sonin spaces form a decreasing chain, so
$L_{a^{-1}}\subset L_a$.  Evaluating the completed Mellin transform is
intrinsic and independent of which member of the chain contains $u$.
More precisely, the evaluator in the smaller space is the orthogonal
projection of the evaluator in the larger space:
\[
 Y^{a^{-1}}_{\rho,k}
 =P_{L_{a^{-1}}}Y^a_{\rho,k}.
\]
Consequently, for $u\in L_{a^{-1}}$,
$\langle u,Y^{a^{-1}}_{\rho,k}\rangle
=\langle u,Y^a_{\rho,k}\rangle$.  Therefore an element of
$P_a\cap L_{a^{-1}}$ is orthogonal, now in $L_{a^{-1}}$, to every
$Y^{a^{-1}}_{\rho,k}$.  Since $a^{-1}>1$, Burnol's
completeness theorem says that these evaluation vectors span
$L_{a^{-1}}$.  Thus $u=0$.  The projection statement is part of the
evaluator construction preceding Proposition~2.2, and the
orthogonality/completeness assertions are Theorem~3.1 of J.-F. Burnol,
\emph{J. Th\'eor. Nombres Bordeaux} 16 (2004), 65--94;
their proofs use co-Poisson summation, Nevanlinna--Krein theory and
Paley--Wiener, not RH.
\end{proof}

\begin{theorem}[Canonical primitive co-Poisson observation --- \PROVED]
Let $F\in C_c^\infty(I_T)\cap\mathcal P_T$, put
\begin{equation}
 g_F(x):=x^{-1/2}F(\log x),\qquad x>0,                \tag{352}
\end{equation}
and define
\begin{equation}
 (\mathfrak S_T^0F)(x)
 :=\sum_{n\ge1}\frac1n g_F(x/n).                    \tag{353}
\end{equation}
Then $\mathfrak S_T^0F\in K_{e^{-T}}\cap P_{e^{-T}}$, where $K_a$
is the cosine--Sonin space.  More precisely, with $\mathcal F_+$ denoting
the cosine transform,
\begin{align}
 \mathfrak S_T^0F(x)&=0 &&(0<x<e^{-T}),\notag\\
 \mathcal F_+\mathfrak S_T^0F(x)
 &=\frac1x\sum_{n\ge1}g_F(n/x)=0 &&(0<x<e^{-T}).      \tag{354}
\end{align}
The map $F\mapsto\mathfrak S_T^0F$ is injective.  If
$F(-t)=\varepsilon F(t)$, $\varepsilon\in\{1,-1\}$, then
\begin{equation}
 \mathcal F_+\mathfrak S_T^0F
 =\varepsilon\mathfrak S_T^0F.                       \tag{355}
\end{equation}
Thus, in a fixed parity sector, reciprocal gap amplification is a
one-component problem.
\end{theorem}

\begin{proof}
The right Mellin moments used in co-Poisson summation are
\[
 \widehat g_F(1)=\int_0^\infty g_F(x)x^{-1}\,dx=M_-F,
 \qquad
 \widehat g_F(0)=\int_0^\infty g_F(x)\,dx=M_+F,
\]
so both constants in the co-Poisson formula vanish.  The support of $g_F$
is contained in $[e^{-T},e^T]$.  If $x<e^{-T}$, every $x/n$ lies below
that support and every $n/x$ lies above it.  Burnol's co-Poisson identity
therefore gives (354).  Smooth compact support makes the completed Mellin
transform rapidly decreasing up to the standard polynomial zeta factor, so
both sides belong to $L^2$ and $\mathfrak S_T^0F\in K_{e^{-T}}$.  Since
(355) is itself a co-Poisson sum of a generator supported in
$[e^{-T},e^T]$, it also belongs to $P_{e^{-T}}$.

Writing $G_F(s)=\int_0^\infty g_F(x)x^{-s}\,dx$, Mellin transformation of
(355) gives
\begin{equation}
 \mathcal M(\mathfrak S_T^0F)(s)=\zeta(s)G_F(s)       \tag{356}
\end{equation}
in an initial half-plane.  If the observation vanishes, (356) and the
identity theorem give $G_F=0$, hence $g_F=F=0$.  Finally inversion parity is
$x^{-1}g_F(1/x)=\varepsilon g_F(x)$.  With $y=x/n$, it gives directly
\[
 \frac1x g_F(n/x)=\frac1x g_F(1/y)
 =\frac{\varepsilon}{n}g_F(x/n).
\]
Summation and (354) prove (355).
\end{proof}

\begin{corollary}[Amplified-gap completion criterion for G --- \PROVED]
Put
\begin{equation}
 a_T=e^{-T},\qquad b_T=e^T=a_T^{-1}.                 \tag{357}
\end{equation}
Condition G follows if the canonical observation (353) extends to every
first-loss state and satisfies the single gap-amplification implication
\begin{equation}
 \boxed{P_TA_TF=0\Longrightarrow
        \mathfrak S_T^0F\in L_{b_T}.}                \tag{358}
\end{equation}
For a parity representative it is enough to prove
$\mathfrak S_T^0F=0$ on $(a_T,b_T)$, because (354)--(355) then place both
the observation and its Fourier transform in $K_{b_T}\subset L_{b_T}$.
Thus (358) is precisely the assertion that the zero-mode Gamma term extends
the natural Tate/co-Poisson gap $(0,a_T)$ across the annulus
$(a_T,b_T)$.
\end{corollary}

\begin{proof}
For a first-loss state, Theorem (352)--(356) and (358) place
$\mathfrak S_T^0F$ in
$P_{a_T}\cap L_{a_T^{-1}}$.  Theorem (351) makes this vector zero, and the
injectivity in Theorem (352)--(356) gives $F=0$, contradicting first loss.
The first-loss theorem then gives condition G.
\end{proof}

\paragraph{Why this is a genuine change of geometry --- \AUDIT.}
The direct logarithmic image of $I_T$ is only the annulus
$(e^{-T},e^T)$, as observed in (334); that fact alone gives no classical
Sonin gap.  Theorem (351) uses information which the direct coordinate
change discarded: a compact co-Poisson generator with vanishing Tate
constants already supplies the adjacent origin gap $(0,e^{-T})$.  The
zero-mode equation is asked to amplify that existing gap up to $e^T$, not to
create an origin gap from an annulus by dilation.  Since $e^T>1$ for every
$T>0$, the amplified vector lands on the complete side of Burnol's sharp
Sonin threshold.

The remaining construction in (358) is narrower than the previous all-source
Hardy estimate.  Theorem (352)--(356) fixes the observation, proves its
natural gap, Fourier parity and injectivity independently of G.  No metric is
left to choose.  What remains is only to derive its enlarged gap from (346).
The required calculation must retain the two unsupertraced Gamma--Poisson
orientations: scalar meromorphic cancellation would make their completed
generators zero everywhere and would not identify the physical observation
on the annulus.

\begin{proposition}[Generator mismatch of the bare co-Poisson observation ---
\PROVED]
Let $G_F(s)$ be as in (356).  The Mellin transform of the canonical
observation is $\zeta(s)G_F(s)$, and therefore
\begin{equation}
 \partial_s\bigl(\zeta(s)G_F(s)\bigr)
 =\zeta'(s)G_F(s)+\zeta(s)G_F'(s).                  \tag{359}
\end{equation}
By contrast, the two arithmetic orientations in the zero-mode generator
have common meromorphic germ
\begin{equation}
 -\left(
   \frac{\zeta'}{\zeta}(s)
   +\frac{\zeta'}{\zeta}(1-s)
  \right)G_F(s),                                    \tag{360}
\end{equation}
as in (321), and the physical operator also contains the affine Gamma and
Tate terms.  Hence neither differentiation in $s$ nor the spatial dilation
generator turns (356) into the zero-mode operator.  In particular, (358)
does not follow by formally integrating $P_TA_TF=0$.
\end{proposition}

\begin{proof}
Equation (359) is the product rule applied to (356).  Equation (360) is the
Mellin calculation (321), with the same compact source.  The first expression
contains $\zeta G_F'$, while the second contains the reflected logarithmic
derivative and no $G_F'$; they cannot be equal as source operators.  The
logarithmic spatial derivative multiplies a Mellin transform by the Mellin
variable and likewise cannot repair this mismatch.  Finally (208)--(213)
show that the missing Gamma and Tate contributions belong to the covariant
derivative of the \emph{completed two-port metric}, not to (359).
\end{proof}

\paragraph{Correct reciprocal-gap orbit --- \AUDIT.}
The reciprocal theorem (351) remains a valid endpoint rigidity theorem, but
the bare observation (353) is only its endpoint candidate.  The physical
interpolating object must be the two-port completed Euler--Gamma orbit whose
positive precritical metric is
\[
 \mathfrak M_t(\tau)
 =\left|\Lambda_{\mathbb Q}
       \left(\frac t2+i\tau\right)\right|^2,
 \qquad t>2,
\]
and whose covariant logarithmic derivative at $t=1$, after localization, is
exactly $a_T$ by (211)--(213).  To prove (358) one must show that its
Kato-compressed, unsupertraced boundary value at equality agrees with
$\mathfrak S_T^0F$ (or with another explicitly injective element of
$P_{e^{-T}}$) and transports the natural gap to $e^T$.  This is the same
physical intertwining absent in (220j), now with a sharp reciprocal-Sonin
endpoint.  Positivity of $\mathfrak M_t$ before continuation does not by
itself prove monotonicity or this boundary identification.

\begin{theorem}[Prime-tower Blaschke--Clark scattering --- \PROVED]
For a prime $p$ put
\begin{equation}
 a_p:=\log p,\qquad r_p:=p^{-1/2},\qquad
 B_p(\tau):=
 \frac{e^{-ia_p\tau}-r_p}{1-r_pe^{-ia_p\tau}}.       \tag{361}
\end{equation}
Then $|B_p(\tau)|=1$ for real $\tau$, and its positive Clark/phase density is
\begin{align}
 d_p(\tau)
 &:=-\frac{d}{d\tau}\arg B_p(\tau)\notag\\
 &=a_p\frac{1-r_p^2}
             {1-2r_p\cos(a_p\tau)+r_p^2}             \tag{362}\\
 &=a_p+2\sum_{k\ge1}
       \frac{\log p}{p^{k/2}}\cos(k\tau\log p)>0.   \tag{363}
\end{align}
Let $\mathcal D_p$ be the positive global Fourier multiplier with symbol
$d_p$.  Zero-extension compression to $I_T$ gives the exact identity
\begin{equation}
 \boxed{
 E_T^*(a_pI-\mathcal D_p)E_T
 =-\sum_{\substack{k\ge1\\k\log p<2T}}
   \frac{\log p}{p^{k/2}}
       (S_{k\log p}+S_{-k\log p}).}                 \tag{364}
\end{equation}
Consequently, with
$\vartheta(e^{2T}):=\sum_{p<e^{2T}}\log p$, the complete localized
primitive operator has the Clark form
\begin{equation}
 \boxed{
 A_T=G_{\Gamma,T}
   +\bigl(\vartheta(e^{2T})-m_0\bigr)I
   -\sum_{p<e^{2T}}E_T^*\mathcal D_pE_T.}            \tag{365}
\end{equation}
Every prime power is contained in one inner scattering factor; no
prime-power port is inserted separately in (365).
\end{theorem}

\begin{proof}
The map $z\mapsto(z-r_p)/(1-r_pz)$ is a disk Blaschke factor.  On
$z=e^{-ia_p\tau}$ it is unimodular, and direct logarithmic differentiation
gives (362).  The Poisson-kernel expansion
\[
 \frac{1-r^2}{1-2r\cos\theta+r^2}
 =1+2\sum_{k\ge1}r^k\cos(k\theta)
\]
gives (363), including positivity.

The Fourier multiplier of
$S_b+S_{-b}$ is $2\cos(b\tau)$.  Hence (363) gives (364) before compression
with all $k\ge1$.  If $k\log p\ge2T$, the translated and original windows
are disjoint up to a null endpoint, so that compressed shift is zero.
Summing (364) over the finitely many primes $p<e^{2T}$ and adding
$G_{\Gamma,T}-m_0I$ gives (365).
\end{proof}

\begin{corollary}[Finite Clark embedding form of G40 --- \PROVED]
At a fixed window $T$, row (d) is equivalent to the source inequality
\begin{equation}
 \boxed{
 \sum_{p<e^{2T}}
   \langle\mathcal D_pE_TF,E_TF\rangle
 \le
 \langle G_\Gamma E_TF,E_TF\rangle
 +\bigl(\vartheta(e^{2T})-m_0\bigr)\|F\|^2,
 \quad F\in\mathcal P_T.}                           \tag{366}
\end{equation}
Thus the missing contraction is equivalently a finite Clark-measure
embedding of the primitive Paley--Wiener/Sonin source into the direct sum of
the prime-tower model spaces, dominated by the Gamma Clark reference plus
the explicit winding budget.
\end{corollary}

\begin{proof}
Pair (365) with $F\in\mathcal P_T$ and rearrange.  No square root of $A_T$
or sign assumption is used, so the equivalence is exact but not a proof of
the inequality.
\end{proof}

\paragraph{Clark-scattering scope --- \AUDIT.}
Equations (361)--(365) supply the missing order-zero inner object at each
\emph{prime tower}, which is stronger than the infinitesimal Julia weights
(221)--(223).  They do not yet prove (366): the winding term $a_pI$ is the
deterministic delay carried by the Blaschke factor and cannot be discarded.
The global bridge must match the direct sum of these finite Clark embeddings
with the Gamma inner/affine boundary and the moving Tate projection.  In
particular, treating (362) merely as a positive density and dropping its
winding constant would change the row-(d) operator.

\begin{theorem}[Global finite-window Gamma--Euler scattering --- \PROVED]
Put $a=1/4$ and define the Archimedean boundary phase
\begin{equation}
 B_\Gamma(\tau):=
 e^{-i\psi(a)\tau}
 \frac{\Gamma(a+i\tau/2)}{\Gamma(a-i\tau/2)}.        \tag{367}
\end{equation}
It is unimodular on the real line and
\begin{equation}
 \frac{d}{d\tau}\arg B_\Gamma(\tau)
 =\Re\psi(a+i\tau/2)-\psi(a)=g_\Gamma(\tau).        \tag{368}
\end{equation}
For a fixed window define the single boundary scattering
\begin{equation}
 \Theta_T(\tau):=
 e^{\,i(\vartheta(e^{2T})-m_0)\tau}
 B_\Gamma(\tau)
 \prod_{p<e^{2T}}B_p(\tau).                         \tag{369}
\end{equation}
Then $|\Theta_T(\tau)|=1$ and its phase derivative
\begin{equation}
 q_T(\tau):=\frac{d}{d\tau}\arg\Theta_T(\tau)
 =g_\Gamma(\tau)+\vartheta(e^{2T})-m_0
   -\sum_{p<e^{2T}}d_p(\tau)                        \tag{370}
\end{equation}
satisfies the exact compressed identity
\begin{equation}
 \boxed{A_T=E_T^*\mathcal F^{-1}M_{q_T}\mathcal FE_T.} \tag{371}
\end{equation}
Thus the complete row-(d) operator is the zero-extension compression of the
phase derivative of one explicit Gamma--Euler scattering.  No zeta zero,
choice of positive spectral part or assumed sign enters (367)--(371).
\end{theorem}

\begin{proof}
The quotient in (367) has modulus one for real $\tau$.  Logarithmic
differentiation gives
\[
 \frac{d}{d\tau}\log
 \frac{\Gamma(a+i\tau/2)}{\Gamma(a-i\tau/2)}
 =\frac i2\psi(a+i\tau/2)+\frac i2\psi(a-i\tau/2)
 =i\Re\psi(a+i\tau/2).
\]
The exponential contributes $-i\psi(a)$, proving (368).  Equations
(361)--(363) give
$\frac d{d\tau}\arg B_p=-d_p$.  Differentiating the finite product (369)
therefore proves (370).  Finally (365) is exactly the Fourier compression of
(370), which proves (371).
\end{proof}

\begin{corollary}[Single-scattering formulation of G40 --- \PROVED]
For every $T>0$, condition G at the window $T$ is equivalent to
\begin{equation}
 \boxed{
 \int_{\mathbb R}q_T(\tau)
       |\widehat{E_TF}(\tau)|^2\,\frac{d\tau}{2\pi}
 \ge0\qquad(F\in\mathcal P_T).}                    \tag{372}
\end{equation}
Equivalently, after removing the two Tate evaluation directions, the
compressed Clark form of $\Theta_T$ is positive.  This is the precise
phase-scattering theorem which a Poisson--Sonin/Hodge construction must now
prove.
\end{corollary}

\begin{proof}
Pair (371) with a primitive vector and use Plancherel.  The statement is an
exact reformulation, not a positivity proof.
\end{proof}

\begin{corollary}[Primitive Wigner--Smith formulation --- \PROVED]
The scalar Wigner--Smith delay of the scattering (369) is
\begin{equation}
 Q_T^{\rm WS}(\tau)
 :=-i\overline{\Theta_T(\tau)}\Theta_T'(\tau)
 =q_T(\tau).                                         \tag{373}
\end{equation}
Consequently
\begin{equation}
 \boxed{
 \langle A_TF,F\rangle
 =\left\langle
 M_{Q_T^{\rm WS}}\widehat{E_TF},
 \widehat{E_TF}\right\rangle_{L^2(d\tau/2\pi)},
 \qquad F\in\mathcal P_T.}                          \tag{374}
\end{equation}
Thus G40 is exactly nonnegative time delay of $\Theta_T$ on the
codimension-two Paley--Wiener subspace cut out by the Tate evaluations.
\end{corollary}

\begin{proof}
Since $|\Theta_T|=1$, differentiation of
$\Theta_T=e^{i\arg\Theta_T}$ gives (373).  Equation (374) is (371) and
Plancherel.
\end{proof}

\begin{theorem}[Lossless Julia realization and neutral prime birth ---
\PROVED]
For $r=r_p$ and $d=(1-r^2)^{1/2}$, the unitary Julia cell
\begin{equation}
 U_p:=\begin{pmatrix}r&d\\d&-r\end{pmatrix}          \tag{375}
\end{equation}
has scalar transfer function
\begin{equation}
 -r+z\,d(1-zr)^{-1}d
 =\frac{z-r}{1-rz}.                                  \tag{376}
\end{equation}
Hence $B_p$ is the boundary transfer of (375) at
$z=e^{-i\tau\log p}$.  Define the winding-corrected birth factor
\begin{equation}
 C_p(\tau):=e^{i\tau\log p}B_p(\tau).               \tag{377}
\end{equation}
Then
\begin{equation}
 \frac d{d\tau}\arg C_p(\tau)
 =a_p-d_p(\tau)
 =-2\sum_{k\ge1}\frac{\log p}{p^{k/2}}
                    \cos(k\tau\log p),             \tag{378}
\end{equation}
and the global scattering factors as
\begin{equation}
 \boxed{
 \Theta_T(\tau)=e^{-im_0\tau}B_\Gamma(\tau)
       \prod_{p<e^{2T}}C_p(\tau).}                  \tag{379}
\end{equation}
At the birth threshold $2T=\log p$, the zero-extension compression of the
Wigner--Smith delay (378) is zero.  Thus adjoining a new prime factor is a
lossless order-zero cascade whose signed physical Gram is neutral at birth;
this is scattering-level old compatibility.
\end{theorem}

\begin{proof}
The matrix in (375) is real orthogonal.  The standard one-state transfer
formula with $A=r$, $B=C=d$, $D=-r$ gives
\[
 D+zC(1-zA)^{-1}B
 =-r+\frac{z(1-r^2)}{1-zr}
 =\frac{z-r}{1-rz},
\]
which proves (376).  Equation (378) follows from (362)--(363), and (379)
is (369) with
$e^{i\vartheta(e^{2T})\tau}=\prod_{p<e^{2T}}e^{ia_p\tau}$.

If $2T=a_p$, then $k a_p\ge2T$ for every $k\ge1$.  Each compressed shift in
(364), equivalently every cosine term in the compressed form of (378), is
zero up to a null endpoint.  The new factor therefore changes the boundary
scattering but contributes zero signed physical Gram at its birth.  Above
threshold its complete tower enters continuously through the growing
overlaps, proving the stated compatibility.
\end{proof}

\begin{theorem}[Gamma all-pass cascade --- \PROVED]
Let $a=1/4$, $b_n=n+a$, and define
\begin{equation}
 C_{\Gamma,n}(\tau):=
 e^{i\tau/b_n}\frac{b_n-i\tau/2}{b_n+i\tau/2}.      \tag{380}
\end{equation}
The product converges locally uniformly on the real line and
\begin{equation}
 \boxed{B_\Gamma(\tau)=\prod_{n=0}^\infty
                  C_{\Gamma,n}(\tau).}              \tag{381}
\end{equation}
Every factor is unimodular and has nonnegative Wigner--Smith delay
\begin{equation}
 -i\overline{C_{\Gamma,n}(\tau)}
       C_{\Gamma,n}'(\tau)
 =\frac1{b_n}-\frac{b_n}{b_n^2+\tau^2/4}\ge0.       \tag{382}
\end{equation}
Consequently
\begin{equation}
 g_\Gamma(\tau)
 =\sum_{n=0}^\infty
 \left(\frac1{b_n}-\frac{b_n}{b_n^2+\tau^2/4}\right), \tag{383}
\end{equation}
and the Archimedean factor in the global scattering has a source-defined
lossless infinite all-pass realization with positive stored-delay density.
\end{theorem}

\begin{proof}
For bounded real $\tau$,
\[
 \log C_{\Gamma,n}(\tau)
 =\frac{i\tau}{b_n}
   -2i\arctan\frac{\tau}{2b_n}
 =O_\tau(b_n^{-3}),
\]
so the product converges locally uniformly.  Its value at $\tau=0$ is one.
Direct differentiation of its phase gives (382).  The standard digamma
series
\[
 \psi(a+z)-\psi(a)
 =\sum_{n=0}^\infty
 \left(\frac1{n+a}-\frac1{n+a+z}\right)
\]
with $z=i\tau/2$, followed by real parts, gives (383).  Therefore the
product in (381) and $B_\Gamma$ have identical logarithmic phase derivative
by (368), are unimodular and take the same value at zero.  Their quotient is
constant one, proving (381).  Each rational quotient in (380) is a
one-state continuous-time all-pass transfer, while the exponential is its
explicit winding correction; cascading the locally convergent factors gives
the final realization.
\end{proof}

\paragraph{Inner-function audit.}
Although every factor in (369) is unimodular on the boundary, it has not been
proved that the completed product belongs to a Schur/inner class whose Clark
form is positive after the Tate compression.  Boundary unimodularity alone
does not control the sign of its phase derivative, and declaring
$\Theta_T$ inner from (369) would assume (372).  The advantage of (369) is
instead that the missing theorem now concerns one explicit scattering and
not an unspecified interconnection of Gamma and prime-power ports.

\begin{proposition}[No scalar half-plane inner shortcut --- \PROVED]
For every window containing at least one prime, the scalar scattering
$\Theta_T$ in (369) is not an inner function of either open half-plane.
Indeed:
\begin{enumerate}
\item every $B_p$ has a pole at $\tau=i/2$, and $B_\Gamma$ has a pole there
      as well, so $\Theta_T$ is not analytic in the upper half-plane;
\item its boundary phase derivative satisfies
\begin{equation}
 q_T(\tau)=\log|\tau|+O_T(1)\qquad(|\tau|\to\infty),   \tag{384}
\end{equation}
whereas an inner function analytic in the lower half-plane with smooth
nonvanishing boundary values has nonpositive boundary phase derivative.
\end{enumerate}
Thus scalar inner-function positivity cannot prove (372).  Any passive
realization which closes G40 must retain at least the two oppositely oriented
Euler/Gamma ports (or an equivalent Pontryagin/Hodge state) before taking
their scalar boundary supertrace.
\end{proposition}

\begin{proof}
At $\tau=i/2$ one has
$e^{-i\tau\log p}=p^{1/2}$, so the denominator
$1-p^{-1/2}e^{-i\tau\log p}$ in (361) vanishes while its numerator does not.
Also $a+i\tau/2=0$ for $a=1/4$, so the numerator Gamma factor in (367) has
a pole.  The remaining factors and exponentials cannot cancel these poles;
this proves the first assertion.

For real $\tau$, each finite prime density $d_p$ is bounded, while the
digamma asymptotic gives
$g_\Gamma(\tau)=\log|\tau|+O(1)$.  Equation (370) proves (384).
For a scalar inner function in the lower half-plane, the canonical
Blaschke--singular factorization gives a nonpositive boundary phase
derivative wherever it is smooth.  The eventually positive function (384)
violates this necessary condition, proving the second assertion.
\end{proof}

\begin{theorem}[Operator-valued Clark two-port colligation --- \PROVED]
Assume $T>\log2$, let $\mathcal P(T)=\{p:p<e^{2T}\}$, and write
$\mathcal Q_T=\operatorname{Dom}G_{\Gamma,T}^{1/2}$ for the common closed
form domain.  Put
\begin{equation}
 h_n(\tau):=\frac1{n+\frac14}
 -\frac{n+\frac14}{(n+\frac14)^2+\tau^2/4},
 \qquad n\ge0.                                       \tag{385}
\end{equation}
For $F\in\mathcal Q_T$ define the positive analysis ports, in Fourier coordinates,
\begin{align}
 \mathcal X_T^{\rm Cl}F
 &:=
 \left(
   \{h_n^{1/2}\widehat{E_TF}\}_{n\ge0},
   \{\sqrt{\log p}\,\widehat{E_TF}\}_{p\in\mathcal P(T)}
 \right),                                           \tag{386}\\
 \mathcal Y_T^{\rm Cl}F
 &:=
 \left(
   \{d_p^{1/2}\widehat{E_TF}\}_{p\in\mathcal P(T)},
   \sqrt{m_0}\,\widehat{E_TF}
 \right).                                           \tag{387}
\end{align}
Their targets carry the ordinary direct-sum $L^2(d\tau/2\pi)$ norm.  They
are source-defined and satisfy, in the sense of closed quadratic forms,
\begin{align}
 (\mathcal X_T^{\rm Cl})^*\mathcal X_T^{\rm Cl}
 &=G_{\Gamma,T}+\vartheta(e^{2T})I,\notag\\
 (\mathcal Y_T^{\rm Cl})^*\mathcal Y_T^{\rm Cl}
 &=m_0I+\sum_{p<e^{2T}}E_T^*\mathcal D_pE_T,         \tag{388}\\
 \boxed{A_T
 &=(\mathcal X_T^{\rm Cl})^*\mathcal X_T^{\rm Cl}
  -(\mathcal Y_T^{\rm Cl})^*\mathcal Y_T^{\rm Cl}.} \tag{389}
\end{align}
The canonical source angular rule
\begin{equation}
 \mathfrak C_T^{\rm Cl}
   (\mathcal X_T^{\rm Cl}F):=\mathcal Y_T^{\rm Cl}F,
 \qquad F\in\mathcal P_T\cap\mathcal Q_T,             \tag{390}
\end{equation}
is well defined on
$\mathcal X_T^{\rm Cl}(\mathcal P_T\cap\mathcal Q_T)$, and
\begin{equation}
 \|\mathfrak C_T^{\rm Cl}\|\le1
 \quad\Longleftrightarrow\quad
 P_TA_TP_T\ge0.                                     \tag{391}
\end{equation}
Thus (391) is exactly G40.
\end{theorem}

\begin{proof}
Equations (382)--(383) give
$\sum_{n\ge0}h_n=g_\Gamma$ pointwise, with nonnegative summands.
Equations (362)--(363) give $d_p>0$.  Tonelli and Plancherel therefore prove
the first two Gram identities in (388), and (365) gives (389).

Since $T>\log2$, the set $\mathcal P(T)$ is nonempty.  The scalar winding
coordinates in (386) give
\[
 \|\mathcal X_T^{\rm Cl}F\|^2
 \ge\vartheta(e^{2T})\|F\|^2,
\]
so $\mathcal X_T^{\rm Cl}$ is injective and (390) is well defined.
Finally (389) gives
\[
 \|\mathcal X_T^{\rm Cl}F\|^2
 -\|\mathcal Y_T^{\rm Cl}F\|^2
 =\langle A_TF,F\rangle.
\]
Thus (390) is contractive exactly when the closed form of $A_T$ is
nonnegative on $\mathcal P_T\cap\mathcal Q_T$, which is
$P_TA_TP_T\ge0$ and hence G40.
\end{proof}

\begin{corollary}[Neutral old-compatible birth in the Clark colligation ---
\PROVED]
When $T$ crosses $T_p=\frac12\log p$, the two-port colligation gains the
positive winding coordinate
$\sqrt{\log p}\,\widehat{E_TF}$ in $\mathcal X_T^{\rm Cl}$ and the Clark
coordinate $d_p^{1/2}\widehat{E_TF}$ in $\mathcal Y_T^{\rm Cl}$.  Their
signed Gram increment is
\begin{equation}
 E_T^*(\log p\,I-\mathcal D_p)E_T,                   \tag{392}
\end{equation}
which is zero at $T=T_p$ and becomes exactly the full active $p^k$ tower
above threshold.  Hence (386)--(390) preserve the old physical defect and
adjoin only a lossless neutral pair at birth.
\end{corollary}

\begin{proof}
This is (364) and the threshold argument following (379), now separated
into its two positive Hilbert ports.
\end{proof}

\begin{proposition}[Clark versus physical delay ports --- \AUDIT]
Let $\mathcal X_T,\mathcal Y_T$ be the authoritative physical factors in
(83), and for each prime put
\begin{equation}
 \mathcal K_{p,T}:=
 \left(\log p-\sum_{\substack{k\ge1\\k\log p<2T}}
       \frac{\log p}{p^{k/2}}\right)I
 +\frac12\sum_{\substack{k\ge1\\k\log p<2T}}
       \frac{\log p}{p^{k/2}}
       (S_{k\log p}+S_{-k\log p}).                   \tag{393}
\end{equation}
Then, as quadratic forms,
\begin{align}
 (\mathcal X_T^{\rm Cl})^*\mathcal X_T^{\rm Cl}
       -\mathcal X_T^*\mathcal X_T
 &=\sum_{p<e^{2T}}\mathcal K_{p,T},\notag\\
 (\mathcal Y_T^{\rm Cl})^*\mathcal Y_T^{\rm Cl}
       -\mathcal Y_T^*\mathcal Y_T
 &=\sum_{p<e^{2T}}\mathcal K_{p,T}.                  \tag{394}
\end{align}
The common correction cancels from the signed defect, explaining (389), but
it is not positive for all windows.  Hence the Clark two-port cannot be
identified with the physical $J_-/J_+$ two-port on the full source by
adjoining a common Hilbert channel.  A primitive-only augmentation would
require a separate proof that the compression of (393) is positive and is
not supplied by this calculation.
\end{proposition}

\begin{proof}
For $b=k\log p$,
\[
 J_{b,-}^*J_{b,-}=I-\tfrac12(S_b+S_{-b}),\qquad
 J_{b,+}^*J_{b,+}=I+\tfrac12(S_b+S_{-b}).
\]
Subtracting the physical minus Gram from the winding Gram $\log p\,I$ gives
(393).  Expanding $\mathcal D_p$ by (363) and subtracting the physical plus
Gram gives the same expression, proving (394).

The full Fourier symbol of the $p=2$ correction, when all powers are active,
at $\tau=\pi/\log2$ tends to
\[
 \log2\left(1-\sum_{k\ge1}2^{-k/2}
                    (1-(-1)^k)\right)
 =\log2\,(1-2\sqrt2)<0.
\]
For sufficiently large finite windows the corresponding finite sum remains
negative.  Long modulated compactly supported test functions concentrate
their Fourier mass near that frequency, so the compressed form is then
negative on some vector.  Thus no window-uniform positive common
augmentation exists.
\end{proof}

\paragraph{Typing consequence.}
The Clark colligation proves a new exact angular realization and preserves
the signed old defect, but it does not finish the physical conservation
identity (220j).  The remaining operator-valued Poisson--Sonin/Hodge
intertwiner must transport the physical delay factors to the Clark/all-pass
ports while also accounting for the indefinite common form (393), rather
than silently treating it as an orthogonal Hilbert summand.

\begin{theorem}[Explicit Krein--Hodge stabilization of the physical ports ---
\PROVED]
For $p<e^{2T}$ put $w_{p,k}=(\log p)p^{-k/2}$ and define the signed analysis
state
\begin{equation}
 \mathcal Z_{p,T}F:=
 \left(
   \sqrt{\log p}\,F,\,
   \{\sqrt{w_{p,k}}J_{k\log p,-}F\}_{k\log p<2T}
 \right)                                             \tag{395}
\end{equation}
in the Krein space with fundamental symmetry
$J_{p,T}=I\oplus(-I)$.  Then
\begin{equation}
 \boxed{\mathcal Z_{p,T}^*J_{p,T}\mathcal Z_{p,T}
       =\mathcal K_{p,T}.}                           \tag{396}
\end{equation}
Let $\mathcal Z_T=\bigoplus_{p<e^{2T}}\mathcal Z_{p,T}$ and
$J_T^K=\bigoplus_pJ_{p,T}$.  Stabilize both authoritative physical ports by
\begin{equation}
 \widetilde{\mathcal X}_TF:=
      (\mathcal X_TF,\mathcal Z_TF),\qquad
 \widetilde{\mathcal Y}_TF:=
      (\mathcal Y_TF,\mathcal Z_TF),                 \tag{397}
\end{equation}
with ambient symmetries $I\oplus J_T^K$.  Their pullback Grams are positive
and exactly the Clark Grams:
\begin{align}
 \widetilde{\mathcal X}_T^*(I\oplus J_T^K)
       \widetilde{\mathcal X}_T
 &=(\mathcal X_T^{\rm Cl})^*\mathcal X_T^{\rm Cl},\notag\\
 \widetilde{\mathcal Y}_T^*(I\oplus J_T^K)
       \widetilde{\mathcal Y}_T
 &=(\mathcal Y_T^{\rm Cl})^*\mathcal Y_T^{\rm Cl}.   \tag{398}
\end{align}
Consequently the source-generated stabilized ranges are positive Hilbert
subspaces of their ambient Krein spaces, and the rules
\begin{align}
 U_{X,T}\widetilde{\mathcal X}_TF&:=\mathcal X_T^{\rm Cl}F,\notag\\
 U_{Y,T}\widetilde{\mathcal Y}_TF&:=\mathcal Y_T^{\rm Cl}F              \tag{399}
\end{align}
extend to isometries of their generated Hilbert completions.  They intertwine
the stabilized physical angular rule with the Clark angular rule:
\begin{equation}
 U_{Y,T}\,
 \bigl[(\mathcal X_TF,\mathcal Z_TF)
       \mapsto(\mathcal Y_TF,\mathcal Z_TF)\bigr]\,
 U_{X,T}^{-1}
 =\mathfrak C_T^{\rm Cl}.                            \tag{400}
\end{equation}
This constructs the physical--Clark Hodge bridge without using the sign of
$A_T$.
\end{theorem}

\begin{proof}
By (395),
\[
 \mathcal Z_{p,T}^*J_{p,T}\mathcal Z_{p,T}
 =\log p\,I-\sum_{k\log p<2T}
        w_{p,k}J_{k\log p,-}^*J_{k\log p,-},
\]
which is (393), proving (396).  Add (396) to each of the two physical Grams
and use (394); this proves (398).

The right sides of (398) are ordinary positive Hilbert Grams.  Hence the
ambient Krein forms restrict to positive definite forms on the generated
ranges; injectivity follows from the winding coordinates as in the proof of
(390) for the $X$ port and from the $m_0$ coordinate for the $Y$ port.
Gram equality makes the rules (399) well defined and isometric, so
they extend to the closures.  Both stabilized ports carry the same
$\mathcal Z_TF$ state.  Applying (399) before and after the source rule
therefore gives (400).
\end{proof}

\begin{corollary}[Exact scope of the Hodge stabilization --- \AUDIT]
The stabilization (395)--(400) solves the type and conservation mismatch
between the physical and Clark ports.  It does not prove contractivity:
\begin{equation}
 \|\widetilde{\mathcal X}_TF\|_{J}^2
 -\|\widetilde{\mathcal Y}_TF\|_{J}^2
 =\|\mathcal X_T^{\rm Cl}F\|^2
 -\|\mathcal Y_T^{\rm Cl}F\|^2
 =\langle A_TF,F\rangle.                            \tag{401}
\end{equation}
Thus positivity of the stabilized angular map on the primitive generated
range remains exactly G40.  Declaring that range contractive merely because
its input and output metrics are individually positive would confuse
positivity of two Grams with domination between them.
\end{corollary}

\begin{proof}
The common signed state cancels in the first difference, and (389) gives the
second.  This proves (401) and the scope statement.
\end{proof}

\paragraph{Panoramic location of the missing theorem --- \AUDIT.}
Rows (a)--(c) do not contradict the certificate above.  They construct the
carrier, the Witt correspondences and the nuclear Lefschetz character on
their stated smooth/nuclear categories.  The hypothetical vector (346) is
forced by (348)--(349) to lie precisely outside the translation-smooth Meyer
range where those nuclear identities already give an $L^2$ causal output.
Thus deriving a contradiction merely from the scalar explicit formula, or
from the already proved local normalizations, would be circular.

There are now two genuinely different, non-circular ways to kill the same
negative certificate.  The analytic way is a boundary-completion theorem
showing that every equality state (346) has the Meyer lift (349), or just the
two Hardy estimates (287).  The geometric way is the extension, requested in
the supplied main manuscript, of periodic Riemann--Roch and effectivity to
the completed mixed class $D_f$: a nonzero effective primitive class would
then have simultaneously positive and zero degree, while Riemann--Roch
applied to $nD_f$ and $-nD_f$ forbids both from remaining ineffective when
$D_f^2>0$.  Either route supplies new information at the completion boundary;
neither is another identity among G39--G42.

The reverse-engineering target is therefore narrowed once more: construct a
\emph{faithful mixed-effectivity/causal completion functor} whose null
objects are exactly the equality states (346), and prove faithfulness before
identifying its Euler--Gamma character with row (c).  The faithfulness step
must rule out (348)--(350); the later character comparison may then use the
already proved amplitude and normalization identities.  Defining this
functor's norm by $-B_{\rm nuc}$ or by $\Delta_T^{1/2}$ would assume the
conclusion and remains forbidden.

\subsection{The panoramic negative: compressed radical versus global radical}

The preceding constructions can be viewed from farther away.  At a first
loss there is no longer an inequality to discover at the equality state:
semidefiniteness has already turned that state into a radical vector for the
\emph{compressed} form.  What is missing is the passage from that compressed
radical to the global radical.  This distinction is invisible in scalar
amplitudes and is the exact place where causality enters.

\begin{theorem}[First-loss radical lifting reduction --- \PROVED]
Let $T_*$ be a first-loss window and let $0\ne F\in\mathcal P_{T_*}\cap
\mathcal Q_{T_*}$ be a corresponding zero mode.  Put $u=E_{T_*}F$ and
\begin{equation}
 w:=L\,\mathcal F^{-1}\!\left(q_{T_*}\widehat u\right),
 \qquad L=\partial_t^2-\tfrac14,
 \qquad
 q_T=g_\Gamma+\vartheta(e^{2T})-m_0
             -\sum_{p<e^{2T}}d_p .                  \tag{402}
\end{equation}
Then
\begin{align}
 \langle A_{T_*}F,K\rangle&=0
       &&(K\in\mathcal P_{T_*}\cap\mathcal Q_{T_*}),\tag{403}\\
 P_{(-T_*,T_*)}w&=0.                                 \tag{404}
\end{align}
If every first-loss state satisfies the radical-lifting implication
\begin{equation}
 P_{(-T,T)}L\mathcal F^{-1}(q_T\widehat{E_TF})=0
 \quad\Longrightarrow\quad
 L\mathcal F^{-1}(q_T\widehat{E_TF})=0,              \tag{405}
\end{equation}
then no first loss exists and condition $G$ holds.  Thus (405), required
only on first-loss states and not on the whole form domain, is a complete
non-circular target for G40.
\end{theorem}

\begin{proof}
At the first-loss parameter the Friedrichs compression of $A_{T_*}$ to the
primitive space is nonnegative and has $F$ in its kernel.  Cauchy--Schwarz
for its nonnegative closed form gives (403).  More precisely, (253) and the
compression identity (371) say on $I_{T_*}$ that
$\mathcal F^{-1}(q_{T_*}\widehat u)$ belongs to
$\operatorname{span}\{e^{-t/2},e^{t/2}\}$.  Applying $L$ kills exactly
these two Tate directions and proves (404), first distributionally on the
form core and then by closure.

Assume (405).  Fourier transformation gives
$(-\tau^2-\tfrac14)q_{T_*}(\tau)\widehat u(\tau)=0$ as a tempered
distribution, hence almost everywhere on the real axis.  The function
$q_{T_*}$ is real analytic on the real axis: $g_\Gamma$ is real analytic
there and the prime sum is finite.  It is not identically zero, since
$q_{T_*}(\tau)=\log|\tau|+O_{T_*}(1)$ as $|\tau|\to\infty$.  Its real zero
set is consequently discrete, hence null; the factor
$-\tau^2-\tfrac14$ has no real zero.  Therefore $\widehat u=0$ almost
everywhere, so $u=0$ and $F=0$, contradicting first loss.  The first-loss
theorem then gives condition $G$.
\end{proof}

\begin{corollary}[Equality-state target, viewed in the negative --- \PROVED]
It is unnecessary to construct a contractive realization on every source
in order to exclude failure.  It is enough to construct, for a hypothetical
first-loss state, a causal conservation law forcing its differentiated
exterior leakage
$w$ in (402) to vanish.  Equivalently, it is enough to prove that a
compressed primitive radical cannot carry a nonzero exterior-only output.
The Gamma--Euler amplitude and the two Tate moments may be used after this
causal statement is established; they cannot replace it.
\end{corollary}

\begin{proof}
The first two assertions are precisely Theorem (402)--(405).  The last
assertion follows because G39 and the two moments determine the interior
cross-amplitude and terminal normalization, whereas (404) already says that
this entire interior observation is zero.  Only the exterior component of
the same global output remains undetermined.
\end{proof}

\begin{proposition}[Why rows (a)--(c) do not by themselves lift the radical
--- \AUDIT]
Hermitian symmetry and all scalar self-pairing identities do not imply that
an isotropic vector is radical.  In particular, in the hyperbolic plane
with Gram matrix
\begin{equation}
 \begin{pmatrix}0&1\\[2pt]1&0\end{pmatrix},           \tag{406}
\end{equation}
the first basis vector has zero self-pairing but pairs nontrivially with the
second.  Consequently a contradiction of a first-loss mode cannot be
deduced solely from the row-(c) character, G39, or the normalizations: one
must additionally prove either the causal radical lifting (405), the
equivalent exterior-leakage uniqueness theorem, or an independent
Hodge/effectivity equality theorem which makes primitive isotropic classes
radical and is faithful on the mixed completion.
\end{proposition}

\begin{proof}
The assertion about (406) is immediate.  Rows (a)--(c) construct the
objects, contacts and Hermitian character, but no theorem there asserts
that every primitive isotropic mixed class lies in the radical of the
completed pairing.  In the analytic realization the missing assertion is
exactly the promotion of (404) to (405).  Theorem (402)--(405) proves that
this additional statement is sufficient, while the hyperbolic example
shows that it is not a formal consequence of symmetry and a vanishing
self-pairing.
\end{proof}

\paragraph{Status of the radical-lifting statement --- \OPEN.}
The implication (405) is the negative, equality-only form of the remaining
problem.  The insertion of $L$ is essential: it removes, rather than
silently discards, the two Tate directions in (253).  The statement is
strictly narrower than constructing the full contraction,
but it has not been proved by the available normalizations.  In the language
of (274)--(287), it says that the exterior Hardy output of a first-loss state
must vanish; in the language of (352)--(360), it is the missing amplification
from the natural co-Poisson gap to the reciprocal Sonin gap.  Treating (404)
as though it were the global equation would erase the projection
$P_{(-T,T)}$ and would be circular.

\begin{proposition}[Tate divisibility does not create scalar positivity ---
\PROVED]
For every $T>\log2$ the Wigner--Smith density (370) satisfies
\begin{equation}
 q_T(0)=-m_0+
 \sum_{p<e^{2T}}\bigl(\log p-d_p(0)\bigr)<0,          \tag{407}
\end{equation}
where
\begin{equation}
 d_p(0)=(\log p)\frac{1+p^{-1/2}}{1-p^{-1/2}}
       >\log p.                                      \tag{408}
\end{equation}
Consequently the clamped-potential factor
$(\tau^2+1/4)^2$ in (93) does not make the scalar multiplier pointwise
nonnegative.  Any proof from (93) must use the Paley--Wiener support geometry,
not a pointwise factorization of its symbol.
\end{proposition}

\begin{proof}
At the origin $g_\Gamma(0)=0$, while
$\vartheta(e^{2T})=\sum_{p<e^{2T}}\log p$.  Substitution in (370) gives
(407), and (408) follows directly from (362).  Since
$(0^2+1/4)^2=1/16>0$, multiplication by the Tate factor preserves the
strictly negative sign at the origin.
\end{proof}

\begin{proposition}[Infinite prime-pole obstruction after the Tate factor ---
\PROVED]
Fix an active prime $p<e^{2T}$.  The meromorphic continuation of its Clark
density has simple poles at
\begin{equation}
 z=\frac{2\pi k}{\log p}+\frac i2,
 \qquad
 z=\frac{2\pi k}{\log p}-\frac i2,
 \qquad k\in\mathbb Z.                              \tag{409}
\end{equation}
The polynomial $(z^2+1/4)^2$ cancels at most the two locations with $k=0$.
All poles with $k\ne0$ remain in
$(z^2+1/4)^2q_T(z)$; they cannot be cancelled by the Gamma term or by a
different prime.  Hence Tate divisibility does not turn the scalar delay
into a finite-index generalized Schur/Nevanlinna object.  In particular a
finite Pontryagin-index scalar colligation cannot prove (405).
\end{proposition}

\begin{proof}
Factor the denominator in (362) as
\[
 1-2r_p\cos((\log p)z)+r_p^2
 =(1-r_pe^{i(\log p)z})(1-r_pe^{-i(\log p)z}).
\]
Since $r_p=p^{-1/2}=e^{-(\log p)/2}$, its zeros are exactly (409), and the
nonzero numerator in (362) makes them simple.  The Gamma poles lie on the
imaginary axis, so they do not meet the points (409) with $k\ne0$.
If poles belonging to distinct primes coincided, then
$k/\log p=\ell/\log q$ for nonzero integers $k,\ell$, which would imply
$p^\ell=q^k$, impossible for distinct primes.  Thus no cross-prime
cancellation occurs.  The Tate polynomial has only the zeros $\pm i/2$.

Finally, the Krein--Langer representation of a scalar generalized Schur
function with finite negative index is a Schur function divided by a finite
Blaschke product; it has only finitely many poles in the chosen half-plane.
The surviving infinite set (409) rules out such a representation.  This
argument concerns only scalar finite-index realizations and does not rule
out the infinite operator-valued two-contour state already constructed.
\end{proof}

\paragraph{Consequence for the next attack --- \AUDIT.}
Equations (407)--(409) close the two most direct scalar exits from (405):
neither the Tate polynomial nor a finite negative-index correction absorbs
the Euler advance.  The remaining viable route is therefore the
operator-valued Wiener--Hopf interconnection of the two prime orientations
with the Gamma extension and its boundary stores.  Its equality case must
show that zero interior output has zero exterior innovation; this is a
causal conservation statement, not a meromorphic pole cancellation.

\begin{theorem}[Divergence of the orthogonal exterior prime reservoir ---
\PROVED]
Let $0\ne F\in L^2(I_T)$ and, for a prime $p\ge e^{2T}$, put
\begin{equation}
 Y_{p,+}(t):=\sum_{k\ge1}p^{-k/2}E_TF(t-k\log p).
                                                               \tag{410}
\end{equation}
Every summand in (410) is supported in the right exterior and the summands
are pairwise orthogonal.  Consequently
\begin{equation}
 \|Y_{p,+}\|_2^2=\frac1{p-1}\|F\|_2^2.              \tag{411}
\end{equation}
In particular, the source-labelled orthogonal prime reservoir has infinite
exterior energy:
\begin{equation}
 \sum_{p\ge e^{2T}}(\log p)\|Y_{p,+}\|_2^2
 =\|F\|_2^2\sum_{p\ge e^{2T}}\frac{\log p}{p-1}
 =+\infty.                                           \tag{412}
\end{equation}
The same assertion holds for the reflected orientation.
\end{theorem}

\begin{proof}
If $a=\log p\ge2T$, the intervals
$I_T+ka$, $k\ge1$, have disjoint interiors and lie to the right of $I_T$.
Translation is isometric, hence orthogonality of the supports gives
\[
 \|Y_{p,+}\|_2^2
 =\sum_{k\ge1}p^{-k}\|F\|_2^2
 =\frac1{p-1}\|F\|_2^2,
\]
which is (411).  To prove divergence, for $s>1$ use the Euler product:
\[
 -\frac{\zeta'}{\zeta}(s)
 =\sum_p\frac{\log p}{p^s-1}.
\]
The left side tends to $+\infty$ as $s\downarrow1$ because $\zeta$ has a
simple pole at one.  Monotone convergence therefore gives
$\sum_p\log p/(p-1)=+\infty$; deleting finitely many primes does not change
this.  Equation (412) follows.  Reflection proves the other orientation.
\end{proof}

\begin{corollary}[The Tate centering must precede the exterior norm ---
\PROVED]
No global G40 observation can take the full family of exterior prime Julia
outputs as mutually orthogonal Hilbert coordinates.  Before completing the
exterior port, it must first synthesize the prime channels and remove their
Lebesgue mean by the two Tate moments.  On the primitive space the resulting
source is exactly the single compensated Chebyshev current
\begin{equation}
 Y_+(t)=\int_{e^{t-T}}^{e^{t+T}}
 x^{-1/2}F(t-\log x)\,dR_\Psi(x),                    \tag{413}
\end{equation}
and its reflection, as in (291).
\end{corollary}

\begin{proof}
Theorem (410)--(412) rules out the orthogonal completion for every nonzero
source.  Equations (289)--(291) show that synthesis followed by the Tate
moment cancellation replaces $d\Psi$ by $dR_\Psi=d\Psi-dx$, proving (413).
No sign assertion is used.
\end{proof}

\paragraph{Revised construction target --- \AUDIT.}
The finite-prime Bayes product (220k) correctly realizes the interior
compression, but (412) proves that its independent innovation coordinates
cannot be sent unchanged to the global exterior boundary.  The missing
Poisson--Sonin/Hodge functor must contain a renormalized, cross-prime
synthesis map
\[
 \{\hbox{prime Julia innovations}\}
 \longrightarrow Y_+\oplus Y_-
 \longrightarrow \hbox{Hardy boundary},
\]
where the first arrow performs Tate centering before completion.  Proving
that this compensated synthesis is bounded on a first-loss equality state
is exactly (287).  Treating the prime innovations as an orthogonal global
direct sum is now ruled out, rather than left as an unspecified option.

\begin{theorem}[Exact exterior jump balance and parity reduction ---
\PROVED]
Let $F\in\mathcal P_T$ be a zero mode and put
\begin{equation}
 R_{+,T}:=1_{(T,\infty)}Y_+,qquad
 R_{-,T}:=1_{(-\infty,-T)}Y_- .                     \tag{414}
\end{equation}
Then, as distributions on the exterior of $I_T$,
\begin{equation}
 1_{\mathbb R\setminus I_T}\mathfrak J_\xi(E_TF)
 =1_{\mathbb R\setminus I_T}G_\Gamma E_TF
  -R_{+,T}-R_{-,T}.                                  \tag{415}
\end{equation}
If $F(-t)=\varepsilon F(t)$, $\varepsilon\in\{1,-1\}$, then
\begin{equation}
 R_{-,T}(t)=\varepsilon R_{+,T}(-t),qquad
 R_{+,T}\in L^2(T,\infty)Longleftrightarrow
 R_{-,T}\in L^2(-\infty,-T).                        \tag{416}
\end{equation}
Consequently, for a parity representative of a nonzero first-loss mode,
both exterior arithmetic tails fail to belong to $L^2$.  Conversely, a
Hardy estimate for either one orientation alone excludes that mode and is
sufficient to prove condition $G$.
\end{theorem}

\begin{proof}
On the right exterior, $Y_-$ vanishes because
$t+\log n>T$ lies outside the support of $E_TF$; symmetrically, $Y_+$
vanishes on the left exterior.  The constant channel $m_0E_TF$ also
vanishes there.  Source identification (279a) therefore gives (415).
Reflection of the defining sums gives
$Y_-(t)=\varepsilon Y_+(-t)$, proving (416).

Suppose one tail of a parity first-loss mode belonged to $L^2$.  By (416)
so would the other.  The exact Laplace/Hardy equivalence (286) then gives
both reflected Hardy memberships.  Hardy pole rigidity (292) forces
$F=0$, a contradiction.  Thus both tails fail.  The same argument shows
directly that one oriented Hardy estimate already excludes every parity
first mode; parity reduction then excludes all first modes and proves $G$.
\end{proof}

\begin{corollary}[One-sided compensated-current target --- \PROVED]
The remaining analytic bridge can be stated using only the right
orientation: for every parity first-loss state, prove
\begin{equation}
 \int_T^\infty\left|
 \int_{e^{t-T}}^{e^{t+T}}
 x^{-1/2}F(t-\log x)\,dR_\Psi(x)
 \right|^2dt<\infty.                                \tag{417}
\end{equation}
Then condition $G$ follows.  Formula (417) is exactly one of the two Hardy
bounds (287), with the prime mean already removed by the Tate moment.
\end{corollary}

\begin{proof}
Equation (413) identifies the integrand with $R_{+,T}$.  Theorem
(414)--(416) supplies the reflected bound and excludes the first-loss mode.
The equivalence with the right Hardy condition is (286).
\end{proof}

\paragraph{Audit of the tempting symmetric closure --- \AUDIT.}
Equation (415) shows why the functional equation and parity do not by
themselves make the arithmetic sum an $L^2$ Gamma output: the omitted term
is precisely the unsupertraced jump $\mathfrak J_\xi(E_TF)$.  Dropping it would
replace the row-(c) source character by its meromorphically cancelled scalar
supercharacter.  The gain from parity is real but limited: it reduces two
Hardy estimates to (417); it does not prove that remaining estimate.

\begin{theorem}[Exact one-sided Riemann--Hilbert decomposition --- \PROVED]
Put $s=\tfrac12+z$ and define the meromorphic continuations
\begin{align}
 \mathscr A_+(z)&:=-\frac{\zeta'}{\zeta}(s)K_F(z)-Q_{+,T}(z),\notag\\
 \mathscr A_-(z)&:=-\frac{\zeta'}{\zeta}(1-s)K_F(z)-Q_{-,T}(z).       \tag{419}
\end{align}
Initially these are respectively the Laplace transforms of the right and
left arithmetic exterior tails in their disjoint convergence half-planes.
Let
\begin{align}
 \mathscr G_+(z)&:=\int_T^\infty
          (G_\Gamma E_TF)(t)e^{-zt}\,dt,\notag\\
 \mathscr G_-(z)&:=\int_{-\infty}^{-T}
          (G_\Gamma E_TF)(t)e^{-zt}\,dt,             \tag{420}
\end{align}
which have a common strip containing the imaginary axis.  For every zero
mode (279), meromorphic continuation gives
\begin{equation}
 \boxed{\mathscr A_+(z)+\mathscr A_-(z)
 =\mathscr G_+(z)+\mathscr G_-(z)
  +P_{\alpha,\beta,T}(z).}                           \tag{421}
\end{equation}
If $F(-t)=\varepsilon F(t)$, then
\begin{equation}
 \mathscr A_-(z)=\varepsilon\mathscr A_+(-z),        \tag{422}
\end{equation}
so (421) fixes the reflected symmetric part of the one-sided arithmetic
tail.  Its boundary on $i\mathbb R$ belongs to $L^2$.  The reflected
antisymmetric part is the only undetermined Riemann--Hilbert datum.
\end{theorem}

\begin{proof}
The first assertion is (285) and its reflected version.  Logarithmic
differentiation of the functional equation gives
\[
 -\frac{\zeta'}{\zeta}(s)
 -\frac{\zeta'}{\zeta}(1-s)
 =\frac12\psi(s/2)+\frac12\psi((1-s)/2)-\log\pi.
\]
The right side is the bilateral Laplace multiplier of
$G_\Gamma-m_0I$.  Hence, in the common Gamma strip,
\[
 \left[-\frac{\zeta'}{\zeta}(s)
 -\frac{\zeta'}{\zeta}(1-s)\right]K_F(z)
 =\Gamma_T(z)-m_0K_F(z)+\mathscr G_+(z)+\mathscr G_-(z).
\]
Subtract (296) and obtain (421).  Equation (422) follows from (297) and
$K_F(-z)=\varepsilon K_F(z)$.  Finally (281) and Plancherel give $L^2$
boundary values for both Gamma exterior transforms; the transform
$P_{\alpha,\beta,T}$ comes from a function on the bounded interval and also
has an $L^2(i\mathbb R)$ boundary.  Thus the symmetric boundary in (421) is
$L^2$.
\end{proof}

\begin{corollary}[The zeta-zero residues live only in the missing
antisymmetric datum --- \PROVED]
Let $z_\rho=\rho-\tfrac12$ for a nontrivial zero $\rho$ of multiplicity
$m_\rho$.  Then the two terms in (419) have opposite residues,
\begin{equation}
 \operatorname*{Res}_{z=z_\rho}\mathscr A_+(z)
   =-m_\rho K_F(z_\rho),\qquad
 \operatorname*{Res}_{z=z_\rho}\mathscr A_-(z)
   =+m_\rho K_F(z_\rho).                             \tag{423}
\end{equation}
Consequently the symmetric identity (421) cancels every zeta-zero pole,
whereas
\begin{equation}
 \operatorname*{Res}_{z=z_\rho}
  (\mathscr A_+-\mathscr A_-)
 =-2m_\rho K_F(z_\rho).                              \tag{424}
\end{equation}
Thus proving the one-sided Hardy estimate (417) is exactly proving that the
antisymmetric jump has no polar or non-$H^2$ boundary component.  Neither
G39 nor the functional equation can see this datum, because both retain
only (421).
\end{corollary}

\begin{proof}
The residue of $-\zeta'/\zeta(s)$ at $s=\rho$ is $-m_\rho$.
At the same $z=z_\rho$, the argument $1-s$ equals $1-\rho$, also a zero of
multiplicity $m_\rho$; the derivative $d(1-s)/dz=-1$ reverses the residue
and gives the second equality in (423).  The entire functions $Q_{\pm,T}$
have no residues.  Addition and subtraction prove (424) and the final
statement follows from the Hardy pole calculation in (292).
\end{proof}

\begin{theorem}[Smirnov symmetry does not select the causal orientation
--- \PROVED]
Fix \(a>0\).  In the even source sector put
\[
 h_{+,a}(z):=\frac{z}{z^2-a^2},
\]
and in the odd source sector put
\[
 h_{-,a}(z):=\frac{1}{z^2-a^2}.
\]
Both functions are meromorphic of bounded type in the right half-plane,
have \(L^2\) boundary values on \(i\mathbb R\), and tend to zero as
\(\operatorname{Re}z\to+\infty\).  Nevertheless
\begin{align}
 h_{+,a}(z)+h_{+,a}(-z)&=0,\notag\\
 h_{-,a}(z)-h_{-,a}(-z)&=0,                           \tag{424a}
\end{align}
while each has a nonremovable pole at \(z=a\) and therefore is not in the
right Hardy space.  Consequently the data
\begin{enumerate}
\item meromorphic Smirnov/Nevanlinna class;
\item the physical decay at \(+\infty\);
\item an \(L^2\) reflected symmetric boundary as in (421);
\end{enumerate}
do not determine \(\mathscr A_+\) and cannot imply (417).  A valid causal
selection theorem must use an additional source-range, passivity or
effectivity property which excludes the gauges (424a).
\end{theorem}

\begin{proof}
The two functions are rational with no pole on \(i\mathbb R\), so their
boundary values are square integrable and they are meromorphic functions of
bounded type in the right half-plane.  Their displayed parity identities
are immediate, as is their decay.  Both have a pole at the interior point
\(a\); hence neither is holomorphic, and a fortiori neither belongs to
\(H^2(\mathbb C_+)\).  Adding \(h_{+,a}\) to an even-sector candidate, or
\(h_{-,a}\) to an odd-sector candidate, preserves the corresponding
reflected identity (421) while changing the causal pole data.  This proves
the nonuniqueness.
\end{proof}

\paragraph{Revised causal target --- \AUDIT.}
The Smirnov class is a correct container for the meromorphic continuation
but not the missing physical law.  The extra datum must distinguish the
actual Paley--Wiener source range from the rational gauges above.  In the
residue realization this is precisely a source-derived invariance or
positivity theorem for the generated range of \(\mathcal E_T\); declaring
that range \(J_{\rm res}\)-nonnegative would be G40 itself and is therefore
not available as a definition.

\paragraph{Operator-valued Hodge target after Riemann--Hilbert splitting ---
\OPEN.}
The new bridge no longer needs to reproduce the symmetric amplitude: that
is already (421), with an $L^2$ boundary.  It must construct a positive
observation of the antisymmetric two-contour datum
$\mathscr A_+-\mathscr A_-$ and prove, on a first-loss equality state, that
this datum lies in the correct one-sided Hardy class.  Equation (424) shows
why this is the substantive equality theorem: it forces
$K_F(\rho-1/2)=0$ at every zeta zero and hence $F=0$.  Defining the norm of
this observation from those residues, from $\Delta_T$, or from the desired
Weil form would remain circular.

\begin{theorem}[Covariant derivative of the one-sided M\"obius lift ---
\PROVED]
Define, initially away from the zeros of the denominator,
\begin{equation}
 \mathcal M_F^+(z):=
 \frac{K_F(z)}{\zeta(\tfrac12+z)}.                   \tag{425}
\end{equation}
Then the right exterior transform in (419) has the exact identity
\begin{equation}
 \boxed{\mathscr A_+(z)
 =\zeta(\tfrac12+z)(\mathcal M_F^+)'(z)
   -K_F'(z)-Q_{+,T}(z).}                             \tag{426}
\end{equation}
The compact counterterms have the opposite elementary support orientation:
\begin{equation}
 e^{-Tz}K_F'(z),\ e^{-Tz}Q_{+,T}(z)\in H^2(\mathbb C_+). \tag{427}
\end{equation}
Consequently the one-sided estimate (417) is equivalent to the affine
M\"obius--Dirichlet condition
\begin{equation}
 \boxed{
 e^{Tz}\left[
 \zeta(\tfrac12+z)(\mathcal M_F^+)'(z)
 -K_F'(z)-Q_{+,T}(z)\right]
 \in H^2(\mathbb C_+).}                              \tag{428}
\end{equation}
This equivalence uses no sign of the Weil form.
\end{theorem}

\begin{proof}
Differentiate (425):
\[
 \zeta(\tfrac12+z)(\mathcal M_F^+)'(z)
 =K_F'(z)-\frac{\zeta'}{\zeta}(\tfrac12+z)K_F(z).
\]
Subtract $K_F'+Q_{+,T}$ and use (419) to obtain (426).

The function $K_F'$ is the Laplace transform of $-tF(t)$ on $[-T,T]$.
After multiplication by $e^{-Tz}$ it is the Laplace transform, in the
variable $x=t+T$, of a function supported in $[0,2T]$, hence belongs to
$H^2(\mathbb C_+)$ by Paley--Wiener.  The same argument applies to
$Q_{+,T}$, whose inverse Laplace transform is supported in
$[\log2-T,T]$.  This proves (427).  Notice that the sign in (427) is
opposite to the causal shift $e^{Tz}$ in (286); the compact prehistory may
not be discarded.  Finally (286) and (426) give exactly (428).
\end{proof}

\begin{definition}[The affine critical M\"obius--Dirichlet boundary space]
For fixed source data $(K,Q)$ let
$\mathscr D_{\zeta,T}^+(K,Q)$ be the affine space of meromorphic germs
$M$ in $\mathbb C_+$ for which
\begin{equation}
 \|M\|_{\mathscr D_{\zeta,T}^+(K,Q)}^2
 :=\left\|e^{Tz}\left[
 \zeta(\tfrac12+z)M'(z)-K'(z)-Q(z)\right]\right\|_{H^2}^2
 <\infty,                                            \tag{429}
\end{equation}
modulo constants whenever the displayed boundary expression exists.  This
is a positive affine Hilbert seminorm defined from the source Zeta
multiplier, differentiation and the physical compact prehistory, before any
Weil sign is chosen.
\end{definition}

\begin{corollary}[Equality-mode M\"obius--Dirichlet criterion --- \PROVED]
Condition $G$ follows if every parity first-loss state satisfies
\begin{equation}
 \mathcal M_F^+=K_F/\zeta(\tfrac12+\cdot)
 \in\mathscr D_{\zeta,T}^+(K_F,Q_{+,T}).             \tag{430}
\end{equation}
Membership forces the polar part at every nontrivial zeta zero to vanish;
equivalently $K_F(\rho-1/2)=0$ for every distinct zero $\rho$, using parity
and the functional-equation reflection for the zeros left of the line.
\end{corollary}

\begin{proof}
Membership (430) is exactly (428), hence (417), and the one-sided criterion proves
$G$.  Alternatively, if $K_F$ failed to vanish to the multiplicity required
at a zero, differentiating $K_F/\zeta$ and multiplying once by $\zeta$
would leave a nonzero simple pole.  Such a pole is incompatible with the
$H^2$ boundary expression in (429), because $K_F'$ and $Q_{+,T}$ are
entire.  This treats $\Re\rho\ge1/2$; parity
$K_F(-z)=\varepsilon K_F(z)$ and the zero reflection
$\rho\mapsto1-\rho$ treat the remaining zeros.  Hardy pole rigidity and the density of distinct
zeros then give $F=0$.
\end{proof}

\paragraph{Status of the new space --- \OPEN.}
Equations (425)--(430) construct the correctly oriented affine target norm
and prove its exact equivalence with the missing one-sided current.  What is
not yet proved is the source map carrying a physical first-loss state into
$\mathscr D_{\zeta,T}^+(K_F,Q_{+,T})$.  The differentiated Poisson
identity/G39 supplies the algebraic covariant derivative (426), but
boundedness of its causal boundary realization, including the compact
prehistory counterterm, is precisely the remaining operator-level
Gamma--Tate normalization.  Defining that source map by completion in the
norm (429) would assume its finiteness and is not permitted.

\subsection{Return to the reciprocal co-Poisson domain gate}

The abstract closability question left below (358) can be separated from
the genuinely arithmetic regularity question.  Write
\[
 \widehat F(\tau)=\int_{-T}^{T}F(t)e^{-i\tau t}\,dt,
 \qquad
 w_\Gamma(\tau)=1+g_\Gamma(\tau).
\]

\begin{theorem}[Maximal physical co-Poisson observation --- \PROVED]
For $F\in C_c^\infty(I_T)\cap\mathcal P_T$, the canonical observation
of (352)--(356) satisfies the exact Plancherel identity
\begin{equation}
 \boxed{
 \|\mathfrak S_T^0F\|_{L^2(0,\infty)}^2
 =\frac1{2\pi}\int_{\mathbb R}
   |\zeta(\tfrac12+i\tau)|^2|\widehat F(\tau)|^2\,d\tau.}       \tag{431}
\end{equation}
Consequently the maximal closed realization is
\begin{align}
 \operatorname{Dom}(\mathfrak S_T^{\max})
 :=\bigl\{F\in L^2(I_T)\cap\mathcal P_T:\
 &\ \zeta(\tfrac12+i\,\cdot)\widehat F
       \in L^2(\mathbb R)\bigr\},                  \tag{432}\\
 \mathcal M(\mathfrak S_T^{\max}F)(\tfrac12+i\tau)
 :=&\ \zeta(\tfrac12+i\tau)\widehat F(\tau).
 \notag
\end{align}
In particular, the smooth-core operator $\mathfrak S_T^0$ is closable in
$L^2$, without any sign assumption.  A physical first-loss vector admits
this observation if and only if it belongs to the explicit domain (432).
\end{theorem}

\begin{proof}
With $x=e^t$, definition (352) gives
\[
 G_F(\tfrac12+i\tau)
 =\int_0^\infty x^{-1/2}F(\log x)x^{-1/2-i\tau}\,dx
 =\widehat F(\tau).
\]
Equation (356) and Mellin--Plancherel now give (431).  Fourier transform on
$I_T$, multiplication by the measurable function
$\zeta(\tfrac12+i\tau)$, and inverse Mellin transform are respectively
bounded, closed, and unitary operations.  Their composition on (432) is
therefore closed.  Since the smooth primitive core is contained in (432),
its restriction is closable.  The last assertion is exactly the definition
of the maximal domain.
\end{proof}

\begin{proposition}[Exact full-form extension criterion --- \PROVED]
The physical form norm of the first-loss operator has the equivalent Gamma
norm
\begin{equation}
 \|F\|_{\Gamma,T}^2
 :=\frac1{2\pi}\int_{\mathbb R}
       w_\Gamma(\tau)|\widehat F(\tau)|^2\,d\tau,
 \qquad
 w_\Gamma(\tau)\asymp 1+\log(2+|\tau|).             \tag{433}
\end{equation}
The observation extends continuously from the smooth primitive core to the
entire physical form domain if and only if there is a finite $C_T$ such that
\begin{equation}
 \boxed{
 \int_{\mathbb R}|\zeta(\tfrac12+i\tau)|^2
          |\widehat F(\tau)|^2\,d\tau
 \le C_T\int_{\mathbb R}w_\Gamma(\tau)
          |\widehat F(\tau)|^2\,d\tau,
 \quad F\in C_c^\infty(I_T)\cap\mathcal P_T.}       \tag{434}
\end{equation}
Thus closability is no longer an open item; what remains open at the domain
level is either (434), or the strictly weaker assertion that every
first-loss eigenvector belongs individually to (432).
\end{proposition}

\begin{proof}
The asymptotic in (433) follows from the standard digamma asymptotic in
(81).  The constant and the finitely many active translated overlaps are
bounded forms, so they do not change the Gamma form domain.  Combining
(431) with (433) proves that (434) is sufficient.  Conversely, a continuous
extension has a finite operator norm between these two Hilbert spaces, and
(431) gives (434) on the core.
\end{proof}

\begin{proposition}[Necessary local zeta bound for full-form extension ---
\PROVED]
If (434) holds for a fixed $T>0$, then there exist
$\delta_T,c_T>0$ such that, for every sufficiently large real $R$,
\begin{equation}
 \boxed{
  \int_{R-\delta_T}^{R+\delta_T}
       |\zeta(\tfrac12+i\tau)|^2\,d\tau
  \le c_T\log(2+R).}                                \tag{435}
\end{equation}
Thus a full-form extension cannot be obtained merely from the asymptotic
second moment on long intervals; it contains a uniform fixed-window
estimate.
\end{proposition}

\begin{proof}
Choose a nonzero $\phi\in C_c^\infty(I_T)$.  After shrinking its support if
necessary, choose $\delta_T>0$ so that
$|\widehat\phi(\tau)|\ge c>0$ for $|\tau|<2\delta_T$.
The modulated function $e^{iRt}\phi(t)$ need not satisfy the two Tate
moments.  Choose fixed $h_-,h_+\in C_c^\infty(I_T)$ for which the matrix
$(M_\sigma h_\nu)_{\sigma,\nu\in\{-,+\}}$ is invertible and subtract the
unique linear combination of $h_-,h_+$ which kills those moments.  Call the
result $F_R\in C_c^\infty(I_T)\cap\mathcal P_T$.

Repeated integration by parts shows that the two moments of
$e^{iRt}\phi$ are $O_N(R^{-N})$ for every $N$; hence the correcting
coefficients are also $O_N(R^{-N})$.  Uniformly for
$|\tau-R|<2\delta_T$,
\[
 \widehat F_R(\tau)=\widehat\phi(\tau-R)+O_N(R^{-N}),
\]
and therefore $|\widehat F_R(\tau)|\ge c/2$ for all sufficiently large
$R$ and $|\tau-R|<\delta_T$.  On the other hand, translation of a Schwartz
function and the rapidly decreasing correction give
\[
 \int_{\mathbb R}w_\Gamma(\tau)|\widehat F_R(\tau)|^2\,d\tau
 =O_T(\log(2+R)).
\]
Apply (434) to $F_R$ and restrict its left side to
$(R-\delta_T,R+\delta_T)$.  The lower bound $c/2$ there proves (435).
\end{proof}

\begin{theorem}[Local physical co-Poisson extension and nilpotent inverse ---
\PROVED]
For every $F\in L^2(I_T)$, extended by zero outside $I_T$, the formula
\begin{equation}
 (\mathfrak S_T^{\rm loc}F)(x)
 :=\sum_{n\ge1}\frac1n x^{-1/2}F(\log(x/n))          \tag{436}
\end{equation}
defines an element of $L^2_{\rm loc}(0,\infty)$; on each compact
multiplicative interval the sum is finite.  Under the unitary logarithmic
coordinate $({\cal U}u)(t)=e^{t/2}u(e^t)$ one has, for $t\in I_T$,
\begin{align}
 ({\cal U}\mathfrak S_T^{\rm loc}F)(t)
 &=\mathcal Z_{T,+}F(t),                             \tag{437}\\
 \mathcal Z_{T,+}
 &:=I+K_{T,+},\qquad
 K_{T,+}F(t):=
 \sum_{2\le n<e^{2T}}n^{-1/2}E_TF(t-\log n).        \tag{438}
\end{align}
If $N\log2>2T$, then $K_{T,+}^{N}=0$ on $I_T$, and hence
\begin{equation}
 \boxed{
 \mathcal Z_{T,+}^{-1}
 =\sum_{j=0}^{N-1}(-K_{T,+})^j.}                    \tag{439}
\end{equation}
No global $L^2$ estimate for $\zeta(\tfrac12+i\tau)\widehat F(\tau)$ is
needed for (436)--(439).
\end{theorem}

\begin{proof}
If $x$ ranges over a compact subset $[c,d]\subset(0,\infty)$ and a summand
in (436) is nonzero, then
$e^{-T}\le x/n\le e^T$.  Thus $ce^{-T}\le n\le de^T$, leaving only
finitely many integers; changes of variables in these finitely many terms
give local $L^2$ boundedness.  Substitution $x=e^t$ gives
\[
 e^{t/2}\frac1n(e^t/n)^{-1/2}F(t-\log n)
 =n^{-1/2}F(t-\log n),
\]
which proves (437)--(438).

Every factor of $K_{T,+}$ delays its argument by at least $\log2$.
Consequently a term in $K_{T,+}^{N}F(t)$ samples $F$ at a point at most
$t-N\log2<-T$ when $t<T$ and $N\log2>2T$.  Zero extension makes every such
term vanish.  Thus $K_{T,+}^{N}=0$, and multiplication verifies the finite
Neumann inverse (439).
\end{proof}

\begin{corollary}[Domain-free annular completion criterion --- \PROVED]
For every physical state $F\in L^2(I_T)$,
\begin{equation}
 \mathfrak S_T^{\rm loc}F=0\ \hbox{ on }(e^{-T},e^T)
 \quad\Longrightarrow\quad F=0.                    \tag{440}
\end{equation}
Consequently condition $G$ follows from the single equality-state
implication
\begin{equation}
 \boxed{
 P_TA_TF=0
 \quad\Longrightarrow\quad
 \mathfrak S_T^{\rm loc}F=0\ \hbox{ on }(e^{-T},e^T),
 \qquad F\in\mathcal P_T.}                          \tag{441}
\end{equation}
This criterion requires neither (434), global $L^2$ membership of the
observation, nor Burnol completeness.  It is strictly the missing local
Gamma--Tate gap-transport identity.
\end{corollary}

\begin{proof}
The annular vanishing in (440), transported by $x=e^t$, says
$\mathcal Z_{T,+}F=0$ on $I_T$.  Apply (439).  If (441) holds, every
first-loss vector furnished by (251)--(253) is zero, a contradiction; the
first-loss theorem then gives $G$.
\end{proof}

\begin{theorem}[Exact local Euler commutator bridge --- \PROVED]
On $L^2(I_T)$ let $QF(t)=tF(t)$ and define the two active von Mangoldt
operators
\begin{equation}
 B_{T,+}F(t):=\sum_{2\le n<e^{2T}}
       \frac{\Lambda(n)}{\sqrt n}E_TF(t-\log n),
 \qquad B_{T,-}:=\mathcal RB_{T,+}\mathcal R.       \tag{442}
\end{equation}
Then the local co-Poisson amplitude of (437) satisfies
\begin{equation}
 \boxed{
 B_{T,+}=\mathcal Z_{T,+}^{-1}[Q,\mathcal Z_{T,+}].} \tag{443}
\end{equation}
If $\mathcal RF=\varepsilon F$, $\varepsilon\in\{1,-1\}$, and
$h=\mathcal Z_{T,+}F$, then the symmetric arithmetic current is
\begin{equation}
 \boxed{
 (B_{T,+}+B_{T,-})F
 =\mathcal Z_{T,+}^{-1}Qh
  +\varepsilon\mathcal R\mathcal Z_{T,+}^{-1}Qh.}  \tag{444}
\end{equation}
Thus both prime orientations are recovered from one observed amplitude;
the two bare $QF$ terms cancel exactly.
\end{theorem}

\begin{proof}
From (438),
\[
 [Q,\mathcal Z_{T,+}]F(t)
 =\sum_{2\le n<e^{2T}}
   \frac{\log n}{\sqrt n}E_TF(t-\log n).
\]
The finite causal inverse (439) has Dirichlet coefficients
$\mu(d)/\sqrt d$.  Hence the coefficient of the delay $\log n$ in the
right side of (443) is
\[
 \frac1{\sqrt n}\sum_{d\mid n}\mu(d)\log(n/d)
 =\frac{\Lambda(n)}{\sqrt n},
\]
which proves (443); truncation is harmless because all contributing delays
are below $2T$.

Equation (443) also gives
\[
 B_{T,+}F=\mathcal Z_{T,+}^{-1}Qh-QF.
\]
Since $B_{T,-}=\mathcal RB_{T,+}\mathcal R$,
$\mathcal RQ=-Q\mathcal R$, and $\mathcal RF=\varepsilon F$, reflection
gives
\[
 B_{T,-}F=QF+\varepsilon
 \mathcal R\mathcal Z_{T,+}^{-1}Qh.
\]
Adding proves (444).
\end{proof}

\begin{corollary}[Observed Gamma--Tate transport equation --- \PROVED]
Let $F$ be a parity zero mode as in (254) or (255), write
$h=\mathcal Z_{T,+}F$, and let
$r_\varepsilon$ denote the corresponding Tate vector
$c\cosh(t/2)$ or $c\sinh(t/2)$.  Then its zero-mode equation is exactly
\begin{equation}
 \boxed{
 \Bigl[(G_{\Gamma,T}-m_0I)\mathcal Z_{T,+}^{-1}
 -\mathcal Z_{T,+}^{-1}Q
 -\varepsilon\mathcal R\mathcal Z_{T,+}^{-1}Q\Bigr]h
 =r_\varepsilon.}                                  \tag{445}
\end{equation}
Conversely, every $h$ satisfying (445), with
$F=\mathcal Z_{T,+}^{-1}h\in\mathcal P_T$ of parity $\varepsilon$, gives a
primitive zero mode.  Therefore the only remaining local bridge is the
non-circular uniqueness statement
\begin{equation}
 \boxed{
 h\in\mathcal Z_{T,+}\mathcal P_T,\quad
 \eqref{445}
 \quad\Longrightarrow\quad h=0.}                   \tag{446}
\end{equation}
By (437), (446) is precisely the annular gap-transport implication (441),
now written with no undefined observation, completion, or global norm.
\end{corollary}

\begin{proof}
On $I_T$ the physical operator is
$A_T=G_{\Gamma,T}-m_0I-B_{T,+}-B_{T,-}$.  Substitute
$F=\mathcal Z_{T,+}^{-1}h$ and (444) into
$A_TF=r_\varepsilon$ to obtain (445).  Reversing these equalities proves
the converse.  Finally $h=0$ is equivalent to $F=0$ by (439), and (437)
identifies $h$ with the logarithmic annular observation.
\end{proof}

\begin{theorem}[Meromorphic Gamma--co-Poisson commutator --- \PROVED]
Let \(\mathbf Z_+\) denote the global co-Poisson operator in logarithmic
coordinates,
\[
 (\mathbf Z_+F)(t)=\sum_{n\ge1}n^{-1/2}F(t-\log n),
\]
initially on \(C_c^\infty(I_T)\cap\mathcal P_T\), and put
\(D_\Gamma:=G_\Gamma-m_0I\).  Then, in the co-Poisson distribution
calculus, equivalently as an identity of common meromorphic Mellin germs on
the primitive smooth core,
\begin{equation}
 \boxed{
 \mathbf Z_+(B_++B_-)=D_\Gamma\mathbf Z_+,\qquad
 \mathbf Z_+\mathcal A_\infty
   =[\mathbf Z_+,D_\Gamma].}                         \tag{447}
\end{equation}
Thus the complete Gamma normalization, the constant \(m_0\), and both
von Mangoldt orientations form one exact regularized commutator identity.
This does not assert that the composed physical boundary
\(\mathbf Z_+\mathcal A_\infty F\) is an ordinary \(L^2\) function.
\end{theorem}

\begin{proof}
Let \(\mathcal C\) be the normalized cosine transform and
\(\mathcal R F(t)=F(-t)\) the logarithmic inversion.  Vanishing of the two
Tate moments removes the constants in co-Poisson summation and gives
\begin{equation}
 \mathcal C\mathbf Z_+=\mathbf Z_+\mathcal R.         \tag{448}
\end{equation}
For \(QF(t)=tF(t)\), the calculation in (443), now before window
restriction, is
\[
 [Q,\mathbf Z_+]=\mathbf Z_+B_+,\qquad
 B_-=\mathcal RB_+\mathcal R.
\]
Differentiate (448) with \(Q\).  Using
\(\mathcal RQ=-Q\mathcal R\) gives
\[
 \{Q,\mathcal C\}\mathbf Z_+
 =\mathbf Z_+(B_+\mathcal R+\mathcal RB_+).
\]
Multiplication on the right by \(\mathcal R\), followed by
\(\mathbf Z_+\mathcal R=\mathcal C\mathbf Z_+\), yields
\begin{equation}
 \mathbf Z_+(B_++B_-)
 =(Q+\mathcal C Q\mathcal C)\mathbf Z_+.             \tag{449}
\end{equation}

For the right Mellin transform, the cosine functional equation is
\[
 \mathcal M(\mathcal Cu)(s)
 =\chi(s)\mathcal Mu(1-s),\qquad
 \chi(s)=\pi^{s-\frac12}
 \frac{\Gamma((1-s)/2)}{\Gamma(s/2)}.
\]
Since multiplication by \(\log x\) becomes
\(-\partial_s\), direct differentiation gives
\[
 \mathcal M\bigl((Q+\mathcal C Q\mathcal C)u\bigr)(s)
 =-\frac{\chi'}{\chi}(s)\mathcal Mu(s).
\]
On \(s=\frac12+i\tau\),
\begin{equation}
 -\frac{\chi'}{\chi}\left(\frac12+i\tau\right)
 =\Re\psi\left(\frac14+\frac{i\tau}{2}\right)-\log\pi
 =g_\Gamma(\tau)-m_0.                               \tag{450}
\end{equation}
Hence \(Q+\mathcal C Q\mathcal C=D_\Gamma\).  Substitution in (449)
proves the first identity in (447).  Since
\(\mathcal A_\infty=D_\Gamma-B_+-B_-\), subtraction gives the commutator
identity.
\end{proof}

\begin{proposition}[Candidate localization defect of the commutator ---
\AUDIT]
Let \(F\in C_c^\infty(I_T)\cap\mathcal P_T\), put
\(h=\mathbf Z_+F\), and let \(P_T\) denote restriction to \(I_T\).
If the regularized composition in (447) has the physical causal boundary
obtained by summing its incoming tail, then
\begin{equation}
 P_T[\mathbf Z_+,D_\Gamma]F
 =\mathcal Z_{T,+}P_T\mathcal A_\infty F
  +\mathcal L_{T,-}(F),                              \tag{451}
\end{equation}
where the candidate localization defect is the past-leakage boundary
\begin{equation}
 \mathcal L_{T,-}(F)(t)
 :=\sum_{\substack{n\ge2\\t-\log n\le -T}}
       n^{-1/2}(\mathcal A_\infty F)(t-\log n),
 \qquad t\in I_T.                                   \tag{452}
\end{equation}
Consequently a zero mode \(P_T\mathcal A_\infty F=r_\varepsilon\) obeys
\begin{equation}
 \boxed{
 P_T[\mathbf Z_+,D_\Gamma]F
 =\mathcal Z_{T,+}r_\varepsilon+\mathcal L_{T,-}(F).} \tag{453}
\end{equation}
Under the stated boundary-realization hypothesis there is no unidentified
interior correction: failure of the global commutator to become the local
gap identity is exactly the incoming Gamma--Euler history (452).
\end{proposition}

\begin{proof}
The first identity in (447) says, at the meromorphic-germ level, that the
left side of (451) is the regularization of
\(P_T\mathbf Z_+\mathcal A_\infty F\).  Under the boundary-realization
hypothesis, split its causal boundary according as
\(t-\log n\in I_T\) or \(t-\log n\le-T\).  The first
part is exactly
\(\mathcal Z_{T,+}P_T\mathcal A_\infty F\); no point
\(t-\log n\ge T\) occurs because all delays are nonnegative.  The second
part is (452), proving the conditional statement (451).  Substitution of
the zero-mode equation gives (453).
\end{proof}

\paragraph{Scope of the commutator bridge --- \AUDIT.}
Equation (447) is the direct meromorphic Gamma--Fourier--Tate identity
sought by the reverse-engineering program: it derives the normalization
rather than postulate a metric.  Equations (451)--(453) identify its only
possible physical localization defect, but remain \AUDIT{} because
\(\mathcal A_\infty F\) is not compactly supported and the subsequent
co-Poisson sum is not automatically locally finite.  To prove (446) one
must first construct that causal boundary and then show, for an equality
state, that the incoming history (452) has the
conservative boundary value which cancels
\(\mathcal Z_{T,+}r_\varepsilon\) and forces the annular amplitude \(h\)
to vanish.  Dropping \(\mathcal L_{T,-}\) would be exactly the forbidden
local/global compression error.  This is now the only term in this route
which is not fixed by an identity.

\paragraph{Arithmetic content of the domain estimate --- \AUDIT.}
Estimate (434) is a weighted Paley--Wiener embedding for the measure
\[
 d\mu_\zeta(\tau)=|\zeta(\tfrac12+i\tau)|^2\,d\tau.
\]
It is not supplied by the classical long-interval second moment alone:
Paley--Wiener functions of type $T$ can concentrate on fixed-size frequency
windows.  Testing (434) on modulated compact sources, with the two Tate
moments removed by a finite-rank correction, requires uniform smoothed local
second-moment bounds at the Gamma scale $\log(2+|\tau|)$.  No such uniform
estimate has been proved in this manuscript.  Hence (431)--(434) genuinely
resolve the operator-theoretic closability issue but do not establish
first-loss membership or the amplified gap (358).

\subsection{Final analytic hygiene of the leakage and Sonin routes}

\begin{theorem}[Global critical-pole rigidity --- \PROVED]
Let \(F\in L^2(I_T)\), extended by zero, and put
\[
 \widehat F(\tau)=\int_{-T}^{T}F(t)e^{i\tau t}\,dt.
\]
If the meromorphic boundary product
\begin{equation}
 \frac{\zeta'}{\zeta}\left(\frac12+i\tau\right)
       \widehat F(\tau)
 \in L^2_{\rm loc}(\mathbb R),                       \tag{454}
\end{equation}
where values at its isolated poles are immaterial, then \(F=0\).
The assertion is unconditional.
\end{theorem}

\begin{proof}
At every simple critical zero
\(\rho=\frac12+i\gamma\),
\[
 \frac{\zeta'}{\zeta}\left(\frac12+i\tau\right)
 =\frac{1}{i(\tau-\gamma)}+O(1).
\]
Condition (454) therefore forces \(\widehat F(\gamma)=0\); otherwise the
squared modulus has a nonintegrable
\(|\tau-\gamma|^{-2}\) singularity.  Heath--Brown's observation applied to
Levinson's theorem gives \(\gg H\log H\) simple zeros on the critical line
with \(0<\gamma\le H\); see D.~R. Heath--Brown,
\emph{Bull. London Math. Soc.} 11 (1979), 17--18,
doi:10.1112/blms/11.1.17.

Paley--Wiener makes \(\widehat F\) an entire Cartwright-class function of
exponential type at most \(T\).  A nonzero such function has only \(O_T(H)\)
zeros, counted with multiplicity, in \(|z|\le H\).  The preceding
\(\gg H\log H\) distinct real zeros are impossible.  Hence
\(\widehat F\equiv0\), and Fourier injectivity gives \(F=0\).
\end{proof}

\begin{proposition}[Why global pole rigidity does not apply to (452) ---
\PROVED]
The incoming localization defect \(\mathcal L_{T,-}\) in (452) is not the
restriction of the translation-invariant von Mangoldt operator whose
meromorphic multiplier is
\(-\zeta'/\zeta(\frac12+i\tau)\).  In particular,
\[
 \mathcal L_{T,-}(F)\in L^2(I_T)
 \quad\not\Longleftrightarrow\quad
 \frac{\zeta'}{\zeta}(\tfrac12+i\,\cdot)\widehat F
 \in L^2(\mathbb R).                                 \tag{455}
\]
Thus Theorem (454) is a sharp exclusion of a stronger global regularity
shortcut, not a proof of the physical boundary statement (451).
\end{proposition}

\begin{proof}
The global oriented von Mangoldt operator is a convolution/translation sum
acting directly on \(F\), so it has the stated meromorphic multiplier.
By contrast, (452) first applies \(\mathcal A_\infty\), then sums with the
co-Poisson weights \(n^{-1/2}\), retains only indices satisfying the moving
condition \(t-\log n\le-T\), and finally observes only \(t\in I_T\).
The indicator depends jointly on \(t\) and \(n\), so the resulting map is
not translation invariant and cannot possess that scalar Fourier
multiplier.  This proves (455).
\end{proof}

\begin{theorem}[Burnol parity audit: reciprocal membership is the missing
gap --- \PROVED]
Let \(F\in C_c^\infty(I_T)\cap\mathcal P_T\) have parity
\(\mathcal RF=\varepsilon F\), let
\[
 u=\mathfrak S_T^0F,\qquad a=e^{-T},\qquad b=a^{-1}=e^T.
\]
Then \(u\in P_a\cap K_a\),
\begin{equation}
 u=\mathcal Cu=0\quad(0<x<a),\qquad
 \mathcal Cu=\varepsilon u,                          \tag{456}
\end{equation}
and the following are equivalent:
\begin{enumerate}
\item \(u\in L_b\);
\item \(u\) and \(\mathcal Cu\) are constant on \((0,b)\);
\item \(u=0\) on \((0,b)\);
\item \(u=0\) on the annulus \((a,b)\).
\end{enumerate}
Consequently the assertion that the zero-mode equation makes
\(\mathcal Cu\) constant on \((a,b)\), with the constant matched to the
inner Tate constant, is already the missing implication (441), not an
additional consequence of co-Poisson intertwining.
\end{theorem}

\begin{proof}
Equations (354)--(355) give (456).  The equivalence of the first two items
is the definition of the extended Sonin space \(L_b\).  Since both functions
already vanish on the nonempty interval \((0,a)\), being constant on the
connected interval \((0,b)\) forces both constants to be zero.  This proves
\((2)\Rightarrow(3)\), while \((3)\Rightarrow(2)\) follows from
\(\mathcal Cu=\varepsilon u\).  Finally (456) makes vanishing on
\((0,b)\) equivalent to vanishing on the remaining annulus \((a,b)\).
\end{proof}

\paragraph{Frozen analytic endpoint --- \OPEN.}
The two primitive Tate moments have already been used in (456) to annihilate
the two inner co-Poisson constants.  They cannot be used again to match an
annular constant.  The analytic route is therefore frozen at exactly
\[
 P_TA_TF=0,\quad F\in\mathcal P_T
 \quad\Longrightarrow\quad
 \mathfrak S_T^{\rm loc}F=0\ \hbox{on }(e^{-T},e^T),
 \tag{457}
\]
equivalently (446).  Neither global \(L^2\) leakage nor Burnol's
intertwining supplies (457).

\subsection{Finite Green audit: the first \(r=2\to3\) compatibility test}

Let \(V_S=\bigoplus_{n\in S}\mathbb Re_n\) and
\[
 E_S=\mathbb RF_1\oplus\mathbb RF_2\oplus V_S,\qquad
 q(F_1,F_1)=q(F_2,F_2)=0,\quad q(F_1,F_2)=1.
\]
For \(x\in V_S\) write
\[
 a_S(x):=q(x,F_1),\qquad b_S(x):=q(x,F_2),\qquad
 C_S(x,y):=q(x,y).
\]

\begin{theorem}[Exact finite mixed-index classification --- \PROVED]
The doubly primitive lift of \(x\in V_S\) is
\begin{equation}
 M_Sx:=x-b_S(x)F_1-a_S(x)F_2,                        \tag{458}
\end{equation}
and its Gram form is
\begin{equation}
 q(M_Sx,M_Sy)
 =C_S(x,y)-a_S(x)b_S(y)-b_S(x)a_S(y).               \tag{459}
\end{equation}
Consequently the finite mixed-index inequality
\[
 q(M_Sx,M_Sx)\le0\qquad(x\in V_S)
\]
holds if and only if
\begin{equation}
 \boxed{
 N_S:=a_Sb_S^*+b_Sa_S^*-C_S\ge0.}                   \tag{460}
\end{equation}
For \(S\subset S'\), compatibility of \(q_{S'}\) with \(q_S\), including
the two degrees, is equivalent to principal restriction compatibility of
\(N_{S'}\) with \(N_S\).
\end{theorem}

\begin{proof}
Pairing (458) with \(F_1,F_2\) gives zero.  Expanding
\(q(M_Sx,M_Sy)\) and using the hyperbolic ruling matrix leaves exactly
(459).  Its negative is the form \(N_S\), proving (460).  Every term in
(460) restricts functorially under \(V_S\hookrightarrow V_{S'}\), proving
the final assertion.
\end{proof}

\begin{corollary}[No-go for the rank-one prime-fusion Green --- \PROVED]
Assume \(\dim V_S\ge2\), the two Tate degree covectors \(a_S,b_S\) are
linearly independent, and a proposed corrected correspondence Gram is
\begin{equation}
 C_S=vv^*,\qquad 0\ne v\in V_S^*.                    \tag{461}
\end{equation}
Then the mixed-index inequality is impossible.  In particular, for two
distinct prime correspondences the candidate
\[
 v_{p^k}=\sqrt{\log p},\qquad
 C_S(e_{p^j},e_{q^k})=\sqrt{\log p\,\log q}
\]
cannot be the required Green completion once both independent Tate degrees
are retained.
\end{corollary}

\begin{proof}
If (460) held, then for every
\(x\in\ker a_S\cap\ker b_S\),
\[
 0\ge q(M_Sx,M_Sx)=|v(x)|^2,
\]
so \(v\in\operatorname{span}\{a_S,b_S\}\).  Write
\(v=\alpha a_S+\beta b_S\).  Since \(a_S,b_S\) are independent, choose
\(x\) with \(b_S(x)=0\), \(a_S(x)\ne0\).  Formula (459) gives
\(|\alpha a_S(x)|^2\le0\), hence \(\alpha=0\).  Interchanging \(a_S,b_S\)
gives \(\beta=0\), contradicting \(v\ne0\).
\end{proof}

\begin{corollary}[Compatibility is not the geometric gate --- \PROVED]
Every abstract finite solution of the mixed-index inequality is of the form
\begin{equation}
 C_S=a_Sb_S^*+b_Sa_S^*-N_S,\qquad N_S\ge0.           \tag{462}
\end{equation}
Compatible positive kernels \(N_S\) produce compatible solutions for all
finite \(S\).  Conversely, if exact comparison with row~(c) is imposed,
then \(C_S\) is fixed by the restriction of \(B_{\rm nuc}\), the Green term
is forced to be \(G_\infty\), and (460) is exactly row~(d) on \(V_S\).
Therefore choosing \(N_S\ge0\), factoring it, or fitting its entries is
circular unless \(N_S\) is independently produced by a mixed
Riemann--Roch/effectivity theory.
\end{corollary}

\begin{proof}
Equation (462) is (460) solved for \(C_S\), and restriction compatibility
is the last assertion of the theorem.  The forced-Green theorem in
\texttt{main.tex} gives
\(B_{\rm nuc}=K+G_\infty\); hence exact comparison removes all freedom in
\(C_S\).  Positivity of the resulting \(N_S\) is then precisely (459) with
the desired sign.
\end{proof}

\paragraph{Outcome of the \(r=2\to3\) test --- \AUDIT.}
The minimal rank-one fusion is ruled out already at \(r=2\), before its
restriction from \(r=3\) is considered.  General compatible matrices do
exist abstractly, so compatibility alone yields neither a construction nor
a no-go for row~(d).  The first genuinely load-bearing datum is provenance:
construct the \emph{forced} kernel
\[
 N_S=a_Sb_S^*+b_Sa_S^*
     -(K_S+G_{\infty,S})
\]
as a section/cohomology Gram whose positivity follows from independently
defined effectivity and quadratic Riemann--Roch.  This is exactly the
mixed-section forcing gate already isolated in \texttt{main.tex}.

\begin{theorem}[The constructed quadratic line is a pairing line, not a
mixed divisor class --- \PROVED]
Let \(\lambda(f,g)\) denote the based normed line assembled in
\texttt{main.tex} from finite contact and the archimedean metric, so that
\begin{equation}
 -\log\operatorname{covol}\lambda(f,g)=B_{\rm nuc}(f,g) \tag{463}
\end{equation}
after polarization.  The assignment called \(D_f=\lambda(f,f)\) there
cannot be an additive Picard-valued mixed divisor functor unless
\(B_{\rm nuc}=0\).  More precisely, in general
\begin{align}
 \lambda(f+g,f+g)
 &\not\simeq
 \lambda(f,f)\otimes\lambda(g,g),                    \tag{464}\\
 \lambda(nf,nf)
 &\not\simeq \lambda(f,f)^{\otimes n}\qquad(n\ge2).  \tag{465}
\end{align}
What has been constructed is the self-pairing line
\(\langle D_f,D_f\rangle\), or equivalently a symmetric biextension
\((f,g)\mapsto\lambda(f,g)\), not the underlying class \(D_f\) whose
sections and tensor powers enter Riemann--Roch.
\end{theorem}

\begin{proof}
If \(f\mapsto\lambda(f,f)\) were additive and the isomorphism in (464)
preserved the assembled metric, taking logarithmic covolumes would give
\[
 B_{\rm nuc}(f+g,f+g)
 =B_{\rm nuc}(f,f)+B_{\rm nuc}(g,g).
\]
Polarization makes the left side equal to the right side plus
\(2\operatorname{Re}B_{\rm nuc}(f,g)\).  Thus the cross form would vanish
for all \(f,g\), forcing the Hermitian form \(B_{\rm nuc}\) to be zero.
Likewise (463) gives quadratic scaling
\[
 -\log\operatorname{covol}\lambda(nf,nf)
 =n^2B_{\rm nuc}(f,f),
\]
whereas the \(n\)-fold tensor power has exponent
\(nB_{\rm nuc}(f,f)\), proving (465) whenever the self-pairing is nonzero.
The bilinearity of \(\lambda(f,g)\), rather than additivity of its diagonal,
is exactly the typing of a pairing/biextension.
\end{proof}

\paragraph{Corrected geometric gate --- \OPEN.}
Before a mixed-section forcing theorem can be applied, one must construct an
additive Picard or stable-category object
\[
 \mathscr D:\mathcal T\longrightarrow\mathsf{Pic}_{\rm mix},
 \qquad
 \mathscr D(f+g)\simeq\mathscr D(f)\otimes\mathscr D(g),          \tag{466}
\]
with genuine section/cohomology objects, together with a comparison
\[
 \langle\mathscr D(f),\mathscr D(g)\rangle_{\rm Del}
 \simeq\lambda(f,g).                                  \tag{467}
\]
Only after (466)--(467) is constructed independently may quadratic
Riemann--Roch and effectivity prove positivity of the forced kernels
\(N_S\).  Defining \(\mathscr D(f)\) from the diagonal line
\(\lambda(f,f)\), or defining its cohomology from the sign of (463), would
be circular or ill-typed.

\subsection{The canonical leakage Gram and the corrected \(\frac12\log5\)
certificate}

Let a self-adjoint operator \(A\) be written relative to an orthogonal
decomposition
\[
 \mathcal H=D\oplus S\oplus Q.
\]
The letters have the endpoint meaning used in
\texttt{row-d-local-analysis.tex}: \(D\) is the last five-dimensional
graph, \(S\) is the already eliminated finite safe block, and \(Q\) is the
whole infinite complement.

\begin{theorem}[Nested shorting and the source-defined leakage Gram ---
\PROVED]
Assume \(A_{SS}>0\), and define
\begin{align}
 K_D&:=A_{DD}-A_{DS}A_{SS}^{-1}A_{SD},\tag{468}\\
 \widehat A_{QQ}&:=A_{QQ}-A_{QS}A_{SS}^{-1}A_{SQ},\tag{469}\\
 C_D&:=A_{DQ}-A_{DS}A_{SS}^{-1}A_{SQ}.\tag{470}
\end{align}
If \(\widehat A_{QQ}>0\), put
\begin{equation}
 \mathscr O_{D\mid S}:=\widehat A_{QQ}^{-1/2}C_D^*,\qquad
 \mathscr L_{D\mid S}:=\mathscr O_{D\mid S}^*\mathscr O_{D\mid S}
 =C_D\widehat A_{QQ}^{-1}C_D^*.                    \tag{471}
\end{equation}
Then
\begin{equation}
 A\ge0
 \quad\Longleftrightarrow\quad
 \widehat A_{QQ}>0,qquad K_D-\mathscr L_{D\mid S}\ge0,          \tag{472}
\end{equation}
with the usual range/Moore--Penrose formulation if a diagonal block is
only semidefinite.  In particular, \(\mathscr L_{D\mid S}\) is a canonical
positive Gram derived from the Gamma--Euler operator and the chosen
orthogonal splitting; it is not a square root defined from the desired
remainder \(K_D-\mathscr L_{D\mid S}\).
\end{theorem}

\begin{proof}
Congruence elimination of the \(S\)-block sends \(A\) to the orthogonal
sum of \(A_{SS}\) and
\[
 \begin{pmatrix}K_D&C_D\\ C_D^*&\widehat A_{QQ}\end{pmatrix}
 \quad\hbox{on }D\oplus Q.
\]
Since \(A_{SS}>0\), this congruence preserves positivity.  Shorting the
remaining positive \(Q\)-block gives exactly
\(K_D-C_D\widehat A_{QQ}^{-1}C_D^*\).  Formula (471) is its canonical
Gram factorization.  This proves (472).
\end{proof}

\begin{corollary}[Minimal directed certificate --- \PROVED]
Suppose
\begin{equation}
 A_{QQ}\ge\delta I,qquad
 \kappa:=\|A_{QS}A_{SS}^{-1/2}\|^2<\delta.           \tag{473}
\end{equation}
Then \(\widehat A_{QQ}\ge(\delta-\kappa)I\), and the source-derived
sufficient test
\begin{equation}
 K_D-(\delta-\kappa)^{-1}C_DC_D^*>0                 \tag{474}
\end{equation}
implies \(A>0\).
For \(T=\frac12\log5\), the certified value in the supplied endpoint
calculation is \(\delta=0.218\).  Thus (473)--(474), with the actual
\(S_{163}\) block and the corrected cross (470), are precisely the two
missing interval enclosures at that endpoint.
\end{corollary}

\begin{proof}
For \(q\in Q\),
\[
 \langle A_{QS}A_{SS}^{-1}A_{SQ}q,q\rangle
 =\|A_{SS}^{-1/2}A_{SQ}q\|^2\le\kappa\|q\|^2.
\]
Together with \(A_{QQ}\ge\delta I\), this proves the first assertion.
Equation (474) bounds the Gram (471) from above and (472) concludes.
\end{proof}

\paragraph{Endpoint status --- \OPEN.}
The existing five-column continuum Gram computes the direct \(D\)-to-\(Q\)
action but does not contain \(A_{SQ}\).  Consequently it determines neither
\(\kappa\) nor the corrected cross \(C_D\), and its positive directed
Gershgorin numbers cannot be substituted for (473)--(474).  No serialized
\(S_{163}\)-to-complement action data were supplied with the manuscript, so
the two enclosures cannot be reconstructed from the printed five-column
data.  This is a finite, reproducible certification task for this endpoint;
it does not by itself prove the all-window statement (457) or row~(d).

\paragraph{Relation to the geometric kernel --- \AUDIT.}
The leakage Gram (471) is the independently defined positive object sought
by reverse engineering at a fixed positive-complement splitting.  It has
the same Schur role as \(N_S\) in (460), but no comparison between them is
currently proved.  Declaring them equal would require exactly the missing
additive mixed class and Deligne-pairing comparison (466)--(467), or an
analytic all-window intertwiner.  Thus (471) supplies the correct object
without supplying the still-open comparison.

\begin{theorem}[Formal Picard lift versus geometric realization ---
\PROVED]
Assume the polarized line \((f,g)\mapsto\lambda(f,g)\) has the biextension
coherences proved in row~(a).  Then equations (466)--(467) admit a tautological
\emph{formal} solution: take the discrete Picard groupoid with objects
\(f\in\mathcal T\), tensor law \(f\otimes g=f+g\), put
\begin{equation}
 \mathscr D_{\rm for}(f):=[f],                       \tag{475}
\end{equation}
and define the pairing of two generators to be \(\lambda(f,g)\).
This formal construction supplies no section functor, effective cone or
Riemann--Roch theorem.

Moreover it cannot be realized faithfully by the factorized external
Picard groupoid already constructed in row~(a).  That groupoid has completed
numerical class space
\begin{equation}
 N^1_{\deg}\simeq\mathbb R^2,\qquad
 [f]\longmapsto\bigl(\widehat f(1),\widehat f(0)\bigr),           \tag{476}
\end{equation}
so every primitive \(f\in\mathcal T^0\) maps to zero.  A pairing descending
through (476) therefore vanishes on
\(\mathcal T^0\times\mathcal T^0\), whereas
\begin{equation}
 -\log\operatorname{covol}\lambda(f,g)=B_{\rm nuc}(f,g)          \tag{477}
\end{equation}
does not vanish identically there.  Hence the required mixed realization
must strictly enlarge the existing external divisor category and must carry
new, independently defined section/effectivity data.
\end{theorem}

\begin{proof}
The first assertion is the universal discrete Picard construction; the
biextension coherence of \(\lambda\) makes its pairing symmetric monoidal.
For the second assertion, Proposition
\texttt{prop:externaleffectivity} of row~(a) identifies the external
numerical category with the two degree coordinates.  By definition both
coordinates vanish on \(\mathcal T^0\).  Any bilinear pairing factoring
through that quotient is therefore zero on the primitive square.  The
localized operator in row~(d), already strictly positive for
\(0<T\le\log2\), supplies primitive test vectors with
\(B_{\rm nuc}(f,f)\ne0\).  Equation (477) rules out such a factorization.
\end{proof}

\paragraph{Sharpened geometric endpoint --- \OPEN.}
The missing datum is not a formal additive symbol and not a formal pairing.
It is a faithful realization of (475) in an enlarged mixed category with
genuine sections or cohomology, an effective cone, quadratic
Riemann--Roch, and comparison (477), all constructed without using the sign
of \(B_{\rm nuc}\).  These hypotheses would activate the already proved
mixed-section forcing theorem; none is present in rows~(a)--(c).

\begin{theorem}[Leakage Grams obey an innovation cocycle, not principal
restriction --- \PROVED]
Let
\(\mathcal H=P\oplus R\oplus Q\), let \(A=A^*\), and assume all blocks
shorted below are strictly positive.  First eliminate \(Q\), writing
\begin{equation}
 \widetilde A_{UV}:=A_{UV}-A_{UQ}A_{QQ}^{-1}A_{QV},
 \qquad U,V\in\{P,R\}.                              \tag{478}
\end{equation}
Let \(\mathscr L_{P\oplus R\mid Q}\) be the leakage Gram from
\(P\oplus R\) to \(Q\), and let
\(\mathscr L_{P\mid R\oplus Q}\) be the leakage Gram from \(P\) to its
whole complement.  Then
\begin{equation}
 \boxed{
 \mathscr L_{P\mid R\oplus Q}
 =P\mathscr L_{P\oplus R\mid Q}P
  +\widetilde A_{PR}\widetilde A_{RR}^{-1}widetilde A_{RP}.}   \tag{479}
\end{equation}
The second summand is positive.  Therefore the family of canonical leakage
Grams has principal restriction compatibility if and only if every added
innovation coupling \(\widetilde A_{RP}\) vanishes.
\end{theorem}

\begin{proof}
Shorting \(Q\) first gives the operator
\(\widetilde A\) on \(P\oplus R\).  Shorting its \(R\)-block gives
\[
 A_{PP}-P\mathscr L_{P\oplus R\mid Q}P
 -\widetilde A_{PR}\widetilde A_{RR}^{-1}\widetilde A_{RP}.
\]
Associativity of Schur complementation says this equals
\(A_{PP}-\mathscr L_{P\mid R\oplus Q}\), proving (479).  The last
summand is
\((\widetilde A_{RR}^{-1/2}\widetilde A_{RP})^*
  (\widetilde A_{RR}^{-1/2}\widetilde A_{RP})\), so it is positive and
vanishes exactly under the stated condition.
\end{proof}

\begin{corollary}[No direct leakage--Green identification --- \PROVED]
The compatible geometric kernels \(N_S\) of (460) cannot be identified,
level by level, with the canonical exterior leakage Grams (471) unless all
successive Schur innovations vanish.  At every threshold where the supplied
old--new block is nonzero, the required object, if it exists, must retain
the positive innovation increments
\begin{equation}
 \mathscr I_{R\mid P,Q}
 :=\widetilde A_{PR}\widetilde A_{RR}^{-1}widetilde A_{RP}\ge0          \tag{480}
\end{equation}
as coherent transition data.  A static identification
\(N_S=\mathscr L_S\) is not the missing comparison.
\end{corollary}

\begin{proof}
Theorem (458)--(460) proves that \(N_{S'}\) restricts exactly to \(N_S\).
Equation (479) proves that leakage restriction differs by (480).  The two
laws agree only when all such increments vanish.  The old--new blocks
\(B_N\) in the threshold-capacity formula of the supplied row~(d) file are
the corresponding innovation couplings; whenever such a block is nonzero,
the extra term cannot be omitted.
\end{proof}

\paragraph{Corrected reverse-engineering target --- \OPEN.}
A viable bridge must be a filtered mixed section object whose transition
maps carry the cocycle (480), while its total Deligne pairing restricts to
the fixed kernels (460).  Positivity of each innovation is automatic from
its source-defined Gram; the unproved statement is the coherent comparison
which makes their accumulated shorting equal to the forced row-(c) pairing
without defining the transition norm from the desired sign.  This is more
specific than asking for an unspecified Green operator.

\begin{theorem}[Minimal positive completion of the innovation cocycle ---
\PROVED]
Let \(P_1\subset P_2\subset\cdots\) be a finite filtration, and suppose the
canonical leakage Grams \(L_n\ge0\) satisfy
\begin{equation}
 L_n=P_nL_{n+1}P_n+I_n,qquad I_n\ge0,               \tag{481}
\end{equation}
as in (479).  Define stored kernels recursively by
\begin{equation}
 E_1=0,qquad
 E_{n+1}:=j_n(E_n+I_n)j_n^*,                        \tag{482}
\end{equation}
where \(j_n:P_n\hookrightarrow P_{n+1}\) and the right side is extended by
zero on \(P_{n+1}\ominus P_n\).  Then
\begin{equation}
 N_n^{\rm inn}:=L_n+E_n\ge0,qquad
 P_nN_{n+1}^{\rm inn}P_n=N_n^{\rm inn}.             \tag{483}
\end{equation}
Consequently the \(N_n^{\rm inn}\) are the finite restrictions of one
positive kernel on the algebraic inductive limit \(\bigcup_nP_n\), and admit
a Kolmogorov/Gram realization obtained by adjoining the square-root feature
of each source-defined innovation \(I_n\).
\end{theorem}

\begin{proof}
Positivity of \(E_n\) follows inductively from (482), so (483) is positive.
Using (481)--(482),
\[
 P_nN_{n+1}^{\rm inn}P_n
 =P_nL_{n+1}P_n+E_n+I_n=L_n+E_n.
\]
This is the compatibility in (483).  A compatible family of positive
finite kernels has its standard Kolmogorov realization on the algebraic
inductive limit.  Equivalently, factor each \(I_n\) positively and take the
orthogonal sum of the resulting innovation feature spaces.
\end{proof}

\begin{corollary}[The minimal innovation completion has the wrong
normalization --- \PROVED]
Let the forced primitive kernel on \(P_n\) be the compression
\(N_{S_n}^{\rm geom}=P_nAP_n\), with signs chosen as in (460).  At the base
level of (482),
\begin{equation}
 N_{S_1}^{\rm geom}-N_1^{\rm inn}
 =A_{P_1P_1}-A_{P_1Q_1}A_{Q_1Q_1}^{-1}A_{Q_1P_1}
 =A/A_{Q_1}.                                         \tag{484}
\end{equation}
Thus the desired normalization
\(N_{S_n}^{\rm geom}=N_n^{\rm inn}\) already fails whenever the base
Schur complement is nonzero.  In particular it fails for the nonzero base
blocks in the directed Feshbach splittings used to certify
\(0<T\le\log2\), where that Schur complement was proved strictly positive.

More generally, initializing (482) with the value needed to repair (484)
requires setting \(E_1=A/A_{Q_1}\).  Positivity of this initial storage is
exactly the Schur inequality being sought.  Hence the innovation completion
does not furnish a non-circular proof of row~(d).
\end{corollary}

\begin{proof}
At level one, (482) gives \(E_1=0\), and the definition of the leakage Gram
gives
\(N_1^{\rm inn}=A_{P_1Q_1}A_{Q_1Q_1}^{-1}A_{Q_1P_1}\).  Subtraction is
the Schur complement (484).  In each supplied certified Feshbach splitting,
the directed lower bound proves this particular complement strictly
positive.  The final assertion follows because any alternative base storage
satisfying the normalization must equal the difference in (484).
\end{proof}

\paragraph{Outcome of the innovation construction --- \AUDIT.}
Equations (481)--(483) construct a legitimate positive compatible kernel,
but (484) proves that it is another positive model with the wrong character.
The calculation is useful as a no-go: neither positive leakage Grams nor
their coherent accumulation determine the forced row-(c) normalization.
The missing input must create the base Schur energy from independent
Gamma--Tate causality or mixed Riemann--Roch/effectivity; it cannot be
recovered from transition compatibility alone.

\subsection{Complete audit of the Abel--Hodge equality-state route}

Let \(F\in\mathcal P_T\) be a parity zero mode, extended by zero, and retain
the notation
\[
 K_F(z)=\int_{-T}^{T}F(u)e^{-zu}\,du,qquad
 Y_+(t)=\sum_{n\ge2}\frac{\Lambda(n)}{\sqrt n}F(t-\log n).
\]
The first-loss theorem, parity reduction and the two Tate moments used below
are already proved in (251)--(255).

\begin{theorem}[Exact Abel energy in its honest convergence half-plane ---
\PROVED]
For \(\sigma>\frac12\), one has
\begin{equation}
 e^{-\sigma t}Y_+(t)\in L^2(\mathbb R),              \tag{485}
\end{equation}
and Laplace--Plancherel gives
\begin{equation}
 \boxed{
 \int_{\mathbb R}e^{-2\sigma t}|Y_+(t)|^2\,dt
 =\frac1{2\pi}\int_{\mathbb R}
 \left|\frac{\zeta'}{\zeta}
   (\tfrac12+\sigma+i\tau)K_F(\sigma+i\tau)\right|^2d\tau.}   \tag{486}
\end{equation}
The Euler series used to derive (486) is absolutely convergent only in this
range.  In particular the frequently suggested assertion that the same
Euler construction is an ordinary positive current for every \(\sigma>0\)
is false.
\end{theorem}

\begin{proof}
For \(\Re z=\sigma>\frac12\), Tonelli gives
\[
 \int_{\mathbb R}Y_+(t)e^{-zt}\,dt
 =K_F(z)\sum_{n\ge2}\frac{\Lambda(n)}{n^{z+1/2}}
 =-\frac{\zeta'}{\zeta}(\tfrac12+z)K_F(z).
\]
The Dirichlet series is uniformly bounded on that vertical line by
\(\sum\Lambda(n)n^{-1/2-\sigma}<\infty\), while
\(K_F(\sigma+i\,\cdot)\in L^2\) by Fourier Plancherel applied to
\(e^{-\sigma u}F(u)\).  This proves (485), and Plancherel gives (486).
The abscissa condition is \(\Re(\frac12+z)>1\), namely
\(\sigma>\frac12\).
\end{proof}

The Tate equality \(K_F(1/2)=0\) cancels the pole of
\(\zeta'/\zeta\) at \(s=1\) in the meromorphic product.  It does not supply
an \(L^2\) continuation through the strip
\(0\le\Re z\le1/2\), where the nontrivial zero poles occur.

\begin{theorem}[Exact critical-pole energy blow-up --- \PROVED]
Let \(\rho=\frac12+i\gamma\) be a zero of \(\zeta\) of multiplicity
\(m_\rho\), and choose \(\delta>0\) so that no other critical-line zero
lies in \((\gamma-2\delta,\gamma+2\delta)\).  Define the meromorphic local
Abel energy
\begin{equation}
 E_{\rho,\sigma}(F):=
 \int_{\gamma-\delta}^{\gamma+\delta}
 \left|\frac{\zeta'}{\zeta}
 (\tfrac12+\sigma+i\tau)K_F(\sigma+i\tau)\right|^2d\tau.       \tag{487}
\end{equation}
If \(K_F(i\gamma)\ne0\), then
\begin{equation}
 \boxed{
 \lim_{\sigma\downarrow0}\sigma E_{\rho,\sigma}(F)
 =\pi m_\rho^2|K_F(i\gamma)|^2.}                    \tag{488}
\end{equation}
Consequently local uniform Abel energy at the critical boundary forces
\(K_F(i\gamma)=0\) at every critical zero.
\end{theorem}

\begin{proof}
Near \(\rho\),
\[
 \frac{\zeta'}{\zeta}(\tfrac12+\sigma+i\tau)
 =\frac{m_\rho}{\sigma+i(\tau-\gamma)}+O(1),
\]
and \(K_F(\sigma+i\tau)=K_F(i\gamma)+O(\sigma+|\tau-\gamma|)\).
The leading squared term integrates to
\[
 m_\rho^2|K_F(i\gamma)|^2
 \int_{-\delta}^{\delta}\frac{dx}{\sigma^2+x^2}
 =\frac{2m_\rho^2|K_F(i\gamma)|^2}{\sigma}
   \arctan(\delta/\sigma).
\]
After multiplication by \(\sigma\), all cross and bounded terms tend to
zero, proving (488).
\end{proof}

\begin{theorem}[The positive Markov reservoir hits the same abscissa barrier
--- \PROVED]
At Abel parameter \(\sigma\), the infinitesimal prime-power defect budget
obtained from (221)--(223) has weights
\begin{equation}
 w_{p,k}(\sigma)=\frac{\log p}{p^{k(1/2+\sigma)}}.    \tag{489}
\end{equation}
For every nonzero \(F\in L^2(I_T)\), the mutually orthogonal exterior
channels belonging to primes \(p\ge e^{2T}\) have total budget bounded below
by
\begin{equation}
 \|F\|_2^2\sum_{p\ge e^{2T}}
       \frac{\log p}{p^{1/2+\sigma}}.                \tag{490}
\end{equation}
It is finite for \(\sigma>1/2\) and infinite for
\(0<\sigma\le1/2\).  Hence the unsupertraced positive Markov defects do not
provide a Hilbert vector across the critical strip.  Tate centering must be
performed before completion, and boundedness of that coherent compensated
synthesis is precisely the open Hardy estimate (287).
\end{theorem}

\begin{proof}
For \(p\ge e^{2T}\), the original and translated supports are disjoint, so
each normalized difference channel has squared norm \(\|F\|_2^2\).  Keeping
only \(k=1\) gives (490).  The prime number theorem, or partial summation
against \(\vartheta(x)\sim x\), shows that
\(\sum_p(\log p)p^{-a}\) converges exactly when \(a>1\).  Here
\(a=1/2+\sigma\), proving the threshold.  Equations (288)--(291) identify
the required pre-completion centering with replacement of \(d\Psi\) by
\(dR_\Psi\); (286)--(287) identify boundedness of its causal synthesis.
\end{proof}

\begin{proposition}[What the equality equation determines --- \PROVED]
For a parity-\(\varepsilon\) zero mode, equations (296)--(298) determine
exactly
\begin{equation}
 Q_{+,T}(z)+\varepsilon Q_{+,T}(-z).                 \tag{491}
\end{equation}
They do not determine the complementary component
\begin{equation}
 Q_{+,T}(z)-\varepsilon Q_{+,T}(-z),                 \tag{492}
\end{equation}
which selects the causal right/left splitting.  G39 and the Tate moments
therefore determine the reflected symmetric amplitude but not the Abel or
Hardy norm of either oriented output.
\end{proposition}

\begin{proof}
This is the direct even/odd decomposition under the involution
\(J_\varepsilon Q(z)=\varepsilon Q(-z)\).  Equation (298) fixes
\((I+J_\varepsilon)Q_{+,T}\); its kernel is the range of
\(I-J_\varepsilon\), so (492) is independent algebraic data.  The right
tail in (285) depends on \(Q_{+,T}\), not only on (491), proving the last
assertion.
\end{proof}

\begin{theorem}[Changing the regularizer cannot remove the energy gate ---
\PROVED]
Let \(Y\in\mathcal S'(\mathbb R)\), and let \(R_\epsilon\) be any family of
Fourier-multiplier contractions on \(L^2\) whose symbols converge pointwise
to one and for which \(R_\epsilon Y\to Y\) in \(\mathcal S'\).  Then
\begin{equation}
 \sup_{0<\epsilon<1}\|R_\epsilon Y\|_2<\infty
 \quad\Longrightarrow\quad Y\in L^2(\mathbb R).     \tag{493}
\end{equation}
Conversely, if \(Y\in L^2\), the displayed bound follows from contractivity.
Thus Gaussian damping, heat damping, finite Euler cutoffs or another Abel
kernel can improve definition at fixed parameter, but a uniform energy
bound is exactly the missing \(L^2\) statement, not a way around it.
\end{theorem}

\begin{proof}
Uniform boundedness gives an \(L^2\)-weakly convergent subnet
\(R_{\epsilon_j}Y\rightharpoonup y\in L^2\).  Weak \(L^2\) convergence
implies distributional convergence to \(y\), while the hypothesis gives
distributional convergence to \(Y\).  Hence \(Y=y\in L^2\).  The converse
is immediate.
\end{proof}

\begin{theorem}[Local Abel equality-state criterion closes G --- \PROVED]
Suppose every parity first-loss mode \(F\) satisfies, for every \(R>0\),
\begin{equation}
 \sup_{0<\sigma<\sigma_R}
 \int_{-R}^{R}
 \left|\frac{\zeta'}{\zeta}
 (\tfrac12+\sigma+i\tau)K_F(\sigma+i\tau)\right|^2d\tau<\infty          \tag{494}
\end{equation}
for some \(\sigma_R>0\).  Then \(F=0\), and condition G and row~(d) follow.
This criterion is weaker than the global Hardy estimate (287), but it is not
supplied by the positive Markov reservoir or by the zero-mode amplitude.
\end{theorem}

\begin{proof}
The blow-up theorem (487)--(488) gives
\(K_F(i\gamma)=0\) at every critical zero.  The unconditional family of
simple critical zeros has counting function \(\gg H\log H\), whereas the
entire Cartwright function \(K_F\), of exponential type at most \(T\), has
only \(O_T(H)\) zeros if nonzero.  Thus \(K_F\equiv0\), and Laplace
injectivity gives \(F=0\).  The first-loss theorem then proves G and row~(d).
The final scope statement is Theorems (489)--(493).
\end{proof}

\begin{theorem}[Critical pole-defect measure --- \PROVED]
Let \(J\subset\mathbb R\) be a compact interval whose endpoints are not
ordinates of critical-line zeros, and define
\begin{equation}
 d\mu_{\sigma,F}(\tau):=\frac{\sigma}{\pi}
 \left|\frac{\zeta'}{\zeta}
 (\tfrac12+\sigma+i\tau)K_F(\sigma+i\tau)\right|^2d\tau.        \tag{495}
\end{equation}
Then, as \(\sigma\downarrow0\),
\begin{equation}
 \mu_{\sigma,F}|_J\ \buildrel *\over\rightharpoonup\
 \sum_{\substack{\rho=1/2+i\gamma\\\gamma\in J}}
 m_\rho^2|K_F(i\gamma)|^2\delta_\gamma.             \tag{496}
\end{equation}
In particular, the weaker condition
\begin{equation}
 \boxed{
 \int_J\left|\frac{\zeta'}{\zeta}
 (\tfrac12+\sigma+i\tau)K_F(\sigma+i\tau)\right|^2d\tau
 =o(\sigma^{-1})}                                    \tag{497}
\end{equation}
for every such \(J\) already forces \(F=0\) and closes G.
\end{theorem}

\begin{proof}
Only finitely many zeros lie above a fixed compact ordinate interval.
Choose disjoint discs around the critical zeros over \(J\).  In each disc,
the proof of (488), with a continuous test function inserted, gives weak
convergence to the stated atom; the normalized Cauchy kernel
\(\pi^{-1}\sigma(\sigma^2+x^2)^{-1}\) is an approximate identity.  On the
complement of those discs, the meromorphic logarithmic derivative and
\(K_F\) remain bounded for all sufficiently small \(\sigma\), because every
off-line zero has a fixed positive horizontal distance from the compact
central segment.  Multiplication by \(\sigma\) sends that part to zero.
This proves (496).  Condition (497) kills every atom, and the
Cartwright/simple-critical-zero argument of (494) gives \(F=0\).
\end{proof}

\paragraph{Reduced equality-state target --- \OPEN.}
The strongest previously isolated target was the Hardy bound (287), and
(494) only asked for local \(O(1)\) Abel energy.  The pole-defect theorem
shows that both can be replaced, for the first-loss contradiction, by the
strictly weaker little-o estimate (497).  A successful conservation theorem
need only show that a zero interior primitive mode carries no
\(1/\sigma\)-order critical Clark/residue mass.  It need not construct an
ordinary \(L^2\) exterior current.  No identity in the supplied manuscript
currently controls this leading residue energy; (491) still leaves the
causal component in which it resides undetermined.

\begin{theorem}[Derived co-Poisson divisibility criterion --- \PROVED]
Put
\begin{equation}
 G_F(s):=\zeta(s)K_F(s-\tfrac12).                    \tag{498}
\end{equation}
For the logarithmic generator \(Q\), the oriented co-Poisson commutator has
the meromorphic Mellin germ
\begin{equation}
 \mathcal M([Q,Z]F)(s)=-\zeta'(s)K_F(s-\tfrac12).    \tag{499}
\end{equation}
At a zero \(\rho\) of multiplicity \(m_\rho\), the following are
equivalent:
\begin{enumerate}
\item \(K_F(\rho-\tfrac12)=0\);
\item \(\operatorname{ord}_\rho G_F\ge m_\rho+1\);
\item the quotient of (499) by \(\zeta(s)\), namely
\begin{equation}
 -\frac{\zeta'}{\zeta}(s)K_F(s-\tfrac12),            \tag{500}
\end{equation}
is holomorphic at \(\rho\).
\end{enumerate}
Consequently G follows if every parity first-loss mode has an oriented
physical source \(U_F\) such that
\[
 [Q,Z]F=ZU_F
\]
and the Mellin transform of \(U_F\) realizes (500) on a connected domain
through the critical line.
\end{theorem}

\begin{proof}
Mellin transformation changes \(Z\) into multiplication by \(\zeta(s)\)
and \(Q\) into differentiation, with the sign fixed by convention.  The
product rule gives (499).  Write
\(\zeta(s)=(s-\rho)^{m_\rho}a(s)\), \(a(\rho)\ne0\).
Then (498) has order at least \(m_\rho+1\) exactly when \(K_F\) vanishes at
\(\rho-1/2\), while \(\zeta'/\zeta\) has a simple pole of residue
\(m_\rho\), proving the three-way equivalence.  A physical factorization
through \(Z\) makes (500) holomorphic at every critical zero.  The density
argument used in (494) then gives \(F=0\), and first loss gives G.
\end{proof}

\paragraph{Why the formal commutator is not yet the theorem --- \AUDIT.}
The algebraic identity \([Q,Z]=ZB_+\) already holds in the common
meromorphic calculus.  Dividing it by \(Z\) there gives (500), which is
meromorphic and may have exactly the forbidden poles.  Thus formal
factorization does not prove physical range membership.  Locally on the
window, \(Z_{T,+}\) is invertible by its finite nilpotent Neumann series,
but this does not construct the exterior causal boundary of \(B_+F\).
G39 and (491) determine the sum of the two reflected derived amplitudes;
the oriented range statement in (500) is their undetermined difference.
This is the second-order divisibility form of the same step~5, not an
additional gap.

\paragraph{Minimal Rellich--Clark construction target --- \OPEN.}
For a compact ordinate interval \(J\), let \(C_{\sigma,J}\) denote the
localized meromorphic analysis map
\[
 C_{\sigma,J}F(\tau)=1_J(\tau)
 \frac{\zeta'}{\zeta}(\tfrac12+\sigma+i\tau)
 K_F(\sigma+i\tau).
\]
It would suffice to construct, directly from the oriented Gamma--Tate and
co-Poisson currents, form-bounded operators \(X_{\sigma,J}\) on the physical
primitive space such that, at a first-contact parameter,
\begin{equation}
 \boxed{
 \frac{\sigma}{\pi}C_{\sigma,J}^*C_{\sigma,J}
 =H_{T_*}X_{\sigma,J}+X_{\sigma,J}^*H_{T_*}
   +R_{\sigma,J},qquad
 R_{\sigma,J}\buildrel{\rm form}\over\longrightarrow0.}       \tag{501}
\end{equation}
Indeed (501) evaluated on \(F_*\in\ker H_{T_*}\) gives (497), hence G.
Constructing \(X_{\sigma,J}\) by a Moore--Penrose solution of (501) would
be equivalent to requiring \(C_{\sigma,J}\ker H_{T_*}=o(\sigma^{-1/2})\)
and would therefore be circular.  The desired theorem is a source formula
for \(X_{\sigma,J}\), a local positive-commutator/flux identity for the
unsupertraced two-port system.  It is strictly weaker than the all-source
factorization (43) and the global Hardy estimate (287).

\begin{theorem}[Sylvester audit of the Rellich--Clark target --- \PROVED]
Let \(H=H^*\ge0\) have compact resolvent, let \(P\) project onto
\(\ker H\), and let \(B_\sigma\ge0\) be bounded.  There are bounded
self-adjoint \(X_\sigma\) and \(R_\sigma\) satisfying
\begin{equation}
 B_\sigma=HX_\sigma+X_\sigma H+R_\sigma,qquad
 R_\sigma=PB_\sigma P.                              \tag{502}
\end{equation}
Consequently
\begin{equation}
 R_\sigma\longrightarrow0
 \quad\Longleftrightarrow\quad
 \|B_\sigma^{1/2}P\|^2\longrightarrow0.             \tag{503}
\end{equation}
For
\(B_\sigma=\frac{\sigma}{\pi}C_{\sigma,J}^*C_{\sigma,J}\), condition
(503) is exactly the residue estimate (497) on the first-loss kernel.
Therefore solving (501) by spectral decomposition, a Sylvester inverse or a
Moore--Penrose formula is circular.
\end{theorem}

\begin{proof}
Compact resolvent gives
\(\mathcal H=\ker H\oplus\bigoplus_{j\ge1}\mathbb Ce_j\), with
\(He_j=\lambda_je_j\) and \(\inf_j\lambda_j>0\).  In this basis set the
matrix entries of \(X_\sigma\) to
\begin{equation}
 (X_\sigma)_{ij}=
 \begin{cases}
 (B_\sigma)_{ij}/(\lambda_i+\lambda_j),
   &\lambda_i+\lambda_j>0,\\
 0,&\lambda_i=\lambda_j=0.
 \end{cases}                                        \tag{504}
\end{equation}
Equivalently on the non-kernel blocks this is the convergent semigroup
integral
\(\int_0^\infty e^{-tH}(B_\sigma-PB_\sigma P)e^{-tH}\,dt\).
The positive spectral gap off the kernel gives boundedness, and entrywise
multiplication proves (502).  Since \(B_\sigma\ge0\),
\(\|PB_\sigma P\|=\|B_\sigma^{1/2}P\|^2\), proving (503).  Substitution of
the displayed \(B_\sigma\) gives the last assertion.
\end{proof}

\paragraph{Connection/Fisher audit --- \AUDIT.}
The logarithmic co-Poisson connection in G39 is linear in
\(\zeta'/\zeta\), whereas \(B_\sigma\) is its positive quadratic score
\(C_{\sigma,J}^*C_{\sigma,J}\).  Vanishing of the primitive first variation
fixes the connection amplitude but does not imply vanishing of its Fisher
energy.  Turning the former into the latter requires a source-derived
flatness or flux theorem.  Equation (502) shows that an abstract
positive-commutator construction cannot provide that theorem: its unremoved
kernel block is precisely the desired residue mass.

\paragraph{First contact is not stationarity --- \AUDIT.}
The equality \(H_{T_*}F_*=0\) says that the lowest eigenvalue vanishes at
the first contact; it does not say that its shape derivative vanishes.
Zero extension makes that eigenvalue nonincreasing, so a transverse first
contact may have \(\lambda_1'(T_*)<0\).  A Rellich identity obtained by
differentiating the moving window therefore contains an outward boundary
flux, naturally proportional to \(-\lambda_1'(T_*)\), and this term cannot
be placed in \(R_{\sigma,J}\to0\) without a new theorem.  Support saturation
(273) rules out deleting it by an empty boundary collar.  Thus a valid proof
of (501) must identify the residue flux with a Gamma--Tate terminal which
vanishes for structural reasons; first contact alone supplies no such
vanishing.

\begin{proposition}[Clamping does not annihilate the Tate terminal ---
\PROVED]
For every \(T>0\), let
\begin{equation}
 \phi_T(t)=(T^2-t^2)^2,qquad
 F_T=(\partial_t^2-\tfrac14)\phi_T.                  \tag{505}
\end{equation}
Then \(\phi_T(\pm T)=\phi_T'(\pm T)=0\), hence
\(F_T\in\mathcal P_T\), but its moving Tate terminal (109) is strictly
positive.  More explicitly,
\begin{equation}
 B_1'F_T=\binom{u_T}{u_T},qquad
 u_T=\frac12\int_{-T}^{T}(T^2-t^2)^2\cosh(t/2)\,dt>0,             \tag{506}
\end{equation}
and
\begin{equation}
 \|\nu_1(F_T)\|^2=\frac{2u_T^2}{2\sinh T+2T}>0.          \tag{507}
\end{equation}
Thus neither the two primitive moments nor the clamped boundary conditions
make the terminal energy vanish.
\end{proposition}

\begin{proof}
The four boundary equalities are immediate.  The primitive-potential
isomorphism (90)--(91) gives \(F_T\in\mathcal P_T\).  Formula (111a), parity
and positivity of \(\phi_T\) in the interior give (506).  Substitution in
the even term of (111), with \(a_1=2\sinh T\) and \(b=2T\), gives (507).
\end{proof}

\begin{theorem}[Residue-signature obstruction to terminal closure ---
\PROVED]
For every primitive \(F\), the already constructed residue boundary
(122)--(125) gives
\begin{align}
 \langle A_TF,F\rangle
 &=\sum_{\rho=1/2+i\gamma}m_\rho|\Phi_F(\gamma)|^2\notag\\
 &\quad+\sum_{\{\rho,1-\bar\rho\},\ \Re\rho\ne1/2}
 2m_\rho\Re\!\left(
   \Phi_F(z_\rho)\overline{\Phi_F(\bar z_\rho)}\right).         \tag{508}
\end{align}
Each off-line summand has signature \((1,1)\).  Consequently
\(H_TF=0\) implies only that the positive critical residue energy in the
first line is balanced by the indefinite off-line fibers; it does not force
either part to vanish.  The Tate terminals remove the polar characters at
\(s=0,1\) and do not alter these nontrivial-zero fibers.

Therefore primitiveness, clamping, the scalar zero-mode equality and the
already computed Tate terminal do not by themselves set the critical
pole-defect measure (496) equal to zero.  A source identity using G39 would
still have to prove that its extra oriented amplitude excludes the negative
directions in every off-line hyperbolic residue plane; that is exactly the
Hodge-index content of row~(d), not a consequence established above.
\end{theorem}

\begin{proof}
Equation (508) is (122)--(125), with multiplicities restored.  At a
first-loss vector its left side is zero.  Formula (125) exhibits one positive
and one negative square for every off-line orbit, so zero total Krein norm
does not imply vanishing of the positive critical coordinates.  The polar
and nontrivial residue fibers are disjoint in the contour decomposition;
the primitive Tate conditions cancel only the former.  This proves every
assertion.
\end{proof}

\begin{theorem}[Canonical negative residue port and logarithmic vertical
flux --- \PROVED]
Choose from every off-line functional-equation orbit the representative
\(\rho=\frac12+\alpha+i\gamma\) with \(\alpha>0\), and define
\begin{equation}
 (\mathcal N_TF)_\rho
 :=\sqrt{\frac{m_\rho}{2}}\left[
   \Phi_F(\gamma-i\alpha)-\Phi_F(\gamma+i\alpha)\right].       \tag{509}
\end{equation}
Then \(\mathcal N_T=P_-\mathcal E_T\) is exactly the negative residue
observation and
\begin{equation}
 \|\mathcal N_TF\|^2
 =\sum_{\rho\ \mathrm{off\ line}}
   \frac{m_\rho}{2}
   |\Phi_F(\gamma-i\alpha)-\Phi_F(\gamma+i\alpha)|^2.          \tag{510}
\end{equation}
Moreover every coordinate is the vertical flux of the continuous
logarithmic generator:
\begin{align}
 \Phi_F(\gamma-i\alpha)-\Phi_F(\gamma+i\alpha)
 &=2\int_{-T}^{T}F(t)e^{i\gamma t}\sinh(\alpha t)\,dt\notag\\
 &=\int_{-\alpha}^{\alpha}\int_{-T}^{T}
      tF(t)e^{i\gamma t-yt}\,dt\,dy .                         \tag{511}
\end{align}
Thus (509) is the canonical candidate for the requested G39 boundary
object: its infinitesimal density is generated by multiplication by \(t\),
the continuous counterpart of the degree operator \(Q\) in (27a)--(27d).
No choice of square root or pseudoinverse enters its construction.
\end{theorem}

\begin{proof}
Diagonalizing the exchange matrix in (124) gives the negative coordinate
\(\sqrt{m_\rho/2}\,[\Phi_F(z_\rho)-\Phi_F(\bar z_\rho)]\), which proves
(509)--(510).  Since
\[
 \Phi_F(\gamma+iy)=\int_{-T}^{T}F(t)e^{i\gamma t-yt}\,dt,
\]
differentiation under the compactly supported integral and the fundamental
theorem of calculus give (511).  The comparison with \(Q\) is the exact
continuous/discrete logarithmic-generator correspondence.
\end{proof}

\begin{theorem}[Exact zero-mode closure criterion --- \PROVED]
Let \(P_{0,T}\) denote the orthogonal projection onto \(\ker H_T\).  On a
first-loss window the residue balance is
\begin{equation}
 \|P_+\mathcal E_TF\|^2=\|\mathcal N_TF\|^2
 \qquad(F\in\ker H_T).                                      \tag{512}
\end{equation}
Consequently the single operator identity
\begin{equation}
 \boxed{\mathcal N_TP_{0,T}=0}                              \tag{513}
\end{equation}
forces every first-loss mode to vanish and hence closes G.  Conversely,
if a positive flux formula has left side
\(\|\mathcal N_TP_{0,T}F\|^2\) and its right side vanishes on zero modes,
then that formula proves precisely (513); it is not an additional weaker
normalization.
\end{theorem}

\begin{proof}
Equation (512) is (508) with zero left side, separated into the positive and
negative eigenspaces of the residue fundamental symmetry.  Under (513),
(512) also gives \(P_+\mathcal E_TF=0\).  In particular \(\Phi_F\) vanishes
at all critical-line zeros.  The proved Paley--Wiener--Cartwright rigidity
criterion then gives \(F=0\), contradicting first loss.  The converse
assertion follows by evaluating the proposed flux identity on
\(\ker H_T\).
\end{proof}

\paragraph{G39-to-negative-port gate --- \OPEN.}
Formula (511) constructs the requested negative boundary port directly from
the logarithmic generator, but the established G39 identity (27g) fixes its
\emph{divergence/cross amplitude}; it supplies no theorem saying that the
vertical integral in (511) vanishes on \(\ker H_T\).  Conservation gives
exactly (512), so replacing that missing implication by a declaration that
the flux is zero would assume (513).  The remaining statement is therefore
the source-derived radical-lifting identity
\(\mathcal N_TP_{0,T}=0\), equivalently the equality-state Hodge theorem.

\begin{theorem}[The negative port is the antisymmetric causal pole jump ---
\PROVED]
Let
\begin{equation}
 \mathscr R_F(s):=-\frac{\zeta'}{\zeta}(s)
                  K_F(s-\tfrac12).                          \tag{514}
\end{equation}
For an off-line pair
\(\rho_+=\frac12+\alpha+i\gamma\) and
\(\rho_-=1-\bar\rho_+=\frac12-\alpha+i\gamma\), with common
multiplicity \(m_\rho\), one has
\begin{align}
 \operatorname*{res}_{s=\rho_+}\mathscr R_F(s)
   &=-m_\rho K_F(\alpha+i\gamma),\notag\\
 \operatorname*{res}_{s=\rho_-}\mathscr R_F(s)
   &=-m_\rho K_F(-\alpha+i\gamma).                         \tag{515}
\end{align}
After the harmless reindexing \(\gamma\mapsto-\gamma\), the negative
coordinate (509) is therefore
\begin{equation}
 (\mathcal N_TF)_\rho
 =\frac{1}{\sqrt{2m_\rho}}\left[
  \operatorname*{res}_{s=\rho_+}\mathscr R_F(s)
 -\operatorname*{res}_{s=\rho_-}\mathscr R_F(s)\right].      \tag{516}
\end{equation}
Thus the desired boundary norm is exactly the squared antisymmetric jump of
the two causal pole residues.  The zero-mode/G39 equation (491) fixes the
reflected symmetric combination; (516) is the complementary datum (492).
\end{theorem}

\begin{proof}
The logarithmic derivative has residue \(m_\rho\) at a zero of
multiplicity \(m_\rho\), which proves (515) with the sign in (514).  The
relations
\[
 K_F(-\alpha+i\gamma)=\Phi_F(-\gamma-i\alpha),\qquad
 K_F( \alpha+i\gamma)=\Phi_F(-\gamma+i\alpha)
\]
turn the difference of (515) into (509), after indexing the conjugate
functional-equation orbit.  Equations (491)--(492) give the final
symmetric/antisymmetric assertion.
\end{proof}

\paragraph{Stable-realization equivalence --- \AUDIT.}
If the two causal outputs belonged to the Hardy classes (286), every pole
in (515) would be removable, so both residues, and hence (516), would vanish.
Conversely, G39 supplies only their reflected symmetric combination and
does not make either meromorphic germ Hardy.  Therefore the proposed
positive-flow proof reduces exactly to establishing the stable causal
realization (286), or the weaker local pole estimate (497), on zero modes.
Calling the finite G39 colligation passive at the critical completion would
insert that missing stability as an assumption.

\begin{theorem}[Causal-splitting gauge obstruction --- \PROVED]
Write, on each off-line orbit,
\begin{equation}
 r_+(F):=\operatorname*{res}_{s=\rho_+}\mathscr R_F(s),qquad
 r_-(F):=\operatorname*{res}_{s=\rho_-}\mathscr R_F(s),       \tag{517}
\end{equation}
and put \(S=r_++r_-\), \(D=r_+-r_-\).  Knowledge of the G39/reflected
symmetric amplitude \(S\) does not determine \(D\): for every scalar
functional \(h\) the transformation
\begin{equation}
 (r_+,r_-)\longmapsto(r_++h,r_--h)                         \tag{518}
\end{equation}
preserves \(S\) and changes \(D\) to \(D+2h\).  Under (516), this is
exactly the freedom to change the negative residue port while leaving all
already proved symmetric G39 data unchanged.

Even after imposing the zero-mode conservation law (512), only the norm of
\(D\) is tied to the positive residue norm; neither its phase nor its
vanishing is determined.  Hence no algebraic consequence of (27g),
(491) and (512) alone can prove
\begin{equation}
 \mathcal N_TP_{0,T}=0.                                  \tag{519}
\end{equation}
A proof of (519) must break the gauge (518) by an independently established
causal selection principle, such as the Hardy realization (286), the local
pole estimate (497), or a faithful geometric Hodge/effectivity theorem.
\end{theorem}

\begin{proof}
Equations (517)--(518) are elementary linear algebra: the sum is invariant
and the difference changes by \(2h\).  Formula (516) identifies that
difference with the negative port.  Equation (512) is one real quadratic
constraint and permits nonzero vectors of equal positive and negative norm,
as already exhibited by the hyperbolic residue plane (125).  Therefore it
cannot select \(D=0\).  Each selection principle listed in the last
sentence removes poles or negative directions by information not contained
in the invariant sum, which proves the final assertion.
\end{proof}

\paragraph{Terminal logical audit --- \AUDIT.}
The demanded positive boundary object has now been constructed explicitly
in three mutually consistent forms: the negative residue coordinate (509),
the logarithmic vertical flux (511), and the antisymmetric causal pole jump
(516).  Its vanishing on a zero mode is (513)/(519).  The gauge theorem
proves that this vanishing cannot be derived from the currently available
G39 amplitude, reflection and scalar conservation identities.  Supplying a
Hardy/passivity or geometric Hodge selection theorem is therefore a genuine
new hypothesis/theorem, not a remaining calculation or normalization.

\begin{proposition}[Physical scope correction of the gauge audit ---
\AUDIT]
The transformation (518) is an algebraic obstruction to recovering the
negative port from the symmetric G39 data; it is not, by itself, a family
of physical zero modes.  A shift \(h\) is physically admissible only if it
is induced by the Paley--Wiener transform of a compact primitive source
which also satisfies the full compressed zero-mode equation.  Therefore
(518) proves insufficiency of the listed scalar identities, but it does not
prove existence of a nonzero first-loss vector.  Any earlier reading of
the gauge argument as a physical counterexample is withdrawn.
\end{proposition}

\begin{theorem}[Finite-orbit physical interpolation survives parity and
Tate --- \PROVED]
Fix an off-line orbit with
\(z_+=\alpha+i\gamma\), \(z_-=-\alpha+i\gamma\), where
\(\alpha\gamma\ne0\), and a parity
\(\varepsilon\in\{+1,-1\}\).  Let
\(\mathcal P_T^\varepsilon\) be the parity-
\(\varepsilon\) primitive subspace.  Then
\begin{equation}
 \mathcal P_T^\varepsilon\longrightarrow\mathbb C^2,qquad
 F\longmapsto\bigl(K_F(z_+),K_F(z_-)\bigr)                 \tag{520}
\end{equation}
is onto.  Hence compact support, Paley--Wiener type, parity and both Tate
moments do not determine the antisymmetric residue coordinate on any fixed
off-line orbit.
\end{theorem}

\begin{proof}
On the even subspace the representing functions for the two evaluations
and for the (single independent) Tate condition are
\begin{equation}
 \cosh(z_+t),\quad\cosh(z_-t),\quad\cosh(t/2),               \tag{521}
\end{equation}
and on the odd subspace they are
\begin{equation}
 \sinh(z_+t),\quad\sinh(z_-t),\quad\sinh(t/2).               \tag{522}
\end{equation}
They are linearly independent on every nonempty interval: exponential
polynomials with distinct exponent pairs are independent, and
\(z_+^2\ne z_-^2\) because \(4i\alpha\gamma\ne0\); neither pair equals
the polar pair \(\{\pm1/2\}\).  Thus the three continuous functionals have
rank three.  Restricting the first two to the kernel of the third leaves
rank two, which proves surjectivity of (520).  In a fixed parity sector the
two Tate moments are proportional, so that third functional is the entire
primitive constraint.
\end{proof}

\paragraph{Corrected physical endpoint --- \OPEN.}
The finite-orbit theorem shows that all local analytic constraints still
leave the negative coordinate free.  The only remaining possible source
of its vanishing is the global zero-mode orthogonality
\begin{equation}
 \sum_{\mathcal O}\left\langle
 J_{\mathcal O}\mathcal E_{\mathcal O}F,
 \mathcal E_{\mathcal O}K\right\rangle=0
 \qquad(K\in\mathcal P_T),                              \tag{523}
\end{equation}
obtained by polarizing the first-loss radical identity.  Passing from the
finite interpolation (520) to a test \(K\) which isolates one orbit in
(523) requires a globally bounded interpolation/synthesis theorem for the
entire zeta-zero sequence.  Such a theorem is not supplied by finite
Paley--Wiener interpolation: the zero sequence has density
\(\asymp R\log R\), whereas a nonzero exponential-type function has only
\(O(R)\) zeros.  This is the precise infinite-dimensional completion gate.

\begin{theorem}[Exact residue-range lifting criterion --- \PROVED]
Let \(F\) be a first-loss vector and let
\[
 \mathscr R_T:=
 \overline{\mathcal E_T(\mathcal P_T\cap\mathcal Q_{\Gamma,T})}
 \subset\mathscr K_{\rm res}
\]
be the closed physical residue range.  If
\begin{equation}
 J_{\rm res}\mathcal E_TF\in\mathscr R_T,             \tag{523a}
\end{equation}
then \(F=0\).  More generally it is enough that there exist
\(K_n\in\mathcal P_T\cap\mathcal Q_{\Gamma,T}\) with
\(\mathcal E_TK_n\to J_{\rm res}\mathcal E_TF\).
Thus a source-defined lift of the residue involution into the closed
physical range is sufficient to prove the causal selection theorem and
exclude first contact.
\end{theorem}

\begin{proof}
By (523),
\[
 \langle J_{\rm res}\mathcal E_TF,\mathcal E_TK_n\rangle=0
 \qquad(n\ge1).
\]
Pass to the limit in the residue Hilbert space, using (549), to obtain
\[
 0=\|J_{\rm res}\mathcal E_TF\|^2
  =\|\mathcal E_TF\|^2.
\]
Hence \(K_F\) vanishes at every nontrivial zeta zero.  The unconditional
distinct-zero density used in (292), together with the Cartwright bound for
the exponential-type function \(K_F\), forces \(K_F\equiv0\).  Transform
injectivity gives \(F=0\).
\end{proof}

\paragraph{Candidate and domain audit --- \AUDIT.}
At the meromorphic-germ level the Fourier/co-Poisson functional equation
interchanges the two members of every functional-equation orbit and hence
has the residue action \(J_{\rm res}\).  To prove (523a), however, its image
must be represented by compact-window primitive sources in the norm (549).
The local operator (436)--(439) is invertible on the window but does not
transport the exterior member of an orbit back into that range.  The
missing statement is exactly the reciprocal gap/domain assertion
(441)/(446), or equivalently control of the incoming history (452).
Therefore (523a) is a genuine source-lifting target, not a consequence of
the meromorphic functional equation alone.

\begin{theorem}[Minimal equality-state range lift --- \PROVED]
Let \(F\in\ker H_T\), put \(e=\mathcal E_TF\), and let
\(\mathscr R_T=\mathcal E_T(\mathcal Q_T^{\rm prim})\), which is closed by
(642)--(644).  Then
\begin{equation}
 \boxed{P_-e\in\mathscr R_T\quad\Longrightarrow\quad F=0.}     \tag{702}
\end{equation}
Equivalently, it is enough to construct sources \(K_n\) such that
\[
 \mathcal E_TK_n\longrightarrow P_-\mathcal E_TF .
\]
On the zero-mode space, the assertion that this membership always holds is
equivalent to (700).  It requires strictly less data than the earlier
sufficient condition (523a), which asks for a lift of the whole vector
\(J_{\rm res}e=P_+e-P_-e\).
\end{theorem}

\begin{proof}
The radical equation \(H_TF=0\), polarized against every source, says
\[
 J_{\rm res}e\perp\mathscr R_T.
\]
If \(P_-e\in\mathscr R_T\), then
\[
 0=\langle J_{\rm res}e,P_-e\rangle
   =-\|P_-e\|^2,
\]
because \(J_{\rm res}P_-=-P_-\).  Hence
\(P_-e=0\), i.e. \(\mathcal N_TF=0\).  The equivalences in
(698)--(700) now give \(F=0\).  The approximation formulation follows
from closedness of \(\mathscr R_T\).

Conversely, on \(\ker H_T\), assertion (700) gives \(F=0\), and then
\(P_-e=0\in\mathscr R_T\).  Thus (700) is equivalent to the assertion that
the membership hypothesis of (702) holds for every equality state.  Since
\(P_-e\) is only one spectral component of \(J_{\rm res}e\), (523a)
implies that membership, while the reverse implication is not asserted.
\end{proof}

\paragraph{Revised Hermitian target --- \OPEN.}
The source construction required from (661) may therefore be restricted to
the equality bundle:
\begin{equation}
 F\in\ker H_T
 \quad\Longrightarrow\quad
 P_-\mathcal E_TF\in\mathscr R_T.                   \tag{703}
\end{equation}
Once (703) is proved, (702) closes the argument without a uniform lift of
\(J_{\rm res}\) on all ordinary sources.  Finite-orbit interpolation proves
the corresponding statement at every finite truncation, but its norms are
not uniform; passing to the closed full range is precisely the remaining
Hermitian graph-tangency estimate.  Membership in (703) may not be imposed
by projecting onto \(\mathscr R_T\), since that would replace the target by
its projection and assume the missing equality.

\begin{theorem}[Exact distance of the negative port from the physical
range --- \PROVED]
Let \(Q_T\) be the orthogonal projection onto \(\mathscr R_T\), let
\(\mathbb H_T=Q_TJ_{\rm res}|_{\mathscr R_T}\), and put
\[
 \mathbb L_T:=(I-Q_T)J_{\rm res}|_{\mathscr R_T}.
\]
For every \(e\in\mathscr R_T\),
\begin{align}
 Q_TP_-e&=\frac12(I-\mathbb H_T)e,\notag\\
 (I-Q_T)P_-e&=-\frac12\mathbb L_Te,                  \tag{704}\\
 \operatorname{dist}(P_-e,\mathscr R_T)^2
 &=\frac14\bigl(\|e\|^2-\|\mathbb H_Te\|^2\bigr).
                                                               \tag{705}
\end{align}
Equivalently,
\begin{equation}
 \boxed{\mathbb H_T^2+\mathbb L_T^*\mathbb L_T=I_{\mathscr R_T}.}
                                                               \tag{706}
\end{equation}
In particular, if \(0\ne e\in\ker\mathbb H_T\), then
\begin{equation}
 \operatorname{dist}(P_-e,\mathscr R_T)=\frac12\|e\|>0.        \tag{707}
\end{equation}
Thus finite interpolation of every finite residue block cannot converge in
the full residue norm to the negative port of a nonzero zero mode.
\end{theorem}

\begin{proof}
Since \(P_-=(I-J_{\rm res})/2\) and \(Q_Te=e\),
\[
 Q_TP_-e=\tfrac12(e-Q_TJ_{\rm res}e)
         =\tfrac12(I-\mathbb H_T)e.
\]
The complementary projection gives the second identity in (704).
Because \(J_{\rm res}\) is unitary and the two orthogonal components of
\(J_{\rm res}e\) are \(\mathbb H_Te\) and \(\mathbb L_Te\),
\[
 \|e\|^2=\|J_{\rm res}e\|^2
        =\|\mathbb H_Te\|^2+\|\mathbb L_Te\|^2.
\]
Polarization gives (706), while (704) and the projection theorem give
(705).  Setting \(\mathbb H_Te=0\) proves (707).
\end{proof}

\paragraph{Range-lift correction --- \AUDIT.}
The conditional implication (702) remains correct, but (707) shows that
(703) is not an intermediate approximation theorem available from finite
interpolation and closedness.  For a hypothetical nonzero equality state,
its proposed target is separated from the physical range by a fixed
positive distance.  Therefore proving (703) already excludes that state;
it is another exact formulation of (700), not a softer construction.
The Hermitian route must establish the product rule (661) directly from a
source conservation identity.  It cannot first manufacture
\(P_-\mathcal E_TF\) by a generic range projection or a limit of the finite
right inverses (606)--(610).

\begin{proposition}[Reality and parity do not implement the residue
involution --- \PROVED]
Let \(F\) be real-valued with parity
\(\mathcal RF=\varepsilon F\), \(\varepsilon\in\{1,-1\}\).  On an off-line
fiber (124), write \(a=\Phi_F(z_\rho)\).  Then
\begin{equation}
 \mathcal E_{\mathcal O}F=(a,\varepsilon\bar a),
 \qquad
 J_{\mathcal O}\mathcal E_{\mathcal O}F
   =(\varepsilon\bar a,a).                            \tag{523b}
\end{equation}
The two vectors agree only when \(a=\varepsilon\bar a\).  This phase
condition is not a consequence of compact support, Tate primitiveness,
reality or parity: on every fixed off-line orbit with
\(\alpha\gamma\ne0\), real primitive parity sources realize arbitrary
\(a\in\mathbb C\).  Hence the real structure and physical reflection do
not prove the range lift (523a).
\end{proposition}

\begin{proof}
For real \(F\),
\[
 \overline{\Phi_F(z)}
 =\Phi_F(-\bar z)
 =\varepsilon\Phi_F(\bar z),
\]
which gives (523b).  Equality with its exchanged vector is exactly the
displayed phase condition.  In the even sector the real and imaginary
parts of \(\cos(zt)\), and in the odd sector those of \(\sin(zt)\), are
linearly independent modulo the single real Tate functional when
\(\alpha\gamma\ne0\).  The same finite-dimensional rank argument as in
(520)--(522), now over \(\mathbb R\), therefore makes
\(F\mapsto\Phi_F(z_\rho)\) onto \(\mathbb C\).  Arbitrary phases occur, so
the phase condition cannot be inferred from the stated symmetries.
\end{proof}

\begin{proposition}[Exact single-orbit isolation is impossible ---
\PROVED]
Let \(F\in L^2(-T,T)\).  If its entire transform \(K_F\) vanishes at all
but finitely many members of any unconditional family of distinct
critical-line zeros having counting function \(\gg R\log R\), then
\begin{equation}
 K_F\equiv0.                                             \tag{524}
\end{equation}
In particular no compact-window test can isolate one residue orbit in
(523) by annihilating every other critical residue coordinate.
\end{proposition}

\begin{proof}
Paley--Wiener places \(K_F\) in the Cartwright class of exponential type at
most \(T\).  A nonzero member has only \(O_T(R)\) zeros in \(|z|\le R\),
counted with multiplicity.  Vanishing on all but finitely many points of a
family with \(\gg R\log R\) distinct points contradicts this bound.  Hence
\(K_F=0\), and transform injectivity also gives \(F=0\).
\end{proof}

\paragraph{Consequence for the radical route --- \AUDIT.}
Equation (523) cannot be diagonalized orbit by orbit with compactly
supported test functions.  A valid use of the full orthogonality must
instead control the entire coupled residue sequence in a completion whose
metric is independently positive or whose causal synthesis is Hardy.  This
returns exactly to (497)/(513), but now without the unjustified assumption
that the algebraic gauge (518) is freely realizable inside the zero-mode
space.

\begin{theorem}[Vertical homotopy flux equals the negative residue energy
--- \PROVED]
For an off-line orbit represented by
\(\rho=\frac12+\alpha+i\gamma\), \(\alpha>0\), set, for
\(0\le y\le\alpha\),
\begin{align}
 D_{\rho,F}(y)&:=\Phi_F(\gamma-iy)-\Phi_F(\gamma+iy),\notag\\
 S_{\rho,QF}(y)&:=\Phi_{QF}(\gamma-iy)+
                         \Phi_{QF}(\gamma+iy),qquad QF(t)=tF(t). \tag{525}
\end{align}
Then
\begin{equation}
 \frac{d}{dy}\frac12|D_{\rho,F}(y)|^2
 =\Re\!\left(S_{\rho,QF}(y)
              \overline{D_{\rho,F}(y)}\right),              \tag{526}
\end{equation}
and, because \(D_{\rho,F}(0)=0\),
\begin{equation}
 \boxed{
 \frac{m_\rho}{2}|D_{\rho,F}(\alpha)|^2
 =m_\rho\int_0^\alpha
   \Re\!\left(S_{\rho,QF}(y)
              \overline{D_{\rho,F}(y)}\right)dy.}           \tag{527}
\end{equation}
The left side is precisely the squared negative-port coordinate in
(509)--(510).  Summing (527) over any finite collection of off-line orbits
gives a source-defined positive boundary energy as an exact logarithmic
bulk flux, without a square-root choice or a sign assumption.
\end{theorem}

\begin{proof}
Differentiation under the compactly supported integral gives
\[
 \frac d{dy}\Phi_F(\gamma-iy)=\Phi_{QF}(\gamma-iy),\qquad
 \frac d{dy}\Phi_F(\gamma+iy)=-\Phi_{QF}(\gamma+iy).
\]
Hence \(D'_{\rho,F}=S_{\rho,QF}\), and differentiating
\(|D|^2/2\) proves (526).  Integrate from zero to \(\alpha\), multiply by
\(m_\rho\), and use (509) to obtain (527).
\end{proof}

\paragraph{Exact comparison still required --- \OPEN.}
Formula (527) supplies the previously missing positive flux identity at the
level of every residue orbit.  To conclude its boundary vanishes for a
zero mode, one must prove that the sum of its bulk integrands is the
physical G39 cross-amplitude after projecting \(QF\) back to the primitive
space and including the moving Tate correction (105)--(111).  This
comparison is not formal: \(QF\) is generally not primitive, and
(505)--(507) prove that its Tate correction has nonzero terminal energy.
Thus discarding that correction would falsely turn (527) into the desired
zero identity.  The next admissible target is an exact formula for the
\emph{sum} of the vertical bulk flux and the Tate terminal, derived from
G39 before taking the first-loss kernel.

\begin{theorem}[Exact primitive correction of the logarithmic generator ---
\PROVED]
Let \(B=B_1\), \(K=BB^*\), \(P=I-B^*K^{-1}B\), and
\(D=\operatorname{diag}(-1,1)\).  For every \(F\in\ker B\),
\begin{align}
 B_1'F&=\frac12D\,B(QF),                              \tag{528}\\
 QF&=PQF+2B^*K^{-1}D\,B_1'F.                         \tag{529}
\end{align}
Consequently
\begin{equation}
 \|(I-P)QF\|^2
 =4\langle DB_1'F,K^{-1}DB_1'F\rangle.               \tag{530}
\end{equation}
This is not the moving Tate terminal (109), whose norm is
\(\langle B_1'F,K^{-1}B_1'F\rangle\).
\end{theorem}

\begin{proof}
Formula (110) is (528).  Orthogonal projection onto
\(\operatorname{Ran}B^*\) gives
\((I-P)QF=B^*K^{-1}B(QF)\); substitution of (528) proves (529).
Taking its norm and using \(BB^*=K\) proves (530).
\end{proof}

Write \(a=2\sinh T\), \(b=2T\).  If \(F\) is even, \(B_1'F\) lies in
the \(a+b\) eigendirection of \(K\), while \(DB_1'F\) lies in the
\(a-b\) eigendirection; for odd \(F\) the roles reverse.  Therefore
\begin{equation}
 \frac{\|(I-P)QF\|^2}{\|\nu_1(F)\|^2}
 =\begin{cases}
 4(a+b)/(a-b),&F\text{ even},\\
 4(a-b)/(a+b),&F\text{ odd},
 \end{cases}                                          \tag{531}
\end{equation}
whenever the denominator is nonzero.

\paragraph{Normalization consequence --- \AUDIT.}
The vertical flux (527) cannot be closed by substituting the moving
terminal square (109): the correct logarithmic normal boundary is (530),
with the Tate orientations exchanged by \(D\).

\begin{theorem}[No conservative transport from the moving Tate terminal to
the logarithmic normal port --- \PROVED]
Define
\begin{equation}
 \mathcal T_TF:=K^{-1/2}B_1'F,\qquad
 \mathcal Q_TF:=2K^{-1/2}DB_1'F.                      \tag{532}
\end{equation}
There is no isometry \(U_T\) on their source-generated ranges satisfying
\begin{equation}
 U_T\mathcal T_TF=\mathcal Q_TF
 \qquad(F\in\mathcal P_T).                            \tag{533}
\end{equation}
\end{theorem}

\begin{proof}
An isometry would preserve the pullback Grams.  On every even primitive
vector with \(B_1'F\ne0\), (531) gives
\begin{equation}
 \frac{\|\mathcal Q_TF\|^2}{\|\mathcal T_TF\|^2}
 =4\frac{a+b}{a-b}>4,                                \tag{534}
\end{equation}
because \(a=2\sinh T>b=2T>0\).  Such vectors exist by (505)--(506).
\end{proof}

\begin{corollary}[What a bare augmented transport would have to pay ---
\AUDIT]
An augmented conservative colligation must compensate the Gram difference
\(\mathcal Q_T^*\mathcal Q_T-\mathcal T_T^*\mathcal T_T\).
Choosing compensating channels abstractly from its square root would be a
KYP completion and would prove no sign.  This does not exclude an
independent source formula: the following theorem constructs one from the
first positive Gamma cell.  The finite doubled G39 port alone still cannot
supply it by an isometric transport.
\end{corollary}

\begin{theorem}[The first Gamma cell pays the complete oriented Tate defect
--- \PROVED]
Let
\begin{equation}
 \mathcal D_T:=\mathcal Q_T^*\mathcal Q_T-
                    \mathcal T_T^*\mathcal T_T.       \tag{535}
\end{equation}
On the even and odd primitive sectors write, respectively,
\(B_1'F=(u,u)\) and \(B_1'F=(-u,u)\).  Then
\begin{equation}
 \langle\mathcal D_TF,F\rangle=
 \begin{cases}
 c_e(T)|u|^2,&F\text{ even},\\
 c_o(T)|u|^2,&F\text{ odd},
 \end{cases}
 \quad
 c_e=\frac8{a-b}-\frac2{a+b},\quad
 c_o=\frac8{a+b}-\frac2{a-b}.                         \tag{536}
\end{equation}
Let \(H_{\Gamma,0}\) be the positive Fourier multiplier
\begin{equation}
 h_0(\tau)=\frac{16\tau^2}{1+4\tau^2}.                \tag{537}
\end{equation}
For every \(T>0\) and every primitive \(F\),
\begin{equation}
 \boxed{
 \langle\mathcal D_TF,F\rangle
 \le\frac34\langle H_{\Gamma,0}F,F\rangle
 \le\frac34\langle G_{\Gamma,T}F,F\rangle.}           \tag{538}
\end{equation}
Consequently
\begin{equation}
 \boxed{G_{\Gamma,T}-\mathcal D_T
        \ge\frac14G_{\Gamma,T}\ge0.}                  \tag{539}
\end{equation}
Thus no Euler channel is needed to compensate this terminal defect.
\end{theorem}

\begin{proof}
Equation (536) follows from (109), (530), and the eigenvalues \(a\pm b\)
of \(K\).  The positive-cell decomposition (385) gives
\[
 g_\Gamma(\tau)=\sum_{n\ge0}h_n(\tau),\qquad
 h_0(\tau)=4-\frac4{1+4\tau^2}
           =\frac{16\tau^2}{1+4\tau^2},
\]
so \(H_{\Gamma,0}\le G_{\Gamma,T}\).

Put \(x=T/2\), \(S=\sinh(2x)\), and \(C=\cosh(2x)\).  In the even sector
\(u\) is represented by
\(r_e(t)=\frac12t\sinh(t/2)\), modulo the primitive annihilator
\(q_e(t)=\cosh(t/2)\).  Choose
\begin{equation}
 \lambda_e=x\coth x-1,\qquad r_e^0=r_e-\lambda_eq_e.  \tag{540}
\end{equation}
Then \(\int_{-T}^Tr_e^0=0\).  If
\(R_e(t)=\int_{-T}^tr_e^0(s)\,ds\), Fourier
Cauchy--Schwarz and
\[
 \frac1{h_0(\tau)}=\frac14+\frac1{16\tau^2}
\]
give
\begin{equation}
 |u|^2\le C_e(x)\langle H_{\Gamma,0}F,F\rangle,\qquad
 C_e=\frac14\|r_e^0\|_2^2+\frac1{16}\|R_e\|_2^2
 =\frac S4+\frac x2-x^2\coth x.                      \tag{541}
\end{equation}
Subtracting \(\lambda_eq_e\) does not change the pairing with primitive
even \(F\).

In the odd sector use
\(r_o(t)=\frac12t\cosh(t/2)\), \(q_o(t)=\sinh(t/2)\), and minimize the
same inverse-\(h_0\) norm over \(r_o-\lambda q_o\).  Direct integration
gives
\begin{align}
 \lambda_o&=\frac{x^2S-2xC-2x+2S}{xC+x-S},\notag\\
 C_o(x)&=\frac{4x^3S-6x^2C-6x^2+xSC+11xS-3S^2}
               {4(xC+x-S)}.                          \tag{542}
\end{align}
Thus \(|u|^2\le C_o\langle H_{\Gamma,0}F,F\rangle\).

The remaining scalar inequalities are
\begin{equation}
 0<c_eC_e\le\frac34,\qquad
 0<C_o\le C_e,\qquad c_o\le c_e.                     \tag{543}
\end{equation}
For the first, after multiplication by \(S^2-4x^2>0\), it suffices that
\begin{equation}
 \bigl(10x^2\coth x+3xS\coth x-8x-4S\bigr)\sinh x
 =\sum_{n\ge2}
 \frac{40n^2-33n+2+(n-2)3^{2n}}{(2n)!}x^{2n}>0.      \tag{544}
\end{equation}
For \(C_o\le C_e\), direct subtraction gives
\[
 C_e-C_o=\frac{M(x)}{2\sinh x\,(xC+x-S)},\qquad
 M(x)=\sum_{n\ge4}\frac{A_n}{(2n+1)!}x^{2n+1},        \tag{545}
\]
where
\begin{align*}
 A_n={}&-2(3^{2n-2}+1)(2n+1)(2n)(2n-1)\\
 &+3(3^{2n-1}+1)(2n+1)(2n)
   -3(3^{2n}-1)(2n+1)
   +\frac{5^{2n+1}-3^{2n+1}-2}{4}.
\end{align*}
Here \(A_4=43008>0\), and
\[
 A_{n+1}-25A_n=\frac89\left[
 9^n(32n^3-126n^2+91n+54)
 +432n^3-378n^2-459n-54\right]>0
\]
for \(n\ge4\).  Also \(xC+x-S>0\) by its power series.  This proves the
middle inequality in (543); the last follows immediately from (536).

If \(c_o\le0\), the odd defect is already nonpositive.  If \(c_o>0\),
(543) gives \(c_oC_o\le c_eC_e\le3/4\).  Hence (538) holds in both
parities.  Finally,
\[
 G_{\Gamma,T}-\frac34H_{\Gamma,0}
 =\frac14G_{\Gamma,T}
  +\frac34(G_{\Gamma,T}-H_{\Gamma,0})
 \ge\frac14G_{\Gamma,T},
\]
which proves (539).
\end{proof}

\begin{corollary}[Explicit noncircular terminal storage --- \PROVED]
Let \(X_0^*X_0=H_{\Gamma,0}\), and let \(v_e,v_o\) be the Riesz vectors
corresponding to (541)--(542).  On a parity sector with \(c_p>0\), put
\(\theta_p=c_p\|v_p\|^2\le3/4\) and \(e_p=v_p/\|v_p\|\).  Then
\begin{equation}
 H_{\Gamma,0}-c_pu_p^*u_p
 =X_0^*(I-\theta_p e_p\otimes e_p)X_0.                \tag{546}
\end{equation}
If \(c_p\le0\), the term \((-c_p)^{1/2}u_p\) is an additional positive
square.  Together with the untouched cells
\(\bigoplus_{n\ge1}h_n^{1/2}\), this explicitly factors
\(G_{\Gamma,T}-\mathcal D_T\).  The rank-one square root is noncircular
because \(0\le\theta_p\le3/4\) was proved independently above.
\end{corollary}

\paragraph{Scope of the Gamma compensation --- \AUDIT.}
Equations (535)--(546) complete the terminal-compensation step and correct
the earlier suggestion that an Euler channel is required merely to pay the
orientation gain.  They do not yet identify the globally summed vertical
residue flux (527) with the physical G39 cross-amplitude.  That source
comparison, including convergence of the off-line orbit sum and the
zero-mode polarization (523), remains necessary before this storage can
imply \(\mathcal N_TP_{0,T}=0\).  Thus this is new positive progress, not
yet a closure of G or row~(d).

\begin{theorem}[Continuous residue Hilbert space and global summability ---
\PROVED]
Let \(\mathcal Z_\xi^{\rm dist}\) be the set of distinct nontrivial
zeros, give an orbit \(\mathcal O\) its common multiplicity
\(m_{\mathcal O}\), and group the zeros into functional-equation orbits as
in (122)--(125).  Define
\begin{equation}
 \mathscr K_{\mathrm{res}}
 :=\bigoplus_{\mathcal O}^{\ell^2}
   (\mathcal K_{\mathcal O},m_{\mathcal O}\langle\cdot,\cdot\rangle),
 \qquad
 J_{\mathrm{res}}:=\bigoplus_{\mathcal O}J_{\mathcal O}.       \tag{547}
\end{equation}
For fixed \(T\), the residue analysis map
\begin{equation}
 \mathcal E_TF:=\{\mathcal E_{\mathcal O}F\}_{\mathcal O}      \tag{548}
\end{equation}
extends boundedly from the smooth primitive core to the logarithmic Gamma
form domain
\[
 \mathcal Q_{\Gamma,T}:=
 \operatorname{Dom}G_{\Gamma,T}^{1/2},\qquad
 \|F\|_{\mathcal Q_{\Gamma,T}}^2
 :=\|F\|_2^2+\langle G_{\Gamma,T}F,F\rangle .
\]
More precisely,
\begin{equation}
 \|\mathcal E_TF\|_{\mathscr K_{\mathrm{res}}}^2
 \le C_T\left(\|F\|_2^2+
              \langle G_{\Gamma,T}F,F\rangle\right).          \tag{549}
\end{equation}
Consequently the negative port
\(\mathcal N_T=P_-\mathcal E_T\) is a genuine bounded Hilbert observation,
and the Krein residue form
\begin{equation}
 \langle\mathcal E_TF,J_{\mathrm{res}}\mathcal E_TG\rangle
 =\sum_{\mathcal O}m_{\mathcal O}
   \langle\mathcal E_{\mathcal O}F,
          J_{\mathcal O}\mathcal E_{\mathcal O}G\rangle        \tag{550}
\end{equation}
converges absolutely for \(F,G\in\mathcal Q_{\Gamma,T}\).
\end{theorem}

\begin{proof}
The Riemann--von Mangoldt formula gives, counting multiplicity,
\[
 \#\{\rho:\ k\le\Im\rho<k+1\}=O(\log(2+|k|)).
\]
The points \(z_\rho=(\rho-\frac12)/i\) lie in the fixed horizontal strip
\(|\Im z|\le\frac12\).  The weighted Plancherel--Pólya sampling estimate
for entire functions of exponential type at most \(T\) therefore gives
\begin{equation}
 \sum_{\rho\in\mathcal Z_\xi^{\rm dist}}
       m_\rho|\Phi_F(z_\rho)|^2
 \le C_T\int_{\mathbb R}\log(2+|\tau|)
                |\widehat{E_TF}(\tau)|^2\,d\tau .             \tag{551}
\end{equation}
For completeness, this estimate follows by dividing the strip into unit
boxes, applying the local Bernstein--Sobolev estimate to the band-limited
function \(\widehat{E_TF}\) in each enlarged box, and using the displayed
\(O(\log(2+|k|))\) zero count; the enlarged boxes have bounded overlap.
Vertical translation by at most \(1/2\) is bounded on the
Paley--Wiener space of type \(T\).

The digamma asymptotic and positivity of (81) give
\[
 \log(2+|\tau|)\le C\bigl(1+g_\Gamma(\tau)\bigr).
\]
Plancherel now turns (551) into (549).  Diagonalizing every off-line
exchange matrix shows that \(P_-\) is an orthogonal projection in the
Euclidean direct sum, proving boundedness of \(\mathcal N_T\).
Finally, Cauchy--Schwarz in \(\mathscr K_{\mathrm{res}}\), together with
\(\|J_{\mathrm{res}}\|=1\), proves absolute convergence of (550).
\end{proof}

\paragraph{What summability does and does not solve --- \AUDIT.}
Equations (547)--(551) put the positive Gamma remainder, the negative
residue port and the Krein explicit-formula boundary in honest continuous
Hilbert spaces; no formal orbit sum remains.  They do not identify the
finite divisibility current (27e)--(27h) with a bounded operator into
\(\mathscr K_{\mathrm{res}}\).  The raw prime channels cannot have such a
strong limit because of (410)--(412); the required limit must first perform
the cross-prime Tate-centered synthesis (413).

\begin{theorem}[Raw finite G39 cannot converge in the residue graph topology
--- \PROVED]
Let
\[
 C_N(s):=\sum_{2\le n\le N}\Lambda(n)n^{-s}.
\]
Every \(C_N(s)K_F(s-\frac12)\) is entire.  There is no locally convex
topology on meromorphic germs which simultaneously
\begin{enumerate}
\item makes every residue functional
      \(f\mapsto\operatorname{res}_{s=\rho}f(s)\) continuous;
\item has
      \[
       C_N(s)K_F(s-\tfrac12)\longrightarrow
       -\frac{\zeta'}{\zeta}(s)K_F(s-\tfrac12)
      \]
      whenever the right side has a nonzero residue at a zeta zero.
\end{enumerate}
In particular the raw finite divisibility/Euler port cannot converge
strongly to the continuous residue port (548).
\end{theorem}

\begin{proof}
Continuity of the residue functional would send the zero residues of all
entire approximants to the residue of the limit.  The latter equals
\(-m_\rho K_F(\rho-\frac12)\), which is nonzero for an admissible source
unless the desired divisibility already holds.  This contradiction proves
the assertion.
\end{proof}

\begin{definition}[Completed continuum G39 graph port --- \AUDIT]
The only type-compatible continuum candidate is therefore the graph pair
\begin{equation}
 \mathbb W_{\mathrm{G39}}F:=
 \left(
  \operatorname{fp}\!\left[
   -\frac{\zeta'}{\zeta}(s)K_F(s-\tfrac12)\right],
  \{-m_\rho K_F(\rho-\tfrac12)\}_\rho
 \right),                                             \tag{552}
\end{equation}
where \(\operatorname{fp}\) denotes the finite-part bulk boundary and the
second component lies in the residue space controlled by (549).  Formula
(552) is source-defined by the completed Gamma--Euler logarithmic
connection and row~(c); it is not asserted to be a passive Hilbert port.
The second coordinate is the defect created by continuation across the
Euler half-plane, and its off-line antisymmetric projection is (516).
\end{definition}

\begin{theorem}[Exact factor-through-\(H_T\) test for the desired
conservation --- \PROVED]
At a first-loss parameter let \(H_T\ge0\) have compact resolvent and let
\(P_{0,T}\) project onto \(\ker H_T\).  Let
\[
 B_T:=\mathcal N_T^*\mathcal N_T+
       \mathcal R_{\Gamma,T}^*\mathcal R_{\Gamma,T}\ge0,
\]
where \(\mathcal R_{\Gamma,T}\) is the explicit compensated Gamma storage
of (546).  A conservation identity of the requested form
\begin{equation}
 B_T=H_TX_T+X_T^*H_T                                  \tag{553}
\end{equation}
with bounded \(X_T\) necessarily satisfies
\begin{equation}
 P_{0,T}B_TP_{0,T}=0
 \quad\Longleftrightarrow\quad
 \mathcal N_TP_{0,T}=0,\quad
 \mathcal R_{\Gamma,T}P_{0,T}=0.                      \tag{554}
\end{equation}
Conversely, if the two vanishings in (554) hold, the Sylvester formula
(504) constructs a bounded \(X_T\) satisfying (553).  Thus existence of
the physical conservation identity is equivalent, on the first-loss
kernel, to the terminal vanishing it is intended to prove.
\end{theorem}

\begin{proof}
Compress (553) on both sides by \(P_{0,T}\).  Since
\(H_TP_{0,T}=0\), its right side vanishes.  Positivity of \(B_T\) gives
\[
 P_{0,T}B_TP_{0,T}
 =(\mathcal N_TP_{0,T})^*(\mathcal N_TP_{0,T})
  +(\mathcal R_{\Gamma,T}P_{0,T})^*
    (\mathcal R_{\Gamma,T}P_{0,T}),
\]
which vanishes exactly under (554).  Conversely, (554) removes the
kernel--kernel block of \(B_T\); the positive spectral gap of \(H_T\) off
its kernel makes the entrywise Sylvester solution (504) bounded and proves
(553).
\end{proof}

\paragraph{Reverse-engineering consequence --- \AUDIT.}
The continuous space and orbit summability are now constructed, and (552)
is the correctly typed completed G39 object.  However, defining \(X_T\) by
the Sylvester solution, a pseudoinverse, or by declaring (553) is circular:
(553) already contains (554).  A noncircular closure must give an
independent source formula for \(X_T\)---for example a Green/virial
operator obtained before first loss---and verify (553) directly from the
Gamma--Fourier--Tate commutator.  Positivity (539) guarantees that such a
formula would have the correct sign; it does not supply the
factor-through-\(H_T\) property.

There is an even sharper obstruction.  By (557),
\(B_T\ge\frac14G_{\Gamma,T}\), and the latter form is strictly positive on
every nonzero compact-window source.  Hence, if (553) were proved from the
Gamma--Fourier--Tate construction for every \(T\), then for
\(F\in\ker H_T\)
\[
 0=\langle(H_TX_T+X_T^*H_T)F,F\rangle
  =\langle B_TF,F\rangle
 \ge\tfrac14\langle G_{\Gamma,T}F,F\rangle
\]
would force \(F=0\).  Thus the requested physical conservation would
indeed prove the no-first-loss theorem and close G40; constructing it is
not a residual convergence or normalization task.  It is the missing
Hodge theorem in operator form.

\begin{theorem}[The logarithmic virial cannot be the conservation operator
--- \PROVED]
Let \(\mathcal R F(t)=F(-t)\), \(QF(t)=tF(t)\), and let \(P_T\) denote
the primitive projection.  The first-loss operator \(H_T\) commutes with
\(\mathcal R\), whereas \(P_TQP_T\) anticommutes with \(\mathcal R\).
Consequently
\begin{equation}
 V_T:=i\bigl[H_T,P_TQP_T\bigr]                        \tag{555}
\end{equation}
anticommutes with parity and
\begin{equation}
 \langle V_TF,F\rangle=0
 \qquad\text{for every even or odd }F.                \tag{556}
\end{equation}
On the other hand, the proposed positive boundary operator satisfies
\begin{equation}
 B_T\ge\mathcal R_{\Gamma,T}^*\mathcal R_{\Gamma,T}
 \ge\frac14G_{\Gamma,T},                              \tag{557}
\end{equation}
and is therefore strictly positive on every nonzero compact-window source.
Thus \(V_T\ne B_T\); the logarithmic generator which creates the vertical
flux (527) cannot by itself furnish the factorization (553).
\end{theorem}

\begin{proof}
The symmetric window, the two Tate constraints and the even multiplier
\(q_T\) imply \([H_T,\mathcal R]=[P_T,\mathcal R]=0\), while
\(\mathcal RQ=-Q\mathcal R\).  Hence \(V_T\mathcal R=-\mathcal RV_T\).
The two parity sectors are orthogonal, which proves (556).  Inequality
(557) is (539).  The Gamma form is zero only when the Fourier transform is
supported at \(\tau=0\); an \(L^2\) compact-window source with that property
is zero.  This proves strict positivity and the no-identification claim.
\end{proof}

\begin{proposition}[Scalar dilation virial also has the wrong sign ---
\PROVED]
On the full-line core let \(\mathsf A\) be the selfadjoint dilation
generator.  For the scattering multiplier \(M_{q_T}\) of (370),
\begin{equation}
 i[M_{q_T},\mathsf A]=M_{\tau q_T'(\tau)}
 \quad\text{up to the fixed sign convention for }\mathsf A.  \tag{558}
\end{equation}
The multiplier \(\tau q_T'(\tau)\) changes sign at arbitrarily large
positive frequencies whenever at least one Euler place is active.  Hence
neither choice of sign in (558) can equal the positive operator \(B_T\),
even after restriction to primitive compact-window wave packets.
\end{proposition}

\begin{proof}
The commutator calculation is the standard identity for the generator of
Fourier dilations.  By (363) and (370), \(q_T'\) is the sum of
\(g_\Gamma'(\tau)=O(1/\tau)\) and a nonzero finite trigonometric
almost-periodic polynomial with zero mean.  Such a real nonconstant
derivative takes both signs recurrently; the \(O(1/\tau)\) term cannot
remove those sign changes at large \(\tau\).  Multiplication by
\(\tau>0\) preserves them.  Long modulated bumps supported in \(I_T\),
followed by the finite-rank correction imposing the two Tate moments,
concentrate their Fourier mass near either sign region.  Therefore the
compressed quadratic form also takes both signs, whereas (557) does not.
\end{proof}

\paragraph{Surviving conservation architecture --- \OPEN.}
The source operator \(X_T\) in (553) cannot be a scalar logarithmic or
dilation virial.  It must be the order-zero, two-contour, cross-place
scattering anticipated in (223p): a nonlocal operator mixing the two Euler
orientations with the Gamma affine state and Tate boundary before
compression.  Its missing physical intertwining (223j) is exactly the
remaining source formula needed to turn the completed graph port (552) into
the conservation (553).

\begin{theorem}[Canonical two-contour scattering and its exact inertia ---
\PROVED]
Let \(\Theta_T=e^{i\varphi_T}\) be the explicit unimodular
Gamma--Euler scattering (369), so that \(\varphi_T'=q_T\).  On
\(L^2(\mathbb R)\oplus L^2(\mathbb R)\) define the multiplier
\begin{equation}
 \mathscr S_T^{(2)}(\tau):=
 \begin{pmatrix}
 0&\Theta_T(\tau)\\
 \overline{\Theta_T(\tau)}&0
 \end{pmatrix}.                                      \tag{559}
\end{equation}
Then \(\mathscr S_T^{(2)}\) is a selfadjoint unitary and its
Wigner--Smith operator is
\begin{equation}
 -i(\mathscr S_T^{(2)})^*(\mathscr S_T^{(2)})'
 =\begin{pmatrix}-q_T&0\\0&q_T\end{pmatrix}.          \tag{560}
\end{equation}
For the second-channel source embedding
\(\iota_2F=(0,\widehat{E_TF})\), compression gives
\begin{equation}
 \iota_2^*
 \left[-i(\mathscr S_T^{(2)})^*(\mathscr S_T^{(2)})'\right]
 \iota_2=A_T.                                        \tag{561}
\end{equation}
Thus the canonical scalar two-contour unitary shadow exists unconditionally,
but its conserved infinitesimal metric has Krein inertia
\((\infty,\infty)\); unitarity alone does not imply row~(d).
\end{theorem}

\begin{proof}
Pointwise multiplication shows
\((\mathscr S_T^{(2)})^*=\mathscr S_T^{(2)}\) and
\((\mathscr S_T^{(2)})^2=I\).  Writing
\(\Theta_T=e^{i\varphi_T}\) and differentiating gives (560).  Equation
(371) says that compression of the multiplier \(q_T\) to the physical
window is exactly \(A_T\), proving (561).
\end{proof}

\begin{corollary}[Hilbertization of the two-contour scattering is exactly
G40 --- \PROVED]
Let
\[
 \mathscr M_T:=
 \overline{\left\{
 \binom{\mathcal X_TP_TF}{\mathcal Y_TP_TF}:F\in L^2(I_T)}
 \right\}}
\]
be the physical graph inside the positive/negative port space with
fundamental symmetry \(J=\operatorname{diag}(I,-I)\).  The following are
equivalent:
\begin{enumerate}
\item the Krein form induced by (560) is nonnegative on the physical
      primitive graph;
\item \(\mathscr M_T\) is the graph of a contraction from the
      \(\mathcal X_T\)-port to the \(\mathcal Y_T\)-port;
\item \(P_TA_TP_T\ge0\);
\item G40 holds at \(T\).
\end{enumerate}
At a first-loss vector the completed realization (552) exhibits
\(\mathcal N_TF\) as the negative residue coordinate whose vanishing would
be required by the desired Hilbert conservation.
\end{corollary}

\begin{proof}
On a source vector the graph Krein form is
\[
 \|\mathcal X_TP_TF\|^2-\|\mathcal Y_TP_TF\|^2
 =\langle P_TA_TP_TF,F\rangle
\]
by (84).  A subspace of a fundamental decomposition is nonnegative exactly
when it is the graph of a contraction, proving the first three
equivalences; the fourth is (84c).  The residue realization
(508)--(516) identifies the negative coordinate at equality.
\end{proof}

\paragraph{Two-contour verdict --- \AUDIT.}
Equation (559) constructs the canonical scalar two-contour \emph{shadow}
of the scattering requested abstractly in (223p), and (560)--(561) give
its exact scalar completed compression.  It does not construct the
source-labelled Euler--Gamma--Tate intertwining required in (223p).
What it supplies is the authoritative Krein conservation, not the positive
conservation (553).
Converting its physical graph into a Hilbert-contractive graph is
equivalent to G40 by the preceding corollary.  Therefore a Potapov,
Julia, graph or polar \emph{Hilbertization} of (559) cannot close row~(d)
unless an independent theorem proves positivity of that graph.

\begin{theorem}[Backward-engineered physical storage and its unavoidable
first-contact pole --- \PROVED]
On every pre-contact window on which \(H_T>0\), the positive boundary form
\(B_T\) has the canonical observability storage defined weakly by
\begin{equation}
 \langle X_T^{\rm obs}f,g\rangle:=\int_0^\infty
 \langle B_T^{1/2}e^{-uH_T}f,
         B_T^{1/2}e^{-uH_T}g\rangle\,du .             \tag{562}
\end{equation}
The integral defines a bounded operator, \(X_T^{\rm obs}\ge0\), and
\begin{equation}
 H_TX_T^{\rm obs}+X_T^{\rm obs}H_T=B_T.              \tag{563}
\end{equation}
More generally, if \(X_T\) is any bounded solution of (563) and
\(H_Tf_T=\lambda_Tf_T\), \(\|f_T\|=1\), then
\begin{equation}
 2\lambda_T\,\Re\langle X_Tf_T,f_T\rangle
 =\langle B_Tf_T,f_T\rangle .                        \tag{564}
\end{equation}
Consequently, if \(T\uparrow T_*\), \(\lambda_T\downarrow0\), and the
lowest eigenvectors converge to a nonzero first-contact vector \(f_*\),
then every such family of solutions satisfies

\[
 \|X_T\|\ \ge\
 \frac{\langle B_Tf_T,f_T\rangle}{2\lambda_T}
 \longrightarrow\infty .
\]
Here the numerator has a strictly positive lower limit because
\(B_T\ge\frac14G_{\Gamma,T}\) and the Gamma form is strictly positive on
the nonzero compact-window vector \(f_*\).
\end{theorem}

\begin{proof}
The established Gamma-form bounds give \(B_T\le C_T(H_T+I)\) as closed
forms.  If \(c_T:=\inf\operatorname{spec}H_T>0\), then
\(H_T+I\le(1+c_T^{-1})H_T\).  The spectral theorem and
\[
 \int_0^\infty
 \|H_T^{1/2}e^{-uH_T}f\|^2du=\tfrac12\|f\|^2
\]
therefore prove that (562) is a bounded positive form.  Applying the same
calculation after polarization, or first on spectral truncations and then
passing to the form limit, proves (563).  Taking the matrix
coefficient of (563) against an eigenvector gives (564).  Form continuity,
compact resolvent, lower semicontinuity of the Gamma form and the lower
bound (557) give the final positive lower limit.
\end{proof}

\paragraph{Meaning of the backward construction --- \AUDIT.}
Formula (562) is the unique positive storage selected by decay at infinite
semigroup time.  It implements the requested engineering specifications on
every interval already known to be positive.  The computation (564) shows
exactly what a source-labelled Gamma--Euler--Tate construction would have
to add: an independent theorem excluding the \(1/\lambda_T\) pole.  Such a
pole cancellation is equivalent to absence of a first-contact vector; it
cannot follow from a different normalization of the same Lyapunov
equation.

\begin{corollary}[Exact uniform-observability closure criterion ---
\PROVED]
Put
\[
 \mathcal O_TF:=
 \bigl(\mathcal N_TF,\mathcal R_{\Gamma,T}F\bigr),
 \qquad \mathcal O_T^*\mathcal O_T=B_T .
\]
The family of physical storages remains uniformly bounded on a pre-contact
interval if and only if there is a constant \(C\), independent of \(T\), such
that
\begin{equation}
 \boxed{\int_0^\infty
 \|\mathcal O_Te^{-uH_T}f\|^2\,du
 \le C\|f\|^2.}                                      \tag{565}
\end{equation}
If (565) is derived from the source-labelled, Tate-centered
Gamma--Euler network, then no finite first contact exists and G40 follows.
Indeed, for a normalized lowest eigenvector,
\begin{equation}
 \int_0^\infty
 \|\mathcal O_Te^{-uH_T}f_T\|^2\,du
 =\frac{\langle B_Tf_T,f_T\rangle}{2\lambda_T}.       \tag{566}
\end{equation}
\end{corollary}

\begin{proof}
The left side of (565) is exactly the quadratic form (562), so uniform
boundedness is equivalent to \(\sup_T\|X_T^{\rm obs}\|<\infty\).
For an eigenvector \(e^{-uH_T}f_T=e^{-u\lambda_T}f_T\); integration gives
(566).  At a hypothetical first contact the numerator has strictly positive
lower limit by (557), while the denominator tends to zero, contradicting
(565).  The first-loss theorem then gives G40.
\end{proof}

\paragraph{Source form of the remaining estimate --- \OPEN.}
The admissibility inequality (565) must be proved after the prime channels
have been synthesized and Tate-centered, not in their divergent orthogonal
sum (410)--(413).  Equivalently, on a parity equality state its arithmetic
component is the one-sided compensated Chebyshev current
\begin{equation}
 t\longmapsto
 \int_{e^{t-T}}^{e^{t+T}}
 x^{-1/2}F(t-\log x)\,dR_\Psi(x),\qquad t>T,          \tag{567}
\end{equation}
whose required \(L^2(T,\infty)\) membership is precisely the open estimate
(417), equivalently the Hardy condition (286).  The Gamma component of
(565) has an explicit positive realization by (535)--(546), while (281)
controls its physical exterior kernel for each fixed source.  Neither fact
by itself gives uniform infinite-semigroup-time admissibility as
\(\lambda_T\downarrow0\).  What remains is one coupled source energy
balance for the explicit Gamma storage and the centered current (567);
that balance must yield (565), rather than estimate \(R_\Psi\)
unconditionally.

\begin{theorem}[Canonical energy-quotient source lift --- \PROVED]
Let \(H_T\ge0\) and put
\[
 \mathscr E_T^0:=\operatorname{Dom}H_T^{1/2}/\ker H_T,
 \qquad
 \|[F]\|_{\mathscr E_T}:=\|H_T^{1/2}F\|_2,
\]
with Hilbert completion \(\mathscr E_T\).  Let
\(\mathcal O_T=(\mathcal N_T,\mathcal R_{\Gamma,T})\).  The source formula
\begin{equation}
 \mathscr L_T[F]:=\mathcal O_TF                         \tag{568}
\end{equation}
defines a bounded operator
\(\mathscr L_T:\mathscr E_T\to
\mathscr K_{\rm res}^{-}\oplus\mathscr K_{\Gamma,T}\)
if and only if there is \(C_T<\infty\) such that
\begin{equation}
 B_T=\mathcal O_T^*\mathcal O_T\le C_TH_T              \tag{569}
\end{equation}
as quadratic forms.  In that case
\begin{equation}
 \mathcal O_T=\widetilde{\mathscr L}_TH_T^{1/2},
 \qquad
 \int_0^\infty\|\mathcal O_Te^{-uH_T}F\|^2du
 \le\frac12\|\widetilde{\mathscr L}_T\|^2\|F\|^2,   \tag{570}
\end{equation}
where \(\widetilde{\mathscr L}_T\) is the same lift transported by the
canonical isometry
\([F]\mapsto H_T^{1/2}F\) from \(\mathscr E_T\) onto
\(\overline{\operatorname{Ran}H_T^{1/2}}\).
Conversely, any bounded source lift satisfying the first identity in
(570) gives (569).  Thus (568) is the unique canonical Hilbert source
object selected by the physical energy; its construction contains no
pseudoinverse or spectral sign choice.
\end{theorem}

\begin{proof}
Formula (568) is well defined on the quotient precisely when
\(\mathcal O_T\ker H_T=0\).  It is bounded for the energy norm precisely
when
\[
 \|\mathcal O_TF\|^2\le C_T\|H_T^{1/2}F\|^2,
\]
which is (569), and then it extends uniquely to the completion.  Conversely
the norm inequality for a bounded extension gives (569).  The factorization
in (570), after the displayed canonical identification, holds first on
\(\operatorname{Dom}H_T^{1/2}\) and then by closure.
Finally the spectral theorem gives
\[
 \int_0^\infty\|H_T^{1/2}e^{-uH_T}F\|^2du
 =\frac12\|(I-P_{0,T})F\|^2\le\frac12\|F\|^2,
\]
and the last assertion of (570) follows.
\end{proof}

\paragraph{Energy-quotient audit --- \AUDIT.}
The existing Gamma and residue estimates prove only
\(B_T\le C_T(H_T+I)\) before contact.  They do not remove the \(I\)-term
uniformly at first contact.  Removing it is exactly (569), and because
\(B_T\ge\frac14G_{\Gamma,T}\) is strictly positive on every nonzero
compact-window source, (569) would itself exclude \(\ker H_T\ne0\).
Accordingly the quotient construction identifies the correct new object,
but its boundedness may not be asserted from closed-graph or completion
alone; it requires a new source estimate at the \(H_T^{1/2}\)-energy level.

\begin{theorem}[Defect telescope for a reachable cascade --- \PROVED]
Let
\(C_j:\mathscr H_{j-1}\to\mathscr H_j\), \(1\le j\le m\), be
contractions, put \(R_0=I\) and
\(R_j=C_j\cdots C_1\), and write
\(D_j=(I-C_j^*C_j)^{1/2}\).  Then
\begin{equation}
 \boxed{I-R_m^*R_m
 =\sum_{j=1}^mR_{j-1}^*D_j^2R_{j-1}.}                \tag{570a}
\end{equation}
Consequently
\begin{equation}
 \|R_mx\|=\|x\|
 \quad\Longleftrightarrow\quad
 D_jR_{j-1}x=0\quad(1\le j\le m).                   \tag{570b}
\end{equation}
In particular, if the reachable observation
\begin{equation}
 x\longmapsto\bigl(D_jR_{j-1}x\bigr)_{j=1}^m        \tag{570c}
\end{equation}
is injective, the cascade has no nonzero one-step isometric vector.  The
same assertions hold for a countable cascade whenever \(R_m\) converges
strongly: the right side of (570a) is then the monotone strong-form sum.
\end{theorem}

\begin{proof}
For each stage,
\[
 R_{j-1}^*R_{j-1}-R_j^*R_j
 =R_{j-1}^*(I-C_j^*C_j)R_{j-1}.
\]
Summing from \(j=1\) to \(m\) telescopes and proves (570a).  Every summand
on its right is the square of \(D_jR_{j-1}x\); hence equality of the two
endpoint norms is equivalent to the simultaneous vanishing in (570b).
This also proves the injectivity assertion.  For an infinite cascade the
partial sums are increasing positive forms and equal
\(I-R_m^*R_m\); strong convergence of \(R_m\) gives the stated limit.
\end{proof}

\begin{proposition}[The compressed Euler port is not a contraction stage
--- \PROVED]
For an active prime \(p\), put \(a=\log p<2T\), \(r=p^{-1/2}\), and
\[
 W_{p,T}F:=\bigl(K_{p,T}F,rB_{p,T}F\bigr).
\]
Then
\begin{equation}
 W_{p,T}^*W_{p,T}=I+C_{p,T}+C_{p,T}^*.               \tag{570d}
\end{equation}
Moreover \(W_{p,T}\) is not a contraction on \(L^2(I_T)\).  Thus the
prime ports in (244)--(246) cannot be inserted as the contraction stages
of Theorem (570a) using their displayed Hilbert norms.
\end{proposition}

\begin{proof}
Identity (570d) is exactly (242).  Take
\(F=1_{I_T}\).  Since
\(C_{p,T}=\sum_{k\ge1}r^kS_p^k\), with a finite sum, and zero-extension
translations preserve positivity,
\[
 \Re\langle C_{p,T}F,F\rangle
 =\sum_{1\le k<2T/a}r^k(2T-ka)>0;
\]
the first term is present because \(a<2T\).  Hence
\[
 \|W_{p,T}F\|^2
 =\|F\|^2+2\Re\langle C_{p,T}F,F\rangle>\|F\|^2.
\]
This proves the claim.
\end{proof}

\begin{theorem}[Complete observability of a compressed prime Julia cell
--- \PROVED]
Let \(a=\log p<2T\), \(r=p^{-1/2}\), \(s=(1-r^2)^{1/2}\), and let
\(S=S_a\) be the zero-extension shift on \(L^2(I_T)\).  Put
\begin{equation}
 \Theta_{p,T}:=(S-rI)(I-rS)^{-1},\qquad
 \mathcal B_{p,T}:=(I-S^*S)^{1/2}(I-rS)^{-1}.        \tag{570e}
\end{equation}
If \(N\) is any integer with \(Na\ge2T\), then
\begin{align}
 I-\Theta_{p,T}^*\Theta_{p,T}
 &=s^2\mathcal B_{p,T}^*\mathcal B_{p,T},            \tag{570f}\\
 \bigcap_{k=0}^{N-1}
 \ker(\mathcal B_{p,T}\Theta_{p,T}^k)&=\{0\}.        \tag{570g}
\end{align}
Thus the compressed Blaschke--Julia transfer is a contraction and its
boundary defect is completely observable in finitely many passes.  In
particular it has no nonzero invariant subspace on which it acts
isometrically.
\end{theorem}

\begin{proof}
Write \(R=(I-rS)^{-1}\) and \(D=(I-S^*S)^{1/2}\).  Multiplication by
\(I-rS^*\) on the left and \(I-rS\) on the right reduces (570f) to
\[
 (I-rS^*)(I-rS)-(S^*-rI)(S-rI)
 =(1-r^2)(I-S^*S),
\]
which proves the defect identity.

It remains to prove observability.  Since \(S^N=0\), all analytic functions
of \(S\) are represented by polynomials modulo \(z^N\).  For
\(0\le k<N\),
\[
 R\Theta_{p,T}^k
 =\frac{(S-rI)^k}{(I-rS)^{k+1}}=:q_k(S).
\]
The \(N\) functions \(q_0,\ldots,q_{N-1}\) form a basis of
\(\mathbb C[z]/(z^N)\).  Indeed, if
\(\sum_{k=0}^{N-1}c_kq_k(z)\) is divisible by \(z^N\), multiplication by
\((1-rz)^N\) gives the polynomial
\[
 \sum_{k=0}^{N-1}c_k(z-r)^k(1-rz)^{N-k-1},
\]
of degree at most \(N-1\), divisible by \(z^N\), hence zero.  The summands
are linearly independent because after division by
\((1-rz)^{N-1}\) they are the distinct powers of the nonconstant M\"obius
coordinate \((z-r)/(1-rz)\).

Now suppose \(Dq_k(S)F=0\) for every \(k<N\).  The basis statement gives
\(DS^jF=0\) for \(0\le j<N\).  The projection \(D\) restricts a function to
the boundary strip of width \(a\), while \(DS^j\) restricts it, after an
isometric translation, to the \(j\)-th successive strip.  Those strips
cover \(I_T\) because \(Na\ge2T\).  Hence \(F=0\), proving (570g).
If \(\Theta_{p,T}\) were isometric on a nonzero invariant subspace, (570f)
would make every observation in (570g) vanish there, a contradiction.
\end{proof}

\begin{corollary}[M\"obius-preconditioned Euler observation --- \PROVED]
With \(W_{p,T}\) from (570d),
\begin{equation}
 W_{p,T}(I-rS)F
 =\bigl(sF,r(I-S^*S)^{1/2}F\bigr),                  \tag{570h}
\end{equation}
and therefore
\begin{equation}
 \|W_{p,T}(I-rS)F\|^2
 =\|F\|^2-r^2\|SF\|^2\le\|F\|^2.                   \tag{570i}
\end{equation}
The preconditioner is precisely the local M\"obius factor occurring in
(150) and (205)--(207).
\end{corollary}

\begin{proof}
Equations (241) give
\(K_{p,T}(I-rS)=sI\) and
\(B_{p,T}(I-rS)=(I-S^*S)^{1/2}\), proving (570h).
Taking norms and using
\(I-S^*S=D^2\) and \(s^2=1-r^2\) gives (570i).
\end{proof}

\begin{theorem}[Winding correction converts the Julia defect into a Krein
boundary flux --- \PROVED]
Let \(z=e^{-i\tau a}\).  The winding-corrected physical factor in
(377) is
\begin{equation}
 z^{-1}\frac{z-r}{1-rz}=\frac{1-rz^{-1}}{1-rz}.      \tag{570j}
\end{equation}
Its zero-extension realization is therefore the source-defined operator
\begin{equation}
 \mathcal C_{p,T}^{\rm wind}
 :=(I-rS^*)(I-rS)^{-1}.                              \tag{570k}
\end{equation}
Put \(D_-=(I-S^*S)^{1/2}\) and
\(D_+=(I-SS^*)^{1/2}\).  For every \(u\in L^2(I_T)\),
\begin{equation}
 \boxed{
 \|(I-rS)u\|^2-
 \|\mathcal C_{p,T}^{\rm wind}(I-rS)u\|^2
 =r^2\bigl(\|D_+u\|^2-\|D_-u\|^2\bigr).}            \tag{570l}
\end{equation}
Thus the winding correction changes the positive Julia defect (570f) into
a two-sided Krein boundary flux.  If \(u\) has either reflection parity,
the two boundary norms in (570l) agree and this local physical flux is zero.
\end{theorem}

\begin{proof}
Equation (570j) is elementary, and replacing the boundary variables
\(z,z^{-1}\) by the zero-extension pair \(S,S^*\) gives (570k) in the
displayed source order.  Set \(A=I-rS\).  Since
\(\mathcal C_{p,T}^{\rm wind}A=A^*\), the left side of (570l) is
\[
 \|Au\|^2-\|A^*u\|^2
 =\langle(A^*A-AA^*)u,u\rangle.
\]
The terms linear in \(r\) cancel and
\[
 A^*A-AA^*=r^2(S^*S-SS^*)=r^2(D_+^2-D_-^2),
\]
which proves (570l).  Reflection interchanges the two boundary strips and
preserves the norm of a parity vector, proving the last assertion.
\end{proof}

\begin{proposition}[An injective Gamma feature is not an injective angular
defect --- \PROVED]
The following three properties, by themselves, do not exclude a nonzero
one-step isometric vector:
\begin{enumerate}
\item a source angular contraction \(C\) with an exact defect law;
\item an injective positive Gamma observation \(G\);
\item injectivity of the input source feature.
\end{enumerate}
Indeed, on \(\mathbb C^2\) take
\begin{equation}
 X=I,\qquad C=Y=\begin{pmatrix}1&0\\0&0\end{pmatrix},
 \qquad G=I.                                         \tag{570m}
\end{equation}
Then \(X\) and \(G\) are injective and
\[
 \|Xf\|^2-\|CXf\|^2
 =\left\langle
   \begin{pmatrix}0&0\\0&1\end{pmatrix}f,f
  \right\rangle\ge0,
\]
but \(e_1\ne0\) is an isometric vector.  Therefore strict Gamma
observability can close the first-contact argument only after a
source-derived factorization places a nonzero Gamma coordinate inside the
\emph{angular defect} (570a), not merely inside the input feature.  In the
present notation that additional assertion is the energy factorization
(569), or an equivalent reachable feedback identity.
\end{proposition}

\begin{proof}
All statements follow by direct matrix multiplication.  In particular
\(Ce_1=e_1\), whereas \(Ge_1=e_1\ne0\).  Thus injectivity of \(G\) has no
logical consequence for equality unless \(G\) is controlled by the defect
of \(C\).
\end{proof}

\begin{theorem}[Scattering norm defect and Wigner--Smith delay are distinct
--- \PROVED]
Let \(a=\log p\), \(r=p^{-1/2}\), \(s=(1-r^2)^{1/2}\), and put
\[
 b_p(z)=\frac{z-r}{1-rz},\qquad |z|=1.
\]
On the bilateral unitary shift \(U\), the operator \(b_p(U)\) is unitary,
so
\begin{equation}
 I-b_p(U)^*b_p(U)=0.                                 \tag{570n}
\end{equation}
Nevertheless its Wigner--Smith delay is the strictly negative operator
\begin{equation}
 Q_p^{\rm WS}
 =-a s^2(I-rU)^{-*}(I-rU)^{-1}<0.                   \tag{570o}
\end{equation}
After winding correction, its delay is
\(aI+Q_p^{\rm WS}\), which is generally nonzero and sign-changing.

For the zero-extension shift \(S=S_a\), equations (242)--(243) give the
two exact finite-window prime-delay forms
\begin{align}
 Q_{p,T}^{\rm Cl}
 &=a\bigl(I-K_{p,T}^*K_{p,T}\bigr),\notag\\
 Q_{p,T}^{\rm phys}
 &=a\bigl(I-K_{p,T}^*K_{p,T}
              -r^2B_{p,T}^*B_{p,T}\bigr).           \tag{570p}
\end{align}
Thus the completely observable norm defect in (570f) and the boundary flux
in (570l) are not the angular defect (248).  They control compression
leakage, whereas (570p) controls stored delay.
\end{theorem}

\begin{proof}
The scalar function \(b_p\) is unimodular on the unit circle, proving
(570n) by functional calculus.  Direct differentiation at
\(z=e^{-i\tau a}\) gives
\[
 -i\overline{b_p(e^{-i\tau a})}
       \frac d{d\tau}b_p(e^{-i\tau a})
 =-\frac{a(1-r^2)}{|1-re^{-i\tau a}|^2},
\]
which is the spectral multiplier of (570o) and is strictly negative.
Multiplication by the winding \(e^{i\tau a}\) adds \(aI\), as in
(377)--(378).

On the compressed interval,
\(K_{p,T}=s(I-rS)^{-1}\), so the first line of (570p) is the compression
of \(aI+Q_p^{\rm WS}\).  The second line is exactly (243).  Their
difference is the positive boundary-storage correction
\(-ar^2B_{p,T}^*B_{p,T}\), proving every assertion.
\end{proof}

\paragraph{Cascade consequence --- \AUDIT.}
Theorem (570a)--(570c) gives a noncircular closing mechanism: if the
Gamma--Euler--Tate source system can be interconnected as contractions and
one of its reachable defects contains the injective residual Gamma square
of (539), then an equality vector must vanish.  Proposition (570d) proves
that the already displayed finite-head ports are not yet such an
interconnection.  Their local Euler supply is signed before the common
Gamma--Tate feedback is imposed.  A positive metric or feedback making the
\emph{reachable combined} stages contractive would prove (249) by (248);
it may therefore not be inferred from the lossless Julia identities alone.
The remaining constructive target is now exact: derive that metric or
feedback from the source-labelled G39 conservation law and prove that its
reachable defect map is (570c), with the residual Gamma coordinate retained.

The new identities (570e)--(570p) show that the arithmetic label network
already has the required finite-pass observability after M\"obius
preconditioning.  They do not yet transfer this property to the physical
angular operator.  The physical Wigner--Smith factor uses the
winding-corrected orientation \(C_p\) in (377)--(379), whose delay is the
\emph{negative} prime contribution.  Replacing it by the Schur orientation
\(\Theta_{p,T}\) would reverse that signed delay.  The remaining
Gamma--Euler--Tate theorem must therefore prove an intertwining of the
reachable, preconditioned Julia defects with the physical orientation while
retaining the positive Gamma residual; simply changing the prime orientation
would change \(A_T\) and is not allowed.  Equation (570l) identifies the
oriented compression flux, but (570n)--(570p) prove that this flux is not
the Wigner--Smith storage whose sign defines \(A_T\).  Complete
observability of the compression boundary therefore cannot be substituted
for strictness of the angular defect.  The still required theorem must place
the injective Gamma observation inside the \emph{stored-delay defect}
(248), not merely inside a scattering norm defect.

\section{Term-by-term source construction of the conservation equation}

Let \(\mathcal R\) denote reflection and put
\(P_\varepsilon=(I+\varepsilon\mathcal R)/2\),
\(\varepsilon\in\{+1,-1\}\).  We now replace every symbolic term in (553)
by an explicit source operator on the common primitive Gamma form domain
\[
 \mathcal Q_T^{\rm prim}:=
 \mathcal P_T\cap\operatorname{Dom}G_{\Gamma,T}^{1/2}.
\]

\begin{theorem}[Explicit construction of all positive terms in (553) ---
\PROVED]
Use the Fourier convention of (385).  For \(n\ge0\), define
\begin{equation}
 (\mathcal X_{\Gamma,n}F)(\tau)
 :=h_n(\tau)^{1/2}\widehat{E_TF}(\tau),              \tag{571}
\end{equation}
where \(h_n\) is (385), so that
\[
 \sum_{n\ge0}\|\mathcal X_{\Gamma,n}F\|_2^2
 =\langle G_{\Gamma,T}F,F\rangle .
\]
For each parity \(\varepsilon\), let \(u_\varepsilon\) be the scalar
functional in (536), let \(c_\varepsilon(T)\) be the corresponding
coefficient, and choose the Riesz vector \(v_\varepsilon\) in the
\(\mathcal X_{\Gamma,0}\)-range as in (541)--(546).  More explicitly, with
\(r_\varepsilon^0\) the mean-zero representatives constructed in
(540)--(542), one may take, up to the fixed Fourier conjugation convention,
\[
 v_\varepsilon(\tau)=
 \frac{\widehat{r_\varepsilon^0}(\tau)}{\sqrt{h_0(\tau)}}.
\]
This belongs to \(L^2(d\tau/2\pi)\) because
\(\widehat{r_\varepsilon^0}(0)=0\) and
\(h_0(\tau)=16\tau^2/(1+4\tau^2)\).  Thus
\begin{equation}
 u_\varepsilon(F)
 =\langle\mathcal X_{\Gamma,0}P_\varepsilon F,
                 v_\varepsilon\rangle,\qquad
 \theta_\varepsilon
 :=c_\varepsilon\|v_\varepsilon\|^2
 \quad(c_\varepsilon>0).                             \tag{572}
\end{equation}
Put \(e_\varepsilon=v_\varepsilon/\|v_\varepsilon\|\), and define
\[
 \mathcal L_{\varepsilon,T}F:=
 \begin{cases}
 (I-\theta_\varepsilon e_\varepsilon\otimes
 e_\varepsilon)^{1/2}\mathcal X_{\Gamma,0}
 P_\varepsilon F,&c_\varepsilon>0,\\[2mm]
 \bigl(\mathcal X_{\Gamma,0}P_\varepsilon F,\,
       (-c_\varepsilon)^{1/2}u_\varepsilon(F)\bigr),
       &c_\varepsilon\le0.
 \end{cases}                                         \tag{573}
\]
Then the compensated Gamma observation is the explicit direct sum
\begin{equation}
 \boxed{\mathcal R_{\Gamma,T}F:=
 \bigoplus_{\varepsilon=\pm1}
 \left(\mathcal L_{\varepsilon,T}F,\,
       \{\mathcal X_{\Gamma,n}P_\varepsilon F\}_{n\ge1}\right),} \tag{574}
\end{equation}
and satisfies
\begin{equation}
 \mathcal R_{\Gamma,T}^*\mathcal R_{\Gamma,T}
 =G_{\Gamma,T}-\mathcal D_T
 \ge\frac14G_{\Gamma,T}.                             \tag{575}
\end{equation}

For every off-line orbit represented by
\(\rho=\frac12+\alpha+i\gamma\), \(\alpha>0\), put
\begin{equation}
 n_\rho(t):=\sqrt{2m_\rho}\,e^{i\gamma t}\sinh(\alpha t).
                                                               \tag{576}
\end{equation}
Then, with the inner-product convention fixed by (509),
\begin{equation}
 (\mathcal N_TF)_\rho=\int_{-T}^{T}F(t)n_\rho(t)\,dt,\qquad
 \langle\mathcal N_TF,\mathcal N_TG\rangle
 =\sum_{\rho}
 \left(\int F n_\rho\right)
 \overline{\left(\int G n_\rho\right)}.              \tag{577}
\end{equation}
The series converges absolutely after polarization on
\(\mathcal Q_T^{\rm prim}\).  Consequently
\begin{align}
 \langle\mathcal N_TF,\mathcal N_TG\rangle
 &=\iint_{I_T^2}F(t)\overline{G(s)}
 \left[\sum_\rho2m_\rho e^{i\gamma(t-s)}
       \sinh(\alpha t)\sinh(\alpha s)\right]dt\,ds,\notag\\
 \langle\mathcal R_{\Gamma,T}F,\mathcal R_{\Gamma,T}G\rangle
 &=\langle G_{\Gamma,T}F,G\rangle
   -\sum_{\varepsilon=\pm1}
       c_\varepsilon u_\varepsilon(F)
       \overline{u_\varepsilon(G)},                  \tag{578a}
\end{align}
where the first bracket is interpreted as the absolutely convergent
polarized sampling form rather than as a pointwise series.  Consequently
\begin{equation}
 \boxed{\mathcal O_TF:=
       (\mathcal N_TF,\mathcal R_{\Gamma,T}F),\qquad
 B_T=\mathcal O_T^*\mathcal O_T}                    \tag{578}
\end{equation}
constructs the entire left side of (553), on one stated domain, without a
spectral square root of \(H_T\) or a sign assumption.
\end{theorem}

\begin{proof}
Tonelli and (383)--(385) give the first Gamma identity.  Equations
(541)--(546), separately on the two parity sectors, prove that (573) has
Gram \(H_{\Gamma,0}-c_\varepsilon u_\varepsilon^*u_\varepsilon\);
adding the untouched cells \(n\ge1\) and both parities gives (574)--(575).

Equation (576) is just (509) expanded using
\[
 \sqrt{\frac{m_\rho}{2}}\,
 e^{i\gamma t}(e^{\alpha t}-e^{-\alpha t})
 =\sqrt{2m_\rho}\,e^{i\gamma t}\sinh(\alpha t).
\]
The bounded sampling theorem (549)--(551) proves convergence of (577) and
its polarized form.  Orthogonal direct sum now proves (578).
\end{proof}

\begin{theorem}[Exact form equation which defines the missing source
operator --- \PROVED]
The authoritative primitive form is
\begin{equation}
 \mathfrak h_T(F,G):=
 \int_{\mathbb R}q_T(\tau)
 \widehat{E_TF}(\tau)
 \overline{\widehat{E_TG}(\tau)}
 \,\frac{d\tau}{2\pi},
 \qquad F,G\in\mathcal Q_T^{\rm prim},               \tag{579}
\end{equation}
where \(q_T\) is the explicit Gamma--Euler delay (370).  Its associated
operator is \(H_T\).

A bounded source operator \(X_T^{\rm src}\), preserving
\(\mathcal Q_T^{\rm prim}\), proves (553) if and only if, for every
\(F,G\) in that domain,
\begin{align}
 &\mathfrak h_T(X_T^{\rm src}F,G)
 +\mathfrak h_T(F,X_T^{\rm src}G)\notag\\
 &\quad=
 \sum_{\rho}
 \left(\int F n_\rho\right)
 \overline{\left(\int G n_\rho\right)}
 +\langle\mathcal R_{\Gamma,T}F,
          \mathcal R_{\Gamma,T}G\rangle .            \tag{580}
\end{align}
If any solution exists, there is one which commutes with reflection.
Hence the construction may be carried out independently on the even and
odd primitive sectors.  Selfadjointness of \(X_T^{\rm src}\) is not
asserted or required.
\end{theorem}

\begin{proof}
Equations (578)--(579) show that (580) is precisely the polarization of
(553), interpreted as an equality of closed forms.  Since \(B_T\) and
\(H_T\) commute with reflection, \(\mathcal RX\mathcal R\) is a solution
whenever \(X\) is.  Thus
\((X+\mathcal RX\mathcal R)/2\) is a parity-preserving solution.
This average preserves boundedness and the common form domain.
\end{proof}

\begin{corollary}[Gamma-only conservation is already sufficient ---
\PROVED]
It is not necessary to solve (580) term by term.  It suffices to construct
a bounded, parity-preserving source operator \(X_{\Gamma,T}^{\rm src}\)
such that
\begin{equation}
 \boxed{
 \langle\mathcal R_{\Gamma,T}F,\mathcal R_{\Gamma,T}G\rangle
 =\mathfrak h_T(X_{\Gamma,T}^{\rm src}F,G)
  +\mathfrak h_T(F,X_{\Gamma,T}^{\rm src}G).}        \tag{581}
\end{equation}
Indeed, at a first contact, \(F\in\ker H_T\) makes the right side vanish,
whereas (575) gives
\[
 \|\mathcal R_{\Gamma,T}F\|^2
 \ge\frac14\langle G_{\Gamma,T}F,F\rangle>0
 \qquad(F\ne0).
\]
Thus (581) alone excludes a nonzero first-contact vector and closes G40.
The residue term in (580) is an optional additional conservation channel,
not part of the minimal construction.
\end{corollary}

\begin{proof}
Evaluate (581) at \(G=F\in\ker H_T\).  Both form terms on the right vanish.
Strictness in (575) forces \(F=0\).  The first-loss theorem then proves the
claim.
\end{proof}

\paragraph{Term-by-term status --- \AUDIT.}
Equations (571)--(580) remove every ambiguity in the proposed identity:
the spaces, kernels, parity blocks and common form domain are explicit.
They also show that attempting to construct separate Sylvester solutions
for every residue coordinate is unnecessary.  The minimal unsolved
identity is (581).  Its left side is already an explicit Gamma ODE/direct
sum observation; its right side couples it to the \emph{total}
Gamma--Euler Wigner--Smith form \(q_T\).  A formula obtained by dividing by
\(q_T(\tau)+q_T(\sigma)\), by \(H_T^{-1}\), or by a semigroup integral
would reproduce the forbidden Sylvester construction.  A valid
\(X_{\Gamma,T}^{\rm src}\) must instead arise from the two-contour
Gamma--Euler state equations before compression.

\begin{theorem}[Cancellation of the unbounded Gamma part and bounded normal
form --- \PROVED]
On \(\mathcal P_T\), set
\begin{align}
 \mathsf G_T&:=P_TG_{\Gamma,T}P_T,\notag\\
 \mathsf W_T&:=P_T\left(m_0I+
   \sum_{2\le n<e^{2T}}w_n(S_{\log n}+S_{-\log n})\right)P_T,\notag\\
 \mathsf D_T&:=P_T\mathcal D_TP_T.                   \tag{582}
\end{align}
Then, as forms,
\begin{equation}
 H_T=\mathsf G_T-\mathsf W_T,\qquad
 \mathcal R_{\Gamma,T}^*\mathcal R_{\Gamma,T}
 =\mathsf G_T-\mathsf D_T.                           \tag{583}
\end{equation}
Consequently, writing
\begin{equation}
 X_{\Gamma,T}^{\rm src}=\frac12I+Y_T,                \tag{584}
\end{equation}
reduces (581) exactly to the bounded-source equation
\begin{equation}
 \boxed{
 \mathsf W_T-\mathsf D_T
 =H_TY_T+Y_T^*H_T.}                                  \tag{585}
\end{equation}
The left side of (585) is a bounded operator: it consists only of the
constant Tate channel, finitely many zero-extension translations and the
finite-rank oriented terminal correction.  Thus the logarithmically
unbounded Gamma multiplier has been matched completely by the universal
choice \(I/2\).
\end{theorem}

\begin{proof}
The first identity in (583) is the authoritative balanced formula
(80)--(84), compressed to the primitive space.  The second is (574)--(575).
Substitution of (584) gives
\[
 H_TX_{\Gamma,T}^{\rm src}
 +(X_{\Gamma,T}^{\rm src})^*H_T
 =H_T+H_TY_T+Y_T^*H_T.
\]
Subtract \(H_T=\mathsf G_T-\mathsf W_T\) from
\(\mathsf G_T-\mathsf D_T\) to obtain (585).  Boundedness follows because
the sum over \(n<e^{2T}\) is finite and every \(S_a\) is a contraction;
\(\mathsf D_T\) is the rank-two parity terminal form (535)--(536).
\end{proof}

\begin{theorem}[Gamma-normalized design equation --- \PROVED]
For fixed \(T\), the strictly positive operator \(\mathsf G_T\) has bounded
inverse on \(\mathcal P_T\).  Define
\begin{equation}
 K_T:=\mathsf G_T^{-1/2}\mathsf W_T\mathsf G_T^{-1/2},
 \qquad
 D_T^\natural:=
 \mathsf G_T^{-1/2}\mathsf D_T\mathsf G_T^{-1/2}.     \tag{586}
\end{equation}
Then \(K_T\) and \(D_T^\natural\) are bounded selfadjoint operators,
\(K_T\) is compact, and
\begin{align}
 H_T&=\mathsf G_T^{1/2}(I-K_T)\mathsf G_T^{1/2},
                                                               \tag{587}\\
 I-D_T^\natural&\ge\frac14I.                         \tag{588}
\end{align}
If
\[
 Z_T:=\mathsf G_T^{1/2}
 X_{\Gamma,T}^{\rm src}\mathsf G_T^{-1/2},
\]
then the Gamma-only conservation (581) is exactly
\begin{equation}
 \boxed{
 I-D_T^\natural
 =(I-K_T)Z_T+Z_T^*(I-K_T).}                          \tag{589}
\end{equation}
Equivalently, with \(Z_T=\frac12I+\mathcal Y_T\),
\begin{equation}
 \boxed{
 K_T-D_T^\natural
 =(I-K_T)\mathcal Y_T+
   \mathcal Y_T^*(I-K_T).}                           \tag{590}
\end{equation}
These are bounded operator equations built solely from the explicitly
constructed terms.
\end{theorem}

\begin{proof}
The Gamma form has compact resolvent on the finite interval and zero
kernel: equality in (82) would make the zero extension translation
invariant, hence zero.  Therefore its lowest primitive eigenvalue is
positive and \(\mathsf G_T^{-1}\) is bounded.  The operator
\(\mathsf W_T\) is bounded, while \(\mathsf G_T^{-1/2}\) is compact;
hence \(K_T\) is compact and selfadjoint.  The same argument, or finite
rank directly, applies to \(D_T^\natural\).

Conjugating (583) by \(\mathsf G_T^{-1/2}\) proves (587).
Inequality (588) is (575) under the same congruence.  Finally substitute
the displayed definition of \(Z_T\) into (581), cancel the outer
\(\mathsf G_T^{1/2}\) factors, and obtain (589).  Substitution of
\(Z_T=I/2+\mathcal Y_T\) gives (590).
\end{proof}

\paragraph{Exact remaining design datum --- \OPEN.}
The construction problem is now finite-order and explicit: build
\(\mathcal Y_T\) from the prime Julia, Gamma ODE and Tate terminal source
maps so that (590) holds and so that
\[
 X_{\Gamma,T}^{\rm src}
 =\mathsf G_T^{-1/2}
   \left(\tfrac12I+\mathcal Y_T\right)\mathsf G_T^{1/2}
\]
extends boundedly and preserves \(\mathcal Q_T^{\rm prim}\).
The formal quotient
\((K_T-D_T^\natural)/(\,2-K_T\otimes I-I\otimes K_T\,)\)
is the normalized Sylvester inverse and is forbidden.  What must be
constructed is a source factorization of the numerator through the two
copies of \(I-K_T\).  Equations (582)--(590) show precisely which local
ports must be coupled: \(K_T\) is the Gamma-normalized constant-plus-Euler
delay, while \(D_T^\natural\) is the Gamma-normalized oriented Tate
terminal.  No residue coordinate remains in the minimal equation.

\begin{theorem}[Coefficient-level engineering specification --- \PROVED]
Because \(K_T\) is compact selfadjoint, choose an orthonormal eigenbasis
\(K_Te_j=\lambda_je_j\), including a basis of its zero eigenspace.  Any
solution of (590) necessarily satisfies
\begin{equation}
 \langle(K_T-D_T^\natural)e_j,e_i\rangle
 =(1-\lambda_i)y_{ij}
  +(1-\lambda_j)\overline{y_{ji}},                    \tag{591}
\end{equation}
where \(y_{ij}=\langle\mathcal Y_Te_j,e_i\rangle\).
Writing
\[
 h_{ij}:=\frac{y_{ij}+\overline{y_{ji}}}{2},\qquad
 a_{ij}:=\frac{y_{ij}-\overline{y_{ji}}}{2},
\]
this is equivalently
\begin{equation}
 \langle(K_T-D_T^\natural)e_j,e_i\rangle
 =(2-\lambda_i-\lambda_j)h_{ij}
  +(\lambda_j-\lambda_i)a_{ij}.                       \tag{592}
\end{equation}
In particular, on a unit eigenvector with \(K_Te=e\),
\begin{equation}
 0=\langle(I-D_T^\natural)e,e\rangle\ge\frac14,       \tag{593}
\end{equation}
which is impossible.  For the selfadjoint ansatz \(a_{ij}=0\), away from
\(\lambda_i+\lambda_j=2\), (592) gives the usual abstract Sylvester
coefficient.
\end{theorem}

\begin{proof}
Take the \(e_i,e_j\) matrix coefficient of (590) to obtain (591).
Substitution of the definitions of \(h_{ij}\) and \(a_{ij}\) gives
(592).  For \(i=j\) and \(\lambda_i=1\), the left side of
(590) equals
\(\langle(I-D_T^\natural)e_i,e_i\rangle\), while its right side is zero.
Inequality (588) gives (593).  The last statement follows by division when
the coefficient is nonzero.
\end{proof}

\paragraph{Interpretation of (591)--(593) --- \AUDIT.}
This is a design specification, not a construction: using (591) to define
\(\mathcal Y_T\) is exactly the forbidden spectral Sylvester inverse.
Its value is that it identifies what a source formula must prove.  The
Gamma--Euler--Tate network must produce the coefficients in (592) with a
uniform bound while never dividing by
\(2-\lambda_i-\lambda_j\).  Such a bounded formula automatically excludes
\(\lambda=1\) by (593), which is why no further inequality would remain.

\begin{theorem}[Hermitian metric connection versus skew Kato connection ---
\PROVED]
Decompose any proposed normalized source connection uniquely as
\[
 \mathcal Y_T=S_T+A_T,\qquad S_T^*=S_T,\quad A_T^*=-A_T.
\]
Then (590) is exactly
\begin{equation}
 \boxed{
 K_T-D_T^\natural
 =\{I-K_T,S_T\}-[K_T,A_T],}                           \tag{594}
\end{equation}
where \(\{U,V\}=UV+VU\).  On every unit eigenvector
\(K_Te=\lambda e\),
\begin{equation}
 2(1-\lambda)\langle S_Te,e\rangle
 =\lambda-\langle D_T^\natural e,e\rangle .          \tag{595}
\end{equation}
The skew connection \(A_T\) makes no contribution to (595).  In
particular, at a hypothetical first contact \(\lambda=1\), (588) gives
\begin{equation}
 0=1-\langle D_T^\natural e,e\rangle\ge\frac14,       \tag{596}
\end{equation}
independently of the choice of Kato transport.
\end{theorem}

\begin{proof}
Substitute \(\mathcal Y_T=S_T+A_T\) into the right side of (590):
\[
 (I-K)(S+A)+(S-A)(I-K)
 =\{I-K,S\}+[I-K,A].
\]
Since \([I-K,A]=-[K,A]\), this is (594).  Pair with an eigenvector of
\(K\).  The diagonal matrix coefficient of a commutator with \(K\)
vanishes, while the anticommutator contributes
\(2(1-\lambda)\langle Se,e\rangle\), proving (595).
Equation (596) follows from (588).
\end{proof}

\paragraph{Connection design consequence --- \AUDIT.}
The already constructed Kato generator (112)--(116) supplies the skew
transport of the moving primitive subspaces, but no skew connection can
produce the diagonal term in (595).  The indispensable new object is the
Hermitian metric connection \(S_T\), obtained by differentiating the
\emph{completed} Euler--Gamma metric before critical compression.  The two
source calculations must verify the single rule
\begin{equation}
 K_T-D_T^\natural
 =\{I-K_T,S_T^{E\Gamma}\}
   -[K_T,A_T^{\rm Kato}] .                            \tag{597}
\end{equation}
Formula (597) is the term-by-term version of the proposed product rule:
the anticommutator is bulk metric variation and the commutator is moving
Tate/cutoff transport.  Equations (101)--(116) construct both raw
derivatives, but their critical completed identification with the left side
of (597) is precisely the still open tangent intertwining (220j)/(223j).
Defining \(S_T^{E\Gamma}\) from (595) would again be spectral division and
is not allowed.

\begin{proposition}[Completed metric connection and the compulsory residue
graph --- \PROVED/\AUDIT]
For \(t>2\), let
\(\mathfrak M_t(\tau)=
 |\Lambda_{\mathbb Q}(t/2+i\tau)|^2\) as in (117), and let
\(\mathcal J_t=M_{\mathfrak M_t^{1/2}}\) be its positive boundary
embedding.  If \(q_t\) denotes the right side of (119), then
\begin{equation}
 2\partial_t\log\mathfrak M_t=q_t,\qquad
 \mathcal J_t^{-1}\mathcal J_t'
 =\frac14M_{q_t}.                                     \tag{598}
\end{equation}
Thus the Hermitian metric connection required in (597) is explicit in the
absolutely convergent half-plane.

Continuation of (598) to \(t=1\) on an arbitrary primitive
Paley--Wiener source cannot be performed as a scalar \(L^2\) multiplier.
Crossing a zero \(\rho\) of \(\xi\) produces the residue
\begin{equation}
 -m_\rho K_F(\rho-\tfrac12),                          \tag{599}
\end{equation}
and the two functional-equation orientations have antisymmetric coordinate
exactly \(\mathcal N_TF\) by (514)--(516).  Hence the only source-typed
critical candidate supplied by (598) is the completed graph connection
\begin{equation}
 F\longmapsto
 \left(
  \operatorname{fp}\!\left[
   -\frac{\zeta'}{\zeta}(s)K_F(s-\tfrac12)\right],
  \{-m_\rho K_F(\rho-\tfrac12)\}_\rho
 \right),                                             \tag{600}
\end{equation}
namely (552), together with the Gamma connection.
\end{proposition}

\begin{proof}
Since \(\mathcal J_t\) is multiplication by
\(\mathfrak M_t^{1/2}\), its logarithmic derivative is
\(\frac12\partial_t\log\mathfrak M_t\).  Equations (118)--(119) therefore
give (598).  The logarithmic derivative of \(\zeta\) has residue
\(-m_\rho\) at \(\rho\), proving (599).  The residue calculation
(514)--(516) identifies its antisymmetric functional-equation coordinate
with \(\mathcal N_TF\).  Separating finite part and residues gives (600).
\end{proof}

\paragraph{Scope of the completed connection --- \AUDIT.}
Proposition (598)--(600) changes the construction strategy in one precise
way.  Although (581) shows that the Gamma square alone would be sufficient
\emph{after} a conservation law is available, its natural metric connection
cannot be continued critically on the full primitive source without the
residue graph.  The full left side \(B_T\) in (578) is therefore the correct
object for constructing \(X_T^{\rm src}\).  What remains is to prove that
the graph connection (600), augmented by the explicit Gamma and Kato
components, is bounded in the norm (549) and satisfies the product rule
(580).  The finite-part source formula is constructed; its Hilbert
conservation and bounded critical pullback remain \OPEN{}.

\begin{theorem}[Finite Hadamard graph connection without division by
\(\xi\) --- \PROVED]
Let
\[
 \Xi_R(s):=b_R+\sum_{\substack{\rho\\|\Im\rho|\le R}}
 m_\rho\left(\frac1{s-\rho}+\frac1{\rho}\right)
\]
be a functional-equation-symmetric finite truncation of the logarithmic
Hadamard expansion of \(\xi'/\xi\), with \(b_R\) containing the fixed
normalization and the omitted regular constants.  For a compact source
\(F\), put \(K_F(s)=\int_{-T}^{T}F(t)e^{st}\,dt\).  Define
\begin{align}
 \mathcal C_R^{\rm bulk}F(s)
 &:=
 b_RK_F(s)+
 \sum_{\substack{\rho\\|\Im\rho|\le R}}m_\rho
 \left[
  \frac{K_F(s)-K_F(\rho)}{s-\rho}
  +\frac{K_F(s)}{\rho}\right],\notag\\
 \mathcal C_R^{\rm res}F
 &:=
 \{-m_\rho K_F(\rho)\}_{|\Im\rho|\le R}.             \tag{601}
\end{align}
Then \(\mathcal C_R^{\rm bulk}F\) is entire and
\begin{equation}
 \Xi_R(s)K_F(s)
 =\mathcal C_R^{\rm bulk}F(s)
  -\sum_{\substack{\rho\\|\Im\rho|\le R}}
       \frac{(\mathcal C_R^{\rm res}F)_\rho}{s-\rho}.
                                                               \tag{602}
\end{equation}
Thus (601) is an exact finite graph decomposition of the completed
logarithmic connection; no quotient by \(\xi\), zero-free region or
pseudoinverse is used.
\end{theorem}

\begin{proof}
For each zero,
\[
 \frac{K_F(s)}{s-\rho}
 =\frac{K_F(s)-K_F(\rho)}{s-\rho}
  +\frac{K_F(\rho)}{s-\rho}.
\]
The first quotient has a removable singularity and is entire.  Substitute
this identity term by term in the finite Hadamard sum.  The second term is
the negative of the residue coordinate in (601), proving (602).
\end{proof}

\begin{theorem}[Finite-strip boundedness of the graph connection ---
\PROVED]
Fix \(T,R,L<\infty\) and a vertical strip
\(|\Re s-\frac12|\le\frac12\).  There is
\(C_{T,R,L}<\infty\) such that
\begin{align}
 &\int_{-L}^{L}
 \left|\mathcal C_R^{\rm bulk}F
       (\tfrac12+i\tau)\right|^2d\tau
 +\sum_{\substack{\rho\\|\Im\rho|\le R}}
       m_\rho|K_F(\rho)|^2
 \le C_{T,R,L}\|F\|_{L^2(I_T)}^2 .                  \tag{603}
\end{align}
Moreover every divided-difference summand in (601) has the source formula
\begin{equation}
 \frac{K_F(s)-K_F(\rho)}{s-\rho}
 =\int_{-T}^{T}tF(t)
   \int_0^1e^{(\rho+u(s-\rho))t}\,du\,dt .            \tag{604}
\end{equation}
Hence the finite graph connection is an explicit bounded integral operator,
not merely a meromorphic formal expression.
\end{theorem}

\begin{proof}
The identity
\[
 \frac{e^{st}-e^{\rho t}}{s-\rho}
 =t\int_0^1e^{(\rho+u(s-\rho))t}\,du
\]
gives (604).  On the stated strip and for a finite set of zeros,
the exponential is uniformly bounded by \(e^{T/2}\), while
\(|t|\le T\).  Cauchy--Schwarz in \(t\), followed by integration over the
finite interval in \(\tau\), bounds every divided difference.  The
ordinary \(K_F(s)\) terms obey the same estimate.  Finally the finite
collection of evaluations \(K_F(\rho)\) is bounded by
\(\|F\|_2\|e^{\rho\cdot}\|_2\).  Summing finitely many bounds proves (603).
\end{proof}

\begin{corollary}[One uniform estimate now controls the critical graph ---
\PROVED]
Suppose there is a Hilbert bulk norm \(\|\cdot\|_{\rm bulk,T}\), fixed
independently of \(R\), for which
\begin{equation}
 \sup_{R\ge1}\left(
 \|\mathcal C_R^{\rm bulk}F\|_{\rm bulk,T}^2+
 \sum_{\substack{\rho\\|\Im\rho|\le R}}
       m_\rho|K_F(\rho)|^2\right)
 \le C_T\|F\|_{\mathcal Q_{\Gamma,T}}^2.              \tag{605}
\end{equation}
Then the finite graph connections have a weakly convergent subnet defining
a bounded completed connection with residue component (548).  If, in
addition, their finite source product rules converge to (580), the limit is
the required bounded \(X_T^{\rm src}\), and G40 follows.
\end{corollary}

\begin{proof}
Bounded sets in Hilbert space are weakly relatively compact.  The residue
coordinates are compatible under truncation and their full norm is already
controlled by (549).  Hence any weak limit has residue component
\(\mathcal E_TF\).  Bounded bilinear forms pass to weak limits, so convergence
of the finite product rules gives (580).  The first-contact argument after
(581), now with the full positive observation (578), proves G40.
\end{proof}

\paragraph{Uniform graph estimate --- \OPEN.}
Equations (601)--(604) construct every finite approximation and prove its
boundedness.  The remaining analytic estimate is no longer a vague
continuation statement: choose a source-defined bulk norm for the finite
part and prove the uniform bound (605), together with convergence of the
finite product rule.  Defining that norm by the left side of (553) would be
circular.  The natural candidate is the two-contour Hardy norm with the
local principal parts removed; establishing its uniform \(R\)-bound is the
same causal boundary control isolated in (417), now with all residue
coordinates retained rather than discarded.

\begin{theorem}[Explicit finite residue interpolation without an operator
inverse --- \PROVED]
Let \(z_1,\ldots,z_M\) be the distinct spectral points
\((\rho-\frac12)/i\) belonging to a finite union of functional-equation
orbits, none equal to the two Tate points \(\pm i/2\).  Choose
\(\epsilon>0\) with \(\epsilon<T\), and let
\(\varphi_\epsilon\in C_c^\infty(-\epsilon,\epsilon)\) be a standard
nonnegative approximate identity of mass one, chosen so narrowly that
\[
 b_\epsilon(z):=
 \int\varphi_\epsilon(t)e^{izt}\,dt
\]
does not vanish at any \(z_j\).  Define the cardinal functions
\begin{equation}
 L_j(z):=
 \frac{b_\epsilon(z)(z^2+\frac14)
       \prod_{k\ne j}(z-z_k)}
      {b_\epsilon(z_j)(z_j^2+\frac14)
       \prod_{k\ne j}(z_j-z_k)}.                     \tag{606}
\end{equation}
Then \(L_j\) is a Paley--Wiener function of type at most \(\epsilon<T\),
belongs to \(L^2(\mathbb R)\), vanishes at \(\pm i/2\), and
\begin{equation}
 L_j(z_k)=\delta_{jk}.                                \tag{607}
\end{equation}
Consequently there are explicit sources
\(\ell_j\in\mathcal P_T\cap C_c^\infty(I_T)\), obtained by inverse Fourier
transform of \(L_j\), and
\begin{equation}
 \mathscr L_R(c_1,\ldots,c_M):=
 \sum_{j=1}^Mc_j\ell_j                                \tag{608}
\end{equation}
is a finite-rank right inverse of the truncated residue evaluation
\(\mathcal E_{T,R}\):
\(\mathcal E_{T,R}\mathscr L_R=I\).
No inverse or pseudoinverse of \(H_T\), \(\mathsf G_T\), or an evaluation
Gram is used.
\end{theorem}

\begin{proof}
As the support of the approximate identity shrinks,
\(b_\epsilon(z_j)\to1\) at every point of the finite set, so all values
are nonzero for a sufficiently narrow choice.  The function
\(b_\epsilon\) has exponential type at most \(\epsilon\) and decays faster
than every power on the real line.  Multiplication by the polynomial in
(606) preserves rapid decay and does not increase exponential type.
Paley--Wiener therefore gives a compactly supported smooth source inside
\((-\epsilon,\epsilon)\subset I_T\).  The factor
\(z^2+\frac14\) gives the two Tate zeros, so the source is primitive.
The Lagrange factors prove (607), and linearity proves (608).
\end{proof}

\begin{theorem}[Exact finite-residue source conservation --- \PROVED]
Let \(\mathscr K_R\) be the finite residue space, \(J_R\) its fundamental
symmetry from (123)--(125), and \(P_{-,R}\) its negative spectral
projection.  Put
\begin{align}
 H_{T,R}&:=\mathcal E_{T,R}^*J_R\mathcal E_{T,R},\notag\\
 \mathcal N_{T,R}&:=P_{-,R}\mathcal E_{T,R},\notag\\
 Y_R&:=\frac12J_RP_{-,R},\notag\\
 X_{T,R}^{\rm res}&:=\mathscr L_RY_R\mathcal E_{T,R}. \tag{609}
\end{align}
Then \(X_{T,R}^{\rm res}\) is an explicit bounded finite-rank operator on
\(\mathcal P_T\), preserves smooth compact support, and satisfies
\begin{equation}
 \boxed{
 H_{T,R}X_{T,R}^{\rm res}
 +(X_{T,R}^{\rm res})^*H_{T,R}
 =\mathcal N_{T,R}^*\mathcal N_{T,R}.}               \tag{610}
\end{equation}
This is the residue part of (553), proved exactly at every finite level
without a Sylvester solution.
\end{theorem}

\begin{proof}
Equation (608) gives
\(\mathcal E_{T,R}X_{T,R}^{\rm res}
 =Y_R\mathcal E_{T,R}\).  Since \(P_{-,R}\) is a spectral projection of
\(J_R\), it commutes with \(J_R\), and \(J_R^2=I\).  Therefore
\[
 J_RY_R+Y_R^*J_R
 =\tfrac12P_{-,R}+\tfrac12P_{-,R}=P_{-,R}.
\]
Pull this identity back by \(\mathcal E_{T,R}\) to obtain (610).
Boundedness and preservation of the smooth core follow from the finite sum
of the explicit sources (606)--(608).
\end{proof}

\begin{proposition}[Exact tail error against the full row-(d) form ---
\PROVED]
Decompose the full residue analysis as
\(\mathcal E_T=(\mathcal E_{T,R},\mathcal E_{T,>R})\) and write
\[
 H_T=H_{T,R}+H_{T,>R}
\]
in the absolutely convergent residue form (550).  Then
\begin{align}
 &H_TX_{T,R}^{\rm res}
 +(X_{T,R}^{\rm res})^*H_T
 -\mathcal N_{T,R}^*\mathcal N_{T,R}\notag\\
 &\qquad=
 H_{T,>R}X_{T,R}^{\rm res}
 +(X_{T,R}^{\rm res})^*H_{T,>R}.                     \tag{611}
\end{align}
Hence the finite construction converges to the residue part of (553)
provided
\begin{equation}
 \sup_R\|X_{T,R}^{\rm res}\|_{\mathcal Q_{\Gamma,T}\to
 \mathcal Q_{\Gamma,T}}<\infty
 \quad\text{and}\quad
 \|\mathcal E_{T,>R}X_{T,R}^{\rm res}\|\longrightarrow0.       \tag{612}
\end{equation}
\end{proposition}

\begin{proof}
Insert \(H_T=H_{T,R}+H_{T,>R}\) and use (610); the finite terms cancel and
give (611).  Estimate the remaining polarized tail by Cauchy--Schwarz in
the residue Hilbert space (549)--(550).  Conditions (612), together with
strong convergence of the residue tail on each fixed source, make both
terms tend to zero.
\end{proof}

\paragraph{Finite-to-global residue gate --- \AUDIT.}
The construction (606)--(610) is the first explicit source formula for a
nontrivial part of \(X_T^{\rm src}\) that uses neither \(H_T^{-1}\) nor a
Moore--Penrose lift.  Its global obstruction is now quantitative:
Lagrange interpolation at a zero set of density \(R\log R\) need not have
uniform norm, and (612) is exactly the stable residue-range lifting problem
(523a).  Thus the residue component is completely solved at every finite
level; proving uniform stability, rather than finding another algebraic
formula, is the remaining infinite-dimensional task.

\begin{theorem}[No global fiberwise lift of the isolated negative residue
connection --- \PROVED]
Assume there is at least one off-line functional-equation orbit.  There is
no bounded operator \(X:\mathcal P_T\to\mathcal P_T\) satisfying
\begin{equation}
 \mathcal E_TX=\frac12J_{\rm res}P_-\mathcal E_T
 \qquad\text{on }\mathcal P_T\cap\mathcal Q_{\Gamma,T}.       \tag{613}
\end{equation}
Consequently the exact finite constructions \(X_{T,R}^{\rm res}\) in
(609) cannot converge to a same-window source operator implementing the
negative connection orbit by orbit.  Any global \(X_T^{\rm src}\) proving
(580) must mix the finite-part/Gamma channel with the residue channel before
the critical limit.
\end{theorem}

\begin{proof}
On every critical-line orbit, \(P_-=0\).  Thus (613) would imply that the
Paley--Wiener transform of \(XF\) vanishes at every distinct critical-line
zero.  The unconditional critical-zero density used in the rigidity theorem
(292), together with the Cartwright bound for an entire function of fixed
exponential type, forces \(XF=0\).

Now choose an off-line orbit and, by finite-orbit interpolation
(520)--(522), choose \(F\) for which
\(P_-\mathcal E_TF\ne0\).  The right side of (613) is then nonzero, whereas
the left side is \(\mathcal E_T(XF)=0\), a contradiction.
\end{proof}

\paragraph{Correction to the uniform residue strategy --- \AUDIT.}
Condition (612) is a sufficient formal limit criterion, but the preceding
theorem proves that it cannot hold for the isolated fiberwise construction
(609) when an off-line orbit is present.  Therefore the residue and bulk
limits may not be taken separately.  The viable object is the joint
Hadamard graph (601)--(602): its divided-difference bulk is exactly what
prevents the critical coordinates from being forced to zero while the
negative off-line coordinate remains nonzero.  The remaining estimate must
be imposed on this \emph{joint} graph and its joint product rule, not on a
standalone residue interpolation operator.

\begin{theorem}[Source realization of the finite Hadamard bulk by Volterra
operators --- \PROVED]
For a zero \(\rho\) define, on \(I_T\),
\begin{equation}
 (\mathcal V_\rho F)(u):=
 \begin{cases}
 \displaystyle\int_u^T F(t)e^{\rho(t-u)}\,dt,&u>0,\\[2mm]
 \displaystyle-\int_{-T}^uF(t)e^{\rho(t-u)}\,dt,&u<0.
 \end{cases}                                         \tag{614}
\end{equation}
Then \(\mathcal V_\rho\) is a bounded Volterra operator on \(L^2(I_T)\),
preserves support in \(I_T\), and
\begin{equation}
 K_{\mathcal V_\rho F}(s)
 =\frac{K_F(s)-K_F(\rho)}{s-\rho}.                   \tag{615}
\end{equation}
Consequently the finite Hadamard bulk in (601) is the transform of the
explicit source operator
\begin{equation}
 \mathcal V_R^{\rm Had}
 :=b_RI+
 \sum_{\substack{\rho\\|\Im\rho|\le R}}m_\rho
 \left(\mathcal V_\rho+\frac1{\rho}I\right),          \tag{616}
\end{equation}
where the scalar \(b_R\) has the same normalization as in (601):
\begin{equation}
 K_{\mathcal V_R^{\rm Had}F}
 =\mathcal C_R^{\rm bulk}F.                          \tag{617}
\end{equation}
Thus both components of the finite joint graph are source-realized:
the bulk by (614)--(617), and the residues by direct evaluations in (601).
\end{theorem}

\begin{proof}
For \(t>0\),
\[
 \frac{e^{st}-e^{\rho t}}{s-\rho}
 =\int_0^te^{su}e^{\rho(t-u)}\,du,
\]
and the same oriented integral identity holds for \(t<0\).  Multiply by
\(F(t)\), integrate in \(t\), and exchange the order on the two triangles
\(0<u<t<T\) and \(-T<t<u<0\).  The result is (614)--(615).
Schur's test gives, uniformly for \(0\le\Re\rho\le1\),
\[
 \|\mathcal V_\rho\|\le 2Te^T.
\]
Summing the finite family and comparing with (601) proves
(616)--(617).
\end{proof}

\paragraph{Joint source colligation still required --- \OPEN.}
The operator \(\mathcal V_R^{\rm Had}\) alone is not
\(X_{T,R}^{\rm src}\): the meromorphic identity (602) also has external
residue states.  Conversely, Theorem (613) excludes lifting those states
separately.  The finite source colligation must therefore feed the residue
vector back into the Volterra bulk before taking \(R\to\infty\), so that
its transfer is (602).  Constructing that feedback and checking its
polarized product rule against (578a) is the next exact algebraic
calculation.

\begin{theorem}[Closed source conservation on one off-line orbit ---
\PROVED]
For an off-line orbit with spectral points
\(z_\pm=\gamma\mp i\alpha\), define the modulated reflection
\begin{equation}
 (\mathcal U_\gamma F)(t):=e^{-2i\gamma t}F(-t).
                                                               \tag{618}
\end{equation}
Then \(\mathcal U_\gamma\) is a selfadjoint unitary on \(L^2(I_T)\) and
\begin{equation}
 \binom{\Phi_{\mathcal U_\gamma F}(z_+)}
       {\Phi_{\mathcal U_\gamma F}(z_-)}
 =J_{\mathcal O}
  \binom{\Phi_F(z_+)}{\Phi_F(z_-)}.                  \tag{619}
\end{equation}
Put
\begin{equation}
 X_{\mathcal O}:=\frac14(\mathcal U_\gamma-I).        \tag{620}
\end{equation}
If
\(H_{\mathcal O}:=
 \mathcal E_{\mathcal O}^*J_{\mathcal O}
 \mathcal E_{\mathcal O}\) and
\(\mathcal N_{\mathcal O}:=P_{-,\mathcal O}
 \mathcal E_{\mathcal O}\), then
\begin{equation}
 \boxed{
 H_{\mathcal O}X_{\mathcal O}
 +X_{\mathcal O}^*H_{\mathcal O}
 =\mathcal N_{\mathcal O}^*\mathcal N_{\mathcal O}.} \tag{621}
\end{equation}
Thus every individual negative residue fiber has a bounded source
conservation formula of norm at most \(1/2\), independent of \(\alpha\).
\end{theorem}

\begin{proof}
A change of variables gives
\[
 \Phi_{\mathcal U_\gamma F}(z)
 =\Phi_F(2\gamma-z),
\]
which swaps \(z_+\) and \(z_-\), proving (619).  Applying (618) twice gives
the identity, and direct adjunction gives
\(\mathcal U_\gamma^*=\mathcal U_\gamma\).
On the residue coordinates, (620) therefore acts by
\[
 Y_{\mathcal O}=\frac14(J_{\mathcal O}-I)
 =\frac12J_{\mathcal O}P_{-,\mathcal O}.
\]
As in the proof of (610),
\[
 J_{\mathcal O}Y_{\mathcal O}
 +Y_{\mathcal O}^*J_{\mathcal O}=P_{-,\mathcal O}.
\]
Pullback by \(\mathcal E_{\mathcal O}\) proves (621).
\end{proof}

\paragraph{Primitive scope of the one-orbit formula --- \AUDIT.}
The modulation in (618) does not preserve the two Tate moments.  Therefore
\(X_{\mathcal O}\) is an ambient source solution, not yet an operator on
\(\mathcal P_T\).  Compressing it by \(P_T\) introduces the normal component
\((I-P_T)\mathcal U_\gamma P_T\), whose differentiated two-moment energy is
of the same rank-two type as (109) and (535).  Thus the primitive correction
cannot be discarded; it is exactly the terminal which the Gamma factor
(574)--(575) was designed to pay.

\begin{proposition}[Exact primitive compression error of one orbit ---
\PROVED]
Let \(P=P_T\), \(H_{\mathcal O}^P=PH_{\mathcal O}P\), and
\[
 X_{\mathcal O}^P:=PX_{\mathcal O}P,\qquad
 C_{\mathcal O}:=(I-P)X_{\mathcal O}P.
\]
Then
\begin{align}
 &H_{\mathcal O}^PX_{\mathcal O}^P+
   (X_{\mathcal O}^P)^*H_{\mathcal O}^P\notag\\
 &\quad=
 P\mathcal N_{\mathcal O}^*\mathcal N_{\mathcal O}P
 -PH_{\mathcal O}C_{\mathcal O}
 -C_{\mathcal O}^*H_{\mathcal O}P.                  \tag{623}
\end{align}
The range of \(C_{\mathcal O}\) lies in the two-dimensional Tate normal
space \(\operatorname{Ran}B_1^*\).  Hence every failure of the ambient
identity (621) after primitive compression is a completely explicit
rank-two cross terminal.
\end{proposition}

\begin{proof}
Compress (621) by \(P\), insert
\(X_{\mathcal O}P=PX_{\mathcal O}P+
(I-P)X_{\mathcal O}P\), and perform the same decomposition in the adjoint
term.  Moving the two normal contributions to the right gives (623).
The formula \(I-P=B_1^*(B_1B_1^*)^{-1}B_1\) proves the range assertion.
\end{proof}

\begin{proposition}[Why the one-orbit solutions cannot simply be summed ---
\PROVED]
For two distinct ordinates \(\gamma\ne\gamma'\),
\(\mathcal U_\gamma\) does not preserve the
\(\gamma'\)-orbit evaluation fiber:
\begin{equation}
 \Phi_{\mathcal U_\gamma F}(\gamma'\mp i\alpha')
 =\Phi_F(2\gamma-\gamma'\pm i\alpha'),               \tag{622}
\end{equation}
which is generally not a zeta-zero coordinate.  Consequently
\(\sum_{\mathcal O}X_{\mathcal O}\) produces uncontrolled cross-orbit
terms in the full residue form and is not the desired
\(X_T^{\rm src}\).  The Hadamard divided-difference bulk (601) is precisely
the source object capable of coupling these different ordinates before
summation.
\end{proposition}

\begin{proof}
Equation (622) is (619) evaluated at the second orbit.  Unless
\(2\gamma-\gamma'\) happens to be another prescribed zero ordinate with
the required real displacement, the image leaves that residue fiber.
Therefore the pullback computation proving (621) does not diagonalize the
sum of two orbit forms, and cross terms remain.
\end{proof}

\begin{theorem}[Universal extended-boundary solution of the conservation
equation --- \PROVED]
Let
\begin{align}
 \mathscr K_T^{\rm ext}
 &:=\mathscr K_{\rm res}\oplus
     \mathscr K_{\Gamma,T}\oplus\mathscr K_{\Gamma,T},\notag\\
 J_T^{\rm ext}
 &:=J_{\rm res}\oplus I\oplus(-I),\notag\\
 \mathscr E_T^{\rm ext}F
 &:=\bigl(\mathcal E_TF,\mathcal R_{\Gamma,T}F,
                    \mathcal R_{\Gamma,T}F\bigr),     \tag{624}\\
 P_T^{\rm obs}&:=P_-\oplus I\oplus0 .
\end{align}
Then, on the common form domain,
\begin{align}
 (\mathscr E_T^{\rm ext})^*J_T^{\rm ext}
       \mathscr E_T^{\rm ext}&=H_T,\notag\\
 (\mathscr E_T^{\rm ext})^*P_T^{\rm obs}
       \mathscr E_T^{\rm ext}&=B_T.                  \tag{625}
\end{align}
The bounded boundary operator
\begin{equation}
 Y_{0,T}:=\frac12J_T^{\rm ext}P_T^{\rm obs}          \tag{626}
\end{equation}
satisfies the universal Krein--Lyapunov identity
\begin{equation}
 \boxed{
 J_T^{\rm ext}Y_{0,T}
 +Y_{0,T}^*J_T^{\rm ext}=P_T^{\rm obs}.}             \tag{627}
\end{equation}
Thus the entire conservation equation has an explicit solution before
pullback to the physical source.
\end{theorem}

\begin{proof}
The two identical Gamma coordinates in (624) have opposite Krein signs and
cancel, leaving the residue representation (550) of \(H_T\).  Projection
by \(P_T^{\rm obs}\) retains exactly the negative residue coordinate and
one Gamma copy, giving (578) and the second identity in (625).
Since \(J_T^{\rm ext}\) commutes with \(P_T^{\rm obs}\) and
\((J_T^{\rm ext})^2=I\),
\[
 J_T^{\rm ext}Y_{0,T}
 =\tfrac12P_T^{\rm obs}
 =Y_{0,T}^*J_T^{\rm ext},
\]
which proves (627).
\end{proof}

\begin{theorem}[Exact tangency/gauge criterion for
\(X_T^{\rm src}\) --- \PROVED]
Let
\[
 \mathscr R_T^{\rm ext}:=
 \overline{\mathscr E_T^{\rm ext}
  (\mathcal Q_T^{\rm prim})}
 \subset\mathscr K_T^{\rm ext}.
\]
Suppose there is a bounded operator \(A_T\) on
\(\mathscr K_T^{\rm ext}\) such that
\begin{align}
 (J_T^{\rm ext}A_T)^*&=-J_T^{\rm ext}A_T,            \tag{628}\\
 (Y_{0,T}+A_T)\mathscr R_T^{\rm ext}
 &\subset\mathscr R_T^{\rm ext},                     \tag{629}
\end{align}
and the induced source rule
\begin{equation}
 \mathscr E_T^{\rm ext}X_T^{\rm src}F
 :=(Y_{0,T}+A_T)\mathscr E_T^{\rm ext}F              \tag{630}
\end{equation}
is represented by a bounded operator preserving
\(\mathcal Q_T^{\rm prim}\).  Then
\begin{equation}
 \boxed{B_T=H_TX_T^{\rm src}
 +(X_T^{\rm src})^*H_T.}                             \tag{631}
\end{equation}
Conversely, every bounded source solution of (631) determines on the
generated range a correction \(A_T\) satisfying (628)--(630), modulo a
Krein-null operator invisible under pullback.
\end{theorem}

\begin{proof}
Condition (628) says
\[
 J_T^{\rm ext}A_T+A_T^*J_T^{\rm ext}=0.
\]
Add this to (627), pull back by \(\mathscr E_T^{\rm ext}\), and use
(625) and (630).  This gives (631).

Conversely, let \(X_T^{\rm src}\) solve (631) and define, on the generated
source vectors,
\[
 A_T\mathscr E_T^{\rm ext}F
 :=\mathscr E_T^{\rm ext}X_T^{\rm src}F
   -Y_{0,T}\mathscr E_T^{\rm ext}F.
\]
The conservation identity and (627), after polarization, say that the
pullback of \(J A+A^*J\) is zero.  Quotienting the generated boundary by
the pullback-null subspace makes this a \(J\)-skew operator.  Any bounded
extension differs only by an invisible Krein-null gauge, proving the final
assertion.
\end{proof}

\begin{proposition}[Why the universal solution needs a nontrivial gauge ---
\PROVED]
The bare operator \(Y_{0,T}\) does not preserve the generated extended
range when \(\mathcal R_{\Gamma,T}\ne0\).  Indeed, on a source vector,
\begin{equation}
 Y_{0,T}\mathscr E_T^{\rm ext}F
 =\frac12\bigl(J_{\rm res}P_-\mathcal E_TF,\,
               \mathcal R_{\Gamma,T}F,\,0\bigr).     \tag{632}
\end{equation}
Every vector in \(\mathscr R_T^{\rm ext}\) has equal second and third Gamma
coordinates, whereas (632) does not unless
\(\mathcal R_{\Gamma,T}F=0\).  Thus the missing mathematics is exactly a
\(J_T^{\rm ext}\)-skew feedback \(A_T\) which restores this diagonal Gamma
constraint while simultaneously coupling the residue finite part.
\end{proposition}

\begin{proof}
Equation (632) follows from (624)--(626).  Equality of the two Gamma
coordinates characterizes the generated diagonal Gamma subspace.  Strict
Gamma observability (575) supplies sources for which the second coordinate
in (632) is nonzero, proving failure of invariance.
\end{proof}

\paragraph{New exact construction target --- \OPEN.}
Equations (624)--(632) construct every term of the desired conservation on
one common boundary and reduce \(X_T^{\rm src}\) to a single geometric
problem: find a bounded \(J_T^{\rm ext}\)-skew gauge making the physical
generated range invariant.  This is not a Sylvester inverse.  The
one-orbit reflections (618) give local pieces of the gauge, the Volterra
Hadamard operator (614)--(617) supplies their cross-orbit bulk, and the
Kato connection supplies the rank-two Tate correction (623).  What remains
is to assemble these three explicit pieces and prove (628)--(630) with a
uniform bound.

\begin{theorem}[Block tangency equation and no Gamma-only gauge ---
\PROVED]
Write a \(J_T^{\rm ext}\)-skew gauge in blocks relative to
\(\mathscr K_{\rm res}\oplus\mathscr K_\Gamma\oplus
\mathscr K_\Gamma\):
\[
 A=\begin{pmatrix}
 A_{rr}&A_{r+}&A_{r-}\\
 A_{+r}&A_{++}&A_{+-}\\
 A_{-r}&A_{-+}&A_{--}
 \end{pmatrix}.
\]
Condition (628) is equivalent to
\begin{align}
 &(J_{\rm res}A_{rr})^*=-J_{\rm res}A_{rr},\notag\\
 &A_{+r}^*=-J_{\rm res}A_{r+},\qquad
   A_{-r}^*= J_{\rm res}A_{r-},\notag\\
 &A_{++}^*=-A_{++},\qquad A_{--}^*=-A_{--},\qquad
   A_{-+}=A_{+-}^*.                                  \tag{633}
\end{align}
For \(e=\mathcal E_TF\) and \(r=\mathcal R_{\Gamma,T}F\), equality of the
two Gamma coordinates in the tangent vector
\((Y_{0,T}+A)(e,r,r)\) is exactly
\begin{equation}
 \boxed{
 (A_{+r}-A_{-r})e+
 \left(\frac12I+C_\Gamma\right)r=0,}                 \tag{634}
\end{equation}
where
\begin{equation}
 C_\Gamma:=A_{++}+A_{+-}-A_{+-}^*-A_{--},
 \qquad C_\Gamma^*=-C_\Gamma.                        \tag{635}
\end{equation}
In particular, no \(J_T^{\rm ext}\)-skew gauge with
\(A_{+r}=A_{-r}=0\) can satisfy tangency on any source with \(r\ne0\).
The residue-to-Gamma block is compulsory.
\end{theorem}

\begin{proof}
Expanding \(JA+A^*J=0\) gives (633).  The plus Gamma coordinate of
\((Y_0+A)(e,r,r)\) is
\[
 \tfrac12r+A_{+r}e+(A_{++}+A_{+-})r,
\]
and the minus coordinate is
\[
 A_{-r}e+(A_{-+}+A_{--})r.
\]
Subtract them, use \(A_{-+}=A_{+-}^*\), and obtain (634)--(635).
If the residue blocks vanish, (634) says
\((\frac12I+C_\Gamma)r=0\).  Taking the real part of its inner product
with \(r\), the skew term disappears and gives
\(\|r\|^2/2=0\), a contradiction for \(r\ne0\).
\end{proof}

\begin{corollary}[Exact residue-to-Gamma synthesis required for closure ---
\PROVED]
Every bounded gauge satisfying (628)--(630) must contain a bounded
residue-to-Gamma synthesis
\begin{equation}
 \mathscr S_T:=A_{+r}-A_{-r}:
 \mathscr K_{\rm res}\longrightarrow\mathscr K_{\Gamma,T}   \tag{636}
\end{equation}
such that, on the physical generated range,
\begin{equation}
 \boxed{
 \mathscr S_T\mathcal E_TF
 =-\left(\frac12I+C_\Gamma\right)
       \mathcal R_{\Gamma,T}F.}                      \tag{637}
\end{equation}
Since \(\frac12I+C_\Gamma\) is bounded below by \(1/2\),
(637) implies
\begin{equation}
 \|\mathcal R_{\Gamma,T}F\|
 \le2\|\mathscr S_T\|\,\|\mathcal E_TF\|.             \tag{638}
\end{equation}
Thus a source conservation law necessarily controls the strict Gamma
observation by the complete residue analysis.
\end{corollary}

\begin{proof}
Equations (636)--(637) are (634).  For skew-adjoint \(C_\Gamma\),
\[
 \|(\tfrac12I+C_\Gamma)r\|\,\|r\|
 \ge\Re\langle(\tfrac12I+C_\Gamma)r,r\rangle
 =\tfrac12\|r\|^2.
\]
Apply this to (637) and use boundedness of \(\mathscr S_T\) to obtain
(638).
\end{proof}

\paragraph{Synthesis audit --- \AUDIT.}
Equation (637), not a free normalization, is the concrete coupling that
the Hadamard/Volterra bulk must produce.  It is stronger than mere
summability (549): it requires a bounded synthesis from the full residue
range into the Gamma storage whose restriction equals the explicit
right-hand side.  On a first-loss vector the polarized radical relation
(523) supplies orthogonality in the residue Krein metric, while (637)
would transport it to the strictly positive Gamma metric.  Constructing
\(\mathscr S_T\) from (614)--(617) and verifying the remaining residue
coordinate tangency in (629) is now the exact final source calculation.

\begin{theorem}[Reverse residue sampling is the exact Gamma tangency
criterion --- \PROVED]
Fix a bounded skew-adjoint \(C_\Gamma\) and put
\[
 L_\Gamma:=\frac12I+C_\Gamma .
\]
There exists a bounded synthesis operator
\(\mathscr S_T:\mathscr K_{\rm res}\to\mathscr K_{\Gamma,T}\)
satisfying
\begin{equation}
 \mathscr S_T\mathcal E_T
 =-L_\Gamma\mathcal R_{\Gamma,T}                     \tag{639}
\end{equation}
on the common source domain if and only if there is \(C_T<\infty\) such
that
\begin{equation}
 \|L_\Gamma\mathcal R_{\Gamma,T}F\|^2
 \le C_T\|\mathcal E_TF\|_{\mathscr K_{\rm res}}^2.  \tag{640}
\end{equation}
When these conditions hold, the rule
\[
 \mathscr S_T(\mathcal E_TF)
 :=-L_\Gamma\mathcal R_{\Gamma,T}F
\]
is well defined and bounded on the physical residue range and extends by
zero from its closure to the full residue Hilbert space.

Every such synthesis necessarily implies the reverse logarithmic sampling
estimate
\begin{equation}
 \boxed{
 \langle G_{\Gamma,T}F,F\rangle
 \le16\|\mathscr S_T\|^2
       \|\mathcal E_TF\|_{\mathscr K_{\rm res}}^2.}  \tag{641}
\end{equation}
\end{theorem}

\begin{proof}
If \(\mathscr S_T\) exists, (640) follows with
\(C_T=\|\mathscr S_T\|^2\).  Conversely, (640) makes the displayed rule
independent of the representative: if \(\mathcal E_TF=0\), its right side
vanishes.  The same inequality proves boundedness on the range, hence
continuous extension to its closure; composing with the orthogonal
projection onto that closure gives an operator on all of
\(\mathscr K_{\rm res}\).

Since \(C_\Gamma^*=-C_\Gamma\),
\[
 \|L_\Gamma r\|\ge\frac12\|r\|.
\]
Equation (639) therefore gives
\[
 \|\mathcal R_{\Gamma,T}F\|
 \le2\|\mathscr S_T\|\,\|\mathcal E_TF\|.
\]
Square this inequality and use (575),
\(\|\mathcal R_{\Gamma,T}F\|^2\ge
\frac14\langle G_{\Gamma,T}F,F\rangle\), to obtain (641).
\end{proof}

\begin{corollary}[What remains beyond an abstract synthesis ---
\AUDIT]
Estimate (640) would construct the residue-to-Gamma block as a bounded
operator without using \(H_T^{-1}\).  It would not yet provide the demanded
source formula for the full gauge: one must additionally
\begin{enumerate}
\item realize \(\mathscr S_T\) by the Hadamard--Volterra transfer
      (601)--(617), rather than by extension from the range;
\item choose the remaining blocks so that (633) holds;
\item verify tangency of the residue coordinate as well as the Gamma
      equality (634).
\end{enumerate}
Thus (641) is the exact analytic estimate behind the construction, while
(614)--(617) give the candidate formula that must prove it.
\end{corollary}

\begin{theorem}[Unconditional reverse residue sampling on every fixed
window --- \PROVED]
For every fixed \(T>0\), there are constants \(a_T,b_T>0\) such that
\begin{align}
 \|F\|_2^2
 &\le a_T\|\mathcal E_TF\|_{\mathscr K_{\rm res}}^2,
                                                               \tag{642}\\
 \langle G_{\Gamma,T}F,F\rangle
 &\le b_T\|\mathcal E_TF\|_{\mathscr K_{\rm res}}^2             \tag{643}
\end{align}
for every \(F\in\mathcal Q_T^{\rm prim}\).  Consequently the physical
residue range is closed in the topology needed to control both the source
norm and Gamma energy.
\end{theorem}

\begin{proof}
Suppose (642) fails.  There are \(F_n\in\mathcal Q_T^{\rm prim}\) with
\(\|F_n\|_2=1\) and \(\|\mathcal E_TF_n\|\to0\).  The residue identity
(550) gives
\[
 |\langle H_TF_n,F_n\rangle|
 \le\|\mathcal E_TF_n\|^2\longrightarrow0.
\]
By (582)--(583),
\[
 \langle G_{\Gamma,T}F_n,F_n\rangle
 =\langle H_TF_n,F_n\rangle+
   \langle\mathsf W_TF_n,F_n\rangle.
\]
The operator \(\mathsf W_T\) is bounded, so the sequence is bounded in the
Gamma form norm.  The compact embedding of that form domain into
\(L^2(I_T)\), proved in the compact-resolvent argument for (251), gives a
subsequence converging strongly in \(L^2\) to some \(F\) with
\(\|F\|_2=1\), and weakly in the Gamma form domain.

The bounded map (548)--(549) sends this weak convergence to
\(\mathcal E_TF_n\rightharpoonup\mathcal E_TF\).  Its left side converges
strongly to zero, so \(\mathcal E_TF=0\).  Hence the Paley--Wiener
transform of \(F\) vanishes at every nontrivial zero.  The unconditional
distinct-zero density and Cartwright rigidity used in (292) force
\(F=0\), contradicting \(\|F\|_2=1\).  This proves (642).

For an arbitrary \(F\),
\[
 \langle G_{\Gamma,T}F,F\rangle
 =\langle H_TF,F\rangle+\langle\mathsf W_TF,F\rangle
 \le\|\mathcal E_TF\|^2+\|\mathsf W_T\|\,\|F\|_2^2.
\]
Apply (642) to obtain (643).
\end{proof}

\begin{corollary}[The mandatory residue-to-Gamma block exists boundedly ---
\PROVED]
For every fixed \(T\) and every bounded skew \(C_\Gamma\), condition (640)
holds.  Hence there is a bounded synthesis \(\mathscr S_T\) satisfying
(639), and in particular the Gamma-coordinate tangency equation (634) can
be solved on the physical residue range.
\end{corollary}

\begin{proof}
The finite-rank form \(\mathcal D_T\) is bounded in the Gamma norm by
(536), (541)--(543); hence
\[
 \|\mathcal R_{\Gamma,T}F\|^2
 =\langle(G_{\Gamma,T}-\mathcal D_T)F,F\rangle
 \le c_T\langle G_{\Gamma,T}F,F\rangle
\]
for some fixed-window \(c_T\).  Combine this with (643) and boundedness of
\(L_\Gamma=\frac12I+C_\Gamma\) to get (640).  Theorem (639)--(641) then
constructs \(\mathscr S_T\).
\end{proof}

\paragraph{What this closes and what it does not --- \AUDIT.}
The previously open norm estimate behind the Gamma equality in the extended
boundary is now proved: the required residue-to-Gamma block exists and is
bounded for every finite window.  This does not yet prove (631), because
one must choose the adjoint blocks mandated by (633) and verify that the
\emph{residue coordinate} of the corrected vector also belongs to the
physical generated range.  The remaining problem is therefore no longer
an estimate or a Gamma normalization; it is the simultaneous graph-tangency
identity for the finite-part residue coordinate.

\begin{theorem}[Closed residue-graph normal form of the full conservation
problem --- \PROVED]
Let
\[
 \mathscr R_T:=\mathcal E_T(\mathcal Q_T^{\rm prim})
 \subset\mathscr K_{\rm res}.
\]
Equations (549), (642) and (643) imply that \(\mathscr R_T\) is closed and
that
\begin{equation}
 \Gamma_T(\mathcal E_TF):=\mathcal R_{\Gamma,T}F     \tag{644}
\end{equation}
defines a bounded operator
\(\Gamma_T:\mathscr R_T\to\mathscr K_{\Gamma,T}\).
Consequently the generated extended range is exactly
\begin{equation}
 \mathscr R_T^{\rm ext}
 =\{(e,\Gamma_Te,\Gamma_Te):e\in\mathscr R_T\}.       \tag{645}
\end{equation}

A bounded source conservation operator exists if and only if there is a
bounded \(T_T:\mathscr R_T\to\mathscr R_T\), admitting a source realization
on \(\mathcal Q_T^{\rm prim}\), such that
\begin{align}
 &\langle J_{\rm res}T_Te,f\rangle+
  \langle J_{\rm res}e,T_Tf\rangle\notag\\
 &\qquad=
 \langle P_-e,f\rangle+
 \langle\Gamma_Te,\Gamma_Tf\rangle
 \qquad(e,f\in\mathscr R_T).                         \tag{646}
\end{align}
If \(T_T\) is source-realized by
\begin{equation}
 T_T\mathcal E_TF=\mathcal E_TX_T^{\rm src}F,        \tag{647}
\end{equation}
then (646) pulls back exactly to (553).
\end{theorem}

\begin{proof}
The upper estimate (549) and reverse estimates (642)--(643) make
\(\mathcal E_T\) a topological isomorphism of the common form domain onto
its range, so the latter is closed.  Equation (575), together with (643),
makes (644) bounded.  Formula (645) is then (624).

For a source operator \(X\), put
\(T_T(\mathcal E_TF)=\mathcal E_TXF\).  Pulling (646) back by
\(\mathcal E_T\) and using (550), (577)--(578) gives
\[
 H_TX+X^*H_T
 =\mathcal N_T^*\mathcal N_T+
  \mathcal R_{\Gamma,T}^*\mathcal R_{\Gamma,T}=B_T.
\]
The converse is the same polarized calculation in reverse.
\end{proof}

\begin{corollary}[A source lift of the residue involution gives an explicit
solution --- \PROVED]
Assume there is a bounded source operator
\(\mathcal C_T:\mathcal Q_T^{\rm prim}\to
\mathcal Q_T^{\rm prim}\) such that
\begin{equation}
 \mathcal E_T\mathcal C_TF
 =J_{\rm res}\mathcal E_TF.                          \tag{648}
\end{equation}
Then \(J_{\rm res}\mathscr R_T=\mathscr R_T\).  On this closed range put
\[
 \mathcal P_T^{\rm obs,res}:=
 P_{\mathscr R_T}
 (P_-+\Gamma_T^*\Gamma_T)\big|_{\mathscr R_T}
\]
and
\begin{equation}
 T_T:=\frac12
 (J_{\rm res}|_{\mathscr R_T})
 \mathcal P_T^{\rm obs,res}.                         \tag{649}
\end{equation}
This \(T_T\) satisfies (646).  If the two factors in (649) are realized by
the co-Poisson/Hardy source transform and its explicit adjoint analysis,
their composition gives \(X_T^{\rm src}\) without using \(H_T^{-1}\).
\end{corollary}

\begin{proof}
Equation (648) gives invariance; applying it twice gives equality of the
range because \(J_{\rm res}^2=I\).  Hence the restricted \(J_{\rm res}\)
is a selfadjoint unitary.  The compressed observation operator in (649) is
selfadjoint on \(\mathscr R_T\).  Therefore
\[
 J_{\rm res}T_T+T_T^*J_{\rm res}
 =\mathcal P_T^{\rm obs,res},
\]
whose polarized form is (646).  The last statement is the source
realization (647).
\end{proof}

\begin{theorem}[The full same-window residue involution lift is equivalent
to RH --- \PROVED]
There exists a bounded complex-linear operator
\(\mathcal C_T:\mathcal Q_T^{\rm prim}\to
\mathcal Q_T^{\rm prim}\) satisfying (648) for one nontrivial window \(T\)
if and only if every nontrivial zero of \(\zeta\) lies on the critical
line.  Under that condition the only possible lift is
\(\mathcal C_T=I\).
\end{theorem}

\begin{proof}
On every critical-line fiber, \(J_{\rm res}=I\).  Hence (648) implies
\[
 \Phi_{\mathcal C_TF-F}(\gamma)=0
\]
at every distinct critical-line zero.  The unconditional family of
distinct critical zeros used in (292) has counting function
\(\gg R\log R\), while a nonzero Paley--Wiener transform of type \(T\)
has only \(O_T(R)\) zeros.  Therefore
\(\mathcal C_TF=F\) for every source \(F\).

If an off-line orbit exists, (648) would now require
\[
 \bigl(\Phi_F(z_\rho),\Phi_F(\bar z_\rho)\bigr)
 =\bigl(\Phi_F(\bar z_\rho),\Phi_F(z_\rho)\bigr)
\]
for every primitive source.  Finite-orbit interpolation (520)--(522)
produces sources with arbitrary unequal coordinates, a contradiction.
Thus every zero is critical.  Conversely, if all zeros are critical, every
residue fiber is one-dimensional with \(J_{\rm res}=I\), and
\(\mathcal C_T=I\) satisfies (648).
\end{proof}

\paragraph{Final graph target --- \OPEN.}
All analytic norms needed for the residue/Gamma graph are now proved.
The full involution lift (648) may not be used as an intermediate lemma:
the preceding theorem proves that it is already equivalent to RH.  The
genuinely weaker remaining source statement is the mixed tangency equation
(646), where \(P_-\) and the positive Gamma graph occur together.
Meromorphic functional-equation symmetry alone supplies only \(J_{\rm res}\)
and therefore overshoots the required target.  A valid co-Poisson/Hardy
construction must retain the Gamma finite part and prove (646) directly;
this is the incoming-history/reciprocal-gap gate with all norm estimates
now supplied.

\begin{theorem}[Fredholm graph form and exact no-first-contact equivalence
--- \PROVED]
Let \(Q_T\) be the orthogonal projection of
\(\mathscr K_{\rm res}\) onto the closed range \(\mathscr R_T\), and define
on \(\mathscr R_T\)
\begin{align}
 \mathbb H_T&:=Q_TJ_{\rm res}\big|_{\mathscr R_T},\notag\\
 \mathbb B_T&:=Q_T(P_-+\Gamma_T^*\Gamma_T)
                    \big|_{\mathscr R_T}.            \tag{650}
\end{align}
Then \(\mathbb H_T\) is bounded selfadjoint Fredholm,
\(\mathbb B_T\) is bounded selfadjoint, and for some \(c_T>0\),
\begin{equation}
 \mathbb B_T\ge c_TI_{\mathscr R_T}.                 \tag{651}
\end{equation}
Moreover the following are equivalent:
\begin{enumerate}
\item there is a bounded \(T_T:\mathscr R_T\to\mathscr R_T\)
      satisfying (646);
\item \(\ker\mathbb H_T=\{0\}\);
\item \(H_T\) has no zero mode.
\end{enumerate}
When these conditions hold, the abstract solution
\begin{equation}
 T_T=\frac12\mathbb H_T^{-1}\mathbb B_T              \tag{652}
\end{equation}
exists.  Formula (652) is recorded only for equivalence auditing and is
forbidden as the requested source construction.
\end{theorem}

\begin{proof}
For \(e=\mathcal E_TF\) and \(f=\mathcal E_TG\),
\[
 \langle\mathbb H_Te,f\rangle
 =\langle J_{\rm res}\mathcal E_TF,\mathcal E_TG\rangle
 =\mathfrak h_T(F,G),
\]
while the corresponding formula for \(\mathbb B_T\) is the right side of
(646).  The equivalence of the residue norm with the Gamma form norm,
from (549) and (642)--(643), transports the normalized representation
\(H_T=\mathsf G_T^{1/2}(I-K_T)\mathsf G_T^{1/2}\) in (587) to
\(\mathbb H_T\).  Since \(K_T\) is compact, \(\mathbb H_T\) is Fredholm.

The first Gamma cell has a positive spectral gap on the fixed compact
window.  Combining (575) with the equivalence of the residue and Gamma
form norms gives
\[
 \langle\mathbb B_Te,e\rangle
 =\|\mathcal N_TF\|^2+\|\mathcal R_{\Gamma,T}F\|^2
 \ge c_T\|e\|^2,
\]
which proves (651).

If (646) holds and \(0\ne e\in\ker\mathbb H_T\), evaluation at
\((e,e)\) gives \(0=\langle\mathbb B_Te,e\rangle\), contradicting
(651).  Hence (1) implies (2).  A selfadjoint Fredholm operator with zero
kernel is boundedly invertible, and (652) gives
\[
 \mathbb H_TT_T+T_T^*\mathbb H_T
 =\tfrac12\mathbb B_T+\tfrac12\mathbb B_T
 =\mathbb B_T.
\]
Thus (2) implies (1).  Finally the first displayed polarized identity and
injectivity of \(\mathcal E_T\) identify
\(\ker\mathbb H_T\) with \(\ker H_T\), proving (2)\(\Leftrightarrow\)(3).
\end{proof}

\paragraph{Meaning of the Fredholm audit --- \AUDIT.}
All existence, boundedness and domain questions surrounding (553) have now
been reduced to a Fredholm equation with a strictly positive right side.
The inverse in (652) proves no new mathematics: it exists precisely after
the first-contact kernel has been excluded.  The only noncircular route is
still to construct \(T_T\) by the joint Hadamard--Volterra,
Gamma-storage and Kato/co-Poisson source formulas and verify (646)
directly.  Success would leave no further analytic or operator-theoretic
gate.

\begin{theorem}[Horizontal opening of the zero orbits and the exact
antisymmetric current --- \PROVED]
For an off-line orbit represented by
\(\rho=\frac12+\alpha+i\gamma\), \(\alpha>0\), and
\(0\le u\le1\), put
\begin{equation}
 \mathcal E_{u,\rho}F:=
 \binom{\Phi_F(\gamma-iu\alpha)}
       {\Phi_F(\gamma+iu\alpha)},\qquad
 D_\rho:=\begin{pmatrix}1&0\\0&-1\end{pmatrix},
 \qquad J_\rho:=\begin{pmatrix}0&1\\1&0\end{pmatrix}. \tag{653}
\end{equation}
Let \(M_tF(t)=tF(t)\).  On every compact window,
\begin{equation}
 \frac d{du}\mathcal E_{u,\rho}F
 =\alpha D_\rho\mathcal E_{u,\rho}(M_tF).            \tag{654}
\end{equation}
If \(P_{-,\rho}=(I-J_\rho)/2\), then the negative residue energy has the
exact integrated-current representation
\begin{align}
 \|\mathcal N_TF\|^2
 &=\sum_{\rho}m_\rho
   \|P_{-,\rho}\mathcal E_{1,\rho}F\|^2\notag\\
 &=2\Re\int_0^1\sum_{\rho}m_\rho\alpha
 \left\langle
 P_{-,\rho}D_\rho\mathcal E_{u,\rho}(M_tF),
 P_{-,\rho}\mathcal E_{u,\rho}F
 \right\rangle du .                                \tag{655}
\end{align}
The sum and the integral in the second line are understood first at finite
height and then as the limit of the nonnegative endpoint forms in the first
line.  Thus (655) introduces no coefficientwise rearrangement of a
conditionally convergent zero sum.
\end{theorem}

\begin{proof}
Since
\(\partial_z\Phi_F(z)=i\Phi_{M_tF}(z)\), differentiation of the two
coordinates in (653) gives respectively
\(+\alpha\Phi_{M_tF}(\gamma-iu\alpha)\) and
\(-\alpha\Phi_{M_tF}(\gamma+iu\alpha)\), proving (654).
At \(u=0\) the two coordinates coincide, hence
\(P_{-,\rho}\mathcal E_{0,\rho}F=0\).  The fundamental theorem of
calculus applied at finite height to
\(\|P_{-,\rho}\mathcal E_{u,\rho}F\|^2\), followed by (654), gives
the second line of (655).  At \(u=1\), diagonalizing the exchange matrix
shows
\[
 m_\rho\|P_{-,\rho}\mathcal E_{1,\rho}F\|^2
 =\frac{m_\rho}{2}
 |\Phi_F(\gamma-i\alpha)-\Phi_F(\gamma+i\alpha)|^2,
\]
which is (510).  Monotone convergence of the endpoint sums, or equivalently
the bounded sampling theorem (549), proves the asserted limiting
interpretation.
\end{proof}

\begin{proposition}[The precise source-lift gate for horizontal opening ---
\PROVED]
Let \(\mathcal E_{T,u}\) be the direct sum of (653), together with the
unchanged critical-line fibers, and set
\[
 \mathscr A\{\mathcal E_{T,u}F\}_\rho
 :=\{\alpha D_\rho\mathcal E_{u,\rho}(M_tF)\}_\rho,
\]
with zero component on a critical-line fiber.  A source operator
\(C_{T,u}\) realizes the horizontal derivative if and only if
\begin{equation}
 \boxed{\mathcal E_{T,u}C_{T,u}F
       =\mathscr A\mathcal E_{T,u}F}                 \tag{656}
\end{equation}
on the common source domain.  In particular, (654) does not by itself
produce a source conservation law: it produces an ambient residue-space
velocity, and it descends to sources exactly when that velocity is tangent
to the physical sampling range.
\end{proposition}

\begin{proof}
Equation (656) is precisely the assertion that the displayed ambient
velocity belongs to the range of \(\mathcal E_{T,u}\) and has the stated
source preimage.  Necessity is immediate.  Conversely, (656) substituted
in (654) realizes the derivative by \(C_{T,u}\).  No inverse or
pseudoinverse is used in either direction.
\end{proof}

\paragraph{Horizontal-deformation verdict --- \AUDIT.}
Equations (653)--(655) give the sought term-by-term continuous object behind
the negative residue square: multiplication by \(t\) opens every symmetric
zero orbit and its graded antisymmetric current integrates exactly to
\(\mathcal N_T^*\mathcal N_T\).  The grading \(D_\rho\), however, changes
sign between the two branches and depends on the orbit width \(\alpha\).
It is therefore not the sampling of one already displayed local source
operator.  Equation (656), rather than (654), is the required Hadamard
feedback statement.  Proving its bounded joint version with the Gamma
finite part and the Tate terminal would give the missing tangency in
(646); asserting that the ambient velocity is automatically in the
sampling range would simply assume that tangency.  Hence the horizontal
opening is a genuine new conservation current, but does not yet close G.

\begin{theorem}[Local boundedness and closed range of the horizontal
opening at the physical boundary --- \PROVED]
Regard \(\mathcal E_{T,u}\) as an operator from
\(\mathcal Q_T^{\rm prim}\), equipped with its Gamma form norm, to the
fixed residue Hilbert space.  Then
\begin{equation}
 \|\mathcal E_{T,u}F-\mathcal E_{T,v}F\|_{\rm res}
 \le C_T|u-v|\,\|F\|_{\mathcal Q_{\Gamma,T}}
 \qquad(0\le u,v\le1).                              \tag{657}
\end{equation}
There is \(\delta_T>0\) such that, for
\(1-\delta_T<u\le1\), every \(\mathcal E_{T,u}\) is bounded below by a
constant independent of \(u\) in that interval.  Hence its range is closed,
and the ambient velocity in (654) is a bounded residue-space observation on
the common source domain.  In particular, the final obstruction in (656)
near the physical boundary is range tangency, not convergence or loss of a
norm estimate.
\end{theorem}

\begin{proof}
The points \(\gamma\mp iu\alpha\) remain in the fixed strip
\(|\Im z|\le1/2\), and \(0\le\alpha\le1/2\).  Apply the weighted
Plancherel--P\'olya estimate used in (551) to the difference obtained by
integrating (654) between \(u\) and \(v\).  Multiplication by \(t\) is a
bounded multiplier on the compact-window logarithmic Gamma form domain:
choose a compactly supported smooth function equal to \(t\) on \(I_T\) and
use the standard order-zero multiplier estimate for the Fourier weight
\(1+\log(2+|\tau|)\).  This proves (657).

Equations (642)--(643) give a constant \(c_T>0\) such that
\begin{equation}
 \|\mathcal E_{T,1}F\|_{\rm res}
 \ge c_T\|F\|_{\mathcal Q_{\Gamma,T}}.              \tag{658}
\end{equation}
Choose \(\delta_T\) with \(C_T\delta_T<c_T/2\).  The reverse triangle
inequality and (657)--(658) then give
\[
 \|\mathcal E_{T,u}F\|_{\rm res}
 \ge\frac{c_T}{2}\|F\|_{\mathcal Q_{\Gamma,T}}
 \qquad(1-\delta_T<u\le1).
\]
A bounded-below operator has closed range.  The same sampling estimate
applied directly to the right side of (654) proves boundedness of the
ambient velocity.
\end{proof}

\begin{theorem}[Affine elimination of the negative residue port ---
\PROVED]
On the closed physical residue range \(\mathscr R_T\), let
\(\mathbb H_T\) be (650).  Then
\begin{equation}
 Q_TP_-\big|_{\mathscr R_T}
 =\frac12(I_{\mathscr R_T}-\mathbb H_T).             \tag{659}
\end{equation}
Consequently (646) is equivalent, after writing
\begin{equation}
 T_T=-\frac14I_{\mathscr R_T}+Y_T,                  \tag{660}
\end{equation}
to the single metric-connection equation
\begin{equation}
 \boxed{
 \mathbb H_TY_T+Y_T^*\mathbb H_T
 =\mathbb S_T,qquad
 \mathbb S_T:=\frac12I_{\mathscr R_T}
                 +\Gamma_T^*\Gamma_T
 \ge\frac12I_{\mathscr R_T}.}                      \tag{661}
\end{equation}
Thus the residue square contributes no further unidentified correction:
after the universal affine term \(-I/4\), the entire remaining datum is
the strictly positive Gamma-graph metric \(\mathbb S_T\).
\end{theorem}

\begin{proof}
On \(\mathscr K_{\rm res}\),
\(P_-=(I-J_{\rm res})/2\).  Since \(Q_Te=e\) for
\(e\in\mathscr R_T\), compression gives
\[
 Q_TP_-e=\tfrac12(e-Q_TJ_{\rm res}e)
          =\tfrac12(I-\mathbb H_T)e,
\]
which proves (659).  Hence
\(\mathbb B_T=(I-\mathbb H_T)/2+\Gamma_T^*\Gamma_T\).
Substitute (660) in
\(\mathbb H_TT_T+T_T^*\mathbb H_T=\mathbb B_T\);
the two affine terms contribute \(-\mathbb H_T/2\), and cancellation gives
(661).  Positivity is immediate.
\end{proof}

\begin{corollary}[Sharp connection blow-up at first contact --- \PROVED]
If \(\mathbb H_Te_T=\lambda_Te_T\), \(\|e_T\|=1\), and a bounded
\(Y_T\) satisfies (661), then
\begin{equation}
 2\lambda_T\Re\langle Y_Te_T,e_T\rangle
 =\frac12+\|\Gamma_Te_T\|^2,                       \tag{662}
\end{equation}
and therefore
\[
 \|Y_T\|\ge\frac{1}{4|\lambda_T|}
 \quad(\lambda_T\ne0).
\]
In particular, no family of bounded source connections satisfying (661)
can extend uniformly to a first-contact parameter.
\end{corollary}

\begin{proof}
Evaluate (661) on \((e_T,e_T)\).  This gives (662); taking absolute values
and discarding the nonnegative Gamma term gives the stated bound.
\end{proof}

\paragraph{Final normalization audit --- \AUDIT.}
Equation (661) is the minimal terminal identity after all residue,
Gamma and Tate normalizations have been performed.  Its right side has a
fixed gap \(1/2\), so a zero mode of \(\mathbb H_T\) is algebraically
incompatible with any bounded solution.  Conversely the Fredholm theorem
(650)--(652) shows that, once the zero mode is excluded, a bounded abstract
solution exists.  Therefore there is no smaller missing scalar equality:
the required new theorem is precisely a source-derived, uniformly bounded
metric connection for (661).  Formula (662) also proves that a construction
which is only defined on pre-contact windows and whose norm is allowed to
depend on \(1/\lambda_T\) cannot close G.

\begin{theorem}[Universal dissipative generator and exact normal-matching
gate --- \PROVED]
Write \(\mathscr K_{\rm res}=\mathscr R_T\oplus\mathscr R_T^\perp\), let
\(Q=Q_T\), \(J=J_{\rm res}\), \(H=\mathbb H_T=QJ|_{\mathscr R_T}\), and
put
\[
 L:=(I-Q)J|_{\mathscr R_T},\qquad
 S:=\frac12I_{\mathscr R_T}+\Gamma_T^*\Gamma_T.
\]
Extend the positive observation
\(\mathcal Oe=(2^{-1/2}e,\Gamma_Te)\) to the ambient residue space by
\(\widehat{\mathcal O}:=\mathcal OQ\).  Then the bounded operator
\begin{equation}
 K_0:=-\frac12J\widehat{\mathcal O}^{\,*}
                   \widehat{\mathcal O}
     =-\frac12JQSQ                                \tag{708}
\end{equation}
satisfies the exact ambient dissipation identity
\begin{equation}
 \boxed{K_0^*J+JK_0=-QSQ.}                         \tag{709}
\end{equation}

Let \(A\) be any bounded \(J\)-skew operator,
\begin{equation}
 A^*J+JA=0.                                        \tag{710}
\end{equation}
Then \(K:=K_0+A\) leaves \(\mathscr R_T\) invariant if and only if the
single normal-block identity
\begin{equation}
 \boxed{(I-Q)AQ=\frac12LS}                         \tag{711}
\end{equation}
holds.  Whenever (711) holds, the source-defined operator
\begin{equation}
 Y_T:=-K|_{\mathscr R_T}
     =\frac12HS-QA|_{\mathscr R_T}                 \tag{712}
\end{equation}
is bounded and satisfies exactly
\begin{equation}
 HY_T+Y_T^*H=S,                                    \tag{713}
\end{equation}
which is (661).  Conversely, every invariant dissipative realization of
(661) whose dissipation is (709) has a \(J\)-skew correction satisfying
(711).
\end{theorem}

\begin{proof}
Since \(J=J^*=J^{-1}\), \(Q=Q^*=Q^2\), and \(S=S^*\),
\[
 K_0^*J=-\tfrac12QSQ,\qquad
 JK_0=-\tfrac12QSQ,
\]
which proves (709).  For \(e\in\mathscr R_T\),
\[
 (I-Q)K_0e=-\tfrac12(I-Q)JSe=-\tfrac12LSe.
\]
Thus \((I-Q)(K_0+A)e=0\) for every \(e\) precisely when (711) holds.

If the range is invariant, \(Ke=-Y_Te\).  Compress (709), using (710), to
\(\mathscr R_T\):
\[
 -Y_T^*H-HY_T=-S.
\]
This is (713).  Formula (712) is the tangential block of
\(-K=-K_0-A\).  Conversely, subtracting \(K_0\) from an ambient generator
with the same dissipation gives a \(J\)-skew \(A\), and invariance forces
its normal block to be (711).
\end{proof}

\begin{corollary}[Why the compressed prime--Gamma phase cannot be the
missing correction --- \PROVED]
Let \(\mathcal W:\mathcal H\to\mathscr R_T\) be any unitary
parametrization of the physical range and let \(C^*=-C\) on
\(\mathcal H\).  The compressed phase
\begin{equation}
 A_{\rm comp}:=\mathcal WC\mathcal W^*Q             \tag{714}
\end{equation}
has
\[
 (I-Q)A_{\rm comp}Q=0.
\]
Therefore it cannot satisfy (711), since \(S\ge\frac12I\) and
\begin{equation}
 LS=0\quad\Longrightarrow\quad L=0.                \tag{715}
\end{equation}
At a hypothetical zero mode \(0\ne e\in\ker H\), (706) gives
\(\|Le\|=\|e\|\), so the required normal correction in (711) is
nonzero.  Thus a formula of the form
\(\mathcal W(K_P+K_\Gamma)\mathcal W^*\), with the phase already compressed
to the physical range, cannot close (661).
\end{corollary}

\begin{proof}
The range of (714) is contained in \(\mathscr R_T\), proving the zero
normal block.  Since \(S\) is boundedly invertible, \(LS=0\) implies
\(L=0\).  Finally (706) gives
\(\|Le\|^2=\|e\|^2-\|He\|^2=\|e\|^2\) on \(\ker H\).
\end{proof}

\paragraph{New construction target --- \OPEN.}
Equations (708)--(713) separate the terminal problem into two independent
pieces.  The positive half-identity and Gamma square are already generated,
without \(H^{-1}\), by the universal dissipative block \(K_0\).  The sole
arithmetic task is now to construct an \emph{ambient}, rather than
pre-compressed, Hadamard--Gamma--prime current \(A_T^{\rm amb}\) which is
\(J_{\rm res}\)-skew and whose normal block obeys
\begin{equation}
 \boxed{
 (I-Q_T)A_T^{\rm amb}Q_T
 =\frac12(I-Q_T)J_{\rm res}Q_T
   \left(\frac12I+\Gamma_T^*\Gamma_T\right).}       \tag{716}
\end{equation}
This is an exact term-by-term test for a proposed new current.  Defining its
normal block to be the right side would be tautological; it must be derived
from the uncompressed Hadamard finite-part and prime/Gamma boundary cells.
If (716) is proved, (712)--(713) close G with no further estimate.

\begin{proposition}[Tangential compatibility forced by the normal block ---
\PROVED]
Write the \(J\)-skew correction in blocks relative to
\(\mathscr R_T\oplus\mathscr R_T^\perp\),
\[
 A=\begin{pmatrix}a&b\\c&d\end{pmatrix}.
\]
If its normal block is the required one,
\(c=\frac12LS\), then the \((1,1)\)-block of
\(A^*J+JA=0\) is exactly
\begin{equation}
 \boxed{
 a^*H+Ha
 =-\frac12\left[
 S(I-H^2)+(I-H^2)S\right].}                         \tag{717}
\end{equation}
The remaining block equations impose additional extension conditions on
\(b,d\); equation (717) is the compulsory tangential condition which must
hold before any such extension is possible.

In particular, if \(0\ne e\in\ker H\), evaluation of (717) on \(e\)
would give
\begin{equation}
 0=-\langle Se,e\rangle
   =-\frac12\|e\|^2-\|\Gamma_Te\|^2<0.              \tag{718}
\end{equation}
Hence an explicit arithmetic current satisfying (710) and (716)
automatically rules out first contact before \(Y_T\) is even formed.
\end{proposition}

\begin{proof}
The involution \(J\) has block form
\[
 J=\begin{pmatrix}H&L^*\\L&D\end{pmatrix}.
\]
The upper-left block of \(A^*J+JA\) is
\[
 a^*H+c^*L+Ha+L^*c.
\]
By \(J^2=I\), \(L^*L=I-H^2\).  Substituting
\(c=\frac12LS\) therefore gives (717).  Finally, if \(He=0\), the
left side of (717) has zero diagonal matrix coefficient and its right side
equals \(-\langle Se,e\rangle\), proving (718).
\end{proof}

\paragraph{Meaning of the new gate --- \AUDIT.}
Formula (716) is not merely another spelling of the full Lyapunov equation:
it tells a cell construction exactly which \emph{normal} boundary current
must be produced.  Formula (717), however, proves that no abstract
completion of that normal block can be declared without checking
\(J\)-skewness: its tangential compatibility already contains the
first-contact contradiction.  The new mathematical opportunity is to
derive both (710) and (716) from the uncompressed Gamma cells (675)--(678),
the prime translations (663), and the Hadamard residue graph (601)--(617).
Their compressed sum cannot be used, by (714)--(715).

\begin{theorem}[Corrected ambient phase ansatz and Ward-range identity ---
\PROVED]
Let \(C_T^{\rm amb}\) be a bounded Hilbert-skew operator on
\(\mathscr K_{\rm res}\),
\begin{equation}
 (C_T^{\rm amb})^*=-C_T^{\rm amb},                 \tag{719}
\end{equation}
and define
\begin{equation}
 A_T^{\rm amb}:=J_{\rm res}C_T^{\rm amb}.           \tag{720}
\end{equation}
Then \(A_T^{\rm amb}\) is automatically \(J_{\rm res}\)-skew.  Moreover
its normal-block condition (716) is equivalent to the single Ward-range
identity
\begin{equation}
 \boxed{
 C_T^{\rm amb}e-\frac12Se\in
 J_{\rm res}\mathscr R_T
 \qquad(e\in\mathscr R_T),}                         \tag{721}
\end{equation}
where \(S=\frac12I+\Gamma_T^*\Gamma_T\).
Thus the corrected prime--Gamma proposal is not
\(\mathcal WC\mathcal W^*\), but an ambient skew phase \(C_T^{\rm amb}\)
whose failure to equal \(S/2\) is a reflected physical vector.
\end{theorem}

\begin{proof}
Using \(J_{\rm res}^*=J_{\rm res}^{-1}=J_{\rm res}\) and (719),
\[
 (A_T^{\rm amb})^*J_{\rm res}
 +J_{\rm res}A_T^{\rm amb}
 =C_T^*+C_T=0.
\]
Substitute (720) in (716) and cancel the leftmost involution:
\[
 (I-Q)J_{\rm res}
 \left(C_T^{\rm amb}e-\tfrac12Se\right)=0.
\]
For any ambient vector \(x\),
\[
 (I-Q)J_{\rm res}x=0
 \quad\Longleftrightarrow\quad
 J_{\rm res}x\in\mathscr R_T
 \quad\Longleftrightarrow\quad
 x\in J_{\rm res}\mathscr R_T.
\]
This proves (721).
\end{proof}

\begin{corollary}[Scalar weak form of the remaining cell calculation ---
\PROVED]
Identity (721) holds if and only if for every
\(e\in\mathscr R_T\) and every \(z\in\mathscr K_{\rm res}\)
\begin{equation}
 \boxed{
 \left\langle C_T^{\rm amb}e-\frac12Se,\,
 (I-J_{\rm res}QJ_{\rm res})z\right\rangle=0.}      \tag{722}
\end{equation}
It is enough to take \(e\) in the source-generated dense core, but the
second variable must range over an ambient core whose projection is dense
in \((J_{\rm res}\mathscr R_T)^\perp\).  Physical source vectors alone do
not provide that test.  This is the exact polarized identity which the
Hadamard divided differences, the prime translations and the Gamma cells
must satisfy term by term.
\end{corollary}

\begin{proof}
The orthogonal projection onto
\((J_{\rm res}\mathscr R_T)^\perp\) is
\(I-J_{\rm res}QJ_{\rm res}\).  Therefore (721) is equivalent to
orthogonality of its left side to that complement, which is (722).
Density and boundedness in the first variable give the final statement;
the ambient quantifier in the second variable cannot be discarded.
\end{proof}

\begin{proposition}[Block form of the Ward gate and disappearance of the
external cell at first contact --- \PROVED]
Write
\[
 J_{\rm res}=\begin{pmatrix}H&L^*\\L&D\end{pmatrix},
 \qquad
 C_T^{\rm amb}=\begin{pmatrix}c&-b^*\\b&d\end{pmatrix}
\]
relative to
\(\mathscr K_{\rm res}=\mathscr R_T\oplus\mathscr R_T^\perp\).
Here \(c^*=-c\), \(d^*=-d\), and \(b:\mathscr R_T\to
\mathscr R_T^\perp\).  Then (716), equivalently (721), is exactly
\begin{equation}
 \boxed{Lc+Db=\frac12LS.}                           \tag{723}
\end{equation}
The involution relations also give
\begin{equation}
 L^*L=I-H^2,\qquad DL=-LH.                          \tag{724}
\end{equation}
Consequently, if \(0\ne e\in\ker H\), the external block \(Db\) has
zero real pairing with \(Le\), whereas (723) would require
\[
 0=\operatorname{Re}\langle ce,e\rangle
  =\frac12\langle Se,e\rangle>0.
\]
Thus the exterior Hadamard cell cannot manufacture the positive diagonal
at first contact: proving the Ward identity forces the contradiction
through the physical skew block \(c\).
\end{proposition}

\begin{proof}
The lower-left block of \(J_{\rm res}C_T^{\rm amb}\) is \(Lc+Db\), so
(716) is precisely (723).  The lower-right and upper-left blocks of
\(J_{\rm res}^2=I\) give (724).  If \(He=0\), then
\(L^*Le=e\) and \(DLe=-LHe=0\).  Pair (723) with \(Le\).  Selfadjointness
of \(D\) makes the real part of the external term vanish:
\[
 \operatorname{Re}\langle Dbe,Le\rangle
 =\operatorname{Re}\langle be,DLe\rangle=0.
\]
Moreover
\[
 \operatorname{Re}\langle Lce,Le\rangle
 =\operatorname{Re}\langle ce,L^*Le\rangle
 =\operatorname{Re}\langle ce,e\rangle=0
\]
because \(c\) is skew, while
\[
 \frac12\operatorname{Re}\langle LSe,Le\rangle
 =\frac12\operatorname{Re}\langle Se,L^*Le\rangle
 =\frac12\langle Se,e\rangle>0.
\]
This contradiction proves the last assertion.
\end{proof}

\begin{theorem}[No hidden shortcut: ambient Ward current is equivalent to
the terminal metric connection --- \PROVED]
There exists a bounded Hilbert-skew operator \(C_T^{\rm amb}\) satisfying
(721) if and only if there exists a bounded operator \(Y_T\) satisfying
\begin{equation}
 HY_T+Y_T^*H=S.                                    \tag{725}
\end{equation}
More precisely, a Ward current gives \(Y_T\) by (712).  Conversely, from
any solution of (725) define on \(\mathscr R_T\)
\begin{equation}
 Ve:=\frac12Se-JY_Te,                               \tag{726}
\end{equation}
write \(c=QV\), \(b=(I-Q)V\), and set
\begin{equation}
 C_T^{\rm amb}:=
 \begin{pmatrix}c&-b^*\\b&0\end{pmatrix}.          \tag{727}
\end{equation}
Then \(C_T^{\rm amb}\) is bounded and Hilbert-skew and satisfies (721).
Thus the ambient construction is a correctly typed realization of the
last gap, but not a theorem weaker than that gap.
\end{theorem}

\begin{proof}
The forward implication is (719)--(721) followed by (708)--(713).  For the
converse, (725) gives
\[
 c+c^*=QV+V^*Q=S-HY_T-Y_T^*H=0.
\]
Hence the block operator (727) is skew.  On \(\mathscr R_T\) it agrees
with \(V\), and (726) gives
\[
 C_T^{\rm amb}e-\tfrac12Se=-JY_Te\in J\mathscr R_T.
\]
This is (721).
\end{proof}

\begin{proposition}[Finite residue truncations are blind to the Ward
normalization --- \PROVED]
For every finite union \(R\) of residue orbits, the interpolation theorem
(606)--(608) makes the truncated evaluation
\(\mathcal E_{T,R}\) surjective onto \(\mathscr K_R\).  Hence its physical
range projection is \(Q_R=I_{\mathscr K_R}\), and therefore
\begin{equation}
 L_R:=(I-Q_R)J_R|_{\mathscr K_R}=0.                 \tag{728}
\end{equation}
The finite version of the Ward block equation (723) consequently reduces
to \(0=0\), independently of the prime and Gamma cells.  Thus checking
(723) on every finite set of zeros does not imply the infinite-dimensional
identity.  Any finite-cell construction must additionally prove a uniform
bound and convergence in the full residue Hilbert space which retains the
nontrivial limiting projection \(Q_T\).
\end{proposition}

\begin{proof}
Equation (608) is a right inverse of \(\mathcal E_{T,R}\), so its range is
all of \(\mathscr K_R\).  This gives \(Q_R=I\) and (728).  Substitution in
(723) annihilates both sides.  Strong convergence of finite coordinate
projections alone does not imply convergence of the corresponding range
projections, because every finite range is the whole truncated space.  A
uniform lifting or graph estimate is therefore indispensable.
\end{proof}

\begin{proposition}[Two-dimensional countermodel to an algebraic closure
of the Ward gate --- \PROVED]
The involution, closed-range and positive-observation axioms used in
(708)--(727) do not by themselves imply (721).  Indeed, take
\(\mathscr K=\mathbb C^2\),
\[
 J=\begin{pmatrix}1&0\\0&-1\end{pmatrix},\qquad
 \mathscr R=\operatorname{span}\{2^{-1/2}(1,1)\},
 \qquad \Gamma=0.
\]
Then \(\mathscr R\) is closed, \(J\) is a selfadjoint involution,
\(S=\frac12I_{\mathscr R}>0\), but
\begin{equation}
 H=QJ|_{\mathscr R}=0.                              \tag{729}
\end{equation}
Consequently neither (725) nor an ambient skew current satisfying (721)
exists.  Thus a proof of the terminal identity must use a genuinely
arithmetic global statement not contained in the abstract Krein geometry.
\end{proposition}

\begin{proof}
For \(e=2^{-1/2}(1,1)\), one has
\(Je=2^{-1/2}(1,-1)\perp\mathscr R\), proving (729).  Hence
\(HY+Y^*H=0\) for every operator \(Y\) on the one-dimensional space
\(\mathscr R\), whereas \(S=I/2\).  The equivalence (725)--(727) then also
excludes a Ward current.  All stated involution, projection, closed-range
and strict-observation assumptions nevertheless hold.
\end{proof}

\begin{proposition}[Why the negative Gamma norm does not close the uniform
graph estimate --- \PROVED]
Let \(w(\tau)=c+g_\Gamma(\tau)\), where the fixed constant \(c\) is chosen
large enough that \(w>0\), so that
\(w(\tau)\asymp1+\log(2+|\tau|)\).  Write
\[
 \|f\|_{\mathcal H_w}^2=\int_{\mathbb R}w|\widehat f|^2,\qquad
 \|u\|_{\mathcal H_w^*}^2=\int_{\mathbb R}w^{-1}|\widehat u|^2.
\]
Then the order-\(w\) multiplier obeys
\begin{equation}
 M_w:\mathcal H_w\longrightarrow\mathcal H_w^*
 \quad\hbox{bounded},\qquad
 \|M_wf\|_{\mathcal H_w^*}=\|f\|_{\mathcal H_w},    \tag{730}
\end{equation}
but \(M_w:\mathcal H_w\to\mathcal H_w\) is unbounded.  Therefore choosing
the Hadamard bulk norm to be merely the dual Gamma norm can make a formal
version of (605) true only in a distributional target; it does not produce
the bounded physical endomorphism required in (630), (646), or (661).
The missing global estimate must prove a genuine logarithmic regularity
gain, or an equivalent cancellation, rather than weaken the codomain.
\end{proposition}

\begin{proof}
The equality in (730) is immediate:
\[
 \int w^{-1}|w\widehat f|^2=\int w|\widehat f|^2.
\]
For the second assertion choose unit \(L^2\) Fourier packets supported
where \(w\asymp N\), and normalize them to have \(\mathcal H_w\)-norm one.
Their images under \(M_w\) have squared \(\mathcal H_w\)-norm comparable
to \(N^2\), which tends to infinity.  Thus dual control cannot be promoted
to a same-space connection.
\end{proof}

\begin{proposition}[Exact full-orbit Volterra kernel and absence of
orbitwise decay --- \PROVED]
Let \(\rho=\beta+i\gamma\), \(0<\beta<1\), \(\beta\ne\frac12\), and group
the full off-line functional-equation orbit
\[
 \mathcal O_\rho=\{\beta+i\gamma,\,1-\beta+i\gamma,\,
                   \beta-i\gamma,\,1-\beta-i\gamma\}.
\]
On the positive Volterra triangle, with \(x=t-u>0\), the sum of the
four exponential kernels in (614) is
\begin{equation}
 \boxed{
 k_{\mathcal O_\rho}(x)
 =2m_\rho\cos(\gamma x)
   \left(e^{\beta x}+e^{(1-\beta)x}\right).}         \tag{731}
\end{equation}
On the negative triangle, with \(x=u-t>0\), the oriented kernel is instead
\[
 -2m_\rho\cos(\gamma x)
   \left(e^{-\beta x}+e^{-(1-\beta)x}\right).
\]
In contrast, the scalar canonical-product correction is summable:
\begin{equation}
 \sum_{\lambda\in\mathcal O_\rho}\frac1\lambda
 =2\frac{\beta}{\beta^2+\gamma^2}
  +2\frac{1-\beta}{(1-\beta)^2+\gamma^2}
 =O(\gamma^{-2}).                                   \tag{732}
\end{equation}
The Volterra part has no corresponding orbitwise norm decay.  More
precisely, for a fixed nonzero smooth \(f\) on the positive half-window and
\(f_\gamma(t)=e^{-i\gamma t}f(t)\),
\[
 \mathcal V_{\beta+i\gamma}f_\gamma
 =e^{-i\gamma\cdot}\mathcal V_\beta f,\qquad
 \mathcal V_{1-\beta+i\gamma}f_\gamma
 =e^{-i\gamma\cdot}\mathcal V_{1-\beta}f.
\]
Hence the two positive-ordinate members add with a norm bounded below
independently of \(\gamma\); the negative-ordinate members tend to zero on
this test vector after demodulation.  Functional-equation pairing alone
therefore cannot provide the uniform dyadic estimate.
For a critical orbit the same conclusion holds with the two-point positive
kernel \(2m_\rho e^{x/2}\cos(\gamma x)\).
\end{proposition}

\begin{proof}
Summing \(e^{\lambda x}\) over the four displayed points gives (731).
Pairing conjugates in the reciprocal sum gives (732), whose orbit sum is
absolutely convergent by the Riemann--von Mangoldt count.

The two modulation identities follow directly from (614).  After
multiplication by \(e^{i\gamma u}\), the two negative-ordinate terms contain
the oscillatory factor \(e^{-2i\gamma t}\).  For fixed smooth \(f\), an
integration by parts on the Volterra triangle makes those terms tend to
zero in \(L^2\).  The remaining operator
\(\mathcal V_\beta+\mathcal V_{1-\beta}\) applied to \(f\) is nonzero for a
suitable \(f\), proving the uniform lower bound and the final assertion.
\end{proof}

\begin{theorem}[Uniform square estimate for the full Volterra residue
analysis --- \PROVED]
Let \(\mathcal V_\rho\) be (614), with every distinct nontrivial zero
weighted by its multiplicity.  Then, for every fixed \(T\), there is
\(C_T<\infty\) such that
\begin{equation}
 \boxed{
 \sum_{\rho}m_\rho
 \|\mathcal V_\rho F\|_{L^2(I_T)}^2
 \le C_T\|F\|_{\mathcal Q_{\Gamma,T}}^2.}           \tag{733}
\end{equation}
Consequently the map
\[
 F\longmapsto\{\sqrt{m_\rho}\,\mathcal V_\rho F\}_\rho
\]
is a bounded analysis operator from the Gamma form domain into
\(\ell^2_\rho(L^2(I_T))\), uniformly under finite residue truncation.
\end{theorem}

\begin{proof}
The Gamma form norm is equivalent on a fixed compact window to
\[
 \int_{\mathbb R}
  \bigl(1+\log(2+|\tau|)\bigr)|\widehat{E_TF}(\tau)|^2\,d\tau.
\]
Multiplication in the physical variable by a half-line indicator is
uniformly bounded in this logarithmic Fourier space.  Indeed, after
Fourier transform it is a linear combination of the identity and a
modulated Hilbert transform; the weight
\(1+\log(2+|\tau|)\) is an \(A_2\) weight, with translations changing its
norm by a fixed-window constant.  Hence, uniformly in \(u\in I_T\),
\begin{equation}
 \|1_{[u,T]}F\|_{\mathcal Q_{\Gamma,T}}
 +\|1_{[-T,u]}F\|_{\mathcal Q_{\Gamma,T}}
 \le C_T\|F\|_{\mathcal Q_{\Gamma,T}}.              \tag{734}
\end{equation}

For \(u>0\), put \(F_u=1_{[u,T]}F\).  Formula (614) gives
\[
 (\mathcal V_\rho F)(u)=e^{-\rho u}K_{F_u}(\rho).
\]
Since \(0<\Re\rho<1\), the exponential factor has modulus at most one.
Apply the weighted Plancherel--Pólya estimate (551) to \(F_u\), use (734),
with the harmless fixed-window multiplier \(e^{t/2}\) converting
\(K_{F_u}(\rho)\) into the coordinate \(\Phi_{e^{t/2}F_u}(z_\rho)\), and
integrate over \(0<u<T\):
\[
 \sum_\rho m_\rho
 \int_0^T|(\mathcal V_\rho F)(u)|^2\,du
 \le C_T\|F\|_{\mathcal Q_{\Gamma,T}}^2.
\]
For \(u<0\), use \(F_u=1_{[-T,u]}F\).  The extra factor
\(|e^{-\rho u}|\le e^T\) is harmless, and the same argument proves the
negative-half estimate.  Adding the two halves gives (733).
\end{proof}

\begin{theorem}[One-sided Gamma--Hadamard--prime--Tate kernel identity ---
\PROVED]
Let the zero sum be taken with multiplicity in the symmetric canonical
Hadamard order, and put
\[
 w_\Gamma(x)=\frac{e^{-x/2}}{1-e^{-2x}},\qquad x>0.
\]
As distributions on the open half-line \((0,\infty)\),
\begin{equation}
 \boxed{
 \sum_\rho m_\rho e^{(\rho-\frac12)x}
 +w_\Gamma(x)
 +\sum_{n\ge2}\frac{\Lambda(n)}{\sqrt n}
       \delta_{\log n}(x)
 =2\cosh(x/2).}                                    \tag{735}
\end{equation}
Thus the complete coherent zero synthesis, after Gamma and prime
centering, collapses exactly to the two Tate exponentials.  No estimate of
individual zero blocks is involved.
\end{theorem}

\begin{proof}
For \(\Re s>1\), logarithmic differentiation gives
\[
 \frac{\xi'}{\xi}(s)
 =\frac1s+\frac1{s-1}-\frac12\log\pi
  \frac12\psi(s/2)-\sum_{n\ge2}\Lambda(n)n^{-s}.
\]
On the other hand, the canonical Hadamard product gives
\[
 \frac{\xi'}{\xi}(s)
 =B+\sum_\rho m_\rho
   \left(\frac1{s-\rho}+\frac1\rho\right).
\]
Constants in \(s\) are supported at \(x=0\) under inverse Laplace
transformation and hence disappear on the open half-line.  The integral
formula
\[
 \psi(s/2)=-\gamma+
 2\int_0^\infty
 \frac{e^{-2x}-e^{-sx}}{1-e^{-2x}}\,dx
\]
shows that the nonlocal inverse Laplace kernel of
\(\frac12\psi(s/2)\) is \(-1/(1-e^{-2x})\).  Inverting the remaining terms
therefore yields
\[
 \sum_\rho m_\rho e^{\rho x}
 =1+e^x-\frac1{1-e^{-2x}}
  -\sum_{n\ge2}\Lambda(n)\delta_{\log n}(x).
\]
Multiplication by \(e^{-x/2}\), using
\(e^{-x/2}\delta_{\log n}=n^{-1/2}\delta_{\log n}\), gives (735).
\end{proof}

\begin{theorem}[Canonical endpoint completion and primitive bilateral
cancellation --- \PROVED]
Let \(\mathfrak z_+\) and \(\mathfrak g_+\) be the distributions supported
in \([0,\infty)\) whose Laplace transforms, for \(\Re z>\frac12\), are
\begin{align}
 \mathcal L\mathfrak z_+(z)
 &=\frac{\xi'}{\xi}\left(\frac12+z\right),\notag\\
 \mathcal L\mathfrak g_+(z)
 &=\frac12\log\pi-\frac12
   \psi\left(\frac14+\frac z2\right),                \tag{736}
\end{align}
and let
\[
 \mu_P:=\sum_{n\ge2}\frac{\Lambda(n)}{\sqrt n}
                  \delta_{\log n}.
\]
Then, as distributions on the closed half-line,
\begin{equation}
 \boxed{\mathfrak z_++\mathfrak g_++\mu_P
        =2\cosh(x/2)\,1_{x\ge0}.}                   \tag{737}
\end{equation}
In particular (736) fixes uniquely the finite-part and delta normalization
at \(x=0\); no separate affine constant remains undetermined.

For a compactly supported smooth source define the right and left
triangular actions
\[
 (\mathcal R_kF)(u)=\langle k(x),F(u+x)1_{0\le x\le T-u}\rangle,
\quad
 (\mathcal L_kF)(u)=\langle k(x),F(u-x)1_{0\le x\le T+u}\rangle.
\]
Then
\begin{equation}
 (\mathcal R_{\mathfrak z_++\mathfrak g_++\mu_P}
 +\mathcal L_{\mathfrak z_++\mathfrak g_++\mu_P})F
 =
 e^{-u/2}M_+(F)+e^{u/2}M_-(F).                     \tag{738}
\end{equation}
Consequently the complete bilateral centered synthesis vanishes identically
on the primitive core.
\end{theorem}

\begin{proof}
Put \(s=\frac12+z\).  The logarithmic derivative formula gives
\[
 \frac{\xi'}{\xi}(s)
 =\frac1s+\frac1{s-1}-\frac12\log\pi
  +\frac12\psi(s/2)+\frac{\zeta'}{\zeta}(s).
\]
Since
\[
 \frac{\zeta'}{\zeta}(s)
 =-\sum_{n\ge2}\Lambda(n)n^{-s},
\]
adding the two transforms in (736) and
\(\mathcal L\mu_P(z)=\sum\Lambda(n)n^{-1/2-z}\) leaves
\[
 \frac1{z+1/2}+\frac1{z-1/2},
\]
the Laplace transform of \(2\cosh(x/2)1_{x\ge0}\).  Injectivity of the
Laplace transform for distributions supported in a half-line proves (737),
including the endpoint normalization.

Apply the two triangular operators to (737).  Their sum is
\[
 \int_{-T}^{T}2\cosh((t-u)/2)F(t)\,dt
 =e^{-u/2}\int_{-T}^{T}e^{t/2}F(t)\,dt
  +e^{u/2}\int_{-T}^{T}e^{-t/2}F(t)\,dt,
\]
which is (738) and vanishes when \(M_+(F)=M_-(F)=0\).
\end{proof}

\begin{theorem}[Tate triangular Green operator and exact normal velocity ---
\PROVED]
Define the bounded triangular operator
\begin{equation}
 (\mathcal C_T^{\rm Tate}F)(u)
 :=\int_u^T2\cosh((t-u)/2)F(t)\,dt.                 \tag{739}
\end{equation}
Let \(B=B_1\), \(P=P_1\), \(K=BB^*\), and let
\(J_2=\left(\begin{smallmatrix}0&1\\1&0\end{smallmatrix}\right)\).
Then
\begin{align}
 \mathcal C_T^{\rm Tate}
 +(\mathcal C_T^{\rm Tate})^*
 &=B^*J_2B,                                         \tag{740}\\
 B\mathcal C_T^{\rm Tate}P
 &=BQP=2DB_1',                                      \tag{741}
\end{align}
where \(QF(t)=tF(t)\) and \(D=\operatorname{diag}(-1,1)\).
Consequently
\[
 P\mathcal C_T^{\rm Tate}P
\quad\hbox{is skew-adjoint on }\mathcal P_T,
\]
and
\[
 (I-P)\mathcal C_T^{\rm Tate}P
 =(I-P)QP
 =2B^*K^{-1}DB_1'.
\]
Thus the canonically summed Gamma--Hadamard--prime current has exactly the
logarithmic Tate normal velocity (529), although its tangential compression
is a bounded skew phase.
\end{theorem}

\begin{proof}
The adjoint of the right triangular operator (739) is the left triangular
operator with the same real kernel.  Their sum has kernel
\(2\cosh((t-u)/2)\) on the full square, hence
\[
 ((\mathcal C_T^{\rm Tate})+
   (\mathcal C_T^{\rm Tate})^*)F(u)
 =e^{-u/2}M_+(F)+e^{u/2}M_-(F)
 =(B^*J_2B)F(u),
\]
which proves (740).

For \(F\in\ker B\), Fubini gives
\begin{align*}
 M_-(\mathcal C_T^{\rm Tate}F)
 &=\int_{-T}^Tte^{-t/2}F(t)\,dt,\\
 M_+(\mathcal C_T^{\rm Tate}F)
 &=\int_{-T}^Tte^{t/2}F(t)\,dt.
\end{align*}
All remaining terms are multiples of \(M_-(F)\) or \(M_+(F)\) and vanish.
This proves \(B\mathcal C_T^{\rm Tate}P=BQP\); equation (528) gives the
last expression in (741).  Formula (740) compressed by \(P\) proves
skew-adjointness, and applying
\(I-P=B^*K^{-1}B\) proves the final normal-block identity.
\end{proof}

\begin{corollary}[Clamped regularization of the logarithmic generator ---
\PROVED]
For \(F\in C_c^\infty(I_T)\cap\mathcal P_T\), put
\(C=\mathcal C_T^{\rm Tate}\) and
\(L=\partial_u^2-\frac14\).  Then
\begin{equation}
 \boxed{L(CF)=-2F',\qquad CF\in H_0^2(-T,T).}        \tag{742}
\end{equation}
Equivalently, after zero extension,
\begin{equation}
 \widehat{CF}(\tau)
 =-\frac{2i\tau}{\tau^2+1/4}\widehat F(\tau),        \tag{743}
\end{equation}
where the apparent poles at \(\tau=\pm i/2\) are removable by the two
primitive moments.  Thus \(C\) is a bounded order-\((-1)\) source
regularization of the logarithmic derivative and retains, by (741),
exactly the same normal Tate component as \(Q\).
\end{corollary}

\begin{proof}
Differentiate (739).  Since
\(k(x)=2\cosh(x/2)\) has \(k(0)=2\), \(k'(0)=0\), and
\(k''=k/4\), one obtains
\[
 (CF)''-\tfrac14CF=-2F'.
\]
At \(u=T\), both \(CF\) and its derivative vanish because \(F\) is compactly
supported in the open interval.  At \(u=-T\), the value and derivative are
linear combinations of \(M_-(F)\) and \(M_+(F)\), hence also vanish.
This proves (742).  Fourier transformation of the clamped zero extension
gives (743).  The final assertion is (741).
\end{proof}

\begin{corollary}[Explicit skew correction of every two-component Tate
terminal --- \PROVED]
Let \(L:\mathcal P_T\to\mathbb C^2\) be bounded and define
\begin{equation}
 \mathcal A_L
 :=B^*K^{-1}LP-PL^*K^{-1}B.                         \tag{744}
\end{equation}
Then
\begin{equation}
 \mathcal A_L^*=-\mathcal A_L,\qquad
 B\mathcal A_LP=L,\qquad
 P\mathcal A_LP=0.                                  \tag{745}
\end{equation}
In particular put
\[
 L_\star=\left(D-\frac14KDK^{-1}\right)B_1',
 \qquad L_0=2DB_1',
\]
and
\begin{equation}
 \mathcal C_T^\star
 :=\mathcal C_T^{\rm Tate}+\mathcal A_{L_\star-L_0}. \tag{746}
\end{equation}
Then
\[
 B\mathcal C_T^\star P=L_\star,\qquad
 P\mathcal C_T^\star P=P\mathcal C_T^{\rm Tate}P,
\]
so the two-component target on the right side of (674) is realized by a
bounded source-defined Tate correction, with no inverse of \(H_T\) and
without changing the tangential skew phase.  This does not identify
\(\mathcal C_T^\star\) with the raw GET commutator phase in (671).
\end{corollary}

\begin{proof}
The two summands in (744) are negative adjoints of one another.  Since
\(BB^*=K\) and \(BP=0\),
\[
 B\mathcal A_LP=BB^*K^{-1}LP=L.
\]
Multiplication by \(P\) on the left annihilates the first summand, while
multiplication by \(P\) on the right annihilates the second; hence
\(P\mathcal A_LP=0\).  Equations (741) and (744)--(745) applied to
\(L=L_\star-L_0\) prove all assertions about (746).
\end{proof}

\begin{theorem}[Absolutely convergent global vertical Green flux ---
\PROVED]
For \(F,G\) in the smooth primitive core define, on every off-line orbit
\(\rho=\frac12+\alpha+i\gamma\), \(0<\alpha<\frac12\),
\begin{align*}
 D_{\rho,F}(y)&=\Phi_F(\gamma-iy)-\Phi_F(\gamma+iy),\\
 S_{\rho,QF}(y)&=\Phi_{QF}(\gamma-iy)+\Phi_{QF}(\gamma+iy).
\end{align*}
Then the polarized series
\begin{equation}
 \mathfrak F_T(F,G):=
 \sum_{\rho\ {\rm off}}m_\rho\int_0^\alpha
 \frac12\left[
 S_{\rho,QF}(y)\overline{D_{\rho,G}(y)}
 +D_{\rho,F}(y)\overline{S_{\rho,QG}(y)}
 \right]dy                                          \tag{747}
\end{equation}
converges absolutely and continuously on the smooth core with its
\(\mathcal Q_{\Gamma,T}\)-graph norm augmented by \(QF,QG\).  Moreover
\begin{equation}
 \boxed{
 \mathfrak F_T(F,G)
 =\langle\mathcal N_TF,\mathcal N_TG\rangle.}        \tag{748}
\end{equation}
Thus the finite-orbit Green identity (527) passes to the full residue
boundary without a formal interchange of an infinite zero sum and a
vertical integral.
\end{theorem}

\begin{proof}
The weighted strip sampling argument used in (551), applied uniformly to
the horizontal segments
\(\{\gamma\pm iy:0\le y\le\alpha\}\), gives
\[
 \sum_{\rho\ {\rm off}}m_\rho\int_0^\alpha
 \left(|D_{\rho,F}(y)|^2+|S_{\rho,QF}(y)|^2\right)dy
 \le C_T\left(
 \|F\|_{\mathcal Q_{\Gamma,T}}^2+
 \|QF\|_{\mathcal Q_{\Gamma,T}}^2\right).
\]
Indeed each unit-height box contains \(O(\log(2+|\gamma|))\) zeros,
vertical displacement is at most \(1/2\), and the enlarged boxes have
bounded overlap.  The same estimate holds for \(G\).  Cauchy--Schwarz in
the measure
\(\sum_\rho m_\rho1_{[0,\alpha]}(y)\,dy\) proves absolute convergence and
continuity of (747).

For a finite family of orbits, polarization of (527) identifies the
truncated form with the corresponding negative residue Gram.  The
left-hand side converges by the preceding estimate, while the right-hand
side converges by boundedness of
\(\mathcal N_T=P_-\mathcal E_T\), proved in (549)--(550).  Passing to the
limit proves (748).
\end{proof}

\begin{theorem}[Deformed one-sided identity and signed Tate tangent ---
\PROVED]
For \(0<\sigma<1\), define half-line distributions by
\begin{align}
 \mathcal L\mathfrak z_{\sigma,+}(z)
 &=\frac{\xi'}{\xi}(\sigma+z),\notag\\
 \mathcal L\mathfrak g_{\sigma,+}(z)
 &=\frac12\log\pi-\frac12\psi((\sigma+z)/2),\notag\\
 \mu_{P,\sigma}
 &=\sum_{n\ge2}\Lambda(n)n^{-\sigma}\delta_{\log n}.
                                                               \tag{749}
\end{align}
Then, by meromorphic continuation from \(\Re(\sigma+z)>1\),
\begin{equation}
 \boxed{
 \mathfrak z_{\sigma,+}
 +\mathfrak g_{\sigma,+}
 +\mu_{P,\sigma}
 =
 \left(e^{-\sigma x}+e^{(1-\sigma)x}\right)1_{x\ge0}.} \tag{750}
\end{equation}
Use (750) with parameter \(\sigma\) on the right triangle and with
parameter \(1-\sigma\) on the reflected left triangle.  Their derivative
at \(\sigma=\frac12\) has the full-square kernel
\begin{equation}
 \boxed{
 -2(t-u)\cosh((t-u)/2).}                             \tag{751}
\end{equation}
Consequently its action is a rank-four Tate operator.  On a primitive
source it reduces to
\[
 -e^{-u/2}\int_{-T}^Tt e^{t/2}F(t)\,dt
 -e^{u/2}\int_{-T}^Tt e^{-t/2}F(t)\,dt,
\]
up to the terms proportional to \(M_\pm(F)\), which vanish.  Thus the
oriented critical tangent of the complete Gamma--Hadamard--prime synthesis
is exactly the logarithmic Tate boundary datum in (528), not an
uncontrolled infinite-rank remainder.
\end{theorem}

\begin{proof}
The logarithmic derivative formula at \(s=\sigma+z\) gives
\[
 \frac{\xi'}{\xi}(s)
 +\frac12\log\pi-\frac12\psi(s/2)
 -\frac{\zeta'}{\zeta}(s)
 =\frac1s+\frac1{s-1}.
\]
The Laplace transform of \(\mu_{P,\sigma}\) is
\(-\zeta'/\zeta(\sigma+z)\), proving (750) by injectivity.

Differentiate the right-hand kernel in (750) at \(\sigma=1/2\):
\(-2x\cosh(x/2)\).  Reflection replaces \(\sigma\) by \(1-\sigma\), so its
derivative on the left triangle has the opposite sign.  With
\(x=|t-u|\), the two pieces combine into (751).  Expanding the hyperbolic
cosine and \(t-u\) expresses the resulting full-square integral through
the four functionals
\(M_\pm(F)\) and
\(\int t e^{\pm t/2}F(t)\,dt\).  Dropping the first two on the primitive
space gives the displayed formula and, by (528), the asserted Tate
identification.
\end{proof}

\paragraph{Scope of the bilateral cancellation --- \AUDIT.}
Equations (736)--(738) complete the coherent Gamma--Hadamard--prime
synthesis as a canonically normalized distribution and show that its
bilateral sum is purely Tate.  This is stronger than the diagonal estimate
(733), but it is still a cancellation identity, not the positive product
rule (661).  The remaining calculation must keep the two triangular
orientations as separate boundary ports and prove that their polarized
energy difference is
\(\frac12I+\Gamma_T^*\Gamma_T\).  Collapsing them with (738) before taking
that energy would erase the jump carrier, exactly as in (269)--(271).
Equations (739)--(741) nevertheless recover one indispensable part of that
energy calculation: the complete current has exactly the previously
computed moving Tate normal velocity, with no normalization adjustment.
What remains is the tangential Hermitian Gamma/residue contribution.
The corrector (744)--(746) also proves that the rank-two primitive terminal
is no longer an obstruction.  It does not manufacture (661): its
compression is zero, so it contributes no Hermitian diagonal on a
first-contact vector.
The deformed identity (749)--(751) further proves that the correctly
oriented tangent of the complete coherent synthesis has no hidden bulk
remainder: after the right/left cancellation it is exactly the rank-four
moving Tate datum.  The still missing term is therefore the derivative of
the separated port \emph{norm}, not the derivative of its synthesized
amplitude.

\begin{proposition}[First-jet gauge obstruction: amplitude tangency does
not determine energy tangency --- \PROVED]
Let \(\mathcal K\) be a Hilbert port space and let
\((r_\sigma,\ell_\sigma)\) be differentiable port vectors.  Knowledge of
the synthesized amplitude
\[
 a_\sigma=r_\sigma+\ell_\sigma
\]
and of its derivative at \(\sigma_0\) does not determine the derivative of
the separated energy
\[
 E_\sigma=\|r_\sigma\|^2+\|\ell_\sigma\|^2.
\]
Indeed, for arbitrary \(h\in\mathcal K\), the gauge
\begin{equation}
 \widetilde r_\sigma
 =r_\sigma+(\sigma-\sigma_0)h,\qquad
 \widetilde\ell_\sigma
 =\ell_\sigma-(\sigma-\sigma_0)h                    \tag{752}
\end{equation}
preserves \(a_\sigma\) identically, while
\begin{equation}
 \widetilde E_{\sigma_0}'-E_{\sigma_0}'
 =2\operatorname{Re}
   \langle h,r_{\sigma_0}-\ell_{\sigma_0}\rangle.    \tag{753}
\end{equation}
On the primitive Gamma--Hadamard--prime synthesis,
\(\ell_{1/2}=-r_{1/2}\) by (738), so (753) can take arbitrary real values
unless the port itself vanishes.  Consequently (735)--(751), including the
complete amplitude tangent and Tate normalization, cannot imply the norm
identity (661).  A separately constructed isometric/Markov port embedding
whose gauge is fixed before critical specialization is indispensable.
\end{proposition}

\begin{proof}
The two perturbations in (752) cancel in their sum, so the amplitude and
all its prescribed first-order data are unchanged.  Differentiate the two
squared norms at \(\sigma_0\); the added terms are
\(2\operatorname{Re}\langle h,r_{\sigma_0}\rangle\) and
\(-2\operatorname{Re}\langle h,\ell_{\sigma_0}\rangle\), which gives
(753).  Substitution of \(\ell_{1/2}=-r_{1/2}\) proves the final assertion.
\end{proof}

\paragraph{Terminal independent theorem --- \OPEN.}
After (752)--(753), the remaining assertion can no longer be phrased as
another kernel or normalization calculation.  One must construct the
joint port embedding \(\mathscr P_{\sigma,T}\), independently of its
desired norm, and prove the critical isometric derivative
\begin{equation}
 \boxed{
 \left.\frac d{d\sigma}
 \langle\mathscr P_{\sigma,T}F_\sigma,
        \mathscr P_{\sigma,T}G_\sigma\rangle
 \right|_{\sigma=1/2}
 =
 \left\langle
 \left(\frac12I+\Gamma_T^*\Gamma_T\right)
 \mathcal E_TF,\mathcal E_TG\right\rangle,}          \tag{754}
\end{equation}
where \(F_\sigma,G_\sigma\) use the already proved moving-Tate transport.
Equation (754) is the norm-level content of (220j), (661), and the Ward
gate (716).  Defining the port norm or its gauge from the right-hand side
would be circular.

\paragraph{Ward identity status --- \OPEN.}
The available prime and Gamma phases give explicit skew source kernels
(663) and (675)--(678), while the Hadamard graph (601)--(617) specifies how
to retain their residue principal parts before compression.  To finish the
construction one must assemble them into one bounded ambient
\(C_T^{\rm amb}\) and prove (722).  The commutator identity (672) controls
the derivative of the phase on physical source vectors, but it does not
contain tests against an ambient core of
\((J_{\rm res}\mathscr R_T)^\perp\); hence (722) does not follow from
(672) alone.  Formula (723) further shows that the exterior block cannot
supply the missing positive diagonal on a zero mode.  No definition of
\(C_T^{\rm amb}\) through the desired range inclusion is permitted.
The equivalence (725)--(727) rules out treating the ambient current as an
automatic dilation: its arithmetic construction must prove new content
strong enough to establish the terminal metric identity itself.
Proposition (728) also rules out proving this content by finite-orbit
verification alone: the normal equation is invisible at every finite
level and appears only in the global closed-range limit.
The countermodel (729) proves that abstract Krein geometry cannot fill this
limit; the missing input must distinguish the arithmetic range from a
general closed subspace of a Krein space.
Proposition (730) rules out the other formal escape, namely declaring the
finite part bounded only in the negative Gamma norm.
Equations (731)--(732) show that the canonical scalar subtraction is
harmless but orbit symmetry does not reduce the Volterra norm.  Gamma and
prime centering must therefore occur at the level of a complete dyadic
operator block before norms are taken.
The square estimate (733) supplies the missing diagonal Bessel bound for
the Volterra residue family.  It does not yet bound its coherent sum; the
remaining dyadic problem is precisely the synthesis/cross-block estimate.

\begin{theorem}[Exact logarithmic Euler commutator and primitive boundary
remainder --- \PROVED]
Let \(M_tF(t)=tF(t)\), let \(S_a\) denote the zero-extension shift on
\(I_T\), and define the finite skew-adjoint logarithmic Euler phase
\begin{equation}
 \mathscr C_T:=
 \sum_{2\le n<e^{2T}}
 \frac{\Lambda(n)}{\log n\sqrt n}
       (S_{\log n}-S_{-\log n}).                    \tag{663}
\end{equation}
Then
\begin{equation}
 \boxed{[M_t,\mathscr C_T]=
 \sum_{2\le n<e^{2T}}
 \frac{\Lambda(n)}{\sqrt n}
       (S_{\log n}+S_{-\log n}).}                  \tag{664}
\end{equation}
Let \(P=P_T\) be the primitive projection, \(P^\perp=I-P\), and put
\(M_P=PM_tP\), \(C_P=P\mathscr C_TP\).  The compressed commutator is
\begin{align}
 P[M_t,\mathscr C_T]P
 &=[M_P,C_P]+\mathscr R_T^{\rm Eul},\notag\\
 \mathscr R_T^{\rm Eul}
 &:=PM_tP^\perp\mathscr C_TP
    -P\mathscr C_TP^\perp M_tP,                    \tag{665}
\end{align}
where \(\mathscr R_T^{\rm Eul}\) is selfadjoint and has rank at most four.
Consequently the bounded equation (585) has the exact source rewrite
\begin{equation}
 \mathsf W_T-\mathsf D_T
 =m_0I+[M_P,C_P]+\mathscr R_T^{\rm Eul}-\mathsf D_T.\tag{666}
\end{equation}
Thus every nonconstant Euler coefficient is generated by one source
commutator; only the scalar Gamma--Tate normalization and a finite-rank
primitive boundary mismatch remain outside that commutator.
\end{theorem}

\begin{proof}
For every real \(a\), direct evaluation on the zero-extended source gives
\([M_t,S_a]=aS_a\).  Apply this to the two orientations in (663); the
factor \(\log n\) cancels and proves (664).  Insert \(I=P+P^\perp\) between
the two factors of each product in \(P[M_t,\mathscr C_T]P\).  The
\(P\)-to-\(P\) terms give \([M_P,C_P]\), and the other two terms are
(665).  Both the left side and \([M_P,C_P]\) are selfadjoint because
\(M_t^*=M_t\) and \(\mathscr C_T^*=-\mathscr C_T\); hence so is the
remainder.  Since \(P^\perp\) has rank two, each summand in (665) has rank
at most two.  Equations (582) and (664)--(665) now give (666).
\end{proof}

\begin{proposition}[Two-component test for the Euler/Tate boundary match ---
\PROVED]
Use the notation of (528):
\(B=B_1\), \(K=BB^*\),
\(D=\operatorname{diag}(-1,1)\), and put
\begin{equation}
 L_{C,T}:=B\mathscr C_TP:
 \mathcal P_T\longrightarrow\mathbb C^2.            \tag{667}
\end{equation}
Then
\begin{equation}
 \mathscr R_T^{\rm Eul}
 =2(B_1')^*DK^{-1}L_{C,T}
  +2L_{C,T}^*K^{-1}DB_1'.                           \tag{668}
\end{equation}
In particular the equality
\(\mathscr R_T^{\rm Eul}=\mathsf D_T\) would follow from the explicit
terminal rule
\begin{equation}
 \boxed{
 L_{C,T}=\left(D-\frac14KDK^{-1}\right)B_1'.}       \tag{669}
\end{equation}
Rule (669), however, is false in general.  More precisely, whenever
\(T\ne\log n\) for all \(2\le n<e^{2T}\), there is
\(F\in C_c^\infty(I_T)\cap\mathcal P_T\) such that
\begin{equation}
 L_{C,T}F=0,qquad
 \left(D-\frac14KDK^{-1}\right)B_1'F\ne0.           \tag{670}
\end{equation}
Thus the logarithmic Euler commutator supplies the correct bulk generator
but cannot, by itself, be identified with the Gamma-compensated Tate
connection.
\end{proposition}

\begin{proof}
Since \(P^\perp=B^*K^{-1}B\), equation (528) gives
\[
 BM_tP=2DB_1',\qquad PM_tB^*=2(B_1')^*D.
\]
Moreover \(\mathscr C_T^*=-\mathscr C_T\), so
\(P\mathscr C_TB^*=-L_{C,T}^*\).  Substitute these three identities in
(665) to obtain (668).  Substitution of (669) in (668) gives
\[
 4(B_1')^*DK^{-1}DB_1'-(B_1')^*K^{-1}B_1'
 =\mathcal Q_T^*\mathcal Q_T-
   \mathcal T_T^*\mathcal T_T=\mathsf D_T,
\]
using (532) and (535).

To disprove the rule, choose a nonempty interval \(J\) centred at zero so
small that, for every active \(a=\log n\), either
\(J\pm a\subset I_T\) when \(a<T\), or
\((J\pm a)\cap I_T=\varnothing\) when \(a>T\).  This is possible because
the active set is finite and no \(a\) equals \(T\).  For a primitive source
supported in \(J\), a shift with \(a<T\) is untruncated and its two Tate
moments are scalar multiples of the vanishing moments of \(F\); a shift
with \(a>T\) vanishes after compression.  Hence
\(B\mathscr C_TF=0\).

On \(C_c^\infty(J)\), the four functionals represented by
\(e^{-t/2},e^{t/2},te^{-t/2},te^{t/2}\) are linearly independent.
Therefore, after imposing the first two primitive moments, the vector
\(B_1'F\) can be chosen outside the kernel of the nonzero matrix
\(D-\frac14KDK^{-1}\).  This gives (670).
\end{proof}

\paragraph{Euler-commutator verdict --- \AUDIT.}
Equations (663)--(668) are the desired term-by-term source engineering for
the Euler bulk and its primitive boundary.  They also locate a tempting but
incorrect closure: setting the rank-four compression remainder equal to the
already compensated Tate defect would erase the independent Gamma finite
part.  The countertest (670) proves that this identification is unavailable.
The surviving construction must add the Gamma/Hadamard current before the
primitive compression; only the \emph{sum} of its terminal with
\(\mathscr R_T^{\rm Eul}\), not the Euler remainder alone, can satisfy
(661).

\begin{theorem}[Single Gamma--Euler--Tate phase commutator before primitive
compression --- \PROVED]
Let \(\varphi_\Gamma:\mathbb R\to\mathbb R\) be the continuous phase of
\(B_\Gamma\) in (367), normalized by \(\varphi_\Gamma(0)=0\); thus
\(\varphi_\Gamma'=g_\Gamma\).  On the smooth compact-window core define
the skew operators
\begin{align}
 \mathscr C_{\Gamma,T}
 &:=E_T^*\mathcal F^{-1}M_{i\varphi_\Gamma}\mathcal FE_T,\notag\\
 \mathscr C_{0,T}
 &:=E_T^*\mathcal F^{-1}M_{-im_0\tau}\mathcal FE_T,\notag\\
 \mathscr C_T^{\rm GET}
 &:=\mathscr C_{\Gamma,T}+\mathscr C_{0,T}-\mathscr C_T. \tag{671}
\end{align}
Here \(\mathscr C_T\) is the finite Euler phase (663).  Then, as an
identity of forms on the common smooth primitive core before applying the
primitive projection,
\begin{equation}
 \boxed{[M_t,\mathscr C_T^{\rm GET}]
 =G_{\Gamma,T}-m_0I-
   \sum_{2\le n<e^{2T}}\frac{\Lambda(n)}{\sqrt n}
      (S_{\log n}+S_{-\log n})=A_T.}                \tag{672}
\end{equation}
Thus the whole Gamma--Euler--Tate Wigner--Smith operator is one exact source
commutator before primitive compression, rather than a formal sum of local
virials.

With \(C_P^{\rm GET}=P\mathscr C_T^{\rm GET}P\), primitive compression
gives
\begin{align}
 H_T&=[M_P,C_P^{\rm GET}]+\mathscr R_T^{\rm GET},\notag\\
 \mathscr R_T^{\rm GET}
 &:=PM_tP^\perp\mathscr C_T^{\rm GET}P
   -P\mathscr C_T^{\rm GET}P^\perp M_tP.            \tag{673}
\end{align}
The terminal is a selfadjoint operator of rank at most four.  If
\(L_T^{\rm GET}:=B\mathscr C_T^{\rm GET}P\), then
\[
 \mathscr R_T^{\rm GET}
 =2(B_1')^*DK^{-1}L_T^{\rm GET}
  +2(L_T^{\rm GET})^*K^{-1}DB_1'.
\]
In particular the explicit two-component identity
\begin{equation}
 \boxed{
 L_T^{\rm GET}
 =\left(D-\frac14KDK^{-1}\right)B_1'}              \tag{674}
\end{equation}
would imply \(\mathscr R_T^{\rm GET}=\mathsf D_T\).
\end{theorem}

\begin{proof}
With the Fourier convention of (385), multiplication by \(t\) becomes
\(-i\partial_\tau\).  Hence, for a real differentiable phase \(\varphi\),
\[
 [-i\partial_\tau,iM_\varphi]=M_{\varphi'}.
\]
The support projection commutes with multiplication by \(t\), so
compression to \(I_T\) introduces no extra term in this commutator.
Apply the formula to \(\varphi_\Gamma\) and to \(-m_0\tau\), and combine
it with (664).  This proves (672), which is also the phase-derivative
identity (370)--(371) in source coordinates.

Insert \(I=P+P^\perp\) exactly as in the proof of (665) to get (673).
The rank and adjoint assertions follow in the same way.  Finally repeat the
calculation (668) with \(\mathscr C_T^{\rm GET}\) in place of
\(\mathscr C_T\).  Substitution of (674) gives \(\mathsf D_T\), exactly as
in the proof following (669).
\end{proof}

\paragraph{Scope of the single-phase commutator --- \AUDIT.}
Equation (672) is the first exact source formula in this section which
places Gamma, the affine constant, both Euler orientations and the Tate
projection in their required order.  It repairs the Euler-only failure
(670).  It does not yet assert (674): the Gamma phase has order
\(|\tau|\log(2+|\tau|)\), so its compressed terminal must be taken through
the finite-part/Hadamard graph rather than evaluated as a bounded scalar
multiplier on the Gamma form domain.  Moreover (674), even if proved, would
identify the rank-four primitive terminal of the Wigner--Smith commutator;
the full closure still requires the positive metric product rule (661).
The next noncircular calculation is therefore to regularize
\(L_T^{\rm GET}\) by (601)--(617), prove (674) on that graph, and check that
the resulting Hermitian connection has product rule (661), not merely
commutator rule (672).

\begin{theorem}[Gamma phase resolved into explicit source cells ---
\PROVED]
Put \(b_n=n+\frac14\) and define
\begin{equation}
 \phi_n(\tau):=\frac{\tau}{b_n}
 -2\arctan\!\left(\frac{\tau}{2b_n}\right),
 \qquad n\ge0.                                      \tag{675}
\end{equation}
Then \(\phi_n(0)=0\), \(\phi_n'=h_n\), with \(h_n\) from (385), and
\begin{equation}
 \varphi_\Gamma(\tau)=\sum_{n\ge0}\phi_n(\tau)       \tag{676}
\end{equation}
locally uniformly together with one derivative.  On the compactly
supported smooth core let
\[
 \mathscr C_{\Gamma,n,T}
 :=E_T^*\mathcal F^{-1}M_{i\phi_n}\mathcal FE_T .
\]
Then
\begin{equation}
 \mathscr C_{\Gamma,T}F
 =\sum_{n\ge0}\mathscr C_{\Gamma,n,T}F,\qquad
 [M_t,\mathscr C_{\Gamma,n,T}]
 =\mathcal X_{\Gamma,n}^*\mathcal X_{\Gamma,n},      \tag{677}
\end{equation}
where the sum is interpreted on the smooth core and the second equality is
an identity of forms.  No logarithm of the full Gamma function is required
to define an individual cell.

With the Fourier convention used throughout the manuscript, the
translation-invariant distribution kernel of the uncompressed cell is
\begin{equation}
 c_n(x)=
 -\frac1{b_n}\delta_0'(x)
 -\operatorname{pv}\frac{e^{-2b_n|x|}}{x}.          \tag{678}
\end{equation}
In particular every cell is skew-adjoint and satisfies
\(x\,c_n(x)=b_n^{-1}\delta_0(x)-e^{-2b_n|x|}\), the
source-kernel form of its positive Gamma commutator.
\end{theorem}

\begin{proof}
Direct differentiation of (675) gives
\[
 \phi_n'(\tau)
 =\frac1{b_n}-\frac{b_n}{b_n^2+\tau^2/4}=h_n(\tau).
\]
The estimate
\[
 |x-\arctan x|\le C\min(|x|^3,|x|)
\]
shows local uniform convergence of \(\sum_n\phi_n\), while (385) gives
local uniform convergence of the derivative series.  Both sides of (676)
vanish at zero and their derivatives equal \(g_\Gamma\), proving (676).
Equations (677) follow by Fourier multiplication and the commutator
calculation used in (672).

For (678), the Fourier transform of
\(-b_n^{-1}\delta_0'\) is \(i\tau/b_n\).  The transform of the odd
principal-value kernel vanishes at \(\tau=0\), and differentiation with
respect to \(\tau\) gives
\[
 -i\int_{\mathbb R}e^{-2b_n|x|}e^{i\tau x}\,dx
 =-\frac{4ib_n}{4b_n^2+\tau^2}.
\]
Integration from zero to \(\tau\) gives
\(-2i\arctan(\tau/(2b_n))\), proving that the transform of (678) is
\(i\phi_n\).  Finally \(x\delta_0'=-\delta_0\), which gives the last
kernel identity and also proves skew-adjointness.
\end{proof}

\begin{corollary}[Explicit two-functional form of the Gamma terminal ---
\PROVED]
For \(\varepsilon\in\{-1,1\}\), put
\(g_\varepsilon(t)=e^{\varepsilon t/2}1_{I_T}(t)\).  On the smooth
primitive core the Gamma contribution to \(L_T^{\rm GET}\) is
\begin{align}
 (L_{\Gamma,T}F)_\varepsilon
 &:=(B\mathscr C_{\Gamma,T}F)_\varepsilon\notag\\
 &=\sum_{n\ge0}
   \left\langle
    c_n*(E_TF),g_\varepsilon
   \right\rangle_{L^2(I_T)}.                         \tag{679}
\end{align}
Equivalently, after keeping the derivative and principal-value pieces in
their paired cell combination,
\begin{equation}
 L_T^{\rm GET}
 =L_{\Gamma,T}+B\mathscr C_{0,T}P-L_{C,T}.           \tag{680}
\end{equation}
Consequently (674) is no longer a symbolic Gamma normalization: it is the
pair of explicit scalar identities
\begin{equation}
 \boxed{
 (L_{\Gamma,T}F)_\varepsilon+
 (B\mathscr C_{0,T}F)_\varepsilon-
 (L_{C,T}F)_\varepsilon
 =
 \left[
 \left(D-\frac14KDK^{-1}\right)B_1'F
 \right]_\varepsilon,}                              \tag{681}
\end{equation}
for \(\varepsilon=\pm1\) and every smooth primitive \(F\).
\end{corollary}

\begin{proof}
Equations (676)--(678), followed by compression and pairing with the two
rows of \(B\), give (679).  The cell pairing is retained before summation;
this is essential because its derivative and principal-value components
need not converge separately.  Equation (671) gives (680), and substitution
in (674) gives (681).
\end{proof}

\begin{corollary}[The affine \(m_0\) phase has no primitive terminal ---
\PROVED]
On \(C_c^\infty(I_T)\cap\mathcal P_T\),
\begin{equation}
 B\mathscr C_{0,T}P=0.                               \tag{682}
\end{equation}
Hence the terminal identity (674) is equivalent on this core to
\begin{equation}
 \boxed{
 L_{\Gamma,T}-L_{C,T}
 =\left(D-\frac14KDK^{-1}\right)B_1'.}              \tag{683}
\end{equation}
\end{corollary}

\begin{proof}
The multiplier \(-im_0\tau\) is \(m_0\partial_t\) in the stated Fourier
convention.  For a compactly supported primitive source, integration by
parts gives
\[
 \int_{-T}^Te^{\varepsilon t/2}m_0F'(t)\,dt
 =-\frac{\varepsilon m_0}{2}
   \int_{-T}^Te^{\varepsilon t/2}F(t)\,dt=0.
\]
This proves (682); substitute it in (680)--(681) to obtain (683).
\end{proof}

\begin{theorem}[The raw single-phase terminal identity is false ---
\PROVED]
For every \(T>0\), the Gamma functional \(L_{\Gamma,T}\) in (679) does not
factor, on the primitive smooth core, through the four Tate functionals
\(B\) and \(B_1'\).  In particular (683), hence the sufficient terminal
identity (674), is false for the raw compressed phase
\(\mathscr C_T^{\rm GET}\).
\end{theorem}

\begin{proof}
Fix a nonempty interval \(J\Subset I_T\) avoiding the finite set of
transition points \(\{\pm(T-\log n):2\le n<e^{2T}\}\), and then shrink it
if necessary.  For each oriented Euler shift, the translate of every
source supported in \(J\) is then either wholly retained by \(I_T\) or
wholly discarded.  In the retained case each Tate moment is a fixed scalar
multiple of the corresponding moment of \(F\).  Hence on such sources the
Euler terminal \(L_{C,T}\) factors through \(B\), and therefore vanishes on
primitive sources.

For \(\varepsilon\in\{-1,1\}\), the distribution representing the
\(\varepsilon\)-component of \(L_{\Gamma,T}\) on \(J\) is
\[
 \ell_\varepsilon(s)
 :=\sum_{n\ge0}\int_{-T}^{T}
 e^{\varepsilon t/2}c_n(t-s)\,dt .
\]
Using (678), write
\(\ell_\varepsilon(s)=
e^{\varepsilon s/2}\sum_{n\ge0}R_{n,\varepsilon}(s)\), where
\begin{align}
 R_{n,\varepsilon}(s)
 ={}&\frac{\varepsilon}{2b_n}
 +\int_0^{T+s}
   \frac{e^{-(2b_n+\varepsilon/2)x}}{x}\,dx\notag\\
 &-\int_0^{T-s}
   \frac{e^{-(2b_n-\varepsilon/2)x}}{x}\,dx .        \tag{684}
\end{align}
The two integrals are paired at zero, as prescribed by the principal
value.  Differentiating first at finite cell height and then summing the
normally convergent derivative series gives
\begin{align}
 \frac d{ds}\sum_{n\ge0}R_{n,\varepsilon}(s)
 &=
 \frac{e^{-(1-\varepsilon)(T-s)/2}}
 {(T-s)(1-e^{-2(T-s)})}\notag\\
 &\quad+
 \frac{e^{-(1+\varepsilon)(T+s)/2}}
 {(T+s)(1-e^{-2(T+s)})}.                            \tag{685}
\end{align}
This is real analytic on \((-T,T)\), but as \(s\uparrow T\) its first
term satisfies
\begin{equation}
 \frac{e^{-(1-\varepsilon)(T-s)/2}}
 {(T-s)(1-e^{-2(T-s)})}
 =\frac1{2(T-s)^2}+O\!\left(\frac1{T-s}\right).      \tag{686}
\end{equation}
Consequently \(\ell_\varepsilon\) cannot belong on any open subinterval to
the four-dimensional exponential-polynomial space
\begin{equation}
 \mathscr T_2:=
 \operatorname{span}\{
 e^{-s/2},e^{s/2},se^{-s/2},se^{s/2}\}.             \tag{687}
\end{equation}
Indeed, equality on \(J\) would extend throughout \((-T,T)\) by real
analyticity, whereas every element of \(\mathscr T_2\) is entire and has no
boundary pole.

If \(L_{\Gamma,T}\) factored through \(B\) and \(B_1'\), its two
representing distributions on \(J\) would lie in \(\mathscr T_2\).
Conversely, equality of a functional with a Tate combination on
\(\ker B\cap C_c^\infty(J)\) implies, by the elementary annihilator theorem
for the two independent rows of \(B\), that its representer differs from
that combination by an element of
\(\operatorname{span}\{e^{-s/2},e^{s/2}\}\), and hence still lies in
\(\mathscr T_2\).  The preceding pole contradiction proves that no such
factorization exists.  Since the right side of (683) factors through
\(B_1'\), while \(L_{C,T}\) vanishes on the chosen primitive sources,
(683) is false.
\end{proof}

\paragraph{Gamma-cell terminal status --- \AUDIT.}
Equations (675)--(683) complete the requested term-by-term construction of
the skew Gamma--Euler--Tate terminal on the nuclear smooth core, and
(684)--(687) decide its proposed normalization: (674)/(683) is false.
Therefore the Hadamard finite-part/residue graph may not be treated merely
as a regularization of the raw phase with the same two Tate outputs.  It
must add a genuine boundary state whose terminal cancels the non-Tate
Gamma functionals detected by (685).  Constructing that enlarged
Hermitian state and proving its bounded product rule (661) remains
\OPEN{}; the raw skew phase (671) is now excluded as \(X_T^{\rm src}\).

\begin{theorem}[Exterior completion of the skew Gamma phase has zero
primitive Tate terminal --- \PROVED]
For \(F\in C_c^\infty(I_T)\cap\mathcal P_T\), \(n\ge0\), and
\(\varepsilon\in\{-1,1\}\), define the Abel-completed exponential pairing
\(\mathscr C_{\Gamma,n}:=
\mathcal F^{-1}M_{i\phi_n}\mathcal F\) on the full line and set
\begin{equation}
 \mathcal L_{n,\varepsilon}^{\rm full}(F)
 :=\lim_{\sigma\to\varepsilon/2\atop |\sigma|<1/2}
 \int_{\mathbb R}e^{\sigma t}
   (\mathscr C_{\Gamma,n}E_TF)(t)\,dt .              \tag{688}
\end{equation}
Then the limit exists and
\begin{equation}
 \mathcal L_{n,\varepsilon}^{\rm full}(F)=0
 \qquad(n\ge0,\ \varepsilon=\pm1).                  \tag{689}
\end{equation}
The affine phase and every full-line Euler shift have the same property.
Consequently the Abel-completed full-line terminal of
\(\mathscr C_T^{\rm GET}\) is zero:
\begin{equation}
 \mathcal L_{\varepsilon}^{\rm GET,full}(F)=0.       \tag{690}
\end{equation}
Thus adjoining only the exterior tails of the skew scattering phase
cancels its interior terminal; it does not create
\(\mathsf D_T\) or the positive right side of (661).
\end{theorem}

\begin{proof}
For \(|\sigma|<1/2\), Fourier--Laplace evaluation of the convolution
multiplier gives
\begin{equation}
 \int_{\mathbb R}e^{\sigma t}
   (\mathscr C_{\Gamma,n}E_TF)(t)\,dt
 =i\phi_n(-i\sigma)\,
   M_\sigma(F),\qquad
 M_\sigma(F):=\int_{-T}^{T}e^{\sigma t}F(t)\,dt.     \tag{691}
\end{equation}
For \(n\ge1\), \(\phi_n(-i\sigma)\) is regular at
\(\sigma=\varepsilon/2\), while primitiveness gives
\(M_{\varepsilon/2}(F)=0\).  Hence the limit is zero.

For \(n=0\), the arctangent representation (675), equivalently
\[
 \arctan z=\frac1{2i}\log\frac{1+iz}{1-iz},
\]
shows
\(\phi_0(-i\sigma)=O(\log|\frac12-|\sigma||)\) at either endpoint.
Since \(M_\sigma(F)\) is entire and has a zero at
\(\sigma=\varepsilon/2\),
\[
 M_\sigma(F)=O(|\sigma-\varepsilon/2|).
\]
The product in (691) therefore tends to zero because
\(x\log x\to0\).  This proves (688)--(689).

For the affine phase, the full pairing is a scalar multiple of
\(\sigma M_\sigma(F)\), which also tends to zero.  A full-line shift
multiplies \(M_\sigma(F)\) by \(e^{\pm\sigma\log n}\); hence every Euler
term vanishes at the primitive endpoint as well.  The finite Euler sum and
the cellwise Abel prescription now give (690).
\end{proof}

\paragraph{Skew exterior-completion verdict --- \AUDIT.}
The non-Tate boundary pole in (685) is exactly an interior/exterior splitting
artifact of the skew phase: the Abel exterior tail cancels it.  The
completed skew terminal is zero, not the positive moving-Tate defect.
Therefore neither the raw phase nor its conservative exterior completion
can be the missing \(X_T^{\rm src}\).  The only surviving use of
(671)--(690) is as the skew/Kato component of the construction.  The
load-bearing new object must be the Hermitian finite-part/residue metric
connection whose diagonal contribution is prescribed by (661).

\begin{theorem}[Diagonal no-go for skew Gamma--Euler--Tate candidates ---
\PROVED]
Let \(A_T^{\rm ph}\) be any bounded skew-adjoint operator on
\(\mathscr R_T\) induced by finite linear combinations, block sums or
conservative exterior completions of the phase cells (663)--(691), and
suppose the corresponding source rules preserve the common form domain.
Then
\begin{equation}
 \langle
  (\mathbb H_TA_T^{\rm ph}+
   (A_T^{\rm ph})^*\mathbb H_T)e,e
 \rangle=0                                          \tag{692}
\end{equation}
for every eigenvector \(\mathbb H_Te=\lambda e\).  On the other hand,
\begin{equation}
 \langle\mathbb S_Te,e\rangle
 =\frac12\|e\|^2+\|\Gamma_Te\|^2
 \ge\frac12\|e\|^2.                                 \tag{693}
\end{equation}
Consequently no realization which remains skew-adjoint on the physical
residue range can satisfy (661).  Every valid
connection on that range must contain a nonzero Hermitian component
\(S_T^{\rm met}\), and on a unit eigenvector it must obey
\begin{equation}
 2\lambda\langle S_T^{\rm met}e,e\rangle
 =\frac12+\|\Gamma_Te\|^2.                          \tag{694}
\end{equation}
\end{theorem}

\begin{proof}
Skew-adjointness is preserved by finite block sums, conservative unitary
completion and operator-norm limits.  If
\(\mathbb H_Te=\lambda e\), then
\[
 \langle\mathbb H_TA_T^{\rm ph}e,e\rangle+
 \langle(A_T^{\rm ph})^*\mathbb H_Te,e\rangle
 =\lambda\langle A_T^{\rm ph}e,e\rangle+
  \lambda\overline{\langle A_T^{\rm ph}e,e\rangle}=0,
\]
because the diagonal coefficient of a skew-adjoint operator is purely
imaginary.  This proves (692); (693) is the definition in (661).
Decomposing any prospective \(Y_T=S_T^{\rm met}+A_T^{\rm ph}\) into
Hermitian and skew parts and evaluating (661) on \(e\) gives (694).
\end{proof}

\paragraph{Final construction audit after the skew-phase test ---
\AUDIT.}
The source pieces presently constructed in the manuscript now have a
complete disposition.  The horizontal residue current is bounded but its
descent is the open tangency (656).  The Euler and Gamma phases give the
exact commutator (672), their raw terminal candidate is false by
(684)--(687), and their conservative exterior completion has zero terminal
by (688)--(691).  Theorem (692)--(694) proves that no realization which
remains skew in the physical residue norm can supply the strictly positive
diagonal in (661).  A nonunitary Hadamard pullback could in principle create
the required Hermitian part, but proving that pullback is exactly the open
tangency problem.  The single
missing mathematical object is therefore not another normalization or
phase correction: it is an independently source-defined, uniformly bounded
Hermitian Hadamard--residue metric connection \(S_T^{\rm met}\) satisfying
(694) without division by \(\lambda\).  Constructing it is equivalent, by
(650)--(652), to excluding the first-contact kernel.

\begin{theorem}[(565) is exactly the no-first-contact theorem ---
\PROVED]
Fix a compact window range \(T\in[T_0,T_1]\), assume \(H_{T_0}>0\), and
transport the primitive spaces to one fixed Hilbert space.  Under the
already proved form continuity, the following are equivalent:
\begin{enumerate}
\item \(H_T\) has no zero eigenvalue on the range;
\item \(\inf_{T\in[T_0,T_1]}\inf\operatorname{spec}H_T>0\);
\item the uniform observability estimate (565) holds.
\end{enumerate}
In particular, (565) cannot be proved merely from the positive Gamma
factorization, residue summability and G39 amplitude: proving it across
every finite range is itself a proof of G40.
\end{theorem}

\begin{proof}
Norm-resolvent continuity and compactness of the \(T\)-range make the first
two assertions equivalent.  If \(H_T\ge cI\) uniformly, the established
uniform form bound \(B_T\le C(H_T+I)\le C(1+c^{-1})H_T\), together with
\[
 \int_0^\infty
 \|H_T^{1/2}e^{-uH_T}f\|^2du
 =\tfrac12\|f\|^2,
\]
proves (565).  Conversely, if \(0\ne F\in\ker H_{T_*}\), then
\(e^{-uH_{T_*}}F=F\) and
\[
 \|\mathcal O_{T_*}e^{-uH_{T_*}}F\|^2
 =\langle B_{T_*}F,F\rangle
 \ge\tfrac14\langle G_{\Gamma,T_*}F,F\rangle>0
\]
for every \(u\).  Its integral is infinite, contradicting (565).  The
first-loss theorem identifies absence of such a kernel on every finite
range with G40.
\end{proof}

\paragraph{Consequence for the proposed proof of (565) --- \AUDIT.}
The Bayes--Gamma defect is positive in the ambient boundary, but pulling it
back to the physical source with exactly the observation
\(\mathcal O_T\) is (220j)/(223j).  If that intertwining were inserted,
Pythagoras would prove (565); conversely the preceding theorem shows that
the pullback cannot be inferred from the existing normalizations.  Thus the
only admissible proof of (565) is a genuinely new source-labelled
intertwining or an independent geometric Hodge theorem.  No estimate
currently proved in this manuscript supplies it.
\paragraph{Terminal-construction verdict --- \AUDIT.}
The Gamma ODE terminals telescope as in (230), and the two outer \(j=0\)
states are killed by \(M_\pm F=0\).  The moving-projection terminal (109),
however, is a different positive square and is generally nonzero by
(505)--(507).  The nontrivial residue terminal is Krein rather than Hilbert
by (508).  Promoting it to a positive terminal whose vanishing follows from
the zero-mode equation would remove precisely the off-line negative fibers;
that promotion is not a remaining normalization calculation but the desired
Hodge theorem itself.

\paragraph{Seven-step verdict --- \AUDIT.}
Steps 1, 2 in the honest half-plane, 4 at the symmetric-amplitude level, 6
conditional on (497), and 7 are rigorous.  Step 3 does construct positive
Euler--Gamma cells, but (490) proves that their global unsynthesized energy
diverges before reaching the critical boundary.  Step 5 may now be weakened
to (497); (494) and the causal estimate (287) are stronger sufficient
versions.  The equation of a first-loss mode
removes the Tate mean and fixes (491), but supplies no identity for the
antisymmetric causal datum (492).  Therefore the Abel proposal localizes the
remaining theorem and quantifies its pole energy by (488); it does not yet
prove that theorem.

\section{Requirement-by-requirement completion audit}

\begin{center}
\begin{tabular}{p{0.20\textwidth}p{0.16\textwidth}p{0.54\textwidth}}
\hline
Requirement & Status & Authoritative evidence \\
\hline
Amplitude & \PROVED/\OPEN & (27g)--(27h) prove the finite divisibility
amplitude; identification with the continuum Hodge/old factor is not supplied.
\\
Normalization & \PROVED & (53)--(55) give the restricted Sonin polar
isometry without using row (d). \\
Gamma--Tate Gram & \PROVED & (81)--(84) give the exact source Gram and
balanced factorization; (138)--(149) construct a global gcd--Gamma reference
space, its Gram isometry and its positive midpoint conservation.  Equations
(361)--(371) additionally assemble every prime tower and Gamma into one
explicit finite-window boundary scattering whose compressed phase derivative
is $A_T$.  Equations (535)--(546) prove that the first Gamma cell
independently dominates and factors the complete oriented Tate terminal
defect with a uniform \(3/4\) bound. \\
Typed observation & \PROVED/\OPEN & (84a)--(84c) define the canonical source
angular operator and (138)--(149) give the global cross-place type.  Its
contractive bound remains open.
\\
Continuous reduction & \PROVED & (90)--(94) identify the primitive space
with clamped potentials and give the exact Paley--Wiener defect kernel. \\
Global Poisson mixing & \PROVED/\OPEN & (96)--(98) give the full continuum
Euler divergence; (310)--(327) reduce its physical boundary defect to the
opposite Fourier-parity component.  Its normalized physical Gram, boundary
variation (104), and the required physical realization of that component
remain open. \\
Gamma normalization & \PROVED/\OPEN & (168)--(178a) give Gamma GNS and the
exact ratio--scale identity deriving $m_0-g_\Gamma$ from one Mellin law;
the unweighted deconvolution is excluded on every nonzero compact-window
source by (179)--(180).  Equations (181)--(183) identify the exact completed
Euler--Gamma coherent metric for $t>2$; its joint critical compression remains
open.  Equations (184)--(186) prove why it cannot be an ordinary vector
limit, while (187)--(194) construct its positive form-level and joint
Euler--Gamma replacement.  Equations (208)--(210) prove the exact
non-deconvolved Gamma--Tate metric termination. \\
Tate terminal motion & \PROVED & (105)--(111) give the exact moving
projection, minimum normal lift and positive even/odd terminal square. \\
Conservation & \PROVED/\OPEN & In addition to the exact Krein identities,
(220a)--(220e) construct an ambient joint Euler--Gamma Markov contraction and
its positive conditional-variance defect.  Identification of its critically
renormalized, Tate-compressed observation with $\mathcal Y_TP_T$ remains
open and is the physical Hilbert conservation statement (84c). \\
Positivity & \PROVED/\OPEN & The supplied certified theorem and (248)--(250)
give strict passivity for $0<T\le\log2$.  The estimate for every $T$, namely
(249), is the Weil criterion and remains open.  Equations (468)--(474)
give the exact non-circular leakage Gram and corrected nested certificate at
$T=\frac12\log5$; its two missing interval enclosures have not been computed.
\\
Old-factor compatibility & \PROVED & Restriction of (83) fixes the
authoritative $X_0,Y_0$ columns; (41s)--(41u) prove exact compatibility under
window enlargement by adjoining only unitary neutral birth channels.
Equations (41p)--(41r) concern the stronger impossible identification of the
finite G39 port alone with the full infinite-rank Hodge factor and are not an
old-compatibility requirement. \\
Condition G / row (d) & \OPEN & Equations (330)--(333) identify the exact
source-typed Gamma--Tate bi-gap uniqueness theorem equivalent to the
first-loss closure.  Equations (346)--(350) give the complete negative
certificate: any failure must be a support-saturating nonsmooth state with
no Meyer lift and a non-$L^2$ oriented arithmetic tail.  Exclusion of this
certificate is not yet proved.  Equations (436)--(440) remove the global
co-Poisson domain and Burnol-completeness gates: the local observation is a
finite causal operator with nilpotent inverse.  The remaining statement is
only (441), namely derivation of its annular vanishing from the physical
zero-mode equation.  That Gamma--Tate gap-transport identity remains open,
so the preceding Gram/passivity requirements still prevent a non-circular
closure.  Equations (485)--(501) audit the proposed Abel route completely:
the honest Euler energy starts only at $\sigma>1/2$, while its critical pole
defect has the exact positive limit (496).  It is enough to prove the
      little-o residue estimate (497), equivalently the oriented derived-range
      divisibility (498)--(500); a source-defined Rellich--Clark identity (501)
      would do so.  Equations (653)--(662) give the exact horizontal residue
      current and reduce the final graph equation to a strictly positive
      metric connection, but do not prove its sampling-range tangency or a
      uniform bounded source realization.  Finally (695)--(700) prove that
      the causal and Hermitian routes coincide on the zero-mode space and
      isolate their common missing arrow as
      \(F\in\ker H_T\Rightarrow\mathcal C_T^{\rm simp}F=0\).
      Equations (702)--(707) further prove that the tempting equality-state
      negative-port lift is separated from the physical range by exactly
      half the norm of a hypothetical zero mode, so finite interpolation
      cannot supply that arrow.
      No such identity is yet constructed. \\
\hline
\end{tabular}
\end{center}

\section{Corrected final ledger after the type audit}

\subsection*{PROVED}
\begin{itemize}
\item The norm-level gauge audit (752)--(754): the full synthesized
      amplitude and its critical first derivative leave an arbitrary
      antisymmetric first-jet gauge in the two separated ports.  This gauge
      changes the derivative of their Hilbert norm.  Hence the completed
      Gamma--Hadamard--prime kernel identities do not imply (661); a
      separately normalized critical isometric embedding is logically
      indispensable.
\item The port-level Green package (739)--(751): the canonical triangular
      Tate operator has symmetric part \(B^*J_2B\), and its normal block on
      primitive sources is exactly the logarithmic boundary velocity
      \(2B^*K^{-1}DB_1'\).  The finite-rank skew correction (744) realizes
      the required two-component terminal without changing the tangential
      phase.  Finally the vertical off-line Green flux is absolutely
      summable and equals the complete negative residue Gram.  Hence both
      the Tate normal energy and the global residue energy are now genuine
      bounded forms; their remaining coupling to the Gamma Hermitian port
      is the open product rule.  The deformed synthesis (749)--(751) also
      proves that its correctly reflected critical tangent is exactly the
      rank-four logarithmic Tate datum; no additional amplitude remainder
      survives.
\item The coherent synthesis identity (735)--(738): inverse Laplace
      transformation of the completed Hadamard logarithmic derivative,
      with its canonical endpoint finite part, gives exactly
      \[
      \mathfrak z_++\mathfrak g_++\mu_P
      =2\cosh(x/2)1_{x\ge0}.
      \]
      Its bilateral triangular action is the two-dimensional Tate operator
      and therefore vanishes on primitive sources.  This completes the
      amplitude-level global Gamma--Hadamard--prime synthesis, including
      the formerly ambiguous \(x=0\) normalization.  The two orientations
      must still be kept separate for the positive energy calculation.
\item The first dyadic Hadamard estimates (731)--(734): complete
      functional-equation orbits have the explicit cosine Volterra kernel
      (731); their canonical \(1/\rho\) correction is summable, but their
      Volterra norm does not decay with height.  Nevertheless the full
      residue Volterra family is a uniformly bounded \(\ell^2\)-valued
      analysis operator on the Gamma form domain.  Thus the diagonal
      Bessel estimate is complete; coherent centered synthesis remains the
      cross-block problem.
\item The global-limit audit (728)--(729): every finite residue evaluation
      is surjective, so its Ward normal equation is the tautology \(0=0\).
      The nontrivial normal block exists only in the infinite closed-range
      limit.  A two-dimensional Krein countermodel proves that the abstract
      involution, closed-range and strict-observation axioms still do not
      imply that block identity.  A uniform arithmetic graph theorem is
      therefore indispensable.  Equation (730) further proves that placing
      the finite part only in the dual Gamma space loses exactly the
      same-space regularity needed for the physical connection.
\item The corrected ambient Ward analysis (719)--(727): an ambient
      Hilbert-skew phase must satisfy the range identity (721), tested
      against the full ambient complement as in (722).  In block form this
      is (723); the exterior block has zero real contribution on a
      first-contact vector.  Moreover such a Ward current exists if and
      only if the terminal bounded metric connection (661) exists.  Hence
      the ambient construction is well typed but is not an automatic
      dilation or a weaker substitute for the missing theorem.
\item The ambient dissipative normal form (708)--(718): the positive
      operator \(\frac12I+\Gamma_T^*\Gamma_T\) is generated without any
      inverse by \(K_0=-\frac12JQSQ\).  Closing (661) is equivalent to
      finding a \(J\)-skew ambient arithmetic correction with the one normal
      block (716).  A phase compressed through the physical-range isometry
      has zero normal block and cannot work.  The compulsory tangential
      identity (717) shows explicitly how such a current would exclude a
      zero mode.
\item The causal critical-defect package (695)--(701): the compact
      Wiener--Hopf prehistory cannot cancel any critical Abel atom; the
      one-sided \(L^2\) estimate is an exact zero theorem; and, on the
      zero-mode space, causal Hardy closure, vanishing of the simple
      critical residue port and vanishing of the negative off-line port are
      equivalent.  Moreover the physical compensated current has trivial
      ordinary \(L^2\)-domain, ruling out a generic bounded-output synthesis.
      The only unproved arrow is now the explicitly isolated implication
      (700).
\item The final analytic hygiene (454)--(457): global
      $\zeta'/\zeta$ pole rigidity is valid, but it does not apply to the
      moving-cutoff leakage (452); the proposed Burnol annular-constancy
      step is exactly the missing gap and has been withdrawn.
\item The finite Green classification and type audit (458)--(467): compatible
      abstract Green matrices are classified, the rank-one fusion candidate
      is excluded, and the constructed diagonal normed line is proved to be
      a self-pairing rather than an additive mixed divisor class.
\item The canonical nested leakage construction (468)--(474): eliminating
      the safe finite block and then the whole complement produces the
      source-defined positive Gram $\mathscr L_{D\mid S}$ and the exact
      corrected certificate for the endpoint $T=\frac12\log5$.
\item The formal/geometric separation (475)--(477): a formal additive
      Picard lift always exists, but the row-(a) external divisor category
      factors through the two Tate degrees and collapses the entire primitive
      module.  Therefore the missing object must be a genuinely enlarged
      section/effectivity realization, not another formal line construction.
\item The innovation-cocycle theorem (478)--(480): canonical leakage Grams
      change under enlargement by an explicit positive Schur increment and
      therefore cannot be identified directly with the principal-compatible
      finite Green kernels.  Any geometric bridge must retain these
      increments as coherent transition data.
\item The innovation completion (481)--(483): accumulated source-defined
      Schur increments produce a positive principal-compatible kernel and a
      canonical Gram/Kolmogorov realization without factoring the desired
      row-(d) remainder.  Equation (484) then proves that this kernel has the
      wrong base normalization; repairing it requires inserting exactly the
      sought Schur remainder and is circular.
\item The complete Abel audit (485)--(494): the exact convergence abscissa,
      critical $1/\sigma$ pole-energy coefficient, Markov-reservoir
      divergence, symmetric/causal split and regularizer-invariance theorem.
\item The critical pole-defect measure (495)--(497) and the derived
      co-Poisson divisibility criterion (498)--(500).  These reduce the
      equality-state target from a global Hardy bound to vanishing of the
      leading local residue mass, or equivalently one additional order of
      divisibility by every critical zeta zero.
\item The Sylvester audit (502)--(504): an abstract Rellich decomposition
      always exists, but its residual is exactly the pole-defect compression
      to the first-loss kernel.  Hence spectral, semigroup and
      Moore--Penrose constructions of the multiplier are circular; a valid
      multiplier must come from an independent two-port flux identity.
\item The terminal audit (505)--(508): clamped primitive potentials do not
      annihilate the moving Tate terminal (an explicit even polynomial gives
      a strictly positive terminal square), while the nontrivial-zero residue
      boundary is a Krein sum with one hyperbolic plane for each off-line
      orbit.  Thus zero total residue balance does not imply vanishing of the
      critical positive residue mass.  Any positive terminal closure must
      independently eliminate those negative fibers; that is the remaining
      Hodge-index assertion, not a consequence of the existing
      normalizations.
\item The canonical negative port (509)--(513): its coordinates and norm are
      explicit, and (511) identifies each coordinate with the vertical flux
      of the logarithmic generator.  On a first-loss kernel conservation is
      exactly equality of its norm with the positive residue norm.  Hence the
      sole closure identity is \(\mathcal N_TP_{0,T}=0\); the current G39
      amplitude does not yet imply this radical-lifting statement.
\item The causal residue identification and gauge obstruction (514)--(519):
      the negative port is precisely the antisymmetric jump of the two
      off-line causal residues.  G39 fixes their sum, which is invariant
      under opposite residue shifts, while the jump changes.  Scalar
      conservation fixes only equality of positive and negative norms.
      Therefore a Hardy/passivity or geometric Hodge selection theorem is
      indispensable and cannot be recovered algebraically from the proved
      G39 data.
\item The physical gauge correction and vertical homotopy
      (520)--(527): finite-orbit interpolation remains onto after imposing
      support, parity and Tate moments, but the algebraic gauge is not
      automatically a zero mode.  Exact orbit isolation by a compact test is
      impossible.  The new homotopy formula (527) realizes the negative
      residue norm itself as a logarithmic bulk flux; completing the proof
      now requires identifying the globally summed bulk, including its Tate
      correction, with a vanishing G39 zero-mode amplitude.
\item The exact logarithmic primitive correction (528)--(531): multiplication
      by \(t\) and motion of the Tate exponent differ by the orientation
      matrix \(D=\operatorname{diag}(-1,1)\).  This exchanges the even/odd
      Tate metric eigenchannels and produces the explicit nonunitary ratios
      (531).  Hence the known moving-terminal square cannot be substituted
      for the normal boundary of the vertical flux; an additional oriented
      transport identity is required.
\item The conservative-transport no-go (532)--(534): the required oriented
      map amplifies every nonzero even terminal by the strict factor
      \(4(a+b)/(a-b)>4\), so it cannot be an isometry with \(J_T=I\).
      An abstract compensating square root would be circular, but the next
      theorem supplies an independent Gamma compensation.
\item The Gamma terminal-compensation theorem (535)--(546) corrects the last
      sentence of that no-go at the terminal level: although the bare
      transport is not isometric, the first independently positive Gamma
      cell \(h_0=16\tau^2/(1+4\tau^2)\) dominates its complete Gram defect
      with the uniform bound \(\mathcal D_T\le\frac34H_{\Gamma,0}\).
      Consequently \(G_{\Gamma,T}-\mathcal D_T\ge
      \frac14G_{\Gamma,T}\), with an explicit rank-one Riesz
      factorization.  No Euler channel is needed for this compensation.
      The remaining open issue is the source comparison between the summed
      vertical residue flux and the G39 zero-mode amplitude.
\item The continuous residue completion (547)--(551): the complete
      functional-equation orbit analysis is a bounded map from the Gamma
      form domain into one weighted \(\ell^2\) Hilbert space.  Hence the
      global negative port and all polarized Krein sums converge
      absolutely without RH.
\item The completed-G39 graph and factorization audit (552)--(554): raw
      finite Euler ports cannot converge in a topology with continuous
      residue traces, because they are entire while the completed
      logarithmic derivative has nonzero residues.  The correct limit is a
      finite-part bulk plus an explicit residue-defect coordinate.
      Factoring the resulting positive boundary Gram through \(H_T\) is
      equivalent on the first-loss kernel to the desired terminal
      vanishing, so a Sylvester or pseudoinverse construction is circular.
\item The scalar-virial no-go (555)--(558): parity forces the logarithmic
      commutator to have zero diagonal form on each parity sector, while
      the dilation virial changes sign because of the high-frequency Euler
      oscillations.  Hence neither scalar local virial can be the positive
      conservation source in (553).
\item The canonical two-contour shadow (559)--(561): the completed scalar
      scattering is an explicit selfadjoint unitary and its Wigner--Smith
      metric compresses exactly to \(A_T\).  Its metric is Krein, and
      nonnegativity of its physical graph is equivalent to G40.  Thus the
      full source-labelled Hilbert scattering remains open; a formal
      Hilbertization would assume the desired positivity.
\item The backward observability storage (562)--(564): before first contact
      the Lyapunov identity has a canonical positive semigroup solution.
      Along a lowest eigenvector every solution has diagonal coefficient
      \(\langle B_Tf_T,f_T\rangle/(2\lambda_T)\); Gamma positivity makes
      this diverge at a nonzero first contact.  Therefore a bounded
      source-labelled continuation across contact is itself a no-first-loss
      theorem, not a choice of normalization.
\item The energy quotient and cascade audit (568)--(570p): a bounded lift of
      the physical observation to the canonical \(H_T^{1/2}\)-energy space
      is exactly the form estimate \(B_T\le C_TH_T\), and the defect of any
      genuine contraction cascade telescopes into the sum of its reachable
      local defect squares.  This would exclude equality if the residual
      Gamma coordinate were retained.  However the displayed compressed
      Euler port has Gram \(I+C_{p,T}+C_{p,T}^*\) and is explicitly expansive
      on a positive source, so the existing prime identities are not yet
      such a Hilbert contraction cascade.  The missing object is the
      source-derived combined Gamma--Euler--Tate feedback/metric, not another
      local Julia dilation.
      Independently, the correctly oriented compressed Blaschke--Julia cell
      is proved completely observable in finitely many boundary passes in
      (570e)--(570g), and Möbius preconditioning gives the explicit
      contraction (570h)--(570i).  The physical Wigner--Smith prime factor
      has the opposite signed-delay orientation.  Equations (570j)--(570l)
      prove that its winding correction converts the positive one-sided
      Julia defect into a right-minus-left Krein boundary flux, which
      vanishes identically on parity sources.  Transferring strict
      observability without changing \(A_T\) therefore requires a genuinely
      joint Gamma--Tate source intertwiner.  The finite-dimensional
      countermodel (570m) additionally proves that injectivity of a Gamma
      \emph{input} feature cannot exclude angular equality unless that
      feature is first placed inside the angular defect by a source-derived
      factorization.  Finally (570n)--(570p) separate two notions that may
      not be interchanged: the compressed scattering norm defect and the
      stored Wigner--Smith delay.  The former can be completely observable
      while the latter remains sign-indefinite; \(A_T\) is the latter.
\item The term-by-term conservation package (571)--(694): the compensated
      Gamma observation is now an explicit parity/direct-sum operator, the
      negative residue observation has an explicit orbit kernel, and their
      common domain and polarized Gram are fixed.  The Gamma term alone is
      sufficient for closure.  Writing
      \(X_{\Gamma,T}^{\rm src}=I/2+Y_T\) cancels the entire unbounded Gamma
      multiplier and reduces the missing identity to the bounded equation
      (585).  Gamma normalization gives the order-zero design equation
      (590), whose data are precisely the normalized constant-plus-Euler
      delay \(K_T\) and oriented Tate terminal \(D_T^\natural\).
      Equations (591)--(597) separate the required source connection into
      its Hermitian completed-metric part and skew Kato part.  The skew
      transport has zero diagonal in every \(K_T\)-eigenspace; all
      first-contact exclusion is therefore carried by the Hermitian
      completed Euler--Gamma metric connection.  Constructing that
      connection without the displayed forbidden Sylvester quotient or
      inverse completed metric remains \OPEN{}.  Equations (598)--(600)
      construct its honest half-plane logarithmic connection and prove that
      critical continuation necessarily adds the completed finite-part plus
      residue graph; the antisymmetric residue coordinate is exactly
      \(\mathcal N_T\).  Bounded Hilbert pullback and the product rule (580)
      remain the final construction.  Equations (601)--(605) give an exact
      Hadamard finite-part/residue graph at every truncation.  Equations
      (606)--(610) then construct explicit Paley--Wiener cardinal sources
      and a finite-rank \(X_{T,R}^{\rm res}\), without any operator inverse,
      satisfying the residue conservation law exactly.  The full-form tail
      is (611); uniform interpolation and tail control (612), equivalent to
      the stable residue-range lift, cannot hold for the isolated fiberwise
      lift by the no-go (613).  Equations (614)--(617) source-realize the
      necessary joint Hadamard bulk by Volterra operators.  Equations
      (618)--(621) give a uniformly bounded conservation operator for each
      individual off-line orbit, while (622) proves that distinct orbit
      solutions have unavoidable cross terms.  Primitive compression
      produces exactly the rank-two Tate cross terminal (623).  Joint
      Hadamard/Gamma feedback, rather than a sum of isolated fibers, remains
      \OPEN{}.  The universal neutral extension (624)--(627) solves the
      conservation algebra on the ambient Krein boundary.  Its physical
      pullback is exactly the \(J\)-skew tangency problem (628)--(631);
      (633)--(638) prove that a residue-to-Gamma block is compulsory.
      The reverse sampling theorem (642)--(643) proves unconditionally that
      the complete residue analysis controls both \(L^2\) and Gamma energy,
      so this mandatory synthesis exists boundedly.  Finally
      (644)--(649) identify the extended physical boundary as a closed
      Gamma graph over the residue range and reduce all remaining blocks to
      the mixed same-window graph tangency identity.  The stronger full
      residue-involution lift is proved equivalent to RH and withdrawn as
      an intermediate target.  Finally (650)--(652) identify the graph
      operator as selfadjoint Fredholm and its observation as strictly
      positive: the conservation equation has a bounded solution exactly
      when the first-contact kernel is absent.  Its co-Poisson/Hardy source
      realization, without the forbidden inverse (652), remains \OPEN{}.
      Equations (653)--(656) additionally construct the continuous
      horizontal-opening current of every off-line zero orbit: multiplication
      by \(t\), followed by the orbit grading, integrates exactly to the
      negative residue square.  They also prove that its descent is precisely
      a sampling-range tangency statement, so ambient differentiation alone
      cannot be used as a source lift.  Equations (657)--(658) prove that in
      a full neighbourhood of the physical boundary the horizontal family
      remains bounded below, has closed sampling range and has bounded
      velocity; thus no local convergence or norm-domain gate remains.
      Finally (659)--(662) eliminate the
      negative port affinely on the closed range and reduce the terminal
      equation to
      \(\mathbb H_TY_T+Y_T^*\mathbb H_T=
      \frac12I+\Gamma_T^*\Gamma_T\).
      The right side has a fixed positive gap, and every such connection must
      blow up at least as \(1/(4|\lambda_T|)\) near first contact.  Hence no
      residue normalization remains to be found: the bounded source-derived
      metric connection itself is the final open theorem.  Equations
      (663)--(666) give a further exact source construction: the complete
      nonconstant Euler delay is the commutator of multiplication by \(t\)
      with one finite skew logarithmic Euler phase.  Primitive compression
      leaves only the explicit rank-four remainder (665).  The
      two-component formula (668) compares it with the oriented Tate defect;
      (669) is the exact equality that would identify them, while the central
      compact-source test (670) proves that equality false for the Euler
      phase alone.  Thus the Gamma/Hadamard finite part must enter before
      primitive compression; it cannot be appended after identifying the
      Euler boundary with \(\mathsf D_T\).  Equations (671)--(673) now
      perform that ordered assembly: the Gamma phase, affine \(m_0\) phase
      and both Euler orientations form one skew source generator whose
      commutator with multiplication by \(t\) is exactly \(A_T\).
      Primitive compression again leaves only a rank-four terminal.
      Equation (674) is the explicit two-component Gamma--Euler--Tate
      terminal identity which would identify the raw phase remainder with
      \(\mathsf D_T\).  It is a candidate test, not an asserted
      normalization; even if it held, the positive product rule (661),
      rather than only the commutator rule (672), would remain necessary for
      closure.  The Gamma cell
      phases and their source kernels are resolved explicitly in
      (675)--(680), while (682) removes the affine terminal on primitive
      sources.  The resulting candidate equality is (683).  Equations
      (684)--(687) then decide it negatively: the Gamma terminal has a
      boundary-pole derivative and cannot factor through the
      four-dimensional Tate jet space.  Thus (674)/(683) is false for the
      raw phase.  The Hadamard residue/finite-part graph must contribute a
      genuinely new boundary state; it cannot merely regularize the same
      two Tate coordinates.  Finally (688)--(691) complete the skew phase
      by its Abel exterior tails and prove that its full primitive Tate
      terminal is exactly zero: the boundary pole in (685) cancels between
      interior and exterior.  Hence neither the compressed skew phase nor
      its conservative full-line completion supplies the positive terminal.
      They can contribute only the skew/Kato part; the missing diagonal
      term must come from a genuinely Hermitian finite-part/residue metric
      state.  The diagonal test (692)--(694) makes this precise on the
      physical residue range: any realization which remains skew has zero
      diagonal in every \(\mathbb H_T\)-eigenvector, whereas the required
      metric has diagonal at least \(1/2\).  A valid Hadamard pullback must
      therefore create a nonzero Hermitian component and satisfy (694)
      without dividing by the eigenvalue.
\item The algebraic G39 divergence realization (27e)--(27h) in the finite
      divisibility model.
\item The restricted CCM polar isometry (52)--(55).
\item The unitary midpoint identity (44)--(47) as an abstract Hilbert-space
      theorem.
\item The typed observation criterion (57)--(61).
\item The order-mismatch no-go (62)--(63), excluding the row-(c) operator
      $K=\mathcal F^{-1}X\mathcal F$ as the Hodge energy feature.
\item The oriented-delay Kolmogorov factor and conservation identity
      (67)--(76) for its own positive reference Gram.
\item The exact Gamma delay formula (81)--(82) and the authoritative balanced
      factorization (83)--(84), including explicit factors
      $\mathcal X_T,\mathcal Y_T$.
\item The signed old-factor correction (41a)--(41c), the exact
      reference-harmonic capacity formula (41d)--(41i), and its equivalence
      (41j) with the return load in the supplied row (d).
\item Old compatibility by neutral birth channels (41s)--(41u): enlarging
      the window restricts to the previous physical factors plus unitary
      channels whose two signed Grams cancel exactly.  The finite-port
      joint-Gram obstruction (41p)--(41r) is thereby separated from, and is
      not required for, physical continuum old compatibility.
\item The physical minimal amplitude current (41k)--(41m) and the complete
      joint-Gram criterion (41n)--(41o) for comparison with G39.
\item The finite-port/full-Hodge dimension obstruction: algebraic G39 cannot
      factor the infinite-dimensional physical old block.
\item The equivalence between the remaining sign and the unit Douglas
      contraction $C_T$.
\item The canonical source angular operator
      $\mathfrak C_T^{\rm src}$ and its exact defect pullback
      (84a)--(84b), constructed without a sign assumption.
\item The correctly typed observation theorem (86)--(87): normalized Sonin
      transport is isometric, whereas the physical observation must be
      contractive; an isometric observation would force $A_T=0$.
\item The equivalence of the all-window passivity estimate (84c), the Weil
      sign criterion, and RH.
\item The no-channelwise-contraction theorem (89): any valid $C_T$ must mix
      places and/or Gamma--Tate terminal channels.
\item The Tate-potential isomorphism and exact continuous kernel
      (90)--(94).
\item The differentiated global Poisson identity (96)--(98), the Gram
      homogeneity obstruction (100), and the exact Zeta metric variation
      (101)--(103).
\item The exact Tate boundary velocity and terminal energy (105)--(111).
\item The Kato primitive transport (112)--(116), the completed
      Euler--Gamma logarithmic metric (117)--(119), and the three-lines
      inverse-metric obstruction (120).
\item The exact residue Krein boundary and its fiber signatures
      (121)--(126), and the amenable unitarization theorem (127)--(128).
\item The critical gcd boundary state and explicit product-Poisson GNS
      realization (133)--(137).
\item The gcd--Gamma reference observation, global oriented midpoint
      colligation and exact source Gram (138)--(149).
\item M\"obius whitening, the centered Euler Fisher score, and the exact
      prime-power innovations (150)--(167).
\item The Archimedean Gamma GNS boundary (168)--(173), and the exact
      Gamma ratio--scale identity (174)--(178a), which derives both
      $g_\Gamma$ and $m_0$ from one canonical Mellin probability law.
\item The Gamma deconvolution domain no-go (179)--(180): neither compact
      support nor the two Tate moments permit removal of the Gamma
      characteristic factor as an $L^2$ operator.
\item The coherent completed Euler--Gamma factorization (181)--(183) in the
      absolutely convergent half-plane.
\item The Euler coherent-orbit escape theorem (184)--(186), separating the
      continuous scaling limit from the multiplicative gcd GNS limit.
\item The positive rigged critical Euler form, exact arithmetic diagonal
      observation, and joint Euler--Gamma boundary (187)--(194).
\item The divisor-incidence realization and deformed GNS equivalence
      (195)--(200), identifying the finite Zeta gauge, harmonic-mass
      renormalization, gcd boundary and M\"obius whitening in one model.
\item The linear/affine Gamma separation (202)--(204), identifying
      $g_\Gamma$ as an unbounded affine cocycle energy rather than a bounded
      coherent-orbit distance.
\item The prime Julia colligations and their exact M\"obius-innovation
      cascade (205)--(207), giving a lossless realization of the arithmetic
      label network.
\item The exact affine Gamma--Tate termination identity (208)--(210), deriving
      $g_\Gamma-m_0$ as the logarithmic derivative of one positive coherent
      metric cell.
\item The complete coherent Euler--Gamma bulk normalization (211)--(213),
      recovering the authoritative localized symbol $a_T$ with every prime
      power and the exact constant.
\item The beta--Gamma and oppositely oriented Euler Markov transports
      (214)--(220), and the Bayes-reversed joint contraction
      (220a)--(220e).  The latter removes the arrow mismatch and gives an
      explicit positive conditional-variance defect independently of G;
      its abstract Julia dilation is canonical.
\item The orthogonal joint innovation Gram (220f)--(220i) and the honest
      finite-window critical boundary (220k): at fixed $T$ the ambient Bayes
      defect is an ordinary positive Hilbert object already at $t=1$.
      The Bayes pullback criterion (220j) isolates the remaining physical
      whitening/cutoff/Tate intertwining without defining a square from G.
\item The prime-power Markov defect rates (221)--(223), deriving exactly
      every coefficient $\Lambda(p^k)/p^{k/2}$ from positive Julia defects.
\item The midpoint Julia two-port rotation (223a)--(223e): its exact lossless
      Gram and the real orthogonality of $J_-,J_+$ eliminate spurious
      $\sqrt\varepsilon$ terms.  Equation (223f) verifies the local defect
      coefficient only.  The audit following it records that $J_+$ is already
      an input of this rotation, so the rotation does not by itself construct
      the missing Bayes output.
\item Zeroth-order scattering necessity (223g)--(223k) and the
      inversion/reflection no-go (223l)--(223m): a strongly identity-limiting
      Markov transport cannot change $J_-$ into $J_+$, while the natural gcd
      involution only separates source parity.  Hence the missing scattering
      must mix prime ports with Gamma/Tate/old channels.
\item The complete Markov-defect realization (224)--(227) of the physical
      reference factor $\mathcal X_T$.
\item The Gamma ODE Green identity, rank-one shift commutator, finite-head
      absorption of $m_0$, and polar-free balanced factorization
      (229)--(238).
\item Gamma common-channel cancellation (238a)--(238d): both the positive
      Gamma defect and the affine $m_0$ channel contain the identical
      $D_H$ port.  Cancelling it leaves exactly the physical input tail
      $D_{\rm tail}$ and output head $C_H$, without factoring
      $G_\Gamma-m_0I$ or assuming its sign.
\item The explicit four-state Gamma head and the compressed Euler--Julia
      resolvent/factorization (239)--(246), resumming every active prime power
      into one filter and one positive boundary port per prime.
\item The finite-head physical angular operator, exact defect law, and
      certified strict passivity on $0<T\le\log2$ (247)--(250).
\item The first-loss zero-mode theorem (251)--(253), reducing any global
      failure to a finite-window primitive solution of an explicit
      Gamma--Euler equation with only two Tate forcing terms.
\item The parity and correctly typed homogeneous exterior-leakage reductions
      (254)--(257), identifying the equality case with a clamped potential
      whose locally finite global differentiated Euler--Gamma current
      vanishes throughout its support.
\item The Fourier-multiplier withdrawal (258), which separates the valid
      map $\mathcal E'\to\mathcal D'$ from an unproved tempered multiplier.
\item Gamma exterior antilocality (261)--(264), valid for compactly supported
      distributions, and the exact infinite Gamma--arithmetic local system
      (265)--(268).
\item The completed logarithmic-derivative collapse (269)--(271): bare
      meromorphic symmetrization is identically zero and loses the contour
      jump carrier.
\item Support saturation (273): every first-loss mode reaches both endpoints
      in essential-support hull.
\item Source identification (279a), equating that jump with the locally
      finite Gamma--arithmetic explicit-form operator after the two Tate
      characters vanish.
\item The finite-energy closure criterion (279)--(280), reducing G to the
      $L^2$ exterior leakage of a first-loss mode.
\item Automatic exterior Gamma control (281)--(282) by the Carleman
      inequality, and the exact Hardy-space formulation (283)--(287) of the
      two remaining oriented arithmetic tails.
\item Exact cancellation of the prime mean by the two Tate moments
      (288)--(291), leaving only the Chebyshev fluctuation, and Hardy pole
      rigidity (292).
\item The zero-mode Wiener--Hopf identity (293)--(298), showing that G39
      fixes the reflected symmetric part while G40 requires its bounded
      causal separation.
\item The equality-mode Poisson-lift criterion (299)--(303): existence of a
      lift only on the finite-dimensional zero-mode space suffices to prove
      finite exterior energy and close G; an all-source metric isometry is
      not necessary for this contradiction route.
\item The exact smooth-source range criterion (304)--(306): the Poisson lift
      exists precisely when the two source half-lifts agree.  This removes
      both spectral division and any unspecified range projection.
\item Compact Poisson-range rigidity (306b), the exact defect transport
      (310)--(314) with its critical-conjugation signs, and causal/Laplace
      rigidity (315)--(318).
\item The common Mellin-germ theorem (319)--(320): both Poisson half-lifts
      continue to $G(s)/\zeta(s)$, which is regular at $s=1$ and vanishes
      there at least doubly for a primitive source.  Thus no additional
      scalar pole cancellation is missing at the Tate character.
\item The differentiated meromorphic observation theorem (321): the
      Poisson current and the two Euler orientations have the same common
      germ.  The physical boundary carrier $\mathfrak B_T$ in (322) is
      therefore the exact, zero-symbol object whose vanishing (323), together
      with realization of (319), would close the zero-mode route.
\item The parity--Fourier reduction (324)--(326): for a parity-$\varepsilon$
      mode the whole boundary carrier is
      $-\varepsilon\mathscr JZ\partial d_\varepsilon$, where
      $\mathcal Fd_\varepsilon=-\varepsilon d_\varepsilon$.
\item The exact nuclear Poisson polar formula (327), and finite-pole
      extended-S rigidity (328)--(329).  Thus it is enough to obtain finite
      nontrivial polar data for the opposite Fourier component; zero-density
      rigidity then excludes the compact source.
\item The exact Gamma--Tate bi-gap reformulation (330)--(333): a clamped
      equality mode has transform
      $H_h(s)=s^2(s-1)^2\Phi(s)$ and its source jump is supported in the
      complementary logarithmic window.  Vanishing of the resulting
      jump--Sonin space for every $T$ is equivalent to the desired uniqueness
      theorem and hence sufficient for G.
\item The coordinate audit (334): the logarithmic observation window becomes
      a multiplicative annulus, not a neighbourhood of the additive origin,
      so classical Sonin membership cannot be inferred from the local carrier
      gap without a new propagation theorem.
\item Translation-jet exclusion (335)--(340): a first-contact kernel vector
      cannot be a compactly supported smooth vector of translations.  Global
      translation invariance would put every derivative in the same
      finite-dimensional kernel and force a compactly supported exponential
      polynomial, hence zero.
\item The prime-oriented collar reduction (342)--(345): on each boundary
      strip of width $\log2$, every outward prime orientation vanishes and
      the zero mode obeys a one-way Gamma--renewal equation.  Hence a
      translation-jet completion need only control the two reflected boundary
      layers.
\item The panoramic negative certificate (346)--(350): failure of G would
      require, simultaneously, support saturation, loss of every admissible
      translation jet, failure of the Meyer equality-mode lift, and a
      non-$L^2$ oriented arithmetic tail, while the exterior Gamma tail
      remains in $L^2$.  The associated Hardy obstruction is either a
      nonremovable zeta-zero pole or a pole-free failure of the uniform Hardy
      norm.
\item The endpoint-shape audit: the lowest primitive eigenvalue is
      nonincreasing by zero-extension and min--max, so a Hadamard/Riccati sign
      alone cannot exclude first contact.  Shape differentiation has been
      removed as an independent closure route.
\item Reciprocal Sonin gap amplification (351): for $0<a<1$, Burnol's
      co-Poisson subspace $P_a$ has zero intersection with the reciprocal
      extended Sonin space $L_{a^{-1}}$.  Co-Poisson orthogonality at all zeta
      zeros and unconditional completeness on the $a^{-1}>1$ side prove the
      assertion without RH.  Criterion (357)--(358) turns this into a new
      equality-state route to G.
\item The physical co-Poisson domain audit (431)--(435): Mellin--Plancherel
      constructs the maximal closed global observation and identifies
      full-form boundedness with the explicit weighted Paley--Wiener estimate
      (434); such boundedness would already force the uniform local zeta
      second-moment estimate (435).
\item The domain-free local observation (436)--(440): on the physical window
      it is $I+K_{T,+}$, with $K_{T,+}$ nilpotent because every nontrivial
      arithmetic term has delay at least $\log2$.  Its inverse is the finite
      Neumann polynomial (439).  Therefore annular vanishing alone forces
      $F=0$, with no global co-Poisson norm or completeness theorem.
\item The exact local Euler commutator (442)--(445): Möbius inversion gives
      \(B_{T,+}=\mathcal Z_{T,+}^{-1}[Q,\mathcal Z_{T,+}]\), and parity
      reconstructs both arithmetic orientations from the single observed
      amplitude \(h\), with cancellation of the two bare \(QF\) terms.
      Equation (445) is the fully explicit observed Gamma--Tate transport
      equation.
\item The meromorphic Gamma--co-Poisson commutator (447)--(450): differentiated
      co-Poisson summation and the cosine Mellin factor give
      \(-\chi'/\chi=g_\Gamma-m_0\), hence
      \(\mathbf Z_+\mathcal A_\infty=[\mathbf Z_+,D_\Gamma]\) as a common
      regularized Mellin germ.  Its physical localization is separately
      marked \AUDIT{} in (451)--(453).
\item The bare-observation generator audit (359)--(360): the Mellin
      derivative of $\zeta G_F$ is not the reflected logarithmic
      Euler--Gamma generator in the zero-mode equation.  Thus reciprocal gap
      amplification cannot be obtained by formally integrating the scalar
      logarithmic derivative; the completed two-port orbit is essential.
\item Prime-tower Blaschke--Clark scattering (361)--(366): each complete
      tower $p^k$ is the nonconstant part of one positive Poisson/Clark phase
      density, with its deterministic winding retained.  Summation gives the
      exact finite Clark embedding form of row (d).
\item The global finite-window Gamma--Euler scattering (367)--(371): one
      explicit unimodular boundary function $\Theta_T$ has compressed phase
      derivative exactly $A_T$.  Equation (372) is the resulting
      single-scattering formulation of G40 on the primitive source.
\item The Wigner--Smith identity (373)--(374), the one-state Julia
      realization (375)--(376), and the winding-corrected prime birth factors
      (377)--(379).  Each prime tower is a lossless cascade and its compressed
      signed Gram is exactly neutral at its threshold, proving scattering-level
      old compatibility.
\item The Gamma all-pass cascade (380)--(383): $B_\Gamma$ is the locally
      convergent product of explicit winding-corrected rational all-pass
      cells, and their nonnegative stored delays sum exactly to $g_\Gamma$.
      Hence every factor of $\Theta_T$ now has an explicit conservative
      realization.
\item The scalar-inner no-go (384): $\Theta_T$ has unavoidable upper
      half-plane poles and an eventually positive phase derivative
      incompatible with lower-half-plane inner orientation.  Therefore the
      final passive realization must remain two-port/operator-valued or
      Krein/Hodge before taking its scalar supertrace.
\item The operator-valued Clark colligation (385)--(392): explicit positive
      Gamma/winding and prime/constant ports have signed Gram exactly $A_T$,
      a canonical angular rule, and neutral old-compatible prime birth.
\item The Clark/physical port audit (393)--(394): both Clark Grams differ
      from the authoritative $J_-/J_+$ Grams by the same explicit form, so
      their defects agree, but that common correction is not positive for
      all windows on the full source.  A full-source common orthogonal
      Hilbert augmentation is therefore unavailable.
\item The explicit Krein--Hodge stabilization (395)--(400): a
      source-defined signed state realizes the common correction without
      spectral square roots, makes both stabilized generated ranges positive,
      and constructs isometric physical--Clark intertwiners.  Equation (401)
      records exactly why this solves type/conservation but not domination.
\item The first-loss radical-lifting reduction (402)--(405): after applying
      $L=\partial_t^2-1/4$ to remove exactly the two Tate directions, every
      hypothetical first mode has a differentiated Gamma--Euler output
      supported entirely outside its window.  Global vanishing of that
      output forces the source to vanish because the finite-window
      scattering delay is real analytic and nonzero almost everywhere.
      The hyperbolic-plane audit (406) proves why isotropy and the scalar
      identities alone do not supply the required lifting.
\item The scalar Tate-factor audits (407)--(409): the delay remains strictly
      negative at the origin, and the clamped polynomial cancels only the
      central pair among infinitely many off-axis prime--Clark poles.  Thus
      neither pointwise positivity nor a finite-index scalar
      Krein--Langer correction can supply the radical lifting.
\item The global exterior-reservoir theorem (410)--(413): independent prime
      Julia coordinates have infinite exterior Hilbert energy for every
      nonzero source.  Therefore prime synthesis and exact Tate centering
      must occur before completion; the only admissible exterior arithmetic
      state is the compensated Chebyshev current $dR_\Psi$ and its reflected
      orientation.
\item The exterior jump/parity theorem (414)--(417): the exact term
      preventing a naive symmetric Gamma closure is the unsupertraced
      $\xi$-jump.  For a parity first-loss state the two arithmetic tails
      are reflections, so both are non-$L^2$ if the state is nonzero and a
      single one-sided compensated-current estimate is sufficient to close
      G40.
\item The one-sided Riemann--Hilbert decomposition (419)--(424): the
      functional equation and the zero-mode identity determine an $L^2$
      reflected symmetric boundary.  Every zeta-zero residue cancels in
      that symmetric part and doubles in the antisymmetric part.  Hence the
      new Hodge observation has one precise job: control the antisymmetric
      two-contour datum in the right Hardy space.
\item The Smirnov-selection audit (424a): meromorphic bounded type,
      physical decay, parity and an \(L^2\) reflected symmetric boundary
      still admit explicit rational gauges carrying a causal pole.  Thus
      Smirnov membership alone cannot select the physical orientation.
\item The residue-range criterion (523a) and real-structure audit (523b):
      lifting \(J_{\rm res}\mathcal E_TF\) into the closed physical residue
      range would make the polarized zero-mode identity force
      \(\mathcal E_TF=0\), hence \(F=0\).  Reality and parity do not supply
      this lift because they leave an arbitrary complex phase on every
      off-line fiber.
\item The minimal equality-state lift and its exact distance audit
      (702)--(707): physical range membership of the negative component
      \(P_-\mathcal E_TF\) alone would force its norm to vanish.  However,
      for a hypothetical normalized zero mode that component lies at
      distance exactly \(1/2\) from the physical range.  Hence finite
      interpolation cannot manufacture the lift; (703) is equivalent to
      the terminal implication (700), not a softer approximation theorem.
\end{itemize}

\subsection*{OPEN}
\begin{itemize}
\item Construct the joint port embedding
      \(\mathscr P_{\sigma,T}\) with its Hilbert metric fixed independently
      of row (d), and prove the critical norm derivative (754).  The
      amplitude, endpoint, residue-flux and moving-Tate calculations are
      complete, but (752)--(753) prove that none fixes this norm derivative.
      Equation (754) is the common remaining content of (220j), (661), and
      the Ward normal block (716).
\item Construct the ambient Hadamard--Gamma--prime current
      \(A_T^{\rm amb}\) satisfying both \(J_{\rm res}\)-skewness (710) and
      the exact normal-block identity (716).  This is now the minimal
      constructive form of (661).  It must be obtained before compression;
      the candidate \(\mathcal W(K_P+K_\Gamma)\mathcal W^*\) is ruled out by
      (714)--(715) because it has no normal component.  Equivalently,
      construct the ambient skew phase in (721) and prove the fully
      ambient weak identity (722); testing only physical source vectors is
      insufficient.  By (725)--(727), this task is equivalent to (661), so
      it cannot be discharged by abstract skew extension alone.
\item Construct, from the joint Hadamard--Volterra/Gamma--Tate source
      system, a uniformly bounded realization of the metric connection
      (661):
      \[
      \mathbb H_TY_T+Y_T^*\mathbb H_T
      =\frac12I+\Gamma_T^*\Gamma_T.
      \]
      Equivalently, prove that the graded horizontal velocity in (656),
      augmented by the Gamma finite part and Tate correction, is tangent to
      the closed physical sampling range and pulls back to one bounded source
      operator.  The connection may not be defined with
      \(\mathbb H_T^{-1}\), a Sylvester quotient, or a pre-contact semigroup;
      (662) proves that any such construction whose bound contains the
lowest spectral gap is circular for first-contact exclusion.  The
Euler bulk of this construction must use the exact commutator
      (663)--(666), but its primitive rank-four remainder may not be set equal
      to \(\mathsf D_T\): (667)--(670) disprove that Euler-only terminal
      identification.  The required tangency must be verified after adding
      the Gamma/Hadamard finite-part current.  The raw-phase terminal
      identity (674) is ruled out by (684)--(687).  Concretely, use the
      joint graph (601)--(617) to construct a Hermitian boundary state.
      Merely adding the conservative skew exterior tail is insufficient:
      (688)--(691) prove that it cancels the non-Tate pole but leaves zero
      terminal.  The new state must instead retain the positive
      finite-part/residue energy and satisfy the product rule (661), not only
      the already proved commutator rule (672).
\item Prove the single critical-port implication (700).  By
      (695)--(699), this is equivalent on the zero-mode space both to the
      causal estimate (417) and to \(\mathcal N_TF=0\).  A source-defined
      two-port Rellich--Clark multiplier \(X_{\sigma,J}\) in (501), or the
      weaker residue estimate (497), would prove it; a pseudoinverse
      construction would be circular.  The Krein balance supplies only the
      non-strict bound (699), so the negative off-line port may not simply be
      discarded.
\item At $T=\frac12\log5$, compute the two directed interval enclosures
      (473)--(474), including the previously omitted
      $S_{163}$--to--whole-complement action $A_{SQ}$.  The printed
      five-column data do not determine these quantities.  Even a successful
      endpoint certificate would extend the unconditional interval but would
      not prove the all-window criterion.
\item Construct the faithful mixed section/effectivity realization isolated
      in (475)--(477).  The existing quadratic line (463) cannot serve as an
      additive class, while the tautological formal Picard lift carries no
      geometry and the row-(a) external category collapses primitive classes.
      Only an independently constructed realization and comparison can
      identify the finite geometric kernel (460) with an analytic leakage
      system carrying the innovation cocycle (478)--(480).
\item Construction of a positive operator-valued observation for the
      antisymmetric Riemann--Hilbert datum
      $\mathscr A_+-\mathscr A_-$ in (424), with a source-derived norm that
      places the right causal component in $H^2$.  The symmetric component
      is already $L^2$ by (421); controlling it again would not address G40.
\item Establish (703) only as a consequence of a direct source product
      rule such as (661), not by compact-window interpolation.
      Equations (704)--(707) prove that for a hypothetical nonzero zero mode
      the target \(P_-\mathcal E_TF\) has fixed positive distance from the
      physical range, so finite right inverses cannot converge to it.
      The meromorphic co-Poisson functional equation does not control the
      incoming history (452), and reciprocal-gap completeness does not
      supply the missing product rule.
\item The equality-only causal radical-lifting implication (405): prove that
      the differentiated exterior-only output of a first-loss state must
      vanish.  This is the narrowest negative form of the remaining G40
      problem; applying $L$ is mandatory and the interior projection may not
      be dropped.  It is equivalent to the first-loss exterior Hardy/Sonin
      completion, not to another scalar normalization.
\item The jump-adapted Gamma--Tate bi-gap theorem
      $\mathscr S_T^\xi=\{0\}$ in (333), equivalently the uniqueness
      implication (260).  A proof must retain the two contour/Hardy
      orientations; scalar meromorphic symmetrization cancels the jump.
\item Construction of a faithful mixed-effectivity/causal completion
      functor which rules out the negative certificate (346)--(350) before
      its character is compared with row (c).  Analytically this may be the
      Meyer lift (349) or the Hardy estimates (287); geometrically it may be
      the extension of periodic Riemann--Roch and effectivity to the completed
      mixed classes $D_f$.  Its norm may not be defined from
      $-B_{\rm nuc}$ or $\Delta_T^{1/2}$.
\item The local equality-state uniqueness/gap-transport identity
      (441)/(446).  The
      observation is now defined for every physical $L^2$ state by (436), and
      its logarithmic restriction to $I_T$ has the finite nilpotent inverse
      (437)--(439).  Hence neither the full-form estimate (434), global
      $L^2$ membership, nor Burnol completeness is required for closure.
      What remains is exactly to construct the physical causal boundary of
      the meromorphic commutator (447), control the incoming history (452),
      and derive annular vanishing of
      $\mathfrak S_T^{\rm loc}F$ from $P_TA_TF=0$, retaining the Gamma and
      both arithmetic orientations.  Once this local identity is proved,
      (440) excludes the first-loss state immediately.
\item Proof that the explicit scattering $\Theta_T$ in (369) has a
      Poisson--Sonin/Hodge realization whose Clark form is nonnegative after
      removal of the two Tate evaluation directions.  Boundary
      unimodularity is proved, but a Schur/inner or generalized-inner class
      with the required primitive index has not been established.
      Equivalently, prove that the positive stored delay (383), together with
      the two moving Tate terminal energies (109), dominates the negative
      winding-corrected prime delays (378) on $\mathcal P_T$.
      The realization cannot be a scalar half-plane inner factor by (384);
      it must use the oppositely oriented Bayes Euler and Gamma ports already
      constructed in (214)--(220j).
\item Proof that the now explicit stabilized angular map (400) is
      contractive on the primitive generated range.  The typed
      physical--Clark intertwiner and the common Krein--Hodge state are
      constructed in (395)--(400); (401) proves that their contractivity is
      still the substantive inequality and may not be inferred from separate
      positivity of the input/output Grams.
\item The boundary regularity bootstrap (341), which would upgrade a
      first-loss solution of the logarithmically elliptic Gamma--arithmetic
      equation to the translation-jet class (335).  Powers of logarithmic
      regularity alone do not imply this upgrade.
\item Positive matching of the two one-way collar problems (343)--(344) in
      the affine Tate/prime-Julia storage norm.  This is the boundary-layer
      form of the still missing G40 Hilbert defect.
\item Source interconnection of the finite G39 cross-amplitude port with the
      infinite Gamma--Poisson continuum colligation without adjoining its
      finite old Gram orthogonally.  Direct identification through
      (41p)--(41r) is ruled out by rank and is no longer listed as the
      compatibility target.
\item Identification of the new source-derived gcd--Gamma reference space
      (138)--(149) with the independently defined CCM Sonin target, if that
      particular target is required; the source Gram/type isometry itself is
      now constructed.
\item The general-window Sonin lower-frame estimate
      \texttt{eq:framebound} from the supplied \texttt{main.tex}; its known
      comparison theorem covers only the prime-free window.
\item Construction of the normalized physical Poisson--Sonin Gram in (99)
      and proof of the boundary monotonicity/square identity (104).
\item Proof that the clamped-potential observation is the conditional
      expectation/martingale of the joint Gamma--Euler product boundary,
      while retaining the Gamma characteristic factor in the Hilbert metric;
      the direct deconvolution route is ruled out by (180).
      More precisely after (214)--(223): prove that the Julia-reversed Euler
      transport and beta--Gamma transport, amplified by the physical delay
      channels and Kato-compressed by $P_t$, have main output
      $\mathcal Y_TP_T$ and defect equal to (109) plus the old-compatible
      Schur residual.
\item Construction of the windowed intertwiner from the authoritative
      clamped-potential/old-new source into the positive joint critical space
      (194), proving its Gram is the Schur boundary form and making it
      compatible with the gcd GNS limit and moving Tate projection.
\item Construction of the Hilbert completion of the Poisson quotient and the
      uniform centered-scaling estimate (127), or an equivalent critical
      closed-range theorem in a source RKHS realization.
\item Construction, from source formulas, of the contraction
      $C_T\mathcal X_TP_T=\mathcal Y_TP_T$ with $\|C_T\|\le1$; equivalently,
      the unit-budget factorization of $Q_c$ through $D_0^{1/2}$.
      Equivalently, prove the global bound (84c) for the already constructed
      source angular operator.
      In the finite-head realization it is enough, by (251)--(253), to exclude
      a nonzero first-loss zero mode and prove the corresponding equality case
      of the ODE/Julia storage network is trivial.  Equivalently, prove the
      exterior-leakage uniqueness statement (257) for clamped potentials.
      After (261)--(273), the irreducible form is the arithmetic separation
      lemma: the coupled cancellation (268) must force simultaneous exterior
      vanishing of a Gamma input/output pair, or directly force all exterior
      Gamma moments to vanish, despite support saturation.  Equivalently,
      after (274)--(287), it suffices to prove the finite exterior-energy
      estimate (280), i.e. the two uniform Hardy bounds (287).
\item Construction and strictness of the unsupertraced jump-carrier
      observation (272).  Its two Poisson orientations must be retained
      before supercharacter cancellation and coupled to the Gamma resolvent
      ladder and the prime Julia boundary storages.
      Equivalently, construct the equality-mode lift (303) from the
      Wiener--Hopf equation (298), or prove the two Hardy estimates (287).
\item Extension of the half-lift criterion (304)--(306) from Meyer's nuclear
      smooth domain to the physical logarithmic form domain of a first-loss
      mode, together with the zero-mode identity (309) in that completion.
      In the defect formulation the exact remaining requirements are the
      Gamma--Tate observation identity (314) on the zero-mode space and a
      boundary-realization theorem identifying the physical defect transform
      with the already proved common meromorphic germ (319).  The latter is
      a topology/growth statement, not a further scalar zeta identity.
      Equivalently, prove the single boundary-carrier theorem (323) in a
      completion where its zero germ (321) determines the physical local
      representative and where the transform of the associated defect has
      the connected-domain realization required by (317).
      After the parity reduction, it is enough to construct finite
      extended-S polar data for $d_\varepsilon$ in (325), or prove directly
      that $P_{I_T}\mathscr JZ\partial d_\varepsilon=0$ plus the physical
      boundary conditions force $d_\varepsilon=0$.
\item Identification of the undefined symbol $\mathcal H_{5/4}$ with a
      precise Gamma form, if that notation is to be retained.
\item The positive factorization (43), hence G40, G41, G42, condition $G$, and
      row (d).
\end{itemize}

\subsection*{AUDIT}
\begin{itemize}
\item Equations (48)--(51) are only a candidate application of the midpoint
      theorem and are not a proof of G40.
\item The proposed shortcut
      $\langle(\mathcal U_{\rm PF}-I)Kf,x\rangle=0$ for every old primitive
      $x$ is unavailable: no supplied theorem says that the physical old
      primitive space is fixed pointwise by $\mathcal U_{\rm PF}$, and the
      row-(c) vector $Kf$ has the order mismatch proved in (62)--(63).
      Differentiated Poisson gives the divergence (96)--(98), not the signed
      normal equation for the physical harmonic minimizer.
\item Equation (56) is ill typed without (57), and it is false when its $K$ is
      interpreted as the row-(c) operator (62).
\item The auxiliary direct-sum G39 port cannot be the midpoint residual:
      Proposition (64)--(66) proves the required unitary orbit is impossible
      under that orthogonal embedding.
\item A full-space identification in (78) with the finite G39 port is ruled
      out by the dimension obstruction.  At most (78) can describe compatible
      finite compressions; it cannot close the continuum row (d).
\item The former identification
      $\mathcal B_T^*\mathcal B_T=\cL$ and
      $\|\mathcal B_TRv\|^2=S_M'$ was a redefinition, not a derivation from
      the supplied row-(d) operators; it has been withdrawn.
\item A direct sum of locally normalized Gamma/Euler isometries cannot supply
      the global contraction; theorem (89) gives an explicit primitive
      obstruction in every fixed arithmetic channel for large windows.
\item Positivity of the Sonin Hilbert--Schmidt Gram $M(f)$ does not identify
      it with the Gamma--Tate reference Gram and does not imply its open lower
      frame bound.
\item Three-lines convexity for the natural completed metric carries the
      wrong $|\Lambda_{\mathbb Q}|^2$ weight; cancelling it by division is a
      zero-free assumption.
\item The analogous attempt to cancel the Gamma characteristic factor in
      (177) is impossible on every nonzero compact-window source by (180),
      including on the primitive Tate subspace.
\item The Gamma ratio coherent orbit cannot replace the Gamma energy current:
      (202)--(204) prove that the former is bounded and the latter is an
      unbounded affine cocycle.
\item The gcd state limit and the continuous scaling-orbit limit cannot be
      identified as ordinary vectors: (185) proves escape of mass and an
      uncountable delta kernel.  They must be joined only after a polar/Tate
      quotient at form level.
\item The formal Krein completion and any square root of the desired Schur
      remainder remain equivalent to the unproved sign and cannot close G.
\item The scalar central-line expression (258) is only mnemonic.  The two
      Euler expansions have disjoint half-planes, while their common
      meromorphic continuation cancels by (270).  It may not be used as a
      Fourier multiplier or as a unique-continuation symbol.
\item Logarithmic-Laplacian unique continuation requires simultaneous local
      vanishing of an input and its logarithmic-Laplacian output.  Equation
      (257) supplies only cancellation of Gamma, constant and translated
      arithmetic outputs, so invoking that theorem directly would separate
      a zero sum without proof.
\item The contour argument (274)--(278) is not used as load-bearing closure:
      coefficientwise elimination of an infinite off-axis exponential
      residue expansion requires a Fourier-hyperfunction carrier-uniqueness
      theorem with explicit convergence and growth hypotheses.  The later
      closure criteria instead use the two ordinary Hardy half-planes and
      pole rigidity (283)--(292).
\item Optimal generic boundary regularity for a logarithmic Dirichlet
      problem occurs at the scale $\sqrt{1/|\log d|}$, which is not $H^1$;
      see (418).  This comparison does not constitute a perturbation theorem
      for the Gamma--Euler system, but it rules out citing generic
      logarithmic ellipticity as the missing translation-jet bootstrap.
\item Meyer's cutoff map $\pi_+$ is only a near-section for nuclear trace
      computations: (307)--(308) and Meyer's own statement show that it is
      not a projection onto the Poisson range.  It cannot be used to define
      the missing positive defect.
\item Compact support and logarithmic Gamma-form regularity do not imply
      membership in the nuclear weighted-Schwartz space $\mathcal H_-$.
      The domain gate (306a) must be proved in a physical completion before
      applying the smooth-source range theorem to a zero mode.
\item Equality of the two meromorphic Mellin germs in (319) is not equality
      of their boundary realizations: the initial convergence regions are
      disjoint.  Using the common germ as though it were already the physical
      Laplace transform would erase the jump carrier and assume the missing
      completion theorem.
\item A causal Laplace transform cannot be used in (317) without zero
      prehistory: translating a delayed term produces lower-limit history
      contributions.  The corrected theorem requires a translation-covariant
      bilateral realization, or the already explicit incoming-support
      hypothesis.
\item Co-Poisson/Sonin membership cannot be inferred from compact support:
      for $a<1$ its closed range is characterized by orthogonality to the
      zeta-zero evaluators.  Multiplicative smoothing cannot repair this
      universally because an exponential-type smoothing multiplier has only
      $O(R)$ zeros, not the required $\gg R\log R$.
\end{itemize}

\end{document}
